% --- CLASSE DOCUMENTO ---
\documentclass[a4paper,11pt]{article}

% --- LINGUA E CODIFICA ---
\usepackage[italian]{babel} 
% Nota: Compilare con XeLaTeX o LuaLaTeX per la gestione dei font OpenType

% --- GEOMETRIA E LAYOUT ---
\usepackage{geometry}
\geometry{a4paper, margin=1in}
\usepackage{parskip}      % Spazio tra i paragrafi senza rientro
\usepackage{setspace}     % Per l'interlinea
\usepackage{calc}         % Per calcoli di dimensioni (richiesto da Pandoc)

% --- FONT E MATEMATICA ---
\usepackage{amsmath, amssymb, amsthm}
\usepackage{bm}           % Grassetto matematico (compatibile con pdflatex)
% Nota: unicode-math richiede XeLaTeX o LuaLaTeX, non pdflatex
\usepackage{microtype}    % Micro-tipografia: riduce overflow di testo nei margini

% --- TESTO BARRATO ---
\usepackage{soul}         % Fornisce il comando \st{} per testo barrato
\usepackage[normalem]{ulem} % Fornisce \sout{} per testo barrato; normalem preserva \emph

% --- GRAFICA E COLORI ---
\usepackage[x11names, svgnames, table]{xcolor}
\usepackage{graphicx}
\usepackage{standalone}
\usepackage{fontawesome5}

% --- TIKZ E PGFPLOTS ---
\usepackage{tkz-euclide}
\usepackage{pgfplots}
\pgfplotsset{compat=1.18}
\usetikzlibrary{
    fit, 
    shapes.geometric, 
    decorations.pathmorphing, 
    arrows.meta, 
    positioning, 
    patterns, 
    shapes
}

% Stili TikZ per diagrammi a blocchi
\tikzset{
    input/.style={coordinate},
    output/.style={coordinate}
}

% --- SCATOLE E CORNICI ---
\usepackage[skins, breakable]{tcolorbox}
\usepackage{framed}
\usepackage{minted} % Richiede -shell-escape

% --- TABELLE E LISTE ---
\usepackage{booktabs, longtable, array}
\usepackage{enumitem}
\usepackage{caption}

% --- INTESTAZIONI E PIÈ DI PAGINA ---
\usepackage{fancyhdr}
\setlength{\headheight}{14pt} % Fix fancyhdr warning
\pagestyle{fancy}
\fancyhf{} % Pulisce tutto
\fancyhead[L]{\nouppercase{\rightmark}} 
\fancyhead[R]{\nouppercase{\leftmark}}  
\fancyfoot[C]{\thepage}
\renewcommand{\headrulewidth}{0.4pt}
\renewcommand{\footrulewidth}{0pt}

% Configura lo stile plain (pagine iniziali)
\fancypagestyle{plain}{
    \fancyhf{}
    \fancyfoot[C]{\thepage}
    \renewcommand{\headrulewidth}{0pt}
}

% --- COMPATIBILITÀ PANDOC (Markdown support) ---
\providecommand{\tightlist}{%
  \setlength{\itemsep}{0pt}\setlength{\parskip}{0pt}}

\makeatletter
\newcommand*\pandocbounded[1]{% scales image to fit in text height/width
  \sbox\pandoc@box{#1}%
  \Gscale@div\@tempa{\textheight}{\dimexpr\ht\pandoc@box+\dp\pandoc@box\relax}%
  \Gscale@div\@tempb{\linewidth}{\wd\pandoc@box}%
  \ifdim\@tempb\p@<\@tempa\p@\let\@tempa\@tempb\fi%
  \ifdim\@tempa\p@<\p@\scalebox{\@tempa}{\usebox\pandoc@box}%
  \else\usebox{\pandoc@box}%
  \fi%
}
\newsavebox\pandoc@box
\makeatother

% --- COLLEGAMENTI IPERTESTUALI ---
\usepackage{hyperref}
\hypersetup{
    pdftitle={Corso istruttori FIN 2025},
    hidelinks,
    breaklinks=true
}
\usepackage{bookmark}

% --- FIX UNIVERSALE TABELLE ---
\usepackage{array}
\usepackage{caption}
\usepackage{tabularx}

% 1. Migliora la spaziatura verticale (più respiro alle righe)
\renewcommand{\arraystretch}{1.5}

% 2. Spazio orizzontale tra le colonne (default=6pt, ridotto per adattarsi)
\setlength{\tabcolsep}{4pt}

% 3. Colori per le righe (opzionale, stile "Zebra")
\rowcolors{2}{gray!5}{white}

% Stile per le intestazioni (Grassetto e linea più spessa)
\renewcommand{\headrulewidth}{0.4pt}
\setlength{\heavyrulewidth}{1.2pt} % Linee sopra e sotto più marcate
\setlength{\lightrulewidth}{0.7pt} % Linea separazione header

\title{Corso istruttori FIN 2025}
\author{Alessandro Vulcu}
\date{}

\begin{document}
\maketitle

\section{1. Aspetti Amministrativi}\label{aspetti-amministrativi}

\subsection{Introduzione}\label{introduzione}

Questa lezione, tenuta da \textbf{Antonello Panza}, Segretario Generale
della Federazione Italiana Nuoto, tratta la \textbf{natura giuridica
      delle federazioni sportive italiane} e, in particolare, gli aspetti
fondamentali dello statuto della FIN che regolano la sua attività.
Conoscere questi aspetti è essenziale per ogni tecnico che opera
all\textquotesingle interno del sistema federale, poiché le norme
statutarie definiscono i diritti e gli obblighi di tutti i tesserati e
le società affiliate.

\begin{center}\rule{0.5\linewidth}{0.5pt}\end{center}

\subsection{La Natura Giuridica della
      FIN}\label{la-natura-giuridica-della-fin}

\subsubsection{La Riforma del 1999}\label{la-riforma-del-1999}

Prima del 1999 le federazioni sportive italiane erano \textbf{enti di
      diritto pubblico}, classificati nel cosiddetto \emph{parastato} (enti
non economici). Svolgevano prevalentemente attività pubblica, pur con
alcune attività di natura privatistica (organizzazione di campionati,
manifestazioni).

Con il \textbf{decreto legislativo n. 242 del 1999} (art. 15), il
legislatore ha trasformato le federazioni in:

\begin{quote}
      \textbf{Associazioni riconosciute senza fine di lucro di diritto
            privato}, rette dal codice civile e dalle relative disposizioni sugli
      enti a base associativa.
\end{quote}

Il mutamento è sintetizzabile così:

\begin{longtable}[]{@{}p{0.30\linewidth}p{0.30\linewidth}p{0.30\linewidth}@{}}
      \toprule\noalign{}
      Periodo                  & Natura giuridica                  & Tipo di attività \\
      \midrule\noalign{}
      \endhead
      \bottomrule\noalign{}
      \endlastfoot
      \textbf{Prima del 1999}  & Ente di diritto \textbf{pubblico} & Con alcune
      attività privatistiche (ente misto)                                             \\
      \textbf{Dal 1999 in poi} & Ente di diritto \textbf{privato}  & Con alcune
      attività che mantengono valenza pubblicistica                                   \\
\end{longtable}

\subsubsection{Attività che Mantengono Valenza
      Pubblicistica}\label{attivituxe0-che-mantengono-valenza-pubblicistica}

Nonostante la privatizzazione, alcune attività delle federazioni
conservano carattere pubblico. Le principali sono:

\begin{itemize}
      \tightlist
      \item
            \textbf{Ammissione e affiliazione} delle società sportive (e relativa
            revoca)
      \item
            \textbf{Controllo del regolare svolgimento} delle competizioni e dei
            campionati professionistici
      \item
            \textbf{Lotta al doping} (repressione e prevenzione)
      \item
            \textbf{Preparazione olimpica}
      \item
            \textbf{Gestione di impianti sportivi pubblici}
\end{itemize}

Per queste materie, in caso di contenzioso, la giurisdizione non è del
giudice ordinario ma del \textbf{giudice amministrativo} (\emph{TAR}).
Ad esempio, il diniego di un\textquotesingle affiliazione o la sua
revoca sono di competenza del TAR, non del tribunale civile.

Un ulteriore effetto pratico: le federazioni, per la valenza
pubblicistica nella gestione di impianti sportivi, possono instaurare
rapporti con gli enti locali come \textbf{rapporti tra pubbliche
      amministrazioni}, senza necessità di bando di gara per
l\textquotesingle individuazione del gestore.

\subsubsection{Il Bilancio e la Deroga al Codice
      Civile}\label{il-bilancio-e-la-deroga-al-codice-civile}

Le federazioni presentano un\textquotesingle ulteriore atipicità: a
differenza delle normali associazioni riconosciute, il bilancio non
viene approvato dall\textquotesingle assemblea degli affiliati, ma:

\begin{enumerate}
      \tightlist
      \item
            È approvato annualmente dal \textbf{Consiglio federale}
      \item
            È poi sottoposto all\textquotesingle approvazione della \textbf{Giunta
                  Nazionale del CONI}
      \item
            Solo in caso di mancata approvazione da parte del CONI si torna alle
            assemblee federali
\end{enumerate}

\begin{center}\rule{0.5\linewidth}{0.5pt}\end{center}

\subsection{La FIN: Finalità e
      Statuto}\label{la-fin-finalituxe0-e-statuto}

La FIN si configura come \textbf{associazione sportiva riconosciuta
      senza fine di lucro dilettantistica}. Il suo statuto la rende, di fatto,
una confederazione degli sport acquatici, data la pluralità di
discipline che raggruppa:

\begin{itemize}
      \tightlist
      \item
            Nuoto (compreso il nuoto di fondo)
      \item
            Tuffi (compresi i tuffi dalle grandi altezze)
      \item
            Pallanuoto
      \item
            Nuoto artistico (già nuoto sincronizzato)
      \item
            Nuoto per salvamento
\end{itemize}

\subsubsection{La Sezione Salvamento}\label{la-sezione-salvamento}

La sezione salvamento ha uno scopo statutario specifico e più ampio
rispetto alle altre discipline:

\begin{itemize}
      \tightlist
      \item
            Promuovere e disciplinare l\textquotesingle attività agonistica del
            salvamento
      \item
            Promuovere l\textquotesingle{}\textbf{insegnamento e
                  l\textquotesingle addestramento} alla prevenzione di incidenti in
            acqua
      \item
            Diffondere le pratiche sanitarie di \textbf{primo soccorso}
      \item
            Operare nell\textquotesingle ambito dei programmi del servizio
            sanitario e della protezione civile
      \item
            Tutela ambientale delle acque
\end{itemize}

La FIN è riconosciuta a livello internazionale come unico soggetto
italiano per l\textquotesingle attività di salvamento (le altre
organizzazioni presenti in Italia non hanno riconoscimento
internazionale).

\subsubsection{La Missione: la Cultura
      dell\textquotesingle Acquaticità}\label{la-missione-la-cultura-dellacquaticituxe0}

Una delle mission fondamentali della FIN è la riduzione delle
\textbf{morti per annegamento}, perseguita attraverso:

\begin{itemize}
      \tightlist
      \item
            La formazione di operatori della sicurezza in acqua
            (\textbf{assistenti bagnanti})
      \item
            L\textquotesingle aumento della pratica sportiva di base e della
            \textbf{cultura dell\textquotesingle acquaticità}
      \item
            La convinzione che più persone acquisiscono consapevolezza del
            corretto utilizzo dell\textquotesingle acqua, minore è il rischio di
            incidenti mortali
\end{itemize}

Le statistiche mostrano un \textbf{continuo decremento delle morti per
      annegamento} dal dopoguerra ai giorni nostri, anche grazie
all\textquotesingle attività della FIN.

\begin{center}\rule{0.5\linewidth}{0.5pt}\end{center}

\subsection{Le Organizzazioni Internazionali di
      Riferimento}\label{le-organizzazioni-internazionali-di-riferimento}

La FIN è affiliata alle seguenti organizzazioni:

\begin{longtable}[]{@{}p{0.30\linewidth}p{0.30\linewidth}p{0.30\linewidth}@{}}
      \toprule\noalign{}
      Sigla                                                      & Nome                                          & Ambito              \\
      \midrule\noalign{}
      \endhead
      \bottomrule\noalign{}
      \endlastfoot
      \textbf{LEN} (\emph{Ligue Européenne de Natation})         & Lega europea di
      nuoto                                                      & Sports natatori a livello europeo                                   \\
      \textbf{COMEN}                                             & Confederazione delle federazioni mediterranee & Sport
      acquatici dell\textquotesingle area mediterranea                                                                                 \\
      \textbf{COLAN}                                             & Confederazione dei paesi latini               & Sport acquatici dei
      paesi latini                                                                                                                     \\
      \textbf{ILS} (\emph{International Life Saving Federation}) & Federazione
      mondiale del salvamento                                    & Attività di formazione e sport di salvamento
      (mondiale)                                                                                                                       \\
      \textbf{ILS Europe}                                        & Federazione europea del salvamento            & Attività di
      formazione e sport di salvamento (europea)                                                                                       \\
\end{longtable}

\begin{quote}
      Per statuto internazionale (ILS e ILS Europe), ogni paese può avere un
      solo soggetto riconosciuto per l\textquotesingle attività di salvamento.
      Per l\textquotesingle Italia, questo soggetto è la FIN.
\end{quote}

\begin{center}\rule{0.5\linewidth}{0.5pt}\end{center}

\subsection{I Soggetti della FIN}\label{i-soggetti-della-fin}

\subsubsection{Persone Giuridiche: Società e
      Associazioni}\label{persone-giuridiche-societuxe0-e-associazioni}

Le società e le associazioni sportive si \textbf{affiliano annualmente}
alla FIN. Le federazioni sono definite dal CONI \textbf{organismi
      affilianti}: hanno il compito di verificare i requisiti dei soggetti che
intendono affiliarsi e di iscriverli nel registro delle società sportive
(attualmente tenuto presso il CONI).

Questa funzione di verifica è continua: le federazioni devono anche
controllare le modifiche statutarie che le società operano nel corso del
tempo.

\subsubsection{Persone Fisiche: i
      Tesserati}\label{persone-fisiche-i-tesserati}

I tesserati riconosciuti dallo statuto FIN sono:

\begin{itemize}
      \tightlist
      \item
            \textbf{Atleti}
      \item
            \textbf{Dirigenti federali}
      \item
            \textbf{Dirigenti sociali}
      \item
            \textbf{Tecnici}
      \item
            \textbf{Ufficiali di gara}
      \item
            \textbf{Medici federali e sociali}
\end{itemize}

\subsubsection{Diritti e Obblighi dei
      Tesserati}\label{diritti-e-obblighi-dei-tesserati}

\textbf{Obblighi:}

\begin{itemize}
      \tightlist
      \item
            Rispettare le norme statutarie e regolamentari
      \item
            Rispettare le decisioni degli organi disciplinari
\end{itemize}

\textbf{Diritti:}

\begin{itemize}
      \tightlist
      \item
            Partecipare all\textquotesingle attività federale nel rispetto delle
            norme tecniche
      \item
            Partecipare alle assemblee elettive federali
      \item
            Avvalersi degli organi disciplinari in caso di privazione ingiusta dei
            propri diritti
\end{itemize}

\begin{center}\rule{0.5\linewidth}{0.5pt}\end{center}

\subsection{Gli Organi Federali}\label{gli-organi-federali}

Gli organi federali si dividono in \textbf{organi centrali} e
\textbf{organi periferici}.

\subsubsection{Organi Centrali}\label{organi-centrali}

\begin{itemize}
      \tightlist
      \item
            \textbf{Assemblea nazionale delle società}: si riunisce ogni \textbf{4
                  anni}; funzione principale è il rinnovo degli organi gestionali
            federali; è il consesso più importante della federazione
      \item
            \textbf{Presidente della federazione}
      \item
            \textbf{Consiglio federale}
      \item
            \textbf{Segretario Generale}
      \item
            \textbf{Collegio dei Revisori dei Conti}
\end{itemize}

\subsubsection{Organi Periferici}\label{organi-periferici}

\begin{itemize}
      \tightlist
      \item
            \textbf{Assemblee regionali} (eleggono i comitati territoriali)
      \item
            \textbf{Comitati regionali}
      \item
            \textbf{Collegi regionali dei Revisori dei Conti}
      \item
            \textbf{Presidenti regionali e provinciali}
      \item
            \textbf{Delegati locali}
\end{itemize}

\begin{center}\rule{0.5\linewidth}{0.5pt}\end{center}

\subsection{Il Sistema della Giustizia
      Sportiva}\label{il-sistema-della-giustizia-sportiva}

\subsubsection{Principio di Separazione dei
      Poteri}\label{principio-di-separazione-dei-poteri}

Gli organi di giustizia sportiva operano in \textbf{piena autonomia}
rispetto agli organi politici di gestione federale (Consiglio federale,
Presidente). Questo principio è garantito da una norma fondamentale: i
giudici sportivi \textbf{non sono revocabili} dal Consiglio federale
durante il mandato per il quale sono stati nominati (a differenza di
qualsiasi altra carica federale). Solo la Commissione federale di
garanzia può revocarli in casi di grave e conclamata negligenza.

\subsubsection{Gli Organi di Giustizia
      Sportiva}\label{gli-organi-di-giustizia-sportiva}

\begin{longtable}[]{@{}p{0.47\linewidth}p{0.47\linewidth}@{}}
      \toprule\noalign{}
      Organo                                    & Funzione principale                       \\
      \midrule\noalign{}
      \endhead
      \bottomrule\noalign{}
      \endlastfoot
      \textbf{Commissione federale di garanzia} & Seleziona la lista dei
      giudici idonei; può revocare i giudici per gravi inadempienze;
      interpreta norme statutarie e regolamentari                                           \\
      \textbf{Giudice arbitro}                  & Giudice di gara; prende provvedimenti sul
      campo (\emph{on field})                                                               \\
      \textbf{Giudice sportivo nazionale}       & Omologa risultati; emette
      squalifiche in primo grado su verbali di gara                                         \\
      \textbf{Giudice sportivo regionale}       & Stessa competenza del nazionale, a
      livello regionale                                                                     \\
      \textbf{Corte sportiva di appello}        & Secondo grado: giudica i ricorsi
      contro le decisioni dei giudici sportivi                                              \\
      \textbf{Procura federale}                 & Avvia procedimenti disciplinari per
      violazioni non legate alla pratica sportiva (tesseramento irregolare,
      ecc.); deferito al Tribunale federale                                                 \\
      \textbf{Tribunale federale}               & Primo grado per i procedimenti avviati
      dalla Procura federale                                                                \\
      \textbf{Corte federale di appello}        & Secondo grado per i procedimenti
      del Tribunale federale                                                                \\
\end{longtable}

\subsubsection{La Procura Generale dello Sport
      (CONI)}\label{la-procura-generale-dello-sport-coni}

Negli ultimi anni è stata istituita una \textbf{Procura Generale dello
      Sport} presso il CONI, che affianca le procure federali e ha poteri
autonomi:

\begin{itemize}
      \tightlist
      \item
            Viene obbligatoriamente notiziata di ogni procedimento disciplinare
      \item
            Può richiedere alla Procura federale un\textquotesingle integrazione
            delle indagini anche in caso di archiviazione
      \item
            Può \textbf{avocare l\textquotesingle azione disciplinare} se la
            Procura federale ha archiviato un caso che la Procura del CONI ritiene
            invece perseguibile
\end{itemize}

\subsubsection{Gradi di Giudizio e
      Garanzie}\label{gradi-di-giudizio-e-garanzie}

Il sistema garantisce il \textbf{doppio grado di giudizio} come diritto
inalienabile. Per le sanzioni più gravi è previsto un terzo grado:

\begin{longtable}[]{@{}p{0.47\linewidth}p{0.47\linewidth}@{}}
      \toprule\noalign{}
      Grado                             & Organo                            \\
      \midrule\noalign{}
      \endhead
      \bottomrule\noalign{}
      \endlastfoot
      \textbf{1° grado} (endo-federale) & Giudice sportivo / Tribunale
      federale                                                              \\
      \textbf{2° grado} (endo-federale) & Corte sportiva di appello / Corte
      federale di appello                                                   \\
      \textbf{3° grado} (eso-federale)  & Collegio di garanzia dello sport
      (CONI)                                                                \\
      \textbf{Ultima istanza}           & TAR del Lazio                     \\
\end{longtable}

\textbf{Condizioni per accedere al Collegio di garanzia del CONI:}

\begin{itemize}
      \tightlist
      \item
            Sanzioni tecnico-sportive superiori a \textbf{90 giorni}
      \item
            Sanzioni pecuniarie superiori a \textbf{10.000 euro}
\end{itemize}

Per sanzioni inferiori a queste soglie non è possibile il ricorso al
Collegio di garanzia.

\textbf{Garanzie del giusto processo:}

\begin{itemize}
      \tightlist
      \item
            Piena tutela dei diritti e degli interessi dei tesserati
      \item
            Parità delle parti (accusa e difesa)
      \item
            Diritto al contraddittorio: il deferito può presentare memorie
            difensive
      \item
            Possibilità di \textbf{recusazione del giudice} in caso di
            incompatibilità
      \item
            Possibilità di \textbf{revisione del giudizio} se sono stati violati i
            diritti procedurali
\end{itemize}

\subsubsection{L\textquotesingle Arbitrato
      Federale}\label{larbitrato-federale}

Per controversie non di natura disciplinare tra affiliati, è disponibile
l\textquotesingle{}\textbf{arbitrato federale}: uno strumento di
composizione delle controversie che evita il ricorso agli organi di
giustizia sportiva.

\begin{center}\rule{0.5\linewidth}{0.5pt}\end{center}

\subsection{I Regolamenti Federali}\label{i-regolamenti-federali}

I regolamenti della FIN sono tutti pubblicati sul sito
\textbf{federnuoto.it}, in osservanza delle norme sulla trasparenza
amministrativa. I principali sono:

\begin{itemize}
      \tightlist
      \item
            \textbf{Statuto federale}
      \item
            \textbf{Regolamento organico}
      \item
            \textbf{Regolamento di giustizia}
      \item
            \textbf{Regolamento antidoping}
      \item
            \textbf{Regolamento dei tecnici}
      \item
            Regolamento delle squadre nazionali, del settore arbitrale, degli
            arbitrati federali, amministrativo
\end{itemize}

\begin{quote}
      Quattro regolamenti richiedono l\textquotesingle{}\textbf{approvazione
            obbligatoria del CONI}: Statuto, Regolamento organico, Regolamento di
      giustizia e Regolamento antidoping.
\end{quote}

La FIN aderisce inoltre al sistema europeo \textbf{SNAQ} (\emph{Sistema
      Nazionale di qualificazione}): le qualifiche tecniche rilasciate dalla
federazione hanno riconoscimento automatico sul territorio europeo, in
quanto i programmi formativi sono allineati agli standard
internazionali.

\begin{center}\rule{0.5\linewidth}{0.5pt}\end{center}

\subsection{Conclusione - Punti Chiave per
      l\textquotesingle Esame}\label{conclusione---punti-chiave-per-lesame}

\begin{itemize}
      \tightlist
      \item
            Prima del 1999 le federazioni erano \textbf{enti di diritto pubblico};
            dal 1999 sono \textbf{associazioni riconosciute senza fine di lucro di
                  diritto privato} (D.lgs. 242/1999, art. 15).
      \item
            Le federazioni sono oggi \textbf{enti privati con alcune attività
                  pubblicistiche}: affiliazione/revoca, controllo campionati
            professionistici, anti-doping, preparazione olimpica, gestione
            impianti.
      \item
            Per le materie pubblicistiche, il giudice competente in caso di
            controversia è il \textbf{TAR} (giudice amministrativo), non il
            giudice ordinario.
      \item
            Il \textbf{bilancio} è approvato dal Consiglio federale e poi
            sottoposto all\textquotesingle approvazione della \textbf{Giunta
                  Nazionale del CONI}.
      \item
            La FIN raggruppa: \textbf{nuoto}, \textbf{tuffi}, \textbf{pallanuoto},
            \textbf{nuoto artistico}, \textbf{nuoto di fondo}, \textbf{salvamento}
            (inclusi tuffi dalle grandi altezze).
      \item
            Per l\textquotesingle attività di salvamento, la FIN è
            l\textquotesingle{}\textbf{unico soggetto italiano riconosciuto} da
            ILS e ILS Europe.
      \item
            Organizzazioni internazionali: \textbf{LEN} (Europa), \textbf{COMEN}
            (Mediterraneo), \textbf{COLAN} (paesi latini), \textbf{ILS}
            (mondiale), \textbf{ILS Europe}.
      \item
            Le federazioni sono \textbf{organismi affilianti}: verificano i
            requisiti di società e associazioni sportive per
            l\textquotesingle iscrizione al registro CONI.
      \item
            Tesserati riconosciuti: atleti, dirigenti federali, dirigenti sociali,
            \textbf{tecnici}, ufficiali di gara, medici.
      \item
            \textbf{Organi centrali}: Assemblea nazionale (ogni 4 anni),
            Presidente, Consiglio federale, Segretario Generale, Collegio dei
            Revisori dei Conti.
      \item
            Principio fondamentale: \textbf{separazione tra potere gestionale e
                  potere giudiziario}. I giudici sportivi non sono revocabili dal
            Consiglio federale.
      \item
            Organi di giustizia: Commissione federale di garanzia, Giudice
            arbitro, Giudice sportivo nazionale/regionale, Corte sportiva di
            appello, Procura federale, Tribunale federale, Corte federale di
            appello.
      \item
            Il \textbf{Collegio di garanzia del CONI} (3° grado) è accessibile
            solo per sanzioni \textbf{\textgreater90 giorni} o
            \textbf{\textgreater10.000 euro}.
      \item
            Quattro regolamenti richiedono approvazione CONI: \textbf{Statuto,
                  Regolamento organico, Regolamento di giustizia, Regolamento
                  antidoping}.
      \item
            La FIN aderisce al sistema \textbf{SNAQ}: le qualifiche tecniche sono
            riconosciute in tutta Europa.
\end{itemize}

\section{2. Il Corpo Umano - Parte I}\label{il-corpo-umano---parte-i}

\subsection{Introduzione}\label{introduzione-1}

Questa lezione è tenuta dal \textbf{dott. Lorenzo Baruggo}, medico
federale della Federazione Italiana Nuoto, ex nuotatore delle squadre
nazionali negli anni \textquotesingle70 e poi medico delle squadre
nazionali dagli anni \textquotesingle80. L\textquotesingle argomento è
l\textquotesingle{}\textbf{anatomia e la fisiologia del corpo umano},
con un focus sugli apparati e i sistemi direttamente collegati al
movimento, alla produzione di energia e alla respirazione. La lezione è
la prima di due parti: questa si concentra sul sistema nervoso e
sull\textquotesingle apparato respiratorio; la seconda tratterà
l\textquotesingle apparato muscolare e l\textquotesingle apparato
locomotore.

\begin{center}\rule{0.5\linewidth}{0.5pt}\end{center}

\subsection{Le Basi della Vita: dalla Cellula
      all\textquotesingle Organismo}\label{le-basi-della-vita-dalla-cellula-allorganismo}

\subsubsection{La Vita sulla Terra}\label{la-vita-sulla-terra}

L\textquotesingle universo è composto quasi esclusivamente da materia
inorganica (minerali e gas). La \textbf{materia organica} è
caratteristica della Terra e si distingue per:

\begin{itemize}
      \tightlist
      \item
            Struttura prevalentemente \textbf{proteica}
      \item
            Presenza di \textbf{grassi} e \textbf{zuccheri}
      \item
            Esistenza in un \textbf{ambiente acquatico}
\end{itemize}

Le prime forme di vita sulla Terra erano \textbf{unicellulari} (alghe o
simili): in un unico involucro era contenuto tutto il necessario per
nutrirsi, riprodursi, respirare, produrre energia e muoversi.

\subsubsection{Organismi Multicellulari}\label{organismi-multicellulari}

Nell\textquotesingle evoluzione si è passati agli \textbf{esseri viventi
      multicellulari}, che si formano principalmente
dall\textquotesingle unione di due \textbf{gameti} (maschile e
femminile), da cui origina una prima cellula che si divide
progressivamente. Le cellule si differenziano in \textbf{stipiti
      cellulari} specializzati, ognuno con una funzione diversa, fino a
formare tutti gli organi e gli apparati. L\textquotesingle intera
informazione per questa differenziazione è scritta nel \textbf{DNA},
contenuto nel nucleo cellulare.

\begin{quote}
      Esiste anche la \textbf{partenogenesi} (riproduzione per autoclonazione
      senza unione di gameti), ma in questo caso l\textquotesingle essere
      prodotto è geneticamente identico al genitore: nessuna variazione,
      nessuna progressione della specie.
\end{quote}

\begin{center}\rule{0.5\linewidth}{0.5pt}\end{center}

\subsection{La Cellula}\label{la-cellula}

La cellula è l\textquotesingle unità fondamentale della vita, sia negli
unicellulari che nei multicellulari. Le sue strutture principali sono:

\subsubsection{Membrana Cellulare}\label{membrana-cellulare}

L\textquotesingle involucro esterno della cellula è composto da
\textbf{proteine} e \textbf{fosfolipidi}, con un polo \emph{lipofilo}
(affine ai grassi) e un polo \emph{idrofilo} (affine
all\textquotesingle acqua). È una struttura altamente selettiva: regola
con precisione cosa entra e cosa esce dalla cellula.

\subsubsection{Nucleo}\label{nucleo}

È il "cervello" della cellula. Contiene il \textbf{DNA}, il patrimonio
genetico che determina cosa la cellula è e cosa deve fare.

\subsubsection{Citoplasma e Organuli}\label{citoplasma-e-organuli}

Nel citoplasma sono presenti diversi organuli con funzioni specifiche:

\begin{itemize}
      \tightlist
      \item
            \textbf{Organuli proteici}: producono proteine
      \item
            \textbf{Mitocondri}: producono \textbf{energia}
      \item
            \textbf{Organuli difensivi}: funzione di protezione
\end{itemize}

Le cellule specializzate hanno morfologie e funzioni proprie. Esempio:
le \textbf{cellule muscolari} contengono miofibrille con proteine
contrattili specifiche (actina, miosina) e un numero elevato di
mitocondri per la produzione di energia.

\begin{center}\rule{0.5\linewidth}{0.5pt}\end{center}

\subsection{I Sistemi Connessi con il
      Movimento}\label{i-sistemi-connessi-con-il-movimento}

Quando l\textquotesingle organismo è completamente sviluppato, i sistemi
direttamente coinvolti nel movimento sono quattro:

\begin{enumerate}
      \tightlist
      \item
            \textbf{Sistema nervoso centrale e periferico}: da cui partono gli
            ordini motori e dove arrivano le informazioni dalla periferia
      \item
            \textbf{Apparato respiratorio}: prende l\textquotesingle ossigeno
            dall\textquotesingle esterno e porta via il CO2
      \item
            \textbf{Apparato cardiocircolatorio}: trasporta
            l\textquotesingle ossigeno ai tessuti e rimuove i prodotti di scarto
      \item
            \textbf{Apparato osteoarticolare + apparato muscolare}: gli effettori
            del movimento
\end{enumerate}

\begin{center}\rule{0.5\linewidth}{0.5pt}\end{center}

\subsection{Le Fonti di Energia}\label{le-fonti-di-energia}

Il muscolo per contrarsi ha bisogno di energia. Le riserve energetiche
nel corpo umano sono:

\begin{longtable}[]{@{}p{0.30\linewidth}p{0.30\linewidth}p{0.30\linewidth}@{}}
      \toprule\noalign{}
      Substrato         & Forma di deposito             & Dove                         \\
      \midrule\noalign{}
      \endhead
      \bottomrule\noalign{}
      \endlastfoot
      \textbf{Zuccheri} & \emph{Glicogeno}              & Fegato e muscolo             \\
      \textbf{Grassi}   & Tessuto adiposo               & Prevalentemente sottocutaneo \\
      \textbf{Proteine} & Solo circolanti e strutturali & Nessun deposito
      specifico                                                                        \\
\end{longtable}

\begin{quote}
      Le proteine vengono utilizzate a scopo energetico solo in situazioni
      particolari.
\end{quote}

\subsubsection{I Meccanismi Energetici}\label{i-meccanismi-energetici}

\begin{longtable}[]{@{}p{0.30\linewidth}p{0.30\linewidth}p{0.30\linewidth}@{}}
      \toprule\noalign{}
      Meccanismo                     & Substrato                                        & Caratteristiche \\
      \midrule\noalign{}
      \endhead
      \bottomrule\noalign{}
      \endlastfoot
      \textbf{Anaerobico alattacido} & Fosfocreatina + ATP                              & Molto rapido,
      brevissima durata; non produce acido lattico                                                        \\
      \textbf{Anaerobico lattacido}  & Zuccheri → acido lattico                         & Rapido,
      durata limitata; produce acido lattico                                                              \\
      \textbf{Aerobico}              & Grassi, zuccheri, (proteine) + \textbf{ossigeno} &
      Principale meccanismo a lungo termine; richiede O2 come comburente                                  \\
\end{longtable}

Il meccanismo aerobico è quello predominante, ma richiede un
\textbf{comburente esterno}: l\textquotesingle ossigeno.
L\textquotesingle organismo non ha riserve significative di ossigeno
(anche dopo iperventilazione dura pochi minuti), quindi deve essere
\textbf{continuamente rifornito} dall\textquotesingle esterno.

\begin{center}\rule{0.5\linewidth}{0.5pt}\end{center}

\subsection{Il Sistema Nervoso}\label{il-sistema-nervoso}

\subsubsection{Sistema Nervoso Centrale
      (SNC)}\label{sistema-nervoso-centrale-snc}

Il SNC è contenuto e protetto in due strutture ossee:

\begin{itemize}
      \tightlist
      \item
            \textbf{Scatola cranica}: contiene il cervello e il cervelletto
      \item
            \textbf{Colonna vertebrale}: contiene il \textbf{midollo spinale}, che
            è un prolungamento del SNC
\end{itemize}

\subsubsection{Sistema Nervoso
      Periferico}\label{sistema-nervoso-periferico}

Dal midollo spinale escono ed entrano i \textbf{nervi spinali}, che si
dividono in:

\begin{itemize}
      \tightlist
      \item
            \textbf{Motoneuroni} (vie efferenti): portano gli ordini dal SNC ai
            muscoli per comandare la contrazione
      \item
            \textbf{Nervi della sensibilità} (vie afferenti): riportano al
            cervello informazioni su dolore, temperatura, pressione, posizione
            nello spazio
\end{itemize}

\subsubsection{Informazioni Coscienti e
      Subcoscienti}\label{informazioni-coscienti-e-subcoscienti}

Non tutte le informazioni sensoriali arrivano alla \textbf{corteccia
      cerebrale} (coscienza):

\begin{itemize}
      \tightlist
      \item
            Le informazioni che raggiungono la coscienza permettono \textbf{azioni
                  volontarie} (decidere di muovere un arto, forzare la respirazione)
      \item
            Le informazioni che si fermano al \textbf{cervelletto} vengono
            elaborate per mantenere la postura e l\textquotesingle equilibrio
            \textbf{in modo automatico}, senza coinvolgere la coscienza
\end{itemize}

Esempio pratico: stare in piedi su una superficie di appoggio minima è
un compito biomeccanico complesso, eppure non richiede attenzione
consapevole. Il cervelletto elabora in continuo i segnali provenienti
da:

\begin{itemize}
      \tightlist
      \item
            Orecchio interno (equilibrio della testa)
      \item
            Pressione dei piedi
      \item
            Vista (posizione nello spazio)
\end{itemize}

...e coordina i muscoli di arti inferiori e tronco per mantenere la
postura eretta senza che la persona ci pensi.

\subsubsection{Sistema Nervoso Autonomo}\label{sistema-nervoso-autonomo}

È un sistema nervoso \textbf{periferico che si autoregola} senza passare
attraverso la coscienza. Controlla le funzioni degli organi interni:
intestino, stomaco, reni, bronchi, vasi sanguigni, ghiandole sudoripare,
cuore.

È bilanciato tra due sottosistemi opposti:

\begin{longtable}[]{@{}p{0.47\linewidth}p{0.47\linewidth}@{}}
      \toprule\noalign{}
      Sottosistema                           & Effetti principali             \\
      \midrule\noalign{}
      \endhead
      \bottomrule\noalign{}
      \endlastfoot
      \textbf{Sistema nervoso simpatico}     & Aumenta la frequenza cardiaca;
      dilata i bronchi; stimola la sudorazione; risposta allo stress/paura    \\
      \textbf{Sistema nervoso parasimpatico} & Riduce la frequenza cardiaca;
      contrasta gli effetti simpatici                                         \\
\end{longtable}

\textbf{Stimolanti del sistema simpatico:}

\begin{itemize}
      \tightlist
      \item
            \textbf{Caffeina} (caffè) → aumento dei battiti cardiaci
      \item
            \textbf{Stress, paura, ansia} → gli ormoni delle ghiandole surrenali
            stimolano il simpatico ("batticuore" pre-gara)
      \item
            Alcuni farmaci
\end{itemize}

\textbf{Implicazioni cliniche pratiche:}

\begin{itemize}
      \tightlist
      \item
            Negli \textbf{asmatici}: i recettori bronchiali reagiscono in modo
            eccessivo con restringimento delle vie aeree; si usano
            \textbf{broncodilatatori} che agiscono direttamente o indirettamente
            sulla muscolatura bronchiale
      \item
            Gli \textbf{antiaritmici} regolano il battito cardiaco agendo sul
            sistema autonomo
\end{itemize}

\subsubsection{Il Movimento Volontario e l\textquotesingle Apprendimento
      Motorio}\label{il-movimento-volontario-e-lapprendimento-motorio}

Per imparare un gesto nuovo, lo stimolo deve passare obbligatoriamente
attraverso la \textbf{coscienza} (corteccia cerebrale). Il nuoto è un
esempio tipico: non è un movimento naturale, deve essere imparato.

\begin{quote}
      "Una volta si diceva che il nuoto è un\textquotesingle arte perché
      bisogna imparare a farlo."
\end{quote}

Una volta che il cervello ha \textbf{memorizzato} quel gesto motorio, il
movimento può avvenire in modalità \textbf{semi-automatica}:

\begin{itemize}
      \tightlist
      \item
            Il nuotatore esperto nuota in modo ritmico e continuo senza dover
            pensare a ogni gesto
      \item
            Interviene la coscienza solo nei momenti che richiedono una
            \textbf{decisione non ripetitiva}: es. quando iniziare la virata,
            quando accelerare, come gestire l\textquotesingle affaticamento
\end{itemize}

\begin{center}\rule{0.5\linewidth}{0.5pt}\end{center}

\subsection{L\textquotesingle Apparato
      Respiratorio}\label{lapparato-respiratorio}

\subsubsection{Struttura delle Vie
      Aeree}\label{struttura-delle-vie-aeree}

L\textquotesingle aria può entrare attraverso due vie:

\begin{longtable}[]{@{}p{0.47\linewidth}p{0.47\linewidth}@{}}
      \toprule\noalign{}
      Via                 & Caratteristiche                                   \\
      \midrule\noalign{}
      \endhead
      \bottomrule\noalign{}
      \endlastfoot
      \textbf{Via nasale} & Più lenta; riscalda e umidifica
      l\textquotesingle aria prima dell\textquotesingle ingresso nella cassa
      toracica                                                                \\
      \textbf{Via orale}  & Più veloce; maggior flusso d\textquotesingle aria
      nell\textquotesingle unità di tempo; preferita sotto sforzo             \\
\end{longtable}

Le due vie si uniscono a livello del \textbf{faringe}, poi continuano
attraverso:

\textbf{Faringe → Laringe} (con le \textbf{corde vocali}, organo della
fonazione) \textbf{→ Trachea} (sostenuta da anelli cartilaginei, sempre
aperta) \textbf{→ Bronchi principali} (destro e sinistro) \textbf{→
      Bronchi secondari → Bronchioli → Bronchioli terminali → Alveoli}

\begin{quote}
      La \textbf{trachea} non collabisce mai, a differenza
      dell\textquotesingle esofago (la via digestiva), che è floscio quando
      vuoto.
\end{quote}

Gli \textbf{alveoli} sono piccoli sacchetti a parete sottilissima a
stretto contatto con i capillari sanguigni: è qui che avvengono gli
scambi gassosi:

\begin{itemize}
      \tightlist
      \item
            \textbf{O2} passa dall\textquotesingle aria al sangue
      \item
            \textbf{CO2} passa dal sangue all\textquotesingle aria
\end{itemize}

\begin{quote}
      A differenza dell\textquotesingle apparato cardiovascolare (che ha vie
      separate per andata e ritorno del sangue), l\textquotesingle apparato
      respiratorio usa \textbf{la stessa via} per l\textquotesingle aria in
      entrata e in uscita.
\end{quote}

\subsubsection{Meccanica Respiratoria}\label{meccanica-respiratoria}

Il movimento dell\textquotesingle aria è determinato dalle variazioni di
pressione all\textquotesingle interno della cassa toracica:

\textbf{Inspirazione} (pressione negativa):

\begin{itemize}
      \tightlist
      \item
            Il \textbf{diaframma} si contrae e si abbassa → aumenta il diametro
            longitudinale del polmone
      \item
            I \textbf{muscoli intercostali} si contraggono e alzano le coste →
            aumenta il diametro trasversale
      \item
            Si crea una pressione interna inferiore a quella esterna →
            l\textquotesingle aria entra
\end{itemize}

\textbf{Espirazione a riposo} (passiva):

\begin{itemize}
      \tightlist
      \item
            I muscoli si rilasciano
      \item
            La \textbf{forza elastica} accumulata fa aumentare la pressione
            interna → l\textquotesingle aria esce spontaneamente
\end{itemize}

\textbf{Espirazione forzata} (attiva):

\begin{itemize}
      \tightlist
      \item
            Si reclutano muscoli specifici per espellere l\textquotesingle aria
            più rapidamente e in quantità maggiore
\end{itemize}

\textbf{Respirazione molto affannosa}: entrano in funzione anche
\textbf{muscoli accessori} (collo, tronco superiore, es.
sternocleidomastoideo) per massimizzare l\textquotesingle espansione
toracica.

\subsubsection{Parametri Respiratori}\label{parametri-respiratori}

\begin{longtable}[]{@{}p{0.22\linewidth}p{0.22\linewidth}p{0.22\linewidth}p{0.22\linewidth}@{}}
      \toprule\noalign{}
      Condizione                & Atti respiratori/min & Volume per atto                    & Ventilazione
      totale                                                                                               \\
      \midrule\noalign{}
      \endhead
      \bottomrule\noalign{}
      \endlastfoot
      \textbf{A riposo}         & \textasciitilde12    & \textasciitilde500 cc (0,5 L)      &
      \textasciitilde6 L/min                                                                               \\
      \textbf{Sforzo massimale} & 40-50                & 3-5 L (atleti fino a \textgreater5
      L)                        & 150+ L/min                                                               \\
\end{longtable}

\begin{quote}
      La ventilazione massimale è \textbf{20-30 volte} superiore a quella a
      riposo.
\end{quote}

\subsubsection{Debito di Ossigeno e
      Affanno}\label{debito-di-ossigeno-e-affanno}

\begin{itemize}
      \tightlist
      \item
            Quando si lavora in modalità \textbf{completamente aerobica}, la
            respirazione torna alla normalità quasi immediatamente dopo lo sforzo
      \item
            Quando si lavora anche in modalità \textbf{anaerobica}, si contrae un
            \textbf{debito di ossigeno}: il metabolismo anaerobico produce
            prodotti che devono poi essere "pagati" attraverso il metabolismo
            aerobico
      \item
            Dopo lo sforzo, la respirazione rimane \textbf{affannosa} per smaltire
            questo debito: più è stato intenso il lavoro anaerobico, più a lungo
            dura l\textquotesingle affanno
\end{itemize}

\begin{quote}
      Per un tecnico attento, la \textbf{modalità e la durata
            dell\textquotesingle affanno} sono un indicatore diretto del grado di
      affaticamento dell\textquotesingle atleta sotto sforzo.
\end{quote}

\begin{center}\rule{0.5\linewidth}{0.5pt}\end{center}

\subsection{Conclusione - Punti Chiave per
      l\textquotesingle Esame}\label{conclusione---punti-chiave-per-lesame-1}

\begin{itemize}
      \tightlist
      \item
            La vita sulla Terra è a struttura prevalentemente \textbf{proteica} in
            ambiente acquatico; gli organismi multicellulari si formano
            dall\textquotesingle unione di due gameti e dalla differenziazione
            cellulare guidata dal \textbf{DNA}.
      \item
            La \textbf{membrana cellulare} (proteine + fosfolipidi) è selettiva;
            il \textbf{nucleo} contiene il DNA; i \textbf{mitocondri} producono
            energia.
      \item
            I sistemi connessi con il movimento sono: \textbf{SNC/periferico},
            \textbf{respiratorio}, \textbf{cardiocircolatorio},
            \textbf{osteoarticolare + muscolare}.
      \item
            Riserve energetiche: \textbf{glicogeno} (fegato e muscolo),
            \textbf{grassi} (tessuto adiposo), \textbf{proteine} (nessun deposito
            specifico).
      \item
            I tre meccanismi energetici: \textbf{anaerobico alattacido}
            (fosfocreatina + ATP), \textbf{anaerobico lattacido} (zuccheri → acido
            lattico), \textbf{aerobico} (richiede O2 come comburente, non ha
            riserve significative).
      \item
            Il \textbf{midollo spinale} è contenuto nella colonna vertebrale ed è
            un prolungamento del SNC; i \textbf{motoneuroni} portano ordini ai
            muscoli.
      \item
            Le informazioni che arrivano alla \textbf{corteccia cerebrale} sono
            coscienti (azioni volontarie); quelle che si fermano al
            \textbf{cervelletto} controllano postura ed equilibrio in modo
            automatico.
      \item
            Il \textbf{sistema nervoso autonomo} (simpatico/parasimpatico)
            controlla organi interni senza passare per la coscienza.
      \item
            \textbf{Simpatico}: aumenta frequenza cardiaca, dilata bronchi;
            stimolato da caffeina, stress, ormoni surrenali.
      \item
            \textbf{Parasimpatico}: riduce frequenza cardiaca.
      \item
            Il nuoto è un gesto \textbf{non naturale}: deve passare per la
            coscienza; una volta memorizzato diventa \textbf{semi-automatico}.
      \item
            L\textquotesingle aria inspirata segue il percorso: \textbf{naso/bocca
                  → faringe → laringe → trachea → bronchi → bronchioli → alveoli}.
      \item
            Negli alveoli avvengono gli scambi: \textbf{O2 dal sangue, CO2 al
                  sangue}.
      \item
            \textbf{Ispirazione} = pressione negativa (diaframma + intercostali);
            \textbf{espirazione a riposo} = passiva (forza elastica).
      \item
            A riposo: \textasciitilde{}\textbf{12 atti/min},
            \textasciitilde{}\textbf{0,5 L per atto} = \textasciitilde{}\textbf{6
                  L/min}; sforzo massimale: 40-50 atti/min, fino a 5 L per atto =
            \textbf{150+ L/min} (20-30x).
      \item
            L\textquotesingle{}\textbf{affanno prolungato} dopo lo sforzo indica
            un debito di ossigeno accumulato durante il lavoro anaerobico: è un
            indicatore di affaticamento utile per il tecnico.
\end{itemize}

\section{3. Il Corpo Umano - Parte II}\label{il-corpo-umano---parte-ii}

\subsection{Introduzione}\label{introduzione-2}

Seconda e ultima parte della lezione di anatomia e fisiologia tenuta dal
\textbf{dott. Lorenzo Baruggo}, medico federale FIN. Partendo
dall\textquotesingle apparato respiratorio trattato nella lezione
precedente, questa sessione si concentra su tre sistemi fondamentali per
il movimento: l\textquotesingle{}\textbf{apparato cardiocircolatorio},
l\textquotesingle{}\textbf{apparato locomotore} (ossa, articolazioni,
muscoli) e la \textbf{biomeccanica della contrazione muscolare}. La
conoscenza di questi argomenti è essenziale per comprendere come il
corpo del nuotatore produce, trasmette e controlla il movimento in
acqua.

\begin{center}\rule{0.5\linewidth}{0.5pt}\end{center}

\subsection{L\textquotesingle Apparato
      Cardiocircolatorio}\label{lapparato-cardiocircolatorio}

\subsubsection{Il Sistema Chiuso e il
      Cuore}\label{il-sistema-chiuso-e-il-cuore}

Mentre l\textquotesingle apparato respiratorio lavora con
l\textquotesingle aria, l\textquotesingle apparato cardiocircolatorio è
un \textbf{sistema liquido a circuito chiuso}: è il \textbf{sangue} che
trasporta l\textquotesingle ossigeno dagli alveoli polmonari a tutti i
tessuti dell\textquotesingle organismo. Il motore di questo sistema è il
\textbf{cuore}, una pompa che mette in moto il sangue in modo continuo e
unidirezionale.

Il cuore è diviso in \textbf{quattro cavità}:

\begin{itemize}
      \tightlist
      \item
            \textbf{Atri} (destro e sinistro): ricevono il sangue dalla periferia
      \item
            \textbf{Ventricoli} (destro e sinistro): si contraggono e spingono il
            sangue verso l\textquotesingle esterno
\end{itemize}

Atrio destro e atrio sinistro sono completamente separati da un setto;
comunicano solo in presenza di patologie.

\subsubsection{I Due Circuiti}\label{i-due-circuiti}

Il sistema cardiocircolatorio è composto da due circuiti principali:

\begin{longtable}[]{@{}p{0.30\linewidth}p{0.30\linewidth}p{0.30\linewidth}@{}}
      \toprule\noalign{}
      Circuito                     & Percorso                                      & Funzione \\
      \midrule\noalign{}
      \endhead
      \bottomrule\noalign{}
      \endlastfoot
      \textbf{Polmonare} (piccolo) & Ventricolo destro → arteria polmonare →
      polmoni → atrio sinistro     & Ossigenazione del sangue e rimozione della
      CO2                                                                                     \\
      \textbf{Sistemico} (grande)  & Ventricolo sinistro → aorta → capillari →
      vene cave → atrio destro     & Distribuzione dell\textquotesingle O2 a tutti
      i tessuti                                                                               \\
\end{longtable}

\textbf{Circuiti paralleli importanti:}

\begin{itemize}
      \tightlist
      \item
            \textbf{Circolo portale}: i nutrienti assorbiti
            dall\textquotesingle intestino entrano in questo circuito e vengono
            convogliati al \textbf{fegato} per essere filtrati ed elaborati prima
            di entrare nella circolazione generale
      \item
            \textbf{Circolo renale}: parte del sangue dell\textquotesingle aorta
            passa attraverso i \textbf{reni}, che filtrano le sostanze di scarto
            poi eliminate con le urine
\end{itemize}

\subsubsection{Le Valvole Cardiache}\label{le-valvole-cardiache}

Il sangue scorre sempre nella stessa direzione grazie a \textbf{quattro
      valvole} che impediscono il reflusso:

\begin{longtable}[]{@{}p{0.30\linewidth}p{0.30\linewidth}p{0.30\linewidth}@{}}
      \toprule\noalign{}
      Valvola                             & Posizione                                                    & Funzione \\
      \midrule\noalign{}
      \endhead
      \bottomrule\noalign{}
      \endlastfoot
      \textbf{Mitrale} (\emph{bicuspide}) & Tra atrio sinistro e ventricolo
      sinistro                            & Si apre per far passare il sangue; si chiude quando il
      ventricolo si contrae                                                                                         \\
      \textbf{Tricuspide}                 & Tra atrio destro e ventricolo destro                         & Stessa
      funzione sul lato destro                                                                                      \\
      \textbf{Semilunare aortica}         & Tra ventricolo sinistro e aorta                              &
      Impedisce il ritorno del sangue dall\textquotesingle aorta al
      ventricolo                                                                                                    \\
      \textbf{Semilunare polmonare}       & Tra ventricolo destro e arteria
      polmonare                           & Impedisce il ritorno del sangue dall\textquotesingle arteria
      polmonare al ventricolo                                                                                       \\
\end{longtable}

\begin{quote}
      Una valvola lesionata costringe il cuore a un \textbf{lavoro doppio}: il
      sangue refluisce e deve essere rimesso in circolo, con grave
      compromissione della funzionalità cardiaca.
\end{quote}

\subsubsection{Parametri Cardiaci}\label{parametri-cardiaci}

\textbf{A riposo:}

\begin{itemize}
      \tightlist
      \item
            Frequenza cardiaca: \textasciitilde{}\textbf{60 battiti/minuto}
      \item
            Volume sistolico (sangue per battito): \textasciitilde{}\textbf{60 cc}
      \item
            Portata cardiaca: 60 × 60 = \textasciitilde{}\textbf{3,6 litri/minuto}
\end{itemize}

\textbf{Sotto sforzo massimale:}

\begin{itemize}
      \tightlist
      \item
            Frequenza cardiaca: \textbf{150-200 battiti/minuto} (in funzione
            dell\textquotesingle età)
      \item
            Volume sistolico: \textasciitilde{}\textbf{100 cc} per battito
      \item
            Portata cardiaca: 150-200 × 100 = \textbf{15-20 litri/minuto}
\end{itemize}

\begin{quote}
      La portata cardiaca aumenta di circa \textbf{4-5 volte} tra riposo e
      sforzo massimale. La ventilazione polmonare, invece, aumenta di circa
      \textbf{25 volte}: questo rende l\textquotesingle apparato
      cardiovascolare il sistema con minor capacità proporzionale di
      incremento del trasporto di ossigeno.
\end{quote}

\subsubsection{Frequenza Cardiaca
      Massimale}\label{frequenza-cardiaca-massimale}

La formula empirica della cardiologia:

\begin{quote}
      \textbf{FC massimale = 220 - età}
\end{quote}

Esempio: un ventenne raggiunge circa 200 bpm (con variabilità
individuale tra 192 e 210). Un cinquantenne arriverà mediamente intorno
a 170 bpm. È una stima approssimativa ma utile in pratica.

\subsubsection{Come Rilevare le
      Pulsazioni}\label{come-rilevare-le-pulsazioni}

Si utilizzano i polpastrelli di indice e medio premuti leggermente su
un\textquotesingle arteria superficiale:

\begin{itemize}
      \tightlist
      \item
            \textbf{Polso}: arteria \textbf{radiale}, compressa
            sull\textquotesingle osso radiale
      \item
            \textbf{Collo}: arteria \textbf{carotide}, vicino
            all\textquotesingle angolo della mandibola
\end{itemize}

\subsubsection{Pressione Arteriosa}\label{pressione-arteriosa}

\begin{itemize}
      \tightlist
      \item
            \textbf{Sistolica} (massima): pressione misurata durante la
            contrazione del ventricolo
      \item
            \textbf{Diastolica} (minima): pressione durante il rilassamento del
            ventricolo
\end{itemize}

\textbf{Valori normali per un adulto:} 120/70 mmHg

\begin{longtable}[]{@{}p{0.47\linewidth}p{0.47\linewidth}@{}}
      \toprule\noalign{}
      Zona vascolare     & Pressione tipica \\
      \midrule\noalign{}
      \endhead
      \bottomrule\noalign{}
      \endlastfoot
      Arterie principali & 120/70 mmHg      \\
      Capillari          & 10-30 mmHg       \\
      Vene               & Molto bassa      \\
\end{longtable}

\textbf{Sotto sforzo}: la pressione sistolica può salire a
\textbf{200-220-230 mmHg}, il che non è pericoloso in un soggetto sano.
La stessa pressione a \textbf{riposo} è invece molto pericolosa, e anche
sotto sforzo è rischiosa nei \textbf{cardiopatici}.

\subsubsection{Frequenza Cardiaca e Intensità dello
      Sforzo}\label{frequenza-cardiaca-e-intensituxe0-dello-sforzo}

\begin{itemize}
      \tightlist
      \item
            Fino a circa il \textbf{70-80\% della FC massimale}: regime
            prevalentemente \textbf{aerobico}
      \item
            Oltre questa soglia: regime sempre più \textbf{misto}, poi
            \textbf{prevalentemente anaerobico}
      \item
            Alla FC massimale: regime \textbf{prevalentemente anaerobico}
\end{itemize}

\begin{center}\rule{0.5\linewidth}{0.5pt}\end{center}

\subsection{L\textquotesingle Apparato
      Locomotore}\label{lapparato-locomotore}

L\textquotesingle apparato locomotore comprende tre tessuti principali
che lavorano in sinergia per produrre il movimento: \textbf{ossa},
\textbf{articolazioni}, \textbf{muscoli}.

\subsubsection{Le Ossa e lo Scheletro}\label{le-ossa-e-lo-scheletro}

Le ossa formano lo \textbf{scheletro}, che ha funzione di sostegno per
l\textquotesingle intero organismo. La loro capacità di muoversi
l\textquotesingle una sull\textquotesingle altra, guidata dai muscoli
attraverso le articolazioni, è alla base di ogni movimento corporeo. Per
l\textquotesingle anatomia dettagliata dello scheletro si rimanda ai
libri FIN.

\subsubsection{Le Articolazioni}\label{le-articolazioni}

Le articolazioni \textbf{connettono un osso all\textquotesingle altro} e
possono consentire o limitare il movimento. Le articolazioni più
complesse sono formate da:

\begin{itemize}
      \tightlist
      \item
            \textbf{Capsula articolare}: involucro fibroso
      \item
            \textbf{Liquido sinoviale}: all\textquotesingle interno della capsula,
            facilita lo scorrimento
      \item
            \textbf{Cartilagine}: riveste le superfici ossee di contatto,
            impedendo l\textquotesingle usura del tessuto osseo
\end{itemize}

\textbf{Classificazione per mobilità:}

\begin{longtable}[]{@{}p{0.30\linewidth}p{0.30\linewidth}p{0.30\linewidth}@{}}
      \toprule\noalign{}
      Tipo                      & Esempio         & Mobilità                             \\
      \midrule\noalign{}
      \endhead
      \bottomrule\noalign{}
      \endlastfoot
      \textbf{Fissa}            & Ossa del cranio & Nessun movimento; si fondono
      completamente nell\textquotesingle invecchiamento                                  \\
      \textbf{Semimobile}       & Vertebre        & Limitata: flessione-estensione e
      lateralizzazione                                                                   \\
      \textbf{Uniplanare}       & Ginocchio       & Principalmente flessione-estensione;
      minima rotazione solo in soggetti molto flessibili                                 \\
      \textbf{Multidirezionale} & \textbf{Spalla} & Massima mobilità su tutti
      i piani; la più complessa e delicata                                               \\
\end{longtable}

\begin{quote}
      La \textbf{spalla} è l\textquotesingle articolazione che offre più
      possibilità di movimento nel corpo umano, ed è anche quella che
      \textbf{causa più problemi nei nuotatori}.
\end{quote}

\begin{center}\rule{0.5\linewidth}{0.5pt}\end{center}

\subsection{La Struttura e la Contrazione
      Muscolare}\label{la-struttura-e-la-contrazione-muscolare}

\subsubsection{Composizione del Muscolo}\label{composizione-del-muscolo}

Il muscolo è un insieme di cellule altamente specializzate, composte
prevalentemente da \textbf{proteine contrattili}. Le due proteine
principali sono:

\begin{itemize}
      \tightlist
      \item
            \textbf{Miosina}: filamento più spesso, con protuberanze (\emph{ponti
                  crociati}) che si agganciano all\textquotesingle actina
      \item
            \textbf{Actina}: filamento più sottile; viene spostata dalla miosina
            durante la contrazione
\end{itemize}

\textbf{Meccanismo di contrazione} (il "remo"):

\begin{enumerate}
      \tightlist
      \item
            La miosina agancia con i suoi ponti crociati l\textquotesingle actina
      \item
            I ponti si piegano come remi, spostando l\textquotesingle actina
      \item
            Le fibrille muscolari si accorciano
      \item
            Moltiplicato per i milioni di miofibrille: il muscolo si accorce
\end{enumerate}

\subsubsection{La Catena dello Stimolo
      Nervoso}\label{la-catena-dello-stimolo-nervoso}

Il segnale che porta alla contrazione percorre questo percorso:

\begin{enumerate}
      \tightlist
      \item
            \textbf{Cervello} (corteccia motoria) → stimolo \textbf{elettrico} via
            motoneurone
      \item
            Il segnale elettrico viaggia lungo il nervo
            (depolarizzazione/ripolarizzazione: sodio e potassio)
      \item
            Alla \textbf{sinapsi} (giunzione neuromuscolare):
            l\textquotesingle ultimo passaggio è \textbf{chimico}
            (neurotrasmettitore rilasciato nello spazio sinaptico)
      \item
            Lo stimolo chimico libera ioni \textbf{calcio}
            all\textquotesingle interno della cellula muscolare
      \item
            Il calcio innesca la reazione \textbf{ATP → ADP + energia}
      \item
            L\textquotesingle energia fa scivolare
            \textbf{l\textquotesingle actina sulla miosina} → contrazione
\end{enumerate}

\subsubsection{Struttura Gerarchica del Muscolo (come bambole
      giapponesi)}\label{struttura-gerarchica-del-muscolo-come-bambole-giapponesi}

\begin{verbatim}
Actina + Miosina
    → Miofibrille
        → Cellula muscolare
            → Fascio di fibre muscolari
                → Muscolo scheletrico (intero)
\end{verbatim}

Ogni livello è avvolto da un involucro di \textbf{collagene} o tessuto
\textbf{connettivo} che separa le strutture.

\subsubsection{Origine e Inserzione}\label{origine-e-inserzione}

Per produrre movimento, ogni muscolo ha:

\begin{itemize}
      \tightlist
      \item
            \textbf{Origine}: attacco su un capo osseo \emph{prossimale}
      \item
            \textbf{Inserzione distale}: attacco su un altro capo osseo
            \emph{distale}
\end{itemize}

Accorciandosi, il muscolo \textbf{avvicina i due capi ossei},
determinando il movimento. I muscoli principali della bracciata
natatoria sono il \textbf{pettorale} e il \textbf{gran dorsale}.

\begin{center}\rule{0.5\linewidth}{0.5pt}\end{center}

\subsection{Le Leve nel Corpo Umano}\label{le-leve-nel-corpo-umano}

Il sistema osso-articolazione-muscolo forma sistemi di \textbf{leve}.
Esistono tre specie:

\begin{longtable}[]{@{}p{0.30\linewidth}p{0.30\linewidth}p{0.30\linewidth}@{}}
      \toprule\noalign{}
      Specie                                                              & Caratteristica                               & Esempio nel corpo \\
      \midrule\noalign{}
      \endhead
      \bottomrule\noalign{}
      \endlastfoot
      \textbf{Prima specie}                                               & Fulcro tra potenza e resistenza              & Variabile         \\
      \textbf{Seconda specie}                                             & Sempre \textbf{vantaggiosa}: piccola forza →
      grande resistenza, piccola escursione                               & Alzarsi sulle \textbf{punte dei
            piedi} (con poca forza del polpaccio si solleva tutto il peso del
      corpo)                                                                                                                                 \\
      \textbf{Terza specie}                                               & Sempre \textbf{svantaggiosa}: grande forza →
      piccola resistenza all\textquotesingle estremità, grande escursione & La
      \textbf{bracciata} del nuoto: i muscoli della spalla e del tronco
      generano molta forza, ma la forza applicata sulla mano è molto minore;
      l\textquotesingle escursione della mano è invece amplissima (1,5-1,8
      m)                                                                                                                                     \\
\end{longtable}

\begin{quote}
      La \textbf{maggior parte delle leve del corpo umano} sono di
      \textbf{terza specie} (svantaggiose in termini di forza, vantaggiose in
      termini di escursione e velocità del movimento).
\end{quote}

\begin{center}\rule{0.5\linewidth}{0.5pt}\end{center}

\subsection{Tipi di Contrazione
      Muscolare}\label{tipi-di-contrazione-muscolare}

Il muscolo produce sempre un accorciamento delle miofibrille, ma il
muscolo nel suo complesso può:

\begin{longtable}[]{@{}p{0.30\linewidth}p{0.30\linewidth}p{0.30\linewidth}@{}}
      \toprule\noalign{}
      Tipo                      & Descrizione                                        & Esempio pratico \\
      \midrule\noalign{}
      \endhead
      \bottomrule\noalign{}
      \endlastfoot
      \textbf{Concentrica}      & Il muscolo si contrae \textbf{e si accorcia}       &
      Sollevare un peso flettendo l\textquotesingle avambraccio sul braccio
      (bicipite si accorce)                                                                            \\
      \textbf{Eccentrica}       & Il muscolo si contrae \textbf{ma si allunga}       &
      Ricaduta a terra dopo un salto: i quadricipiti devono frenare il peso
      accelerato dalla caduta, contraendosi mentre si allungano                                        \\
      \textbf{Isometrica}       & Il muscolo si contrae \textbf{ma non si muove} (né
      si allunga né si accorce) & Spingere con la mano su una superficie fissa
      inamovibile                                                                                      \\
\end{longtable}

\begin{quote}
      La contrazione \textbf{eccentrica} è la più stressante per il muscolo: i
      carichi possono superare anche di molto la forza massimale che il
      muscolo può esercitare concentricamente.
\end{quote}

\begin{center}\rule{0.5\linewidth}{0.5pt}\end{center}

\subsection{Tipi di Movimento}\label{tipi-di-movimento}

Oltre al tipo di contrazione, il movimento nel suo insieme può essere
classificato per la costanza del \textbf{carico} o della
\textbf{velocità}:

\begin{longtable}[]{@{}p{0.30\linewidth}p{0.30\linewidth}p{0.30\linewidth}@{}}
      \toprule\noalign{}
      Tipo                                                             & Definizione                                               & Note  \\
      \midrule\noalign{}
      \endhead
      \bottomrule\noalign{}
      \endlastfoot
      \textbf{Isotonico}                                               & Carico (forza resistente) costante durante tutto il
      movimento                                                        & Raro in natura; riprodotto da alcune macchine da palestra
      che compensano le variazioni angolari                                                                                                \\
      \textbf{Isometrico}                                              & Nessuno spostamento dei capi ossei                        & Forza
      applicata contro resistenza fissa inamovibile                                                                                        \\
      \textbf{Isocinetico}                                             & Velocità del movimento costante per tutta
      l\textquotesingle escursione (né accelerazioni né decelerazioni) &
      Raramente realizzabile; il nuoto vi si avvicina grazie alla resistenza
      dell\textquotesingle acqua                                                                                                           \\
\end{longtable}

\subsubsection{Il Nuoto e il Movimento
      Isocinetico}\label{il-nuoto-e-il-movimento-isocinetico}

La resistenza dell\textquotesingle acqua \textbf{smorzando le
      accelerazioni e le decelerazioni} rende il movimento natatoria più
vicino all\textquotesingle isocinetico rispetto alla maggior parte dei
movimenti atletici terrestri, pur senza raggiungerlo perfettamente.

La \textbf{panca biocinetica} (\emph{biokinetic bench}) era un attrezzo
a secco usato nella preparazione dei nuotatori che simulava
approssimativamente il pattern di forza del nuoto, con limitazioni alle
accelerazioni e decelerazioni, avvicinandosi al movimento in acqua senza
però essere un vero movimento isocinetico.

\begin{center}\rule{0.5\linewidth}{0.5pt}\end{center}

\subsection{Conclusione - Punti Chiave per
      l\textquotesingle Esame}\label{conclusione---punti-chiave-per-lesame-2}

\begin{itemize}
      \tightlist
      \item
            L\textquotesingle apparato cardiocircolatorio è un \textbf{sistema
                  chiuso a circuito doppio}: circuito \textbf{polmonare} (piccolo,
            ossigena il sangue) e circuito \textbf{sistemico} (grande,
            distribuisce l\textquotesingle O2).
      \item
            Il cuore ha \textbf{4 cavità}: atrio destro, ventricolo destro, atrio
            sinistro, ventricolo sinistro.
      \item
            Le \textbf{4 valvole cardiache}: \textbf{mitrale} (bicuspide, atrio sx
            → ventricolo sx), \textbf{tricuspide} (atrio dx → ventricolo dx),
            \textbf{semilunare aortica} (ventricolo sx → aorta),
            \textbf{semilunare polmonare} (ventricolo dx → arteria polmonare).
      \item
            A riposo: \textasciitilde{}\textbf{60 bpm},
            \textasciitilde{}\textbf{60 cc/battito} = \textasciitilde{}\textbf{3,6
                  L/min}. Sforzo max: 150-200 bpm, \textasciitilde{}\textbf{100
                  cc/battito} = \textbf{15-20 L/min} (4-5x rispetto a riposo; il sistema
            respiratorio aumenta 25x).
      \item
            \textbf{FC massimale = 220 - età}.
      \item
            Pulsazioni rilevabili all\textquotesingle{}\textbf{arteria radiale}
            (polso) o alla \textbf{carotide} (collo, angolo mandibola).
      \item
            Pressione normale: \textbf{120/70 mmHg}. Sistolica = contrazione;
            Diastolica = rilassamento.
      \item
            Soglia aerobica/anaerobica: \textasciitilde{}\textbf{70-80\% della FC
                  massimale}.
      \item
            Apparato locomotore: tre tessuti: \textbf{ossa} (sostegno),
            \textbf{articolazioni} (connessione e guida), \textbf{muscoli}
            (motori).
      \item
            La \textbf{spalla} è l\textquotesingle articolazione più mobile e la
            principale causa di problemi nei nuotatori.
      \item
            Proteine contrattili: \textbf{actina} (filamento sottile) e
            \textbf{miosina} (filamento spesso con ponti crociati).
      \item
            Catena dello stimolo: cervello → stimolo \textbf{elettrico}
            (motoneurone) → \textbf{sinapsi chimica} → liberazione di
            \textbf{calcio} → \textbf{ATP → ADP + energia} → actina scivola su
            miosina.
      \item
            Struttura gerarchica: actina+miosina → miofibrille → cellula muscolare
            → fascio di fibre → \textbf{muscolo scheletrico}.
      \item
            Il muscolo ha \textbf{origine} (prossimale) e \textbf{inserzione
                  distale} (distale) su due capi ossei diversi.
      \item
            Leve: \textbf{seconda specie} = vantaggiosa (es. punta del piede);
            \textbf{terza specie} = svantaggiosa ma con grande escursione (es.
            bracciata). La maggior parte delle leve del corpo sono di
            \textbf{terza specie}.
      \item
            Contrazioni: \textbf{concentrica} (si accorce), \textbf{eccentrica}
            (si contrae ma si allunga; la più stressante), \textbf{isometrica}
            (nessun movimento).
      \item
            Movimenti: \textbf{isotonico} (carico costante), \textbf{isometrico}
            (zero spostamento), \textbf{isocinetico} (velocità costante). Il nuoto
            si avvicina all\textquotesingle isocinetico grazie alla resistenza
            dell\textquotesingle acqua.
      \item
            La \textbf{panca biocinetica} (\emph{biokinetic bench}) simula a secco
            il pattern di forza del nuoto.
\end{itemize}

\section{4. Teoria del Movimento}\label{teoria-del-movimento}

\subsection{Introduzione}\label{introduzione-3}

Questa lezione è tenuta da \textbf{Filippo Tassara} e tratta le
\textbf{capacità motorie}: l\textquotesingle insieme delle
caratteristiche fisiche che un individuo possiede e che permettono
l\textquotesingle apprendimento e l\textquotesingle esecuzione delle
azioni motorie. Le capacità motorie influenzano
l\textquotesingle intensità e la qualità della risposta motoria
all\textquotesingle ambiente, sono in parte legate
all\textquotesingle ereditarietà e al patrimonio genetico, e possono
modificarsi con l\textquotesingle allenamento.
L\textquotesingle obiettivo è fornire al futuro istruttore un quadro
teorico solido per comprendere come l\textquotesingle individuo impara,
controlla e perfeziona il movimento.

\begin{center}\rule{0.5\linewidth}{0.5pt}\end{center}

\subsection{L\textquotesingle Organismo come Sistema
      Stimolo-Risposta}\label{lorganismo-come-sistema-stimolo-risposta}

L\textquotesingle organismo vive in continuo equilibrio con
l\textquotesingle ambiente, reagendo in modo proporzionato a ogni
stimolo proveniente dall\textquotesingle esterno o dal proprio corpo.
Tutto il comportamento dell\textquotesingle organismo può essere
descritto con il modello \textbf{stimolo → risposta} (\emph{input} →
\emph{output}): i movimenti e le azioni sono
l\textquotesingle espressione dell\textquotesingle attività vitale del
soggetto.

La relazione tra i comandi inviati dal cervello (via nervi efferenti) e
il risultato prodotto è però \textbf{complessa e non univoca}, per due
ragioni fondamentali:

\begin{enumerate}
      \tightlist
      \item
            I muscoli sono strutture \textbf{elastiche}: l\textquotesingle effetto
            della loro azione dipende dalla posizione e dalla velocità del sistema
            all\textquotesingle inizio dell\textquotesingle attività muscolare
      \item
            Ogni azione deve vincere \textbf{forze esterne} non controllabili dal
            soggetto (gravità, attrito, opposizione
            dell\textquotesingle avversario, ostacoli)
\end{enumerate}

Per questo l\textquotesingle organismo ha bisogno di una
\textbf{continua attività di controllo} realizzata attraverso gli organi
di senso (afferenza nervosa). I principali strumenti di segnalazione
sono:

\begin{itemize}
      \tightlist
      \item
            \textbf{Fusi intramuscolari}: nei muscoli
      \item
            \textbf{Corpuscoli tendinei del Golgi}: nei tendini
      \item
            \textbf{Meccanocettori articolari}: nelle articolazioni
      \item
            \textbf{Apparato uditivo}
      \item
            \textbf{Apparato visivo}
      \item
            \textbf{Struttura vestibolare} (labirintica o staticodinamica)
\end{itemize}

\begin{center}\rule{0.5\linewidth}{0.5pt}\end{center}

\subsection{Sviluppo Motorio per Fasce
      d\textquotesingle Età}\label{sviluppo-motorio-per-fasce-detuxe0}

\begin{longtable}[]{@{}p{0.30\linewidth}p{0.30\linewidth}p{0.30\linewidth}@{}}
      \toprule\noalign{}
      Età / Fase                                             & Caratteristiche motorie        & Categoria FIN \\
      \midrule\noalign{}
      \endhead
      \bottomrule\noalign{}
      \endlastfoot
      \textbf{\textasciitilde10 anni}                        & Il patrimonio dei movimenti si
      perfeziona; si accumulano movimenti piccoli e precisi  & Esordiente B                                   \\
      \textbf{Pubertà/adolescenza}                           & Riorganizzazione
      dell\textquotesingle apparato motorio: rigidità, angolosità, carenze
      coordinative, disturbi dell\textquotesingle inibizione & Allenamento di
      costruzione, prime competizioni                                                                         \\
\end{longtable}

Il \textbf{talento motorio} si manifesta spesso in età
pre-adolescenziale e richiede una gestione precisa, intelligente e
consapevole da parte del tecnico.

\begin{center}\rule{0.5\linewidth}{0.5pt}\end{center}

\subsection{Abilità Motorie e Capacità
      Motorie}\label{abilituxe0-motorie-e-capacituxe0-motorie}

\subsubsection{\texorpdfstring{Abilità Motorie (\emph{Motor
                  Skills})}{Abilità Motorie (Motor Skills)}}\label{abilituxe0-motorie-motor-skills}

Le \textbf{abilità motorie} sono le possibilità di movimento acquisite
\emph{ontogeneticamente} durante la vita. Sono:

\begin{itemize}
      \tightlist
      \item
            Elementi \textbf{determinati e conosciuti}
      \item
            \textbf{Automatizzati} con la ripetizione sistematica
      \item
            \textbf{Infiniti}: quante se ne apprendono
      \item
            In ambito sportivo diventano la \textbf{tecnica} e i
            \textbf{fondamentali} della disciplina
\end{itemize}

Si distinguono in:

\begin{longtable}[]{@{}p{0.30\linewidth}p{0.30\linewidth}p{0.30\linewidth}@{}}
      \toprule\noalign{}
      Tipologia                        & Descrizione                               & Esempio \\
      \midrule\noalign{}
      \endhead
      \bottomrule\noalign{}
      \endlastfoot
      \textbf{\emph{Closed skills}}    & Parametri fissi: angolazioni, ritmo,
      precisione millimetrica          & Nuoto (sport ciclico)                               \\
      \textbf{\emph{Open skills}}      & Devono essere modificate continuamente in
      risposta alla situazione tattica & Pallanuoto (sport situazionale)                     \\
\end{longtable}

\subsubsection{Capacità Motorie}\label{capacituxe0-motorie}

Le \textbf{capacità motorie} (definite da Zacciorski \emph{qualità
      fisiche}) sono il presupposto o prerequisito motorio di base su cui si
strutturano le abilità:

\begin{itemize}
      \tightlist
      \item
            Hanno significato più \textbf{generale}: riguardano
            un\textquotesingle intera classe di azioni
      \item
            Non si differenziano chiaramente nei \textbf{primi anni di vita}
      \item
            Possono essere \textbf{migliorate} con l\textquotesingle allenamento
\end{itemize}

\begin{quote}
      Abilità e capacità si integrano: le capacità motorie consentono
      l\textquotesingle apprendimento delle abilità; lo sviluppo delle une
      condiziona l\textquotesingle apprendimento delle altre.
\end{quote}

\begin{center}\rule{0.5\linewidth}{0.5pt}\end{center}

\subsection{La Triade delle Capacità Motorie (Gundlach,
      1968)}\label{la-triade-delle-capacituxe0-motorie-gundlach-1968}

\pandocbounded{\includegraphics[keepaspectratio]{/home/alessandrovulcu/Documenti/Corso Istruttori FIN 25-26/Pasted image 20260309113532.png}}

Gundlach propone la classificazione in due grandi categorie, a cui nelle
discipline natatorie si aggiunge una terza componente:

\begin{verbatim}
Capacità coordinative
        +
Capacità condizionali
        +
Mobilità articolare (flessibilità)
\end{verbatim}

Queste tre componenti formano una \textbf{triade} indissolubile che
definisce la qualità, la precisione e l\textquotesingle efficienza della
risposta motoria del soggetto.

\begin{center}\rule{0.5\linewidth}{0.5pt}\end{center}

\subsection{Capacità Coordinative}\label{capacituxe0-coordinative}

Le capacità coordinative sono \textbf{legate in modo stretto al sistema
      nervoso}: sono i presupposti che determinano
l\textquotesingle apprendimento, l\textquotesingle organizzazione, il
controllo e la trasformazione del movimento.

Alla nascita le forme di movimento sono minime; con il tempo si
acquisiscono \textbf{schemi motori di base} e \textbf{schemi posturali}
(\emph{human skills}), che dipendono dall\textquotesingle apprendimento
e dai processi di maturazione.

La \textbf{coordinazione} si realizza a due livelli:

\begin{itemize}
      \tightlist
      \item
            \textbf{Intramuscolare}: regolazione degli stimoli
            all\textquotesingle interno di un singolo muscolo
      \item
            \textbf{Intermuscolare}: attività sinergica tra più muscoli o gruppi
            muscolari
\end{itemize}

\subsubsection{Capacità Coordinative Generali (fasi
      dell\textquotesingle apprendimento)}\label{capacituxe0-coordinative-generali-fasi-dellapprendimento}

Seguono una progressione in tre tappe, che costituisce il modello
didattico di base:

\begin{enumerate}
      \tightlist
      \item
            \textbf{Capacità di apprendimento/scoperta}: prima fase; il soggetto
            scopre la nuova motricità attraverso l\textquotesingle esperienza
            diretta
      \item
            \textbf{Capacità di controllo}: il gesto viene riprodotto in modo
            sempre più efficace e vicino al modello ideale (che
            l\textquotesingle istruttore deve conoscere perfettamente)
      \item
            \textbf{Capacità di adattamento/trasformazione}: il gesto viene
            modificato in risposta a variazioni ambientali (es. nuotare in mare
            anziché in piscina; stile libero classico vs. testa alta in
            pallanuoto)
\end{enumerate}

\begin{quote}
      \textbf{Pericolo}: gli errori commessi nella fase di scoperta vengono
      \textbf{automatizzati} nella fase di controllo ed entrano nella memoria
      motoria. Un errore non corretto in tempo diventa difficilissimo da
      eliminare.
\end{quote}

\subsubsection{Capacità Coordinative Speciali (Tassonomia di
      Bloom)}\label{capacituxe0-coordinative-speciali-tassonomia-di-bloom}

Le capacità coordinative speciali riguardano fasi evolute del controllo
motorio e hanno un \textbf{sviluppo accentuato tra i 6 e i 12 anni}
(esordienti B ed A).

\begin{longtable}[]{@{}p{0.30\linewidth}p{0.30\linewidth}p{0.30\linewidth}@{}}
      \toprule\noalign{}
      Capacità                                                           & Descrizione                                                          & Applicazione nelle discipline FIN \\
      \midrule\noalign{}
      \endhead
      \bottomrule\noalign{}
      \endlastfoot
      \textbf{Equilibrio}                                                & Mantenere la postura utile
      all\textquotesingle azione; legata all\textquotesingle apparato
      vestibolare                                                        & Virate in tutti gli stili, variazioni di posizione del
      corpo                                                                                                                                                                         \\
      \textbf{Reazione}                                                  & Fase 1: elaborazione rapida dello stimolo
      percettivo; Fase 2: velocità di trasmissione agli organi effettori &
      Partenza al segnale acustico; salto per prendere la palla in
      pallanuoto                                                                                                                                                                    \\
      \textbf{Combinazione motoria}                                      & Sincronizzazione di azioni secondo
      simultaneità, alternanza e successione                             & Sincronizzato, tuffi, tutte le
      discipline                                                                                                                                                                    \\
      \textbf{Orientamento}                                              & Mantenere il rapporto corretto con il campo
      d\textquotesingle azione e i riferimenti spaziali                  & Nuotatori a dorso
      (soffitto/cielo); pallanuotisti (compagni, calottine, avversari);
      nuotatori in acque libere                                                                                                                                                     \\
      \textbf{Differenziazione cinestesica}                              & Dosare finemente la forza per
      massima efficacia                                                  & Variazione di ritmo nel nuoto (frequenza bracciata
      mantenendo la tecnica corretta)                                                                                                                                               \\
      \textbf{Anticipazione motoria}                                     & Previsione probabilistica di situazioni
      non ancora verificatesi; pre-costituzione dell\textquotesingle azione
      successiva                                                         & Fondamentale in pallanuoto; importante in acque libere e
      salvamento; secondaria in nuoto classico e discipline
      tecnico-compositive                                                                                                                                                           \\
      \textbf{Ritmo}                                                     & Scelta temporale degli impulsi per realizzare in modo
      adeguato un\textquotesingle azione finalizzata; può essere adeguato a
      una sequenza esterna o autodeterminato                             & Distribuzione dello sforzo in
      gara; ritmo bracciata                                                                                                                                                         \\
      \textbf{Fantasia motoria}                                          & Usare le abilità immagazzinate nella memoria
      motoria combinandole in modo creativo per risolvere il problema motorio
                                                                         & Sport situazionali (pallanuoto); Teodorescu: "organizzare il proprio
      gioco per disorganizzare quello dell\textquotesingle avversario,
      rimanendo organizzati"                                                                                                                                                        \\
\end{longtable}

\begin{center}\rule{0.5\linewidth}{0.5pt}\end{center}

\subsection{Gli Analizzatori}\label{gli-analizzatori}

Gli \textbf{analizzatori} sono strutture anatomico-funzionali che
partecipano all\textquotesingle{}\textbf{afferenza} (arrivo del
messaggio) e alla \textbf{riafferenza} (ritorno del messaggio dopo una
risposta). Ogni analizzatore è composto da:

\begin{itemize}
      \tightlist
      \item
            \textbf{Recettori specifici}
      \item
            \textbf{Vie nervose afferenti}
      \item
            \textbf{Campo primario di proiezione} della corteccia cerebrale
\end{itemize}

\subsubsection{Analizzatori Esterni}\label{analizzatori-esterni}

\begin{longtable}[]{@{}p{0.30\linewidth}p{0.30\linewidth}p{0.30\linewidth}@{}}
      \toprule\noalign{}
      Analizzatore      & Recettori                                                         & Funzione e applicazione      \\
      \midrule\noalign{}
      \endhead
      \bottomrule\noalign{}
      \endlastfoot
      \textbf{Visivo}   & Telecettori (recettori della distanza)                            & Ricezione del
      modello di apprendimento; controllo parziale del movimento; visione
      periferica; la dimostrazione è il principale strumento didattico;
      fondamentale in sport di squadra (posizione della palla e dei compagni
      in pallanuoto)                                                                                                       \\
      \textbf{Acustico} & Recettori uditivi                                                 & Trasmissione di informazioni
      verbali; ritmo delle bracciate; sequenza di utilizzo della palla in
      pallanuoto                                                                                                           \\
      \textbf{Tattile}  & Pressocettori, termocettori, nocicettori (sulla cute)
                        & Determinante nel nuoto: permette di applicare la forza propulsiva
      sull\textquotesingle acqua mantenendo la posizione idrodinamica; grande
      contrasto tra scivolare nell\textquotesingle acqua e applicare forza
      propulsiva                                                                                                           \\
\end{longtable}

\subsubsection{Analizzatori Interni}\label{analizzatori-interni}

\begin{longtable}[]{@{}p{0.30\linewidth}p{0.30\linewidth}p{0.30\linewidth}@{}}
      \toprule\noalign{}
      Analizzatore                                        & Recettori                                               & Funzione e applicazione \\
      \midrule\noalign{}
      \endhead
      \bottomrule\noalign{}
      \endlastfoot
      \textbf{Cinestesico}                                & Propriocettori in muscoli, tendini e
      articolazioni                                       & Percepisce il movimento e le variazioni di tensione sin
      dall\textquotesingle inizio; vie di conduzione più veloci degli altri
      analizzatori; es. posizione della gamba nel sincro, braccio nel lancio,
      recupero nel dorso                                                                                                                      \\
      \textbf{Statico-dinamico} (vestibolare/labirintico) & Apparato
      vestibolare dell\textquotesingle orecchio interno   & Percepisce la
      posizione della testa nel campo gravitazionale; registra direzione,
      accelerazioni e decelerazioni; regola equilibrio statico e dinamico;
      determinante nelle virate, nelle discipline subacquee (salvamento), nei
      tuffi e nelle posizioni speciali del sincronizzato e della pallanuoto                                                                   \\
\end{longtable}

\begin{center}\rule{0.5\linewidth}{0.5pt}\end{center}

\subsection{Capacità Condizionali}\label{capacituxe0-condizionali}

Le capacità condizionali sono legate ai \textbf{grandi apparati}
(muscolo-scheletrico e cardiorespiratorio) e determinate dai
\textbf{processi energetici} (metabolici e plastici). Hanno uno
\textbf{sviluppo accentuato tra i 12 e i 18 anni} (pubertà/adolescenza),
quando si sviluppano i grandi apparati.

\subsubsection{L\textquotesingle Unità Motoria}\label{lunituxe0-motoria}

L\textquotesingle unità elementare della contrazione volontaria è
l\textquotesingle{}\textbf{unità motoria}, composta da:

\begin{itemize}
      \tightlist
      \item
            \textbf{Neurone} (corpo cellulare + nucleo)
      \item
            \textbf{Assone} (\emph{neurite}): estroflessione del neurone, anche di
            notevole lunghezza
      \item
            \textbf{Placca motrice}: punto di inserzione del nervo sulla fibra
            muscolare
\end{itemize}

\subsubsection{Le Tre Capacità
      Condizionali}\label{le-tre-capacituxe0-condizionali}

\paragraph{1. Velocità / Rapidità}\label{velocituxe0-rapidituxe0}

La capacità di compiere azioni nel \textbf{minor tempo possibile}. Si
articola in:

\begin{itemize}
      \tightlist
      \item
            \textbf{Velocità di azione}: eseguire un gesto autonomo nel minor
            tempo possibile (es. tuffo di partenza volontario)
      \item
            \textbf{Capacità di reazione}: rispondere a uno stimolo prefissato o
            improvviso nel minor tempo (es. partenza al segnale
            dell\textquotesingle allenatore)
      \item
            \textbf{Frequenza dei movimenti}: unione di movimenti rapidi e
            semplici in azioni complesse (es. prime bracciate dopo il tuffo di
            partenza)
\end{itemize}

\paragraph{2. Forza}\label{forza}

La capacità di \textbf{opporsi o vincere una resistenza} con una
contrazione muscolare. Si distingue in:

\begin{itemize}
      \tightlist
      \item
            \textbf{Forza massima}: produrre un\textquotesingle intensità di forza
            massima
      \item
            \textbf{Forza rapida}: produrre forza da elevata a massima nel
            brevissimo tempo (collegata alla rapidità)
      \item
            \textbf{Forza resistente}: produrre una forza di intensità scelta
            mantenuta nel tempo
\end{itemize}

\paragraph{3. Resistenza}\label{resistenza}

La capacità di \textbf{tollerare l\textquotesingle affaticamento} in
esercizi di vario tipo. È la base di tutte le altre capacità
condizionali:

\begin{quote}
      Lo sviluppo di forza rapida, coordinazione e rapidità è possibile con
      successo \textbf{solo sulla base di un adeguato sviluppo della
            resistenza}.
\end{quote}

Si distingue in:

\begin{itemize}
      \tightlist
      \item
            \textbf{Resistenza di breve termine}
      \item
            \textbf{Resistenza di medio termine}
      \item
            \textbf{Resistenza di lungo termine}
\end{itemize}

\begin{center}\rule{0.5\linewidth}{0.5pt}\end{center}

\subsection{Mobilità Articolare (Terzo Vertice della
      Triade)}\label{mobilituxe0-articolare-terzo-vertice-della-triade}

La mobilità articolare ha componenti sia coordinative che condizionali e
dipende da:

\begin{itemize}
      \tightlist
      \item
            \textbf{Componente congenita}: struttura del tessuto connettivo,
            muscoli e articolazioni
      \item
            \textbf{Fattori ormonali}: si modificano nel corso dello sviluppo
      \item
            \textbf{Abitudini motorie}: sollecitano determinate ampiezze
            articolari
\end{itemize}

\subsubsection{Definizioni}\label{definizioni}

\begin{itemize}
      \tightlist
      \item
            \textbf{Mobilità muscolare} (flessibilità): capacità di eseguire
            movimenti con la massima ampiezza nell\textquotesingle escursione
            articolare (es. slancio di una gamba nel sincronizzato)
      \item
            \textbf{Elasticità muscolare}: capacità del muscolo scheletrico di
            lasciarsi stirare recuperando la lunghezza fisiologica senza subire
            traumi
      \item
            \textbf{Mobilità articolare} (anche: flessibilità, articolabilità,
            articolarità, estensibilità): capacità di una o più articolazioni di
            muoversi liberamente per tutto il proprio \emph{range} di mobilità,
            senza limiti e senza dolore
\end{itemize}

\subsubsection{Tipologie}\label{tipologie}

\begin{longtable}[]{@{}p{0.47\linewidth}p{0.47\linewidth}@{}}
      \toprule\noalign{}
      Tipo             & Descrizione                                   \\
      \midrule\noalign{}
      \endhead
      \bottomrule\noalign{}
      \endlastfoot
      \textbf{Attiva}  & I muscoli agonisti distendono gli antagonisti \\
      \textbf{Passiva} & Una forza esterna agisce
      sull\textquotesingle articolazione (gravità, attrezzo, azione del
      tecnico)                                                         \\
      \textbf{Mista}   & Combinazione delle due                        \\
\end{longtable}

\begin{quote}
      Test di valutazione: \textbf{flessibilità scapolomerale} con asta
      graduata; \textbf{flessione della colonna} sul plinto con gradazione
      millimetrica.
\end{quote}

L\textquotesingle articolazione più sollecitata nel nuoto è la
\textbf{scapolomerale}, coinvolta in tutti e quattro gli stili. La
mobilità articolare favorisce: benessere fisico, corretta postura,
economia dei gesti, miglioramento della performance, prevenzione degli
infortuni.

\begin{center}\rule{0.5\linewidth}{0.5pt}\end{center}

\subsection{Abilità Motorie: Fasi di
      Sviluppo}\label{abilituxe0-motorie-fasi-di-sviluppo}

Le abilità motorie entrano nella \textbf{memoria motoria di lungo
      termine} e rappresentano i \emph{fondamentali} della disciplina. Il
processo di acquisizione si articola in tre fasi:

\begin{longtable}[]{@{}p{0.47\linewidth}p{0.47\linewidth}@{}}
      \toprule\noalign{}
      Fase                                                  & Caratteristiche                      \\
      \midrule\noalign{}
      \endhead
      \bottomrule\noalign{}
      \endlastfoot
      \textbf{Coordinazione grezza} (scoperta/comprensione) & Prime esecuzioni
      complete; risposta motoria ancora imprecisa; coinvolgimento affettivo e
      cognitivo; costruzione di un automatismo globale                                             \\
      \textbf{Coordinazione fine}                           & Esecuzione perfezionata; sensibile a
      fattori esterni; padronanza completa dell\textquotesingle azione;
      possibilità di trasformare il gesto                                                          \\
      \textbf{Disponibilità variabile}                      & Tipica del grande campione: tecnica
      perfetta + capacità di adattarla a qualsiasi situazione; risposta
      flessibile e imprevedibile                                                                   \\
\end{longtable}

\begin{center}\rule{0.5\linewidth}{0.5pt}\end{center}

\subsection{Fattori della Prestazione nel
      Nuoto}\label{fattori-della-prestazione-nel-nuoto}

I fattori che concorrono alla prestazione nelle discipline natatorie
sono:

\begin{itemize}
      \tightlist
      \item
            \textbf{Condizione}: forma fisica e rendimento della prova
      \item
            \textbf{Gestione della gara}: distribuzione delle forze
            nell\textquotesingle arco della prova
      \item
            \textbf{Biomeccanica}: rendimento delle tecniche di nuotata
      \item
            \textbf{Elementi specifici}: partenze, virate, subacque; fondamentali
            con la palla (pallanuoto); tecniche del sincronizzato e dei tuffi
      \item
            \textbf{Tattica}: predominante nella pallanuoto; presente nelle acque
            libere e nel salvamento; minore nel nuoto classico
      \item
            \textbf{Adesione emotiva}: come il soggetto affronta la prestazione
            (il "leone" in allenamento può diventare "coniglietto" in gara, e
            viceversa)
\end{itemize}

\begin{center}\rule{0.5\linewidth}{0.5pt}\end{center}

\subsection{Conclusione - Punti Chiave per
      l\textquotesingle Esame}\label{conclusione---punti-chiave-per-lesame-3}

\begin{itemize}
      \tightlist
      \item
            Le \textbf{capacità motorie} sono presupposti generali del movimento,
            in parte genetici, modificabili con l\textquotesingle allenamento. Le
            \textbf{abilità motorie} sono specifiche, automatizzate, infinite.
      \item
            \textbf{Closed skills}: sport ciclici (nuoto) → parametri fissi.
            \textbf{Open skills}: sport situazionali (pallanuoto) → adattamento
            continuo.
      \item
            L\textquotesingle organismo funziona come sistema \textbf{stimolo →
                  risposta}; il controllo avviene tramite \textbf{afferenza nervosa}
            (recettori).
      \item
            \textbf{Triade}: capacità \textbf{coordinative} (SNC) + capacità
            \textbf{condizionali} (grandi apparati) + \textbf{mobilità
                  articolare}.
      \item
            Capacità coordinative generali: \textbf{scoperta → controllo →
                  adattamento}. Il pericolo: gli errori della fase di scoperta si
            automatizzano.
      \item
            Capacità coordinative speciali (sviluppo massimo \textbf{6-12 anni}):
            equilibrio, reazione, combinazione motoria, orientamento,
            differenziazione cinestesica, anticipazione motoria, ritmo, fantasia
            motoria.
      \item
            \textbf{Analizzatori esterni}: visivo (telecettori, dimostrazione),
            acustico (informazioni verbali), tattile (pressocettori, fondamentale
            nel nuoto per applicare la forza propulsiva).
      \item
            \textbf{Analizzatori interni}: cinestesico (propriocettori in
            muscoli/tendini/articolazioni, vie più veloci),
            vestibolare/labirintico (posizione testa, accelerazioni/decelerazioni,
            equilibrio statico e dinamico).
      \item
            \textbf{Unità motoria}: neurone + assone + placca motrice.
      \item
            Capacità condizionali: \textbf{velocità/rapidità}, \textbf{forza}
            (massima, rapida, resistente), \textbf{resistenza}. Sviluppo
            accentuato \textbf{12-18 anni}.
      \item
            La \textbf{resistenza} è la base di tutte le capacità condizionali;
            senza di essa non si sviluppano forza rapida, coordinazione e
            rapidità.
      \item
            Mobilità articolare: \textbf{attiva} (muscoli agonisti),
            \textbf{passiva} (forza esterna), \textbf{mista}. Articolazione più
            sollecitata nel nuoto: \textbf{scapolomerale}.
      \item
            Fasi dell\textquotesingle abilità motoria: \textbf{coordinazione
                  grezza} → \textbf{coordinazione fine} → \textbf{disponibilità
                  variabile} (caratteristica del grande campione).
      \item
            La \textbf{disponibilità variabile} è la capacità di esprimere la
            tecnica perfetta adattandola a qualsiasi situazione.
\end{itemize}

\section{Area Psicopedagogica 1 - Comunicazione, Intelligenza Emotiva e
  Apprendimento}\label{area-psicopedagogica-1---comunicazione-intelligenza-emotiva-e-apprendimento}

\textbf{Docente:} Paola Contardi - Psicologa, psicoterapeuta e psicologa
dello sport\\
\textbf{Piattaforma:} FAD FIN

\begin{center}\rule{0.5\linewidth}{0.5pt}\end{center}

\subsection{Introduzione}\label{introduzione-4}

Questa lezione introduce le basi psicopedagogiche necessarie
all\textquotesingle istruttore di nuoto per svolgere la propria attività
in modo efficace. I tre argomenti trattati -\/- il processo di
comunicazione, l\textquotesingle intelligenza emotiva e il binomio
insegnamento-apprendimento -\/- costituiscono il fondamento teorico su
cui si costruisce la competenza didattica
dell\textquotesingle istruttore.

\begin{center}\rule{0.5\linewidth}{0.5pt}\end{center}

\subsection{1. Il Processo di
      Comunicazione}\label{il-processo-di-comunicazione}

\subsubsection{Definizione e principi
      fondamentali}\label{definizione-e-principi-fondamentali}

Il termine \textbf{comunicazione} deriva dal latino \emph{comunicare},
ovvero "mettere in comune". Comunicare significa far capire agli altri
il proprio messaggio attraverso parole, gesti e immagini. Non si tratta
di un semplice trasferimento di informazioni: comunicare implica lo
scambio di emozioni e sensazioni, e richiede attenzione e il superamento
di \textbf{barriere comunicative}.

La comunicazione è alla base dei rapporti interpersonali. In ambito
didattico, il processo è \textbf{bidirezionale}:
l\textquotesingle istruttore deve essere capace non solo di comunicare,
ma anche di \textbf{saper ascoltare}.

\begin{quote}
      Principio chiave: per produrre un cambiamento
      nell\textquotesingle azione didattica, conta di piu cio che
      \textbf{arriva} rispetto a cio che \textbf{parte}.
      L\textquotesingle effetto ottenuto e piu importante
      dell\textquotesingle intenzione iniziale.
\end{quote}

\begin{center}\rule{0.5\linewidth}{0.5pt}\end{center}

\subsubsection{Elementi del processo
      comunicativo}\label{elementi-del-processo-comunicativo}

Un atto comunicativo si compone di elementi fondamentali:

\begin{longtable}[]{@{}p{0.47\linewidth}p{0.47\linewidth}@{}}
      \toprule\noalign{}
      Elemento           & Descrizione                                         \\
      \midrule\noalign{}
      \endhead
      \bottomrule\noalign{}
      \endlastfoot
      \textbf{Emittente} & La persona che avvia la comunicazione attraverso un
      messaggio                                                                \\
      \textbf{Ricevente} & Colui che accoglie, decodifica, interpreta,
      comprende e mette in atto il messaggio                                   \\
      \textbf{Messaggio} & Il contenuto della comunicazione; cio che viene
      scambiato sotto forma di informazione, e sia il "cosa" che il "come"     \\
      \textbf{Contesto}  & L\textquotesingle ambiente fisico e sociale in cui
      la comunicazione avviene                                                 \\
      \textbf{Empatia}   & La sintonia relazionale che favorisce
      l\textquotesingle apprendimento                                          \\
\end{longtable}

Poiché la comunicazione e bidirezionale, \textbf{insegnante e allievo
      sono al tempo stesso emittenti e riceventi}.\\
% Immagine mancante: Pasted image 20260309105803.png

\begin{center}\rule{0.5\linewidth}{0.5pt}\end{center}

\subsubsection{Il ruolo dell\textquotesingle istruttore come
      emittente}\label{il-ruolo-dellistruttore-come-emittente}

Quando l\textquotesingle istruttore emette un messaggio deve:

\begin{itemize}
      \tightlist
      \item
            \textbf{Individuare il destinatario}: valutare se si tratta di un
            bambino, un adolescente o un adulto, e se e un principiante o ha già
            acquisito un buon livello di acquaticità.
      \item
            \textbf{Scegliere i mezzi con cui comunicare}.
      \item
            \textbf{Accertarsi che il destinatario (l\textquotesingle allievo) sia
                  in grado di raccogliere fedelmente il messaggio}.
      \item
            \textbf{Considerare eventuali interferenze}.
\end{itemize}

\begin{center}\rule{0.5\linewidth}{0.5pt}\end{center}

\subsubsection{Codice verbale e non
      verbale}\label{codice-verbale-e-non-verbale}

\begin{itemize}
      \tightlist
      \item
            Il \textbf{codice verbale} indica \emph{cosa si dice}.
      \item
            Il \textbf{codice non verbale} riguarda \emph{come lo si dice}.
\end{itemize}

L\textquotesingle istruttore deve prestare attenzione a:

\begin{itemize}
      \tightlist
      \item
            Postura, gestualità e movimenti del corpo
      \item
            Espressioni del volto e sguardo
      \item
            Uso dello spazio
      \item
            Tono della voce
      \item
            Espressioni legate alle emozioni
\end{itemize}

\begin{quote}
      Se conta molto ciò che si dice, conta ancora di più \textbf{come lo si
            dice}.
\end{quote}

\begin{center}\rule{0.5\linewidth}{0.5pt}\end{center}

\subsubsection{Il linguaggio: modulazione del
      livello}\label{il-linguaggio-modulazione-del-livello}

La scelta del linguaggio e fondamentale. L\textquotesingle istruttore
deve scegliere il \textbf{livello linguistico piu idoneo} in funzione
degli interlocutori:

\begin{itemize}
      \tightlist
      \item
            \textbf{Livello massimo} (linguaggio tecnico): per chi ha gia
            raggiunto un livello di nuoto avanzato.
      \item
            \textbf{Livello minimo}: per chi e principiante.
\end{itemize}

\begin{center}\rule{0.5\linewidth}{0.5pt}\end{center}

\subsubsection{Il contesto: l\textquotesingle ambiente
      piscina}\label{il-contesto-lambiente-piscina}

L\textquotesingle ambiente piscina presenta caratteristiche specifiche
da considerare:

\begin{itemize}
      \tightlist
      \item
            E un ambiente ricco di \textbf{elementi di disturbo}.
      \item
            Sono presenti diverse \textbf{barriere alla comunicazione}: rumore di
            sottofondo, presenza di molte persone contemporaneamente.
      \item
            E un ambiente sociale che richiede di riconoscere i componenti del
            gruppo.
\end{itemize}

\begin{center}\rule{0.5\linewidth}{0.5pt}\end{center}

\subsubsection{Feedback e comunicazione
      efficace}\label{feedback-e-comunicazione-efficace}

L\textquotesingle istruttore deve verificare che il proprio messaggio
sia stato compreso nel modo corretto, cogliendo il \textbf{feedback}
dell\textquotesingle allievo. Per farlo e necessario:

\begin{itemize}
      \tightlist
      \item
            Mantenere un coinvolgimento positivo nella comunicazione didattica.
      \item
            Rendere le lezioni \textbf{interessanti, coinvolgenti e divertenti}.
\end{itemize}

\begin{center}\rule{0.5\linewidth}{0.5pt}\end{center}

\subsubsection{Indicazioni pratiche per la
      lezione}\label{indicazioni-pratiche-per-la-lezione}

Durante la lezione, per una comunicazione efficace,
l\textquotesingle istruttore dovrà:

\begin{itemize}
      \tightlist
      \item
            Usare \textbf{parole adeguate} e una \textbf{gestualità corretta}.
      \item
            Stabilire \textbf{regole di ingaggio} condivise con gli allievi.
      \item
            Assicurarsi che gli allievi possano \textbf{osservarlo e ascoltarlo}
            quando spiega gli esercizi.
      \item
            Definire un \textbf{linguaggio e una gestualità condivisa} per
            comunicare anche durante l\textquotesingle esecuzione degli esercizi.
      \item
            Fornire costantemente \textbf{feedback di rinforzo positivo},
            sottolineando impegno, attenzione e risultati ottenuti, per mantenere
            alto interesse e rendimento.
\end{itemize}

\begin{center}\rule{0.5\linewidth}{0.5pt}\end{center}

\subsection{2. L\textquotesingle Intelligenza
      Emotiva}\label{lintelligenza-emotiva}

\subsubsection{Origine del concetto}\label{origine-del-concetto}

Il concetto di \textbf{intelligenza emotiva} e stato introdotto nel
\textbf{1990} da \textbf{Salovey e Mayer}. Diventa popolare nel
\textbf{1995} con la pubblicazione del libro \emph{Intelligenza emotiva}
di \textbf{Daniel Goleman}, che la descrive come un insieme di
competenze fondamentali per affrontare con successo la vita, attraverso
autocontrollo, entusiasmo, perseveranza e capacita di automotivarsi.

\begin{center}\rule{0.5\linewidth}{0.5pt}\end{center}

\subsubsection{I cinque ambiti dell\textquotesingle intelligenza emotiva
      (Goleman, Mayer,
      Salovey)}\label{i-cinque-ambiti-dellintelligenza-emotiva-goleman-mayer-salovey}

\begin{longtable}[]{@{}p{0.47\linewidth}p{0.47\linewidth}@{}}
      \toprule\noalign{}
      Ambito                                     & Descrizione                              \\
      \midrule\noalign{}
      \endhead
      \bottomrule\noalign{}
      \endlastfoot
      \textbf{Riconoscere le proprie emozioni}   & La consapevolezza di
      identificare le emozioni nel momento in cui insorgono; e la base
      dell\textquotesingle intelligenza emotiva e incrementa
      l\textquotesingle autoconsapevolezza                                                  \\
      \textbf{Padroneggiare le proprie emozioni} & Accrescere la
      consapevolezza di se e tollerare senza traumi gli eventi positivi e
      negativi della vita                                                                   \\
      \textbf{Motivare se stessi}                & La spinta interiore che ogni soggetto si
      da per raggiungere un obiettivo prefissato                                            \\
      \textbf{Riconoscere le emozioni altrui}    & Capacita empatica fondamentale
      per i tecnici sportivi; permette di fare distinzioni tra soggetti in
      base a umore, temperamento e motivazione                                              \\
      \textbf{Utilizzare le competenze sociali}  & Capacita di analisi della
      situazione sociale; riconoscere e comprendere sentimenti, motivazioni e
      preoccupazioni degli altri                                                            \\
\end{longtable}

\begin{center}\rule{0.5\linewidth}{0.5pt}\end{center}

\subsubsection{Importanza in ambito lavorativo e
      sportivo}\label{importanza-in-ambito-lavorativo-e-sportivo}

L\textquotesingle intelligenza emotiva e essenziale sia in ambito
personale che lavorativo. Permette di:

\begin{itemize}
      \tightlist
      \item
            Entrare in relazione con gli altri effettuando critiche in forma
            costruttiva.
      \item
            Creare un\textquotesingle atmosfera in cui le diversità tra i soggetti
            diventino qualcosa da apprezzare, non una fonte di attrito.
\end{itemize}

Per i tecnici sportivi, \textbf{riconoscere le emozioni altrui} e
particolarmente critico per costruire le competenze sociali necessarie
in fase di formazione.

Per relazionarsi efficacemente -\/- con colleghi e allievi -\/- occorre
essere:

\begin{itemize}
      \tightlist
      \item
            \textbf{Specifici} nei feedback
      \item
            Capaci di \textbf{offrire soluzioni}
      \item
            \textbf{Presenti} e attenti
      \item
            \textbf{Empatici}
\end{itemize}

\begin{center}\rule{0.5\linewidth}{0.5pt}\end{center}

\subsubsection{L\textquotesingle intelligenza emotiva nei bambini: il
      compito
      dell\textquotesingle educatore}\label{lintelligenza-emotiva-nei-bambini-il-compito-delleducatore}

L\textquotesingle intelligenza emotiva puo essere stimolata \textbf{sin
      da bambini}. E compito dell\textquotesingle educatore:

\begin{itemize}
      \tightlist
      \item
            Favorirne uno sviluppo ottimale.
      \item
            Conoscere se stesso per gestire con intelligenza il proprio profilo
            emotivo.
\end{itemize}

I bambini che imparano a riconoscere e gestire le proprie emozioni:

\begin{itemize}
      \tightlist
      \item
            Tollerano in modo piu efficace le situazioni stressanti.
      \item
            Comunicano meglio i propri stati emozionali.
      \item
            Sviluppano relazioni migliori con famiglia e amici.
      \item
            Ottengono risultati migliori a scuola, nello sport e nella vita.
\end{itemize}

\begin{center}\rule{0.5\linewidth}{0.5pt}\end{center}

\subsubsection{La ricerca sul "marshmallow test" (Michel e
      Ebbesen)}\label{la-ricerca-sul-marshmallow-test-michel-e-ebbesen}

Un celebre studio assegno il seguente compito a bambini di quattro anni:
\textbf{attendere il ritorno del ricercatore evitando di mangiare una
      caramella} lasciata davanti a loro. Come premio per
l\textquotesingle attesa, avrebbero ottenuto due caramelle; in caso
contrario, solo quella gia consumata.

I risultati in eta adolescenziale mostrarono che i bambini capaci di
attendere erano:

\begin{itemize}
      \tightlist
      \item
            Piu competenti socialmente, efficienti e sicuri di se.
      \item
            Piu capaci di fronteggiare le frustrazioni.
      \item
            Piu abili nell\textquotesingle accettare e perseguire i propri
            obiettivi senza farsi scoraggiare dalle difficolta.
      \item
            Studenti migliori, piu competenti sia sul piano scolastico che
            educativo.
\end{itemize}

La capacita di \textbf{attendere la gratificazione} e un indicatore
precoce di intelligenza emotiva.

\begin{center}\rule{0.5\linewidth}{0.5pt}\end{center}

\subsubsection{Analfabetismo emotivo e
      prevenzione}\label{analfabetismo-emotivo-e-prevenzione}

Attualmente si registra una diffusa condizione di \textbf{analfabetismo
      emotivo} nei giovani. Gli studi confermano un rapporto direttamente
proporzionale tra analfabetismo emotivo e disturbi del comportamento.

L\textquotesingle{}\textbf{alfabetizzazione emotiva} dei bambini e la
strategia piu efficace per fare prevenzione attiva. In ambito formativo
sono stati sviluppati programmi per insegnare nelle strutture educative
a riconoscere e gestire le emozioni.

Gli istruttori, attraverso un\textquotesingle adeguata educazione,
possono contribuire a:

\begin{itemize}
      \tightlist
      \item
            Sviluppare l\textquotesingle intelligenza emotiva degli allievi.
      \item
            Evitare che i ragazzi vivano in futuro un inaridimento dei rapporti
            interpersonali.
      \item
            Prevenire senso di inadeguatezza e caduta
            dell\textquotesingle autostima.
\end{itemize}

\begin{quote}
      L\textquotesingle intelligenza emotiva e la sintesi tra la razionalita e
      la vita emotiva, e migliora il benessere individuale e collettivo. Se
      sviluppata parallelamente alla preparazione professionale, ottimizza i
      processi di acquisizione, apprendimento e scambio delle informazioni.
\end{quote}

\begin{center}\rule{0.5\linewidth}{0.5pt}\end{center}

\subsection{3. Insegnamento e
      Apprendimento}\label{insegnamento-e-apprendimento}

\subsubsection{Definizioni}\label{definizioni-1}

\textbf{Insegnamento}: l\textquotesingle atto di fornire specifiche
conoscenze, competenze e informazioni agli allievi attraverso una vasta
gamma di metodi e strategie. Gli insegnanti progettano e offrono lezioni
interessanti e coinvolgenti per stimolare l\textquotesingle interesse e
la motivazione, tenendo conto dell\textquotesingle eta e delle
caratteristiche peculiari dei singoli allievi e del gruppo.

In senso piu ampio, l\textquotesingle insegnamento consiste in una
molteplice serie di comportamenti realizzati dagli educatori allo scopo
di far apprendere qualcosa agli allievi. Puo attuarsi sia tra adulti che
tra bambini.

\textbf{Apprendimento}: la capacita di modificare pensieri e
comportamenti a seguito di una situazione esperita ripetutamente. E
definito come processo che coinvolge aspetti cognitivi e affettivi
(motivazione, percezione, emozione, sentimenti).

\begin{center}\rule{0.5\linewidth}{0.5pt}\end{center}

\subsubsection{L\textquotesingle insegnamento deve rispettare i tempi
      dell\textquotesingle individuo}\label{linsegnamento-deve-rispettare-i-tempi-dellindividuo}

Punto fondamentale: l\textquotesingle insegnamento deve seguire le
\textbf{modalita di apprendimento spontaneo} e rispettare i
\textbf{tempi dell\textquotesingle individuo}, che sono in relazione ai
limiti e alle risorse presenti.

L\textquotesingle insegnamento richiede la massima competenza
professionale. Piu gli istruttori sono preparati a livello tecnico,
educativo, psicologico e didattico, piu:

\begin{itemize}
      \tightlist
      \item
            Si instaurera un buon rapporto con gli allievi.
      \item
            Gli allievi si sentiranno coinvolti, interessati e responsabilizzati.
      \item
            Si favorira l\textquotesingle apprendimento.
      \item
            Gli allievi diventeranno protagonisti attivi e motivati.
\end{itemize}

\begin{center}\rule{0.5\linewidth}{0.5pt}\end{center}

\subsubsection{I principali contributi
      teorici}\label{i-principali-contributi-teorici}

\begin{longtable}[]{@{}p{0.47\linewidth}p{0.47\linewidth}@{}}
      \toprule\noalign{}
      Autore                                & Contributo                    \\
      \midrule\noalign{}
      \endhead
      \bottomrule\noalign{}
      \endlastfoot
      \textbf{Piaget}                       & Mise in evidenza il ruolo
      dell\textquotesingle{}\emph{imitazione} nel processo di apprendimento \\
      \textbf{Bandura} (psicologo canadese) & Sottolineo
      l\textquotesingle importanza del \emph{modeling} (apprendimento per
      osservazione e imitazione di un modello)                              \\
      \textbf{Vygotsky} (scienziato russo)  & Spiego come
      l\textquotesingle apprendimento del bambino avvenga con
      l\textquotesingle aiuto degli altri; introdusse il concetto di
      \emph{zona di sviluppo prossimale}                                    \\
\end{longtable}

\textbf{Zona di sviluppo prossimale} (Vygotsky): la differenza tra cio
che il bambino puo imparare da solo e cio che puo apprendere grazie
all\textquotesingle aiuto e alla guida di adulti e coetanei.

\begin{center}\rule{0.5\linewidth}{0.5pt}\end{center}

\subsubsection{Creare un ambiente per
      l\textquotesingle apprendimento}\label{creare-un-ambiente-per-lapprendimento}

Il ruolo dell\textquotesingle insegnamento e quello di \textbf{creare le
      condizioni adatte affinche l\textquotesingle apprendimento possa
      avvenire}. Questo significa lavorare su:

\begin{itemize}
      \tightlist
      \item
            La \textbf{relazione} tra istruttore e allievo.
      \item
            La \textbf{comunicazione efficace}.
      \item
            La creazione di un ambiente ricco di \textbf{esperienze sensoriali,
                  emotive e tecniche} che permettano all\textquotesingle allievo di
            avere un ruolo attivo.
\end{itemize}

\begin{center}\rule{0.5\linewidth}{0.5pt}\end{center}

\subsubsection{Caratteristiche
      dell\textquotesingle apprendimento}\label{caratteristiche-dellapprendimento}

\begin{itemize}
      \tightlist
      \item
            Si compone dell\textquotesingle unione degli aspetti \textbf{cognitivi
                  e affettivi} (motivazione, percezione, emozione, sentimenti).
      \item
            L\textquotesingle{}\textbf{esperienza} e la condizione imprescindibile
            affinche si abbia apprendimento.
      \item
            Non riguarda solo l\textquotesingle acquisizione di nuove competenze,
            ma coinvolge \textbf{tutte le sfere della personalita} con una
            finalita adattativa.
      \item
            Implica la capacita di \textbf{analizzare, valutare e applicare} le
            informazioni acquisite.
      \item
            Non si riferisce solo a conoscenze e capacita, ma anche ad
            \textbf{atteggiamenti, valori e abitudini}.
\end{itemize}

\begin{center}\rule{0.5\linewidth}{0.5pt}\end{center}

\subsubsection{Il profilo dell\textquotesingle istruttore
      ideale}\label{il-profilo-dellistruttore-ideale}

L\textquotesingle istruttore deve essere una persona competente e
professionalmente preparata, in grado di:

\begin{itemize}
      \tightlist
      \item
            Costruire un percorso \textbf{personalizzato} sulle emozioni e sulla
            realta dei propri allievi.
      \item
            Ottenere un cambiamento nel pieno \textbf{rispetto degli allievi}.
      \item
            Agire nella consapevolezza di essere un \textbf{educatore} che tenga
            conto dei principi della \textbf{cultura dell\textquotesingle acqua}.
\end{itemize}

\begin{center}\rule{0.5\linewidth}{0.5pt}\end{center}

\subsection{Conclusione: Punti chiave per
      l\textquotesingle esame}\label{conclusione-punti-chiave-per-lesame}

\begin{enumerate}
      \item
            \textbf{Comunicazione}: processo bidirezionale; conta piu
            l\textquotesingle effetto che arriva che l\textquotesingle intenzione
            di partenza. L\textquotesingle istruttore gestisce codice verbale e
            non verbale, adatta il livello linguistico all\textquotesingle allievo
            e coglie il feedback.
      \item
            \textbf{Emittente e ricevente}: in ambito didattico, istruttore e
            allievo ricoprono entrambi i ruoli contemporaneamente.
      \item
            \textbf{Contesto piscina}: ambiente pieno di barriere (rumore,
            affollamento); richiede strategie comunicative specifiche (gestualita
            condivisa, regole di ingaggio).
      \item
            \textbf{Intelligenza emotiva}: introdotta da Salovey e Mayer (1990),
            resa popolare da Goleman (1995). I cinque ambiti: riconoscere le
            proprie emozioni, padroneggiare le proprie emozioni, motivare se
            stessi, riconoscere le emozioni altrui, utilizzare le competenze
            sociali.
      \item
            \textbf{Marshmallow test}: la capacita di attendere la gratificazione
            nei bambini predice competenza sociale, successo scolastico e capacita
            di gestire le frustrazioni in eta adolescenziale.
      \item
            \textbf{Zona di sviluppo prossimale} (Vygotsky): distanza tra cio che
            il bambino impara da solo e cio che impara con l\textquotesingle aiuto
            degli altri -\/- fondamentale per l\textquotesingle azione didattica
            dell\textquotesingle istruttore.
      \item
            \textbf{Apprendimento}: modifica di pensieri e comportamenti
            attraverso l\textquotesingle esperienza ripetuta; coinvolge aspetti
            cognitivi, affettivi e tutta la sfera della personalita.
      \item
            \textbf{L\textquotesingle istruttore} deve essere competente a livello
            tecnico, educativo, psicologico e didattico, e agire come educatore
            consapevole della cultura dell\textquotesingle acqua.
\end{enumerate}

\section{Area Psicopedagogica 2 - Sviluppo, Autostima e Gestione del
  Gruppo}\label{area-psicopedagogica-2---sviluppo-autostima-e-gestione-del-gruppo}

\textbf{Docente:} Paola Contardi - Psicologa, psicoterapeuta e psicologa
dello sport\\
\textbf{Piattaforma:} FAD FIN

\begin{center}\rule{0.5\linewidth}{0.5pt}\end{center}

\subsection{Introduzione}\label{introduzione-5}

Questa seconda lezione di area psicopedagogica approfondisce tre temi
strettamente interconnessi: lo sviluppo cognitivo, emotivo e affettivo
del bambino; la costruzione dell\textquotesingle autostima; la gestione
del singolo e del gruppo in contesto didattico. La comprensione di
questi argomenti e fondamentale per l\textquotesingle istruttore, che
non opera solo come tecnico ma come figura educativa di riferimento.

\begin{center}\rule{0.5\linewidth}{0.5pt}\end{center}

\subsection{1. Sviluppo Cognitivo, Emotivo e
      Affettivo}\label{sviluppo-cognitivo-emotivo-e-affettivo}

\subsubsection{Il concetto di sviluppo}\label{il-concetto-di-sviluppo}

Lo \textbf{sviluppo} e un processo dinamico costituito da svariati
cambiamenti in ognuna delle fasi della vita, con implicazioni
fondamentali per il futuro adulto. Con questo termine si intende lo
studio dei diversi aspetti della personalità e di come questi si
modificano e si integrano in relazione non solo a fattori
\textbf{biologici} ma anche a quelli \textbf{ambientali}, in tutto
l\textquotesingle arco della vita di un individuo.

Gli aspetti cognitivo, emotivo, affettivo, sociale e relazionale sono
\textbf{strettamente interconnessi e interdipendenti}: non e possibile
considerarli separatamente.

\begin{center}\rule{0.5\linewidth}{0.5pt}\end{center}

\subsubsection{Il ruolo del cervello nei primi anni di
      vita}\label{il-ruolo-del-cervello-nei-primi-anni-di-vita}

Durante i \textbf{primi anni di vita} il cervello umano e alla ricerca
di stimoli esperienziali. Grazie alla sua notevole \textbf{plasticità},
l\textquotesingle organismo in via di sviluppo e
l\textquotesingle ambiente possono adattarsi e armonizzarsi insieme.

Da un punto di vista neuropsicologico, e di fondamentale importanza:

\begin{itemize}
      \tightlist
      \item
            La \textbf{relazione di attaccamento}, inizialmente con la madre e
            successivamente con chi opera in ambito educativo.
      \item
            L\textquotesingle assenza di traumi nelle fasi precoci della vita: i
            traumi precoci sono causa di uno scorretto sviluppo cerebrale e, di
            conseguenza, di comportamenti psichicamente deviati.
\end{itemize}

\begin{center}\rule{0.5\linewidth}{0.5pt}\end{center}

\subsubsection{I principali contributi
      teorici}\label{i-principali-contributi-teorici-1}

\begin{longtable}[]{@{}p{0.47\linewidth}p{0.47\linewidth}@{}}
      \toprule\noalign{}
      Autore             & Contributo principale                              \\
      \midrule\noalign{}
      \endhead
      \bottomrule\noalign{}
      \endlastfoot
      \textbf{Bowlby}    & Per la salute psicologica del bambino e essenziale
      l\textquotesingle esperienza di una relazione continuativa, intima e
      calda con la madre o con chi ne rappresenta la figura.
      L\textquotesingle ambiente precoce e un ambiente sociale, e la madre e
      l\textquotesingle oggetto sociale primario che media
      l\textquotesingle approccio con il mondo esterno                        \\
      \textbf{Winnicott} & Rafforza il concetto dell\textquotesingle ambiente
      precoce come ambiente sociale; sottolinea l\textquotesingle importanza
      della figura materna come mediatore                                     \\
      \textbf{Piaget}    & L\textquotesingle intelligenza e il risultato
      dell\textquotesingle interazione tra l\textquotesingle organismo e
      l\textquotesingle ambiente                                              \\
      \textbf{Vygotsky}  & Le varie funzioni psichiche
      dell\textquotesingle essere umano dipendono dalla rete di fattori
      sociali e culturali in cui il bambino cresce e si sviluppa              \\
      \textbf{Bruner}    & La percezione e un processo cognitivo che proviene
      dall\textquotesingle interno del soggetto, influenzato da valori,
      motivazioni e credenze. L\textquotesingle attivita umana e finalizzata a
      uno scopo; gli strumenti per raggiungerlo vanno costruiti tramite il
      linguaggio e gli insegnamenti su come fare le cose                      \\
\end{longtable}

\begin{quote}
      Le esperienze del bambino con la figura genitoriale e con le persone
      significative del suo mondo portano alla formazione di \textbf{modelli
            mentali del se} che faranno da filtro nei futuri processi di
      elaborazione delle informazioni.
\end{quote}

\begin{center}\rule{0.5\linewidth}{0.5pt}\end{center}

\subsubsection{Sviluppo affettivo ed
      emotivo}\label{sviluppo-affettivo-ed-emotivo}

Lo sviluppo affettivo e quello emotivo sono aspetti del funzionamento
umano \textbf{inscindibili} l\textquotesingle uno
dall\textquotesingle altro, e a loro volta si intrecciano con lo
sviluppo cognitivo. Lo sviluppo emotivo si associa alle tappe formulate
da Piaget per lo sviluppo cognitivo.

Le emozioni svolgono più funzioni:

\begin{itemize}
      \tightlist
      \item
            Comunicare agli altri uno stato interno.
      \item
            Preparare l\textquotesingle organismo all\textquotesingle azione.
      \item
            Contribuire al \textbf{processo di negoziazione} tra individui, che
            avviene all\textquotesingle interno dell\textquotesingle interazione
            sociale bidirezionale.
\end{itemize}

Le emozioni e la loro regolazione sono l\textquotesingle esito
dell\textquotesingle intreccio tra le componenti cognitive ed emotive
che riguardano due o più persone coinvolte in una relazione.

Inoltre, le emozioni del bambino sono influenzate dalle reazioni della
figura significativa agli stessi eventi: il comportamento verbale del
genitore trasmette modelli emotivi al bambino.

\begin{center}\rule{0.5\linewidth}{0.5pt}\end{center}

\subsubsection{Il ruolo dello sport e della cultura
      dell\textquotesingle acqua}\label{il-ruolo-dello-sport-e-della-cultura-dellacqua}

Lavorare con tutti gli stimoli sensoriali previsti dalla nuova
\textbf{cultura dell\textquotesingle acqua} favorisce lo sviluppo delle
capacita motorie e consente al bambino di acquisire con gradualità la
competenza natatoria necessaria per sviluppare un \textbf{senso di
      sicurezza e benessere in acqua}, e solo successivamente, se lo desidera,
di diventare un buon nuotatore o un atleta agonista.

Il miglioramento delle \textbf{capacita sensopercettive} si ottiene
proponendo attività che richiedono l\textquotesingle uso degli
analizzatori.

Attraverso le proposte didattiche l\textquotesingle istruttore dovrà:

\begin{itemize}
      \tightlist
      \item
            \textbf{Soddisfare la curiosità} per favorire lo sviluppo cognitivo.
      \item
            \textbf{Stimolare la capacita di iniziativa e di scelta} per
            promuovere l\textquotesingle aspetto emotivo.
      \item
            \textbf{Proporre le attività in un clima di partecipazione e fiducia}
            per sviluppare gli aspetti sociali e affettivi.
      \item
            Strutturare le attività in forma \textbf{adeguata, gratificante e
                  ludica} per favorire lo sviluppo dell\textquotesingle autostima.
\end{itemize}

\begin{center}\rule{0.5\linewidth}{0.5pt}\end{center}

\subsection{2. La Costruzione
      dell\textquotesingle Autostima}\label{la-costruzione-dellautostima}

\subsubsection{Definizione e componenti}\label{definizione-e-componenti}

L\textquotesingle{}\textbf{autostima} e l\textquotesingle insieme dei
giudizi valutativi che l\textquotesingle individuo da di se stesso. Può
essere costruita giorno dopo giorno attraverso una serie di strategie.

Dalla letteratura scientifica emergono \textbf{tre elementi
      fondamentali} ricorrenti in tutte le definizioni:

\begin{enumerate}
      \tightlist
      \item
            La presenza di un sistema che consente all\textquotesingle individuo
            di \textbf{auto-osservarsi e auto-conoscersi}.
      \item
            L\textquotesingle{}\textbf{aspetto valutativo}: il giudizio generale
            di se stessi.
      \item
            L\textquotesingle{}\textbf{aspetto affettivo}: la capacita di valutare
            e considerare in modo positivo o negativo gli elementi descritti.
\end{enumerate}

L\textquotesingle autostima non emerge solo da fattori interni
individuali, ma anche dai confronti che l\textquotesingle individuo fa
-\/- più o meno consapevolmente -\/- con l\textquotesingle ambiente in
cui vive.

% Immagine mancante: Pasted image 20260309105909.png

\begin{center}\rule{0.5\linewidth}{0.5pt}\end{center}

\subsubsection{Se reale e se ideale}\label{se-reale-e-se-ideale}

La formazione dell\textquotesingle autostima e determinata
dall\textquotesingle interazione tra due componenti:

\begin{longtable}[]{@{}p{0.47\linewidth}p{0.47\linewidth}@{}}
      \toprule\noalign{}
      Componente         & Descrizione             \\
      \midrule\noalign{}
      \endhead
      \bottomrule\noalign{}
      \endlastfoot
      \textbf{Se reale}  & Cio che si e realmente  \\
      \textbf{Se ideale} & Come si vorrebbe essere \\
\end{longtable}

\begin{quote}
      Maggiore sara la discrepanza tra cio che si e e cio che si vorrebbe
      essere, minore sara la stima in se stessi.
\end{quote}

\begin{center}\rule{0.5\linewidth}{0.5pt}\end{center}

\subsubsection{Alta autostima vs bassa
      autostima}\label{alta-autostima-vs-bassa-autostima}

\begin{longtable}[]{@{}p{0.47\linewidth}p{0.47\linewidth}@{}}
      \toprule\noalign{}
      Alta autostima                               & Bassa autostima                  \\
      \midrule\noalign{}
      \endhead
      \bottomrule\noalign{}
      \endlastfoot
      Riconosce in modo realistico pregi e difetti & Riconosce esclusivamente
      le proprie debolezze                                                            \\
      Apprezza i propri punti di forza             & Trascura i propri punti di forza \\
      Si impegna a migliorare i propri limiti      & Si arrende facilmente di
      fronte alle difficolta                                                          \\
      Persevera nel raggiungere gli obiettivi      & Cede di fronte a pareri
      contrari o difficolta                                                           \\
      Sa relativizzare gli insuccessi              & Subisce gli insuccessi senza
      elaborarli                                                                      \\
\end{longtable}

L\textquotesingle autostima e il fondamento su cui ogni bambino
costruisce il proprio benessere ed e considerata la \textbf{chiave dei
      successi in eta adulta}.

\begin{center}\rule{0.5\linewidth}{0.5pt}\end{center}

\subsubsection{Il ruolo dell\textquotesingle istruttore nella
      costruzione
      dell\textquotesingle autostima}\label{il-ruolo-dellistruttore-nella-costruzione-dellautostima}

L\textquotesingle istruttore ha una responsabilita diretta nello
sviluppo dell\textquotesingle autostima dei propri allievi. E importante
che abbia piena \textbf{consapevolezza del proprio ruolo educativo}.

L\textquotesingle autostima e messa in relazione al grado di
\textbf{autonomia} e al \textbf{senso di autoefficacia}
dell\textquotesingle individuo. Pertanto l\textquotesingle istruttore
dovrebbe:

\begin{itemize}
      \tightlist
      \item
            Valutare per ciascun allievo il suo grado di autonomia (che varia con
            l\textquotesingle eta).
      \item
            Proporre attivita che facciano riferimento
            all\textquotesingle autonomia stessa.
      \item
            Creare un clima educativo in cui sia costantemente presente la
            possibilita di dare \textbf{rinforzi positivi}.
\end{itemize}

Esempio: quando un allievo riesce per la prima volta a eseguire un
esercizio in acqua, sottolineare il successo rafforza il senso di
autoefficacia del bambino. Il \textbf{rinforzo} aiuta il bambino a
rendersi consapevole dei propri punti di forza e a considerare i limiti
come margini di miglioramento.

L\textquotesingle{}\textbf{incoraggiamento} e il modo migliore per
aiutare il bambino a sviluppare il \textbf{dialogo interno positivo}:
una modalita di autoincoraggiamento (es. "ce la posso fare", "io valgo",
"posso affrontare questa difficolta") che lo supportera nel percorso di
crescita.

\begin{center}\rule{0.5\linewidth}{0.5pt}\end{center}

\subsubsection{Didattica ludica e convinzioni
      razionali}\label{didattica-ludica-e-convinzioni-razionali}

Con i bambini piu piccoli e molto importante il \textbf{momento del
      gioco}. Una didattica ludica consente di costruire convinzioni razionali
quali:

\begin{itemize}
      \tightlist
      \item
            "Io valgo anche se a qualcuno non piace cio che io faccio."
      \item
            "Si puo essere amici anche quando non si vogliono fare le stesse
            cose."
      \item
            "Alcuni sono gentili con me, altri a volte si comportano male."
\end{itemize}

\begin{center}\rule{0.5\linewidth}{0.5pt}\end{center}

\subsubsection{Comportamenti utili dell\textquotesingle istruttore per
      l\textquotesingle autostima}\label{comportamenti-utili-dellistruttore-per-lautostima}

Per migliorare la crescita dell\textquotesingle autostima degli allievi,
l\textquotesingle istruttore deve:

\begin{itemize}
      \tightlist
      \item
            Favorire un \textbf{clima di accettazione} dei propri limiti e delle
            difficolta.
      \item
            \textbf{Incentivare gli allievi ad aiutarsi} tra di loro.
      \item
            \textbf{Comprendere le peculiarita di ciascuno} mostrando un
            atteggiamento di accettazione.
      \item
            Usare la \textbf{critica costruttiva} per correggere gli errori.
      \item
            \textbf{Valorizzare le esperienze positive}.
      \item
            \textbf{Supportare} l\textquotesingle allievo nei momenti di
            difficolta, mostrando interessamento.
      \item
            \textbf{Criticare il comportamento e mai la persona}.
      \item
            Costruire un \textbf{rapporto di collaborazione con i genitori} per
            sensibilizzarli ad adottare comportamenti utili a costruire
            l\textquotesingle autostima dei propri figli.
\end{itemize}

\begin{center}\rule{0.5\linewidth}{0.5pt}\end{center}

\subsubsection{L\textquotesingle equilibrio tra spinta e
      protezione}\label{lequilibrio-tra-spinta-e-protezione}

Aiutare i bambini a sviluppare talenti e acquisire competenze richiede
un equilibrio tra due opposte necessita:

\begin{itemize}
      \tightlist
      \item
            \textbf{Spinta}: incentivare il bambino a cimentarsi con nuove
            attivita.
      \item
            \textbf{Protezione}: preservarlo da aspettative non realistiche che
            potrebbero compromettere la sua autostima.
\end{itemize}

\begin{quote}
      Gli obiettivi didattici dovrebbero essere costruiti in modo che
      l\textquotesingle allievo possa considerarli \textbf{raggiungibili},
      seppur con il giusto grado di fatica e impegno.
\end{quote}

\begin{center}\rule{0.5\linewidth}{0.5pt}\end{center}

\subsection{3. Gestione del Singolo e del
      Gruppo}\label{gestione-del-singolo-e-del-gruppo}

\subsubsection{Competenze relazionali
      dell\textquotesingle istruttore}\label{competenze-relazionali-dellistruttore}

Il lavoro dell\textquotesingle istruttore non si esaurisce nel possesso
di competenze tecniche: implica una \textbf{formazione personale} mirata
allo sviluppo delle capacita relazionali. Lo \textbf{stile
      dell\textquotesingle istruttore} incide in modo diretto sulla gestione
efficiente della corsia.

L\textquotesingle istruttore deve impostare una \textbf{didattica attiva
      e individualizzata} all\textquotesingle interno del gruppo, dimostrando
adeguate capacita di:

\begin{itemize}
      \tightlist
      \item
            Attenzione e osservazione degli allievi.
      \item
            Identificazione di problemi di natura tecnica o relazionale.
      \item
            Organizzazione delle attivita con messaggi chiari e in sequenza
            logica.
      \item
            Valorizzazione del principio che \textbf{tutti gli allievi possono
                  imparare}.
\end{itemize}

\begin{center}\rule{0.5\linewidth}{0.5pt}\end{center}

\subsubsection{La diversita degli
      allievi}\label{la-diversita-degli-allievi}

Le tipologie di allievi che l\textquotesingle istruttore incontra in una
giornata di lavoro sono molto varie: bambini, adolescenti, adulti,
anziani. Tutti condividono il desiderio di apprendere o migliorare la
propria capacita di muoversi in sicurezza in acqua.

Chi si iscrive a un corso ha una propria \textbf{motivazione
      individuale}, che deve pero trovare realizzazione
all\textquotesingle interno di un \textbf{gruppo dinamico}, collegato
alla necessita del momento, alle capacita del singolo e in continua
evoluzione.

L\textquotesingle istruttore deve:

\begin{itemize}
      \tightlist
      \item
            Riconoscere e assecondare le \textbf{differenze individuali}.
      \item
            Adottare una \textbf{didattica varia} che garantisca
            un\textquotesingle alternanza di stimoli.
      \item
            Assicurare a ciascuno un adattamento al proprio \textbf{stile di
                  apprendimento personale}.
      \item
            Prevedere, quando necessario, \textbf{percorsi didattici
                  differenziati} sulla base delle diverse difficolta degli allievi.
      \item
            Essere \textbf{tempestivo e specifico} nel colmare le lacune che
            potrebbero compromettere il percorso di apprendimento.
\end{itemize}

\begin{center}\rule{0.5\linewidth}{0.5pt}\end{center}

\subsubsection{Il gruppo: definizione e
      complessita}\label{il-gruppo-definizione-e-complessita}

Il \textbf{gruppo} e un insieme di persone unite tra loro da vincoli
naturali, rapporti di interesse e obiettivi comuni, che stabiliscono tra
loro relazioni interpersonali e di interdipendenza.

\begin{quote}
      Un gruppo e qualcosa di diverso dalla semplice somma dei suoi singoli
      componenti.
\end{quote}

Il gruppo rappresenta un \textbf{sistema complesso} legato a:

\begin{itemize}
      \tightlist
      \item
            Il numero degli elementi.
      \item
            La tipologia degli elementi stessi.
      \item
            Le relazioni che si instaurano al suo interno.
\end{itemize}

L\textquotesingle apprendimento e un\textquotesingle azione
\textbf{individuale ma non solitaria}: l\textquotesingle efficacia
dell\textquotesingle insegnamento e i risultati
dell\textquotesingle apprendimento sono strettamente correlati alla
capacita dell\textquotesingle istruttore di gestire educativamente e
didatticamente sia il singolo che il gruppo.

\begin{center}\rule{0.5\linewidth}{0.5pt}\end{center}

\subsubsection{La conduzione del gruppo}\label{la-conduzione-del-gruppo}

Il \textbf{nodo centrale} della gestione del gruppo e la capacita
dell\textquotesingle istruttore di guidare e incanalare correttamente
gli allievi considerando i loro bisogni, le loro particolarita e il loro
grado di maturita.

La conduzione del gruppo deve essere:

\begin{itemize}
      \tightlist
      \item
            \textbf{Intenzionale}: fondata su scopi educativi chiari.
      \item
            \textbf{Sistematica}: organizzata e coerente nel tempo.
      \item
            \textbf{Flessibile}: capace di adattarsi alle situazioni che emergono.
\end{itemize}

Essere in sintonia con gli allievi significa \textbf{calibrare
      costantemente} modalita e atteggiamenti comunicativi. Cio richiede
attenzione continua e una costante rimodulazione degli interventi.

\begin{center}\rule{0.5\linewidth}{0.5pt}\end{center}

\subsubsection{Le difficolta
      dell\textquotesingle istruttore}\label{le-difficolta-dellistruttore}

Anche l\textquotesingle istruttore puo incontrare difficolta nella
gestione quotidiana, tra cui:

\begin{itemize}
      \tightlist
      \item
            La \textbf{paura di essere criticato}.
      \item
            La difficolta nel trovare il giusto equilibrio tra \textbf{severita e
                  flessibilita}, per concedere il grado di autonomia e liberta
            sufficienti a favorire l\textquotesingle apprendimento.
\end{itemize}

La risposta a queste difficolta passa attraverso la consapevolezza del
proprio ruolo, la formazione continua e la riflessione sulla propria
pratica educativa.

\begin{center}\rule{0.5\linewidth}{0.5pt}\end{center}

\subsection{Conclusione: Punti chiave per
      l\textquotesingle esame}\label{conclusione-punti-chiave-per-lesame-1}

\begin{enumerate}
      \item
            \textbf{Sviluppo}: processo dinamico che coinvolge aspetti cognitivi,
            emotivi, affettivi, sociali e relazionali, strettamente interconnessi
            e influenzati sia da fattori biologici che ambientali.
      \item
            \textbf{Plasticita cerebrale}: nei primi anni di vita il cervello e
            molto ricettivo agli stimoli esperienziali; le relazioni di
            attaccamento precoci sono fondamentali per lo sviluppo sano.
      \item
            \textbf{Contributi teorici chiave}: Bowlby e Winnicott (attaccamento e
            ambiente precoce), Piaget (sviluppo cognitivo e interazione
            organismo-ambiente), Vygotsky (fattori socioculturali nello sviluppo),
            Bruner (percezione come processo influenzato da valori e motivazioni).
      \item
            \textbf{Autostima}: giudizi valutativi che l\textquotesingle individuo
            da di se stesso; determinata dall\textquotesingle interazione tra
            \emph{se reale} e \emph{se ideale}. Piu grande e la discrepanza, piu
            bassa e l\textquotesingle autostima.
      \item
            \textbf{Rinforzo positivo e autoefficacia}: sottolineare i successi
            costruisce il senso di autoefficacia;
            l\textquotesingle incoraggiamento sviluppa il dialogo interno
            positivo.
      \item
            \textbf{Critica costruttiva}: criticare sempre il comportamento, mai
            la persona.
      \item
            \textbf{Obiettivi didattici}: devono essere raggiungibili con il
            giusto grado di fatica, per non compromettere
            l\textquotesingle autostima dell\textquotesingle allievo.
      \item
            \textbf{Gruppo}: sistema complesso e qualcosa di piu della somma dei
            suoi componenti; la sua gestione richiede una conduzione intenzionale,
            sistematica e flessibile.
      \item
            \textbf{Didattica individualizzata}: l\textquotesingle istruttore deve
            adattare la propria didattica alle differenze individuali, agli stili
            di apprendimento e al grado di autonomia di ciascun allievo.
      \item
            \textbf{Equilibrio spinta-protezione}: incentivare il bambino a
            cimentarsi con nuove attivita, ma proteggerlo da aspettative
            irrealistiche.
\end{enumerate}

\section{Area Psicopedagogica - Parte 3: Strategie Didattiche, Clima
  Psicologico, Allievo Adolescente e
  Adulto}\label{area-psicopedagogica---parte-3-strategie-didattiche-clima-psicologico-allievo-adolescente-e-adulto}

\textbf{Docente:} Paola Contardi, psicologa, psicoterapeuta e psicologa
dello sport

\begin{center}\rule{0.5\linewidth}{0.5pt}\end{center}

\subsection{Introduzione}\label{introduzione-6}

Questa terza parte dell\textquotesingle area psicopedagogica affronta
tre temi strettamente connessi alla pratica
dell\textquotesingle istruttore di nuoto: le \textbf{strategie
      didattiche} e il \textbf{clima psicologico} da costruire in vasca, le
specificità dell\textquotesingle{}\textbf{allievo adolescente} e quelle
dell\textquotesingle{}\textbf{allievo adulto e anziano}. Conoscere le
caratteristiche di ciascuna fascia di età consente
all\textquotesingle istruttore di calibrare il proprio intervento in
modo professionale ed efficace.

\begin{center}\rule{0.5\linewidth}{0.5pt}\end{center}

\subsection{1. Strategie Didattiche}\label{strategie-didattiche}

\subsubsection{1.1 Elementi da
      Considerare}\label{elementi-da-considerare}

Le strategie didattiche devono tenere conto di:

\begin{itemize}
      \tightlist
      \item
            \textbf{Età dell\textquotesingle allievo}
      \item
            \textbf{Caratteristiche psicofisiche}
      \item
            \textbf{Motivazioni} verso lo sport
      \item
            \textbf{Atteggiamento} verso l\textquotesingle attività sportiva
\end{itemize}

Queste informazioni permettono di individuare \textbf{risorse e limiti}
di ogni allievo.

\subsubsection{1.2 Elementi Universali della
      Didattica}\label{elementi-universali-della-didattica}

Due principi valgono per qualsiasi contesto e fascia di età:

\begin{enumerate}
      \tightlist
      \item
            \textbf{Rispetto per l\textquotesingle allievo:} mettersi nei suoi
            panni, riconoscerne il valore come persona indipendentemente dai
            risultati e proporre attività coerenti con questa visione.
      \item
            \textbf{Aspetto realizzativo:} fare leva sulla \emph{base
                  motivazionale} di ciascun individuo, ovvero il desiderio universale di
            vedere riconosciuto il proprio valore e di essere gratificati. Alla
            base di questo principio si trovano gli studi di \textbf{Abraham
                  Maslow} sulla motivazione.
\end{enumerate}

\subsubsection{1.3 Il Ruolo Morale e Professionale
      dell\textquotesingle Istruttore}\label{il-ruolo-morale-e-professionale-dellistruttore}

L\textquotesingle istruttore ha il dovere morale e professionale di:

\begin{itemize}
      \tightlist
      \item
            Considerare risorse e limiti del singolo allievo e del gruppo
      \item
            Riconoscere i motivi che hanno spinto l\textquotesingle allievo a
            frequentare la piscina
      \item
            Immedesimarsi nel vissuto dell\textquotesingle allievo, cogliendo
            desideri e rispettando i tempi
      \item
            Proporre \textbf{obiettivi raggiungibili}: obiettivi troppo lontani
            dalle risorse percepite generano atteggiamenti rinunciatari, perché
            valorizzano i limiti anziché le risorse
      \item
            Apprezzare non solo le doti tecniche, ma anche: impegno, puntualità,
            determinazione, miglioramenti progressivi
\end{itemize}

\begin{quote}
      La pratica sportiva insegna che per ottenere risultati bisogna
      impegnarsi, faticare e saper aspettare: un insegnamento di grande valore
      etico e morale, in contrasto con la filosofia del "tutto e subito".
\end{quote}

\subsubsection{1.4 Linee Guida per una Comunicazione Didattica
      Efficace}\label{linee-guida-per-una-comunicazione-didattica-efficace}

La competenza comunicativa dell\textquotesingle istruttore è in grado di
valorizzare --- o invalidare --- anche il miglior curriculum
professionale. Le regole fondamentali:

\begin{itemize}
      \tightlist
      \item
            Rivolgersi agli allievi \textbf{chiamandoli sempre per nome}
      \item
            Mettere in evidenza le \textbf{doti e l\textquotesingle impegno}
      \item
            Tradurre l\textquotesingle errore in un \textbf{consiglio pratico}
            (\emph{critica costruttiva})
      \item
            Essere \textbf{specifici e chiari} sia nello spiegare gli esercizi che
            nel definire i tempi di esecuzione
      \item
            Parlare sempre in forma \textbf{semplice e affermativa}
      \item
            \textbf{Criticare solo l\textquotesingle azione} (il compito), mai la
            persona: ogni essere umano ha il diritto di essere corretto
            nell\textquotesingle azione, non mortificato nella sua dignità
\end{itemize}

\subsubsection{1.5 La Programmazione delle
      Attività}\label{la-programmazione-delle-attivituxe0}

L\textquotesingle attività didattica non può essere frutto di
improvvisazione. In ogni piscina esistono appuntamenti fissi come gare
sociali e brevetti tecnici: in linea di massima, \textbf{ogni trimestre
      ha un obiettivo specifico}.

Il processo di programmazione prevede:

\begin{enumerate}
      \tightlist
      \item
            Analisi attenta del \textbf{livello di partenza}
      \item
            Definizione di \textbf{obiettivi intermedi} progressivi
      \item
            Raggiungimento dell\textquotesingle{}\textbf{obiettivo finale}
      \item
            \textbf{Coinvolgimento degli allievi}, promuovendo autonomia e senso
            di \emph{autoefficacia}
\end{enumerate}

\begin{center}\rule{0.5\linewidth}{0.5pt}\end{center}

\subsection{2. Il Clima Psicologico}\label{il-clima-psicologico}

\subsubsection{2.1 Clima per Allievi Adolescenti e
      Adulti}\label{clima-per-allievi-adolescenti-e-adulti}

Un buon clima psicologico si costruisce attraverso:

\begin{itemize}
      \tightlist
      \item
            \textbf{Progettualità} e \textbf{divertimento}
      \item
            \textbf{Apertura}, collaborazione e dialogo
      \item
            Conoscenza reciproca tra i membri del gruppo
      \item
            Chiarire le aspettative personali degli allievi e i motivi che li
            hanno spinti a scegliere il nuoto
\end{itemize}

L\textquotesingle istruttore deve:

\begin{itemize}
      \tightlist
      \item
            Distribuire attenzioni e tempo in modo \textbf{equo ed equilibrato}
            tra tutti gli allievi
      \item
            Utilizzare \textbf{spiegazioni chiare} con rinforzo positivo
      \item
            Sottolineare il progresso tra la fase iniziale e il momento presente
      \item
            Usare la \textbf{terminologia tecnica} con gli allievi più grandi,
            facendoli sentire parte integrante dello sport scelto
      \item
            Essere aggiornato sugli \textbf{avvenimenti sportivi} (gare, record,
            personaggi) e sulle revisioni del regolamento (ogni quadriennio
            olimpico)
      \item
            Proporre esercizi a \textbf{piccoli gruppi} con finalità tecniche
            diversificate, favorendo la collaborazione (es. lavori di
            senso-percezione)
      \item
            Organizzare, laddove possibile, \textbf{attività collaterali} (terme,
            musei, teatri), specialmente con gli anziani a rischio di ritiro
            sociale
\end{itemize}

\subsubsection{2.2 Clima per i Bambini}\label{clima-per-i-bambini}

Per i bambini l\textquotesingle interesse prioritario è giocare, quindi
la lezione in piscina deve essere un \textbf{grande gioco} fatto di
attività facili e un po\textquotesingle{} più difficili, ma sempre
divertenti.

Il \textbf{divertimento} è l\textquotesingle elemento chiave: riesce a
coinvolgere anche i bambini più riservati e insicuri.

\subsubsection{2.3 Le Regole come Strumento
      Educativo}\label{le-regole-come-strumento-educativo}

L\textquotesingle attività educativa richiede
un\textquotesingle impostazione disciplinare. Lo sport è fatto di
regole, e anche l\textquotesingle ambiente piscina ne impone alcune
(spazi assegnati, doccia obbligatoria, lavapiedi, ecc.).
L\textquotesingle istruttore deve definire regole specifiche, ad
esempio:

\begin{itemize}
      \tightlist
      \item
            Si entra in acqua solo in presenza dell\textquotesingle istruttore
      \item
            Il materiale didattico va ricollocato negli spazi dedicati
      \item
            Lo spazio tra un allievo e l\textquotesingle altro durante gli
            esercizi
      \item
            Quando e come uscire dall\textquotesingle acqua
\end{itemize}

\textbf{Poche regole, ma sostanziali}: evitano interpretazioni
soggettive, favoriscono collaborazione e impegno, rendono
l\textquotesingle ambiente sicuro e fonte di serenità.

\paragraph{Effetti dei due climi educativi estremi sui
      bambini}\label{effetti-dei-due-climi-educativi-estremi-sui-bambini}

\begin{longtable}[]{@{}p{0.30\linewidth}p{0.30\linewidth}p{0.30\linewidth}@{}}
      \toprule\noalign{}
      Clima educativo                    & Caratteristiche                      & Conseguenze \\
      \midrule\noalign{}
      \endhead
      \bottomrule\noalign{}
      \endlastfoot
      \textbf{Eccessivamente protettivo} & Troppe regole, adulti che
      impediscono esplorazione           & Dipendenza dagli adulti, mancanza di
      autostima e autoefficacia                                                               \\
      \textbf{Lassista}                  & Assenza di regole o regole basate
      sull\textquotesingle emotività     & Forte insicurezza, stato di allerta
      permanente, difficoltà nell\textquotesingle orientarsi                                  \\
\end{longtable}

\begin{quote}
      Le regole, se stabili e uguali per tutti, diventano \textbf{amiche con
            le quali crescere}: permettono al bambino di percepire confini che
      amplificano spazi e fantasia. \textbf{Autonomia e autostima camminano di
            pari passo.}
\end{quote}

\begin{center}\rule{0.5\linewidth}{0.5pt}\end{center}

\subsection{3. L\textquotesingle Allievo
      Adolescente}\label{lallievo-adolescente}

\subsubsection{3.1 Caratteristiche Generali
      dell\textquotesingle Adolescenza}\label{caratteristiche-generali-delladolescenza}

L\textquotesingle adolescenza è il periodo caratterizzato da grande
\textbf{instabilità e turbolenza}: difficoltà e malesseri si affiancano
a entusiasmo e vitalità.

\begin{itemize}
      \tightlist
      \item
            Inizio della \textbf{pubertà}: femmine tra i 10 e i 14 anni, maschi
            tra gli 11 e i 15
      \item
            Tarda adolescenza e acquisizione di autogestione sociale: oltre i 20
            anni
      \item
            Gli adolescenti non sono più bambini, ma non sono ancora adulti:
            lottano per l\textquotesingle indipendenza pur dipendendo ancora dai
            genitori
\end{itemize}

\subsubsection{3.2 Cambiamenti Biologici}\label{cambiamenti-biologici}

La cosiddetta \emph{proceritas secunda} scatena una \textbf{tempesta
      ormonale} con enormi influenze sulle trasformazioni biologiche e
psicologiche:

\begin{itemize}
      \tightlist
      \item
            Gli ormoni determinano la \textbf{vagotonia} tipica di questa età
      \item
            L\textquotesingle{}\textbf{irregolarità endocrina} è responsabile
            dello squilibrio emozionale
      \item
            Forte crescita della muscolatura degli arti inferiori (allungamento
            delle fibre muscolari)
      \item
            Gli apparati cardiocircolatorio e respiratorio aumentano la capacità
            funzionale
      \item
            Il \textbf{sistema nervoso centrale} completa la maturazione,
            permettendo il perfezionamento della coordinazione e lo sviluppo del
            \textbf{ragionamento astratto} (capacità logico-deduttiva)
      \item
            Aumenta la capacità \textbf{analitica e introspettiva}
\end{itemize}

\subsubsection{3.3 Conseguenze per lo Schema
      Corporeo}\label{conseguenze-per-lo-schema-corporeo}

Lo schema corporeo del preadolescente diventa \textbf{disarmonico e
      sgraziato}: i movimenti sono goffi e imprecisi, si rompe
l\textquotesingle equilibrio staturo-ponderale. Questo può determinare
la \textbf{spiacevole sensazione di perdita di efficacia della nuotata},
che può essere causa di abbandono precoce.

L\textquotesingle istruttore deve \textbf{informare e rendere
      consapevole} l\textquotesingle allievo adolescente delle cause di questa
perdita di efficacia.

\subsubsection{3.4 Caratteristiche
      Psicologiche}\label{caratteristiche-psicologiche}

\begin{itemize}
      \tightlist
      \item
            L\textquotesingle aspetto fisico diventa un interesse prioritario
            (vestiti, accessori, taglio di capelli, piercing)
      \item
            Nel primo periodo dell\textquotesingle adolescenza: maggiore interesse
            per le persone dello \textbf{stesso sesso} (per rinforzare
            l\textquotesingle identità di genere)
      \item
            Successivamente: interesse verso il \textbf{sesso opposto}, primi
            intensi innamoramenti e prime delusioni
      \item
            La sofferenza per la fine di un innamoramento adolescenziale va
            \textbf{presa molto sul serio}, poiché può innescare un malessere
            profondo ed esteso
      \item
            Il desiderio di accettazione e approvazione aumenta la sensibilità al
            rifiuto
\end{itemize}

\subsubsection{3.5 Benefici dello Sport durante
      l\textquotesingle Adolescenza}\label{benefici-dello-sport-durante-ladolescenza}

Lo sport in questa fase:

\begin{itemize}
      \tightlist
      \item
            Favorisce la costruzione del \textbf{gruppo dei pari} e la
            condivisione di valori, norme e obiettivi
      \item
            Promuove l\textquotesingle ampliamento del gruppo di riferimento delle
            figure adulte
      \item
            Permette la \textbf{riscoperta del proprio corpo} attraverso
            l\textquotesingle acqua (attività natatoria)
      \item
            Favorisce la percezione di risorse e limiti
\end{itemize}

\subsubsection{3.6 Strategie Educative con gli
      Adolescenti}\label{strategie-educative-con-gli-adolescenti}

La strada maestra per instaurare un corretto rapporto educativo si basa
su due elementi:

\begin{enumerate}
      \tightlist
      \item
            \textbf{Accettare l\textquotesingle esuberanza} e aprire il dialogo
            per modulare e correggere comportamenti non idonei
      \item
            \textbf{Mostrare apprezzamento} per i comportamenti rispettosi e
            diligenti, che spesso nessuno nota o sottolinea
\end{enumerate}

\paragraph{Il Locus of Control}\label{il-locus-of-control}

Il \emph{locus of control} indica dove l\textquotesingle individuo
colloca il centro di responsabilità della propria vita:

\begin{itemize}
      \tightlist
      \item
            \textbf{Esterno} (fortuna, destino, caso): sensazione di impotenza,
            scarsa efficacia nel poter incidere sugli eventi
      \item
            \textbf{Interno}: consapevolezza della propria capacità di influenzare
            le situazioni
\end{itemize}

Attraverso la pratica sportiva, l\textquotesingle istruttore può
facilitare l\textquotesingle{}\textbf{empowerment} degli allievi
adolescenti, promuovendo la consapevolezza di limiti e risorse e
favorendo lo sviluppo dell\textquotesingle{}\emph{autoefficacia}.

\paragraph{La Resilienza}\label{la-resilienza}

La \textbf{resilienza} è la capacità di affrontare e superare le
avversità della vita. Chi è resiliente:

\begin{itemize}
      \tightlist
      \item
            È ottimista
      \item
            È consapevole delle proprie risorse
      \item
            Vede nei cambiamenti delle opportunità
      \item
            Mantiene la speranza per il futuro di fronte alle frustrazioni
\end{itemize}

\begin{quote}
      Il mondo dello sport è l\textquotesingle unico a poter promuovere in
      modo strutturato la resilienza. Lo sport agonistico è una vera
      \emph{metafora della vita}: \textbf{l\textquotesingle atteggiamento vale
            quanto i muscoli}.
\end{quote}

\begin{center}\rule{0.5\linewidth}{0.5pt}\end{center}

\subsection{4. L\textquotesingle Allievo Adulto e
      Anziano}\label{lallievo-adulto-e-anziano}

\subsubsection{4.1 Il Contesto
      Demografico}\label{il-contesto-demografico}

Il numero crescente della popolazione anziana richiederà alla società
interventi specifici. Gli istruttori, insieme agli impianti natatori,
devono essere pronti ad accogliere questa tipologia di utenza.

\subsubsection{4.2 Processi di
      Invecchiamento}\label{processi-di-invecchiamento}

I principali cambiamenti fisiologici legati all\textquotesingle età:

\begin{itemize}
      \tightlist
      \item
            \textbf{Apparato sensoriale:} perde efficienza
      \item
            \textbf{Sistema cardiovascolare:} le pareti vascolari diventano più
            rigide, la portata del cuore diminuisce
      \item
            \textbf{Apparato respiratorio:} riduzione della capacità funzionale
      \item
            \textbf{Sistema muscolo-scheletrico:} le articolazioni perdono
            elasticità, si riduce la massa ossea (in particolare nelle donne, con
            rischio di osteoporosi)
      \item
            \textbf{Capacità cognitive:} decadimento progressivo
\end{itemize}

\subsubsection{4.3 Benefici dell\textquotesingle Esercizio Fisico
      nell\textquotesingle Anziano}\label{benefici-dellesercizio-fisico-nellanziano}

Lo sport \textbf{può fare veramente la differenza}. I benefici
documentati:

\begin{itemize}
      \tightlist
      \item
            Favorisce il \textbf{metabolismo osseo} e lo sviluppo della massa
            scheletrica, riducendo il rischio di fratture e osteoporosi
      \item
            Fa diminuire l\textquotesingle uso di \textbf{analgesici}
      \item
            Aumenta la funzionalità nelle \textbf{abilità quotidiane}
      \item
            Migliora il \textbf{tono muscolare}, prevenendo il rischio di cadute
      \item
            Aumenta il consumo calorico e il \textbf{metabolismo basale}
      \item
            Migliora la \textbf{circolazione} e la funzionalità cardiorespiratoria
      \item
            Aumenta la produzione endogena di \textbf{beta-endorfine}, che:

            \begin{itemize}
                  \tightlist
                  \item
                        Aumentano la tolleranza al dolore
                  \item
                        Migliorano l\textquotesingle appetito
                  \item
                        Riducono ansietà e tensione nervosa
                  \item
                        Riducono la reattività emozionale
            \end{itemize}
      \item
            Contribuisce a una generale \textbf{sensazione di benessere}
      \item
            Offre benefici \textbf{sociali}: possibilità di intrecciare e
            consolidare relazioni
\end{itemize}

\subsubsection{4.4 La Dimensione Psicologica
      dell\textquotesingle Anziano}\label{la-dimensione-psicologica-dellanziano}

Il \textbf{pensionamento} rappresenta un momento critico
dell\textquotesingle esistenza: era il lavoro a garantire socialità e
autostima. La pensione può diventare l\textquotesingle inizio di un
\textbf{isolamento sociale}, una situazione potenzialmente traumatica.

Spesso l\textquotesingle anziano ha insicurezze che lo portano
inconsapevolmente a una condizione di \textbf{dipendenza}, con rischio
di apatia e perdita di libertà. Tuttavia, \textbf{essere anziani non è
      sinonimo di patologia}.

La \textbf{resilienza} --- legata alla consapevolezza di limiti e
risorse propri --- è anche per l\textquotesingle anziano una
caratteristica fondamentale per adattarsi alla nuova condizione.

\subsubsection{4.5 Strategie Didattiche per
      l\textquotesingle Anziano}\label{strategie-didattiche-per-lanziano}

\textbf{Obiettivo primario:} mantenere un funzionamento ottimale.

Per impostare una seduta di attività acquatica per la terza età:

\begin{itemize}
      \tightlist
      \item
            Proporre attività \textbf{semplici e motivanti}
      \item
            Lasciare ampio spazio alla \textbf{fantasia}
      \item
            Evitare attività ad alto rischio traumatico
\end{itemize}

\textbf{Elementi chiave} per l\textquotesingle attività sportiva di
adulti e anziani:

\begin{longtable}[]{@{}p{0.47\linewidth}p{0.47\linewidth}@{}}
      \toprule\noalign{}
      Capacità                          & Indicazione                    \\
      \midrule\noalign{}
      \endhead
      \bottomrule\noalign{}
      \endlastfoot
      Coordinazione dinamica generale   & Prioritaria                    \\
      Forza                             & Con moderazione                \\
      Equilibrio miotensivo e posturale & Fondamentale                   \\
      Resistenza                        & Da sviluppare progressivamente \\
\end{longtable}

\subsubsection{4.6 La Socializzazione
      Trasversale}\label{la-socializzazione-trasversale}

L\textquotesingle istruttore farà un grande servizio favorendo la
\textbf{socializzazione trasversale}, cioè
l\textquotesingle integrazione tra diverse fasce di età:

\begin{itemize}
      \tightlist
      \item
            L\textquotesingle anziano trae vantaggio dalla vitalità e
            dall\textquotesingle entusiasmo dei più giovani
      \item
            I giovani adulti possono abbandonare una visione esclusivamente
            competitiva e scoprire lo sport come veicolo di benessere e scambio
            umano
\end{itemize}

Decidere di fare sport in età avanzata significa:

\begin{itemize}
      \tightlist
      \item
            Assumere un atteggiamento diverso nei confronti del proprio corpo
      \item
            Accettare e imparare i sistemi migliori per un invecchiamento
            \textbf{equilibrato e più lento possibile}
      \item
            Prendere in mano la propria vita, manifestare apertura verso
            l\textquotesingle ambiente e le persone
\end{itemize}

\begin{center}\rule{0.5\linewidth}{0.5pt}\end{center}

\subsection{Conclusione: Punti Chiave per
      l\textquotesingle Esame}\label{conclusione-punti-chiave-per-lesame-2}

\begin{enumerate}
      \item
            \textbf{Strategie didattiche universali:} rispetto per
            l\textquotesingle allievo e valorizzazione
            dell\textquotesingle aspetto realizzativo (Maslow). Gli obiettivi
            devono essere percepiti come raggiungibili.
      \item
            \textbf{Comunicazione efficace:} chiamare per nome, critica
            costruttiva, linguaggio affermativo, criticare
            l\textquotesingle azione mai la persona.
      \item
            \textbf{Programmazione:} trimestrale, con analisi del livello di
            partenza, obiettivi intermedi e coinvolgimento degli allievi.
      \item
            \textbf{Clima psicologico:} divertimento per i bambini; progettualità,
            dialogo e terminologia tecnica per adolescenti e adulti. Le regole ---
            poche e stabili --- rendono l\textquotesingle ambiente sicuro e
            favoriscono l\textquotesingle autostima.
      \item
            \textbf{Adolescente:} caratterizzato da tempesta ormonale
            (\emph{vagotonia}), schema corporeo disarmonico, sviluppo del
            ragionamento astratto. Può abbandonare lo sport per perdita di
            efficacia della nuotata. Strumenti chiave: locus of control interno,
            empowerment, resilienza.
      \item
            \textbf{Adulto e anziano:} i cambiamenti fisiologici sono reali ma non
            sono patologia. L\textquotesingle esercizio fisico contrasta il
            decadimento, le beta-endorfine migliorano il benessere psicofisico. Il
            pensionamento è un momento di rischio sociale: il gruppo sportivo è un
            antidoto fondamentale.
      \item
            \textbf{Socializzazione trasversale:} l\textquotesingle integrazione
            tra fasce di età diverse è un valore aggiunto che
            l\textquotesingle istruttore dovrebbe attivamente promuovere.
\end{enumerate}

\section{Elementi di Primo Soccorso - Parte
  I}\label{elementi-di-primo-soccorso---parte-i}

\begin{center}\rule{0.5\linewidth}{0.5pt}\end{center}

\subsection{Introduzione}\label{introduzione-7}

Questa lezione introduce i principi fondamentali del \textbf{primo
      soccorso} applicati al contesto degli impianti natatori.
L\textquotesingle istruttore di nuoto non è un operatore sanitario, ma
ha il dovere morale e legale di saper riconoscere
un\textquotesingle emergenza, attivare correttamente il sistema di
soccorso e intervenire nelle fasi iniziali con le manovre appropriate.
La lezione affronta il sistema di emergenza territoriale, le regole
generali di comportamento, i disturbi della coscienza,
l\textquotesingle epilessia e le emorragie.

\begin{center}\rule{0.5\linewidth}{0.5pt}\end{center}

\subsection{1. Il Sistema di Emergenza
      Territoriale}\label{il-sistema-di-emergenza-territoriale}

\subsubsection{1.1 Il Numero Unico di Emergenza:
      112}\label{il-numero-unico-di-emergenza-112}

Il \textbf{112} è il \emph{Numero Unico di Emergenza} (NUE), attivo in
tutta Italia dal 1996. È un numero a tre cifre, di facile
memorizzazione, che può essere chiamato:

\begin{itemize}
      \tightlist
      \item
            gratuitamente da qualsiasi telefono, inclusi i cellulari \textbf{senza
                  credito residuo}
      \item
            in qualsiasi regione italiana (in alcune aree è ancora attivo il 118:
            se si chiama il 112 dove non è ancora attivo, rispondono i carabinieri
            e trasferiscono la chiamata al numero dedicato
            all\textquotesingle emergenza sanitaria)
\end{itemize}

Esiste una \textbf{rete di backup} su ponti radio dedicata alla sanità:
anche in caso di blackout delle centraline dei gestori telefonici, il
sistema di emergenza continua a funzionare.

\subsubsection{1.2 L\textquotesingle Applicazione
      WhereAreU}\label{lapplicazione-whereareu}

L\textquotesingle app \textbf{WhereAreU} permette di chiamare il 112 con
funzioni avanzate:

\begin{itemize}
      \tightlist
      \item
            \textbf{Geolocalizzazione automatica} tramite GPS: la centrale
            operativa conosce la posizione del chiamante senza che questi debba
            spiegare dove si trova
      \item
            Scelta del tipo di soccorso: sanitario, forze
            dell\textquotesingle ordine, antincendio
      \item
            Chiamata \textbf{vocale}, chiamata in formato \textbf{testo} oppure in
            \textbf{modalità muta} (microfono ambientale attivo: la centrale sente
            l\textquotesingle ambiente circostante senza che il chiamante parli,
            utile in situazioni di pericolo)
\end{itemize}

\subsubsection{1.3 Come Comunicare con il
      112}\label{come-comunicare-con-il-112}

Chiamare il 112 chiedendo genericamente "mandami
un\textquotesingle ambulanza" non è sufficiente.
L\textquotesingle operatore ha bisogno di informazioni precise per
decidere quante e quali risorse inviare. Il soccorritore deve essere
\textbf{gli occhi della centrale} e riferire:

\begin{itemize}
      \tightlist
      \item
            dove si trova (indirizzo preciso)
      \item
            cosa sta accadendo
      \item
            quante persone sono coinvolte
      \item
            se il soggetto è cosciente o no
      \item
            se sta respirando
\end{itemize}

\subsubsection{1.4 Organizzazione del Sistema
      118/112}\label{organizzazione-del-sistema-118112}

Sul territorio sono presenti le \textbf{PET} (\emph{Postazioni di
      Emergenza Territoriale}), con medici e infermieri, che coprono il
territorio in modo omogeneo. I tempi medi di intervento sono:

\begin{longtable}[]{@{}p{0.47\linewidth}p{0.47\linewidth}@{}}
      \toprule\noalign{}
      Area         & Tempo massimo di intervento \\
      \midrule\noalign{}
      \endhead
      \bottomrule\noalign{}
      \endlastfoot
      Urbana       & 8-10 minuti                 \\
      Extra-urbana & 20 minuti                   \\
\end{longtable}

Il sistema è organizzato a \textbf{cerchi concentrici} per coprire
l\textquotesingle intera penisola. Per eventi gravi (infarto, ictus) la
centrale operativa non indirizza al pronto soccorso più vicino, ma
all\textquotesingle{}\textbf{ospedale più idoneo} (ad esempio con
emodinamica o rianimazione), anche ricorrendo a mezzi più veloci come
l\textquotesingle elicottero.

\textbf{Obiettivi del sistema:}

\begin{itemize}
      \tightlist
      \item
            Ridurre le morti evitabili
      \item
            Diminuire i tassi di invalidità (intervenendo rapidamente si previene
            il peggioramento)
\end{itemize}

\begin{center}\rule{0.5\linewidth}{0.5pt}\end{center}

\subsection{2. Regole Generali del Primo
      Soccorso}\label{regole-generali-del-primo-soccorso}

\subsubsection{2.1 Cosa Fare Sempre}\label{cosa-fare-sempre}

\begin{itemize}
      \tightlist
      \item
            \textbf{Allertare il 112}: è il primo atto in qualsiasi emergenza.
            L\textquotesingle articolo \textbf{593 del Codice Penale} punisce il
            reato di \emph{omissione di soccorso}: chi trova un soggetto inanimato
            e non interviene né chiama i soccorsi può incorrere in responsabilità
            penale
      \item
            \textbf{Mettersi in sicurezza}: nessuno ha la maglietta di Superman.
            Intervenire in una situazione pericolosa senza i presidi adeguati
            rischia di aumentare il numero di vittime e rendere inefficace il
            soccorso
      \item
            \textbf{Raccogliere informazioni}: capire se si tratta di un
            \textbf{trauma} (es. caduta, scivolata) o di un \textbf{evento
                  medico/clinico} (es. infarto, ictus) è fondamentale perché cambia
            completamente il percorso di soccorso
      \item
            \textbf{Mantenere la calma}: in presenza di curiosi, assegnare loro un
            compito specifico (es. "vai ad aspettare l\textquotesingle ambulanza
            all\textquotesingle ingresso") permette di allontanarli senza entrare
            in conflitto, liberando il campo e preservando la concentrazione
      \item
            \textbf{Praticare il primo soccorso} se si è formati e capaci
      \item
            \textbf{Indossare i guanti in lattice}: negli impianti sportivi
            l\textquotesingle accesso è aperto a tutti, inclusi soggetti affetti
            da epatite B, epatite C o HIV, patologie che non richiedono
            certificazione preventiva per l\textquotesingle iscrizione. I guanti
            sono indispensabili ogni volta che si è in presenza di sangue o fluidi
            corporei
\end{itemize}

\subsubsection{2.2 Cosa Non Fare Mai}\label{cosa-non-fare-mai}

\begin{itemize}
      \tightlist
      \item
            \textbf{Non mettere mai a rischio la propria sicurezza}
      \item
            \textbf{Non spostare mai il paziente vittima di trauma}: un soggetto
            con frattura a collo, schiena o arti spostato senza presidi adeguati
            (stecobende, asse spinale, collare cervicale) rischia la paralisi.
            L\textquotesingle unica eccezione è il rischio imminente di esplosione
            o pericolo letale immediato
      \item
            \textbf{Non somministrare nulla per bocca}: né acqua, né acqua e
            zucchero, né sostanze alcoliche. Se il soggetto è privo di sensi,
            qualsiasi liquido può finire nelle vie aeree causando una
            \emph{polmonite ab ingestis}. Nota: le sostanze alcoliche,
            contrariamente alla credenza popolare, \textbf{non riscaldano}: danno
            solo la sensazione di calore a livello del sistema nervoso centrale,
            provocando in realtà una vasodilatazione periferica e un ulteriore
            raffreddamento del corpo
      \item
            \textbf{Non caricare il soggetto sulla prima macchina} e trasportarlo
            in ospedale: il sistema 112 deve essere sempre attivato; farlo
            autonomamente priva il paziente di cure appropriate durante il
            trasporto
\end{itemize}

\begin{center}\rule{0.5\linewidth}{0.5pt}\end{center}

\subsection{3. Disturbi della Coscienza}\label{disturbi-della-coscienza}

La \textbf{coscienza} è la consapevolezza di noi stessi e
dell\textquotesingle ambiente che ci circonda. Le principali alterazioni
sono due.

\subsubsection{3.1 Lipotemia (Presincope)}\label{lipotemia-presincope}

La \emph{lipotemia} è un malessere passeggero: sensazione di giramento
di testa, visione scura, debolezza. Il soggetto se ne accorge in tempo e
si siede. La sensazione passa spontaneamente.

\subsubsection{3.2 Sincope (Svenimento)}\label{sincope-svenimento}

La \emph{sincope} è una perdita di coscienza improvvisa con caduta a
terra per debolezza generalizzata e incapacità di mantenere la posizione
eretta.

\textbf{Segni e sintomi:}

\begin{itemize}
      \tightlist
      \item
            Pallore cutaneo
      \item
            Sudorazione fredda
      \item
            Visione confusa o annebbiata
      \item
            Senso di vertigine
      \item
            Polso carotideo o radiale debole
      \item
            Organi di senso appannati
\end{itemize}

\textbf{Cause possibili:}

\begin{longtable}[]{@{}p{0.47\linewidth}p{0.47\linewidth}@{}}
      \toprule\noalign{}
      Cause "soft"                                                          & Cause gravi                              \\
      \midrule\noalign{}
      \endhead
      \bottomrule\noalign{}
      \endlastfoot
      Affaticamento, digiuno prolungato, mancanza di sonno                  & Malattie
      cardiache (fibrillazione ventricolare, arresto cardiaco)                                                         \\
      Ambiente troppo caldo (es. genitori in piscina vestiti invernalmente) &
      Patologie del sistema nervoso centrale (ictus, epilessia)                                                        \\
      Dolore intenso (es. frattura)                                         & Alterazioni metaboliche (ipoglicemia nel
      diabetico, iperglicemia)                                                                                         \\
      Emozioni intense o stress                                             & Patologie polmonari                      \\
      Stare troppo a lungo in piedi (\emph{sindrome ortostatica})           &                                          \\
\end{longtable}

\textbf{Cosa fare:}

\begin{itemize}
      \tightlist
      \item
            Posizionare il soggetto supino e \textbf{sollevare gli arti inferiori
                  di circa 30 cm} (non a 90°: un angolo eccessivo stringe i vasi come un
            tubo piegato e impedisce il flusso sanguigno). Il sangue scende per
            gravità verso il torace e la testa, aumentando la pressione cerebrale
            e favorendo il recupero della coscienza
      \item
            \textbf{Slacciare} tutto ciò che è stretto: cravatte, cinture,
            cinturini. Nota storica: nei costumi gommati usati nel nuoto
            agonistico tra il 2009 e il 2010, molti atleti svenivano perché i
            costumi, indossati di taglia ridotta per migliorare le prestazioni,
            comprimevano il torace e impedivano la respirazione
      \item
            \textbf{Non somministrare nulla}
      \item
            Allontanare i curiosi
      \item
            In caso di \textbf{arresto cardiocircolatorio} (soggetto non respira,
            non ha attività circolatoria, non è cosciente): avviare immediatamente
            le manovre di \textbf{BLS} (\emph{Basic Life Support}) e, come
            previsto dalla legge per tutti gli impianti sportivi, utilizzare il
            \textbf{defibrillatore} (BLSD)
\end{itemize}

\begin{center}\rule{0.5\linewidth}{0.5pt}\end{center}

\subsection{4. Epilessia}\label{epilessia}

\subsubsection{4.1 Definizione e
      Diffusione}\label{definizione-e-diffusione}

L\textquotesingle{}\textbf{epilessia} è una malattia del sistema nervoso
centrale che interessa circa il \textbf{2\% della popolazione}, in tutte
le fasce di età. I soggetti epilettici possono ottenere
l\textquotesingle idoneità medico-sportiva sia agonistica che non
agonistica: \textbf{negli impianti natatori possono quindi essere
      presenti atleti epilettici}.

\subsubsection{4.2 Cause}\label{cause}

\begin{itemize}
      \tightlist
      \item
            \textbf{Congenite}: il soggetto nasce con questa condizione
      \item
            \textbf{Acquisite}: esito di trauma cerebrale importante, ictus,
            neoplasie cerebrali (l\textquotesingle attacco epilettico può essere
            il primo segnale di un tumore cerebrale), incremento termico
            importante (febbre alta)
\end{itemize}

\subsubsection{4.3 Forme di Epilessia}\label{forme-di-epilessia}

\paragraph{Grande Male (forma
      classica)}\label{grande-male-forma-classica}

Caratterizzato da \textbf{scosse tonico-cloniche} (contrazioni muscolari
involontarie violente). Può essere preceduto da una fase prodromica
chiamata \textbf{aura}: segnali premonitori come irritabilità, cefalea,
ansia, percezioni sensoriali anomale (odori strani, visioni alterate). I
soggetti epilettici esperti riconoscono l\textquotesingle aura e possono
assumere in anticipo la terapia.

\textbf{Fasi dell\textquotesingle attacco di grande male:}

\begin{enumerate}
      \tightlist
      \item
            Perdita di coscienza con caduta a terra
      \item
            Fase \textbf{tonica-clonica}: scosse violente, contrazione dei muscoli
            della faccia. Durata: \textbf{1-2 minuti} (sembra molto più lungo per
            chi assiste)
      \item
            Possibili: schiuma alla bocca, morso della lingua, perdita
            involontaria di feci e urine (rilascio degli sfinteri)
      \item
            Fase di \textbf{narcolessia}: sonno profondo post-ictale
      \item
            Al risveglio il soggetto \textbf{non ricorda nulla} di quanto accaduto
\end{enumerate}

\paragraph{Piccolo Male (assenze)}\label{piccolo-male-assenze}

Il soggetto perde il contatto con l\textquotesingle ambiente esterno,
non risponde alle domande, ha uno sguardo fisso nel vuoto. Può evolvere
in attacco tonico-clonico.

\subsubsection{4.4 Cosa Fare durante un Attacco
      Epilettico}\label{cosa-fare-durante-un-attacco-epilettico}

L\textquotesingle attacco epilettico ha un \textbf{alto impatto emotivo}
per chi non è abituato, ma se gestito correttamente \textbf{si risolve
      da solo} senza conseguenze gravi, lasciando solo stordimento e
stanchezza.

\begin{itemize}
      \tightlist
      \item
            \textbf{Non cercare di bloccare le scosse}: nascono direttamente dal
            sistema nervoso centrale, i movimenti degli arti sono riflessi.
            Tenerlo fermo è inutile
      \item
            \textbf{Evitare che il soggetto cada} se è ancora in piedi:
            accompagnarlo delicatamente a terra
      \item
            Posizionarlo su un \textbf{fianco} nella \emph{posizione laterale di
                  sicurezza}
      \item
            Mettere sotto la testa del \textbf{materiale morbido}
      \item
            \textbf{Slacciare} cravatte, cinture e indumenti stretti
      \item
            \textbf{Annotare il tempo} di inizio della fase convulsiva
      \item
            \textbf{Non lasciare mai il soggetto da solo}
      \item
            Controllare che siano presenti \textbf{polso} e \textbf{attività
                  respiratoria}
      \item
            Verificare che le \textbf{vie aeree siano libere}: se la base della
            lingua cade all\textquotesingle indietro ostruendo le vie aeree,
            eseguire la \textbf{manovra di iperestensione della testa} (portare la
            testa in estensione verso il basso, spostando la base della lingua)
      \item
            Se si tratta di un bambino e il genitore ha fornito il farmaco
            prescritto: somministrarlo nel \textbf{dosaggio e con le modalità
                  indicate dal medico}
\end{itemize}

\subsubsection{4.5 Cosa NON Fare durante un Attacco
      Epilettico}\label{cosa-non-fare-durante-un-attacco-epilettico}

\begin{itemize}
      \tightlist
      \item
            \textbf{Non aprire la bocca con la forza}: rischio elevato di essere
            morsi
      \item
            \textbf{Non inserire fazzoletti o materiale morbido in bocca}: inutile
            e pericoloso
      \item
            \textbf{Non usare il cuneo apribocca} presente in alcuni kit di primo
            soccorso: può provocare la frattura degli incisivi o dei denti, con
            conseguente \textbf{responsabilità civile} per il soccorritore.
            Semmai, si può inserire lateralmente (non al centro) qualcosa che
            impedisca la chiusura completa della bocca
\end{itemize}

\begin{center}\rule{0.5\linewidth}{0.5pt}\end{center}

\subsection{5. Emorragie}\label{emorragie}

\subsubsection{5.1 Definizione}\label{definizione}

Un\textquotesingle{}\textbf{emorragia} si verifica quando uno dei vasi
sanguigni (arterie, vene, capillari) si rompe e il sangue fuoriesce dal
sistema circolatorio chiuso. La gravità dipende dalla quantità di sangue
perso: oltre il \textbf{10\% del volume ematico totale} si può
instaurare uno \textbf{shock ipovolemico}, con calo pressorio e rischio
di arresto cardiaco.

\subsubsection{5.2 Classificazione per
      Sede}\label{classificazione-per-sede}

\begin{longtable}[]{@{}p{0.47\linewidth}p{0.47\linewidth}@{}}
      \toprule\noalign{}
      Tipo                                      & Descrizione                      \\
      \midrule\noalign{}
      \endhead
      \bottomrule\noalign{}
      \endlastfoot
      \textbf{Emorragia esterna}                & Il sangue fuoriesce visibilmente
      all\textquotesingle esterno del corpo                                        \\
      \textbf{Emorragia interna}                & Il sangue si raccoglie
      all\textquotesingle interno di cavità naturali o infiltra i tessuti; non
      si vede sangue rosso. Esempi: ematoma, ecchimosi, lesioni muscolari          \\
      \textbf{Emorragia interna esteriorizzata} & Il sangue si raccoglie
      internamente ma fuoriesce da orifizi naturali (bocca, orecchio, naso).
      \textbf{Non va tamponata} a differenza delle emorragie esterne               \\
\end{longtable}

\subsubsection{5.3 Classificazione per Tipo di
      Vaso}\label{classificazione-per-tipo-di-vaso}

\begin{longtable}[]{@{}p{0.30\linewidth}p{0.30\linewidth}p{0.30\linewidth}@{}}
      \toprule\noalign{}
      Tipo                                                     & Caratteristiche                                     & Gravità \\
      \midrule\noalign{}
      \endhead
      \bottomrule\noalign{}
      \endlastfoot
      \textbf{Arteriosa}                                       & Sangue rosso vivo, fuoriesce a \textbf{schizzi
      pulsanti} (risente dei battiti cardiaci), pressione alta & Molto
      grave                                                                                                                    \\
      \textbf{Venosa}                                          & Sangue rosso scuro/brunastro, fuoriesce in modo
      \textbf{continuo e lento}, pressione bassa               & Meno grave, più frequente
      (le vene sono più superficiali)                                                                                          \\
      \textbf{Capillare}                                       & Sangue rosso vivo che \textbf{stilla} attorno alla
      ferita, si arresta spesso da sola                        & Lieve                                                         \\
      \textbf{Mista}                                           & Arterie e vene scorrono in parallelo: un trauma può
      coinvolgerle entrambe contemporaneamente                 & Variabile                                                     \\
\end{longtable}

\subsubsection{5.4 Cosa Fare per le Emorragie
      Esterne}\label{cosa-fare-per-le-emorragie-esterne}

\begin{enumerate}
      \tightlist
      \item
            \textbf{Indossare i guanti in lattice} prima di intervenire (rischio
            epatite B, C, HIV)
      \item
            \textbf{Comprimere con forza} la ferita usando garze sterili: una
            pressione superiore a quella del sangue ne arresta il flusso
      \item
            Se l\textquotesingle emorragia non si arresta, \textbf{non rimuovere
                  il tampone}: aggiungere altro materiale sopra
      \item
            \textbf{Non usare il laccio emostatico} (pratica abbandonata)
      \item
            Quando il flusso rallenta, applicare un \textbf{bendaggio compressivo}
            (garze o fazzoletti puliti)
      \item
            Applicare \textbf{ghiaccio} sulla ferita: ha effetto vasocostrittore e
            riduce l\textquotesingle afflusso di sangue
      \item
            \textbf{Tranquillizzare il soggetto} e monitorare i segni di shock
            (pallore, sudorazione fredda, polso debole, respiro rapido e
            superficiale)
      \item
            Inviare il soggetto in \textbf{ospedale}
\end{enumerate}

\begin{quote}
      La frequenza respiratoria normale è di \textbf{12-16 atti/minuto}. Se il
      soggetto inizia a respirare a 25-40 atti/minuto in modo superficiale, è
      un segnale di shock imminente.
\end{quote}

\subsubsection{5.5 Cosa Fare per le Emorragie
      Interne}\label{cosa-fare-per-le-emorragie-interne}

Le emorragie interne non presentano sangue visibile
all\textquotesingle esterno. I campanelli d\textquotesingle allarme
sono:

\begin{itemize}
      \tightlist
      \item
            Pallore cutaneo
      \item
            Pelle fredda e sudata
      \item
            Polso debole
      \item
            Respiro rapido e superficiale
\end{itemize}

Per le forme lievi: applicare \textbf{ghiaccio}
sull\textquotesingle area interessata.\\
Per le forme gravi (es. trauma addominale, frattura del bacino):
\textbf{chiamare immediatamente il 112} e predisporre il trasporto
urgente in ospedale. Un\textquotesingle emorragia interna lenta ma
costante riempie progressivamente le cavità corporee, come una cannella
che perde riempie un secchio nel corso delle ore.

\begin{center}\rule{0.5\linewidth}{0.5pt}\end{center}

\subsection{Conclusione: Punti Chiave per
      l\textquotesingle Esame}\label{conclusione-punti-chiave-per-lesame-3}

\begin{enumerate}
      \item
            \textbf{112 e WhereAreU}: il numero unico di emergenza si chiama con o
            senza credito, da tutta Italia. L\textquotesingle app fornisce
            geolocalizzazione automatica e permette la chiamata in modalità muta.
      \item
            \textbf{Comunicare correttamente al 112}: riferire luogo, numero di
            vittime, stato di coscienza e respirazione. Essere gli "occhi" della
            centrale.
      \item
            \textbf{Art. 593 Codice Penale}: omissione di soccorso è un reato
            penale. Anche chi non è formato ha l\textquotesingle obbligo minimo di
            attivare il 112.
      \item
            \textbf{Sicurezza del soccorritore}: non intervenire mai in situazioni
            di pericolo senza presidi. Non spostare mai il traumatizzato (rischio
            paralisi), salvo pericolo di vita immediato.
      \item
            \textbf{Svenimento}: sollevare le gambe di circa 30 cm, slacciare gli
            indumenti stretti, non somministrare nulla. No acqua, no zucchero, no
            alcolici.
      \item
            \textbf{Epilessia}: interessa il 2\% della popolazione, anche atleti
            agonisti. L\textquotesingle attacco di grande male dura 1-2 minuti.
            Non bloccare le scosse, non aprire la bocca con la forza, non usare il
            cuneo apribocca. Posizione laterale di sicurezza, materiale morbido
            sotto la testa, verificare polso e respiro.
      \item
            \textbf{Emorragie}: arterioso = rosso vivo, pulsante, gravissimo;
            venoso = rosso scuro, continuo, più frequente; capillare = lieve, si
            arresta da sola. Emorragie interne esteriorizzate non vanno tamponate.
            Per le esterne: compressione con garze + guanti in lattice + bendaggio
            compressivo + ghiaccio.
      \item
            \textbf{Guanti in lattice}: obbligatori sempre in presenza di sangue,
            poiché negli impianti sportivi possono accedere soggetti con epatite
            B, C o HIV senza che sia richiesta alcuna certificazione preventiva.
\end{enumerate}

\section{Elementi di Primo Soccorso - Parte
  II}\label{elementi-di-primo-soccorso---parte-ii}

\begin{center}\rule{0.5\linewidth}{0.5pt}\end{center}

\subsection{Introduzione}\label{introduzione-8}

Questa seconda parte approfondisce le principali situazioni
traumatologiche e mediche che un istruttore di nuoto può trovarsi ad
affrontare sul bordo vasca o in palestra. Vengono trattati i traumi
muscolari e le relative fasi di trattamento, l\textquotesingle uso
corretto della crioterapia, le lesioni articolari (contusioni,
distorsioni, lussazioni), le fratture, il trauma cranico e le manovre di
\textbf{BLS} (\emph{Basic Life Support}) incluse la defibrillazione e la
manovra di Heimlich per l\textquotesingle ostruzione delle vie aeree.

\begin{center}\rule{0.5\linewidth}{0.5pt}\end{center}

\subsection{1. Lesioni Muscolari}\label{lesioni-muscolari}

\subsubsection{1.1 Principi Generali}\label{principi-generali}

Qualsiasi tipo di trauma produce un effetto distruttivo totale o
parziale delle fibre muscolari. Il tessuto muscolare ha una
\textbf{capacità di rigenerazione limitata}:

\begin{itemize}
      \tightlist
      \item
            Le lesioni \textbf{minori} si rigenerano quasi totalmente
      \item
            Le lesioni \textbf{gravi} portano alla formazione di un
            \textbf{tessuto cicatriziale fibroso} che sostituisce le fibre
            danneggiate
\end{itemize}

\subsubsection{1.2 Lesioni Indirette: Fattori
      Predisponenti}\label{lesioni-indirette-fattori-predisponenti}

Le lesioni da sovraccarico o da \emph{overuse} sono frequenti anche nel
nuoto. I principali fattori di rischio sono:

\begin{itemize}
      \tightlist
      \item
            \textbf{Riscaldamento inadeguato}: troppo breve, troppo corto o non
            idoneo
      \item
            \textbf{Condizioni climatiche avverse}: freddo, umidità eccessiva
      \item
            \textbf{Deficit di forza o flessibilità} muscolare
      \item
            \textbf{Ripresa precoce} dell\textquotesingle attività prima del
            completo recupero
      \item
            \textbf{Dismorfismi muscolari}: muscolo non in condizioni ottimali
      \item
            \textbf{Disidratazione}: perdita significativa di liquidi che
            compromette la funzionalità muscolare
      \item
            \textbf{Squilibri idrosalini}: carenza di sali minerali o deficit
            nutrizionali
      \item
            \textbf{Stanchezza}: sia locale che generale (nel calcio, ad esempio,
            la maggior parte delle lesioni muscolari avviene a fine partita)
      \item
            \textbf{Alterata coordinazione neuromuscolare} o deficit
            propriocettivi
\end{itemize}

\subsubsection{1.3 Lesioni da Dolore a Insorgenza Tardiva:
      DOMS}\label{lesioni-da-dolore-a-insorgenza-tardiva-doms}

I \textbf{DOMS} (\emph{Delayed Onset Muscle Soreness}) sono dolori
muscolari che compaiono 3-4 giorni dopo uno sforzo insolito o nuovo. Si
verificano tipicamente a inizio preparazione o quando si introducono
esercizi mai svolti in precedenza.

\subsubsection{1.4 Le Tre Fasi del
      Trattamento}\label{le-tre-fasi-del-trattamento}

\paragraph{Fase 1 --- Prime 48 ore (competenza
      dell\textquotesingle istruttore)}\label{fase-1-prime-48-ore-competenza-dellistruttore}

Il metodo di riferimento è l\textquotesingle acronimo \textbf{RICE}:

\begin{longtable}[]{@{}p{0.30\linewidth}p{0.30\linewidth}p{0.30\linewidth}@{}}
      \toprule\noalign{}
      Lettera    & Termine            & Significato pratico                    \\
      \midrule\noalign{}
      \endhead
      \bottomrule\noalign{}
      \endlastfoot
      \textbf{R} & \emph{Rest}        & Riposo immediato, sospensione
      dell\textquotesingle attività                                            \\
      \textbf{I} & \emph{Ice}         & Applicazione del ghiaccio              \\
      \textbf{C} & \emph{Compression} & Bendaggio compressivo per ridurre la
      diffusione dell\textquotesingle edema                                    \\
      \textbf{E} & \emph{Elevation}   & Sollevamento dell\textquotesingle arto
      infortunato                                                              \\
\end{longtable}

La somministrazione di farmaci (FANS, antinfiammatori, antiedematosi) è
di \textbf{esclusiva competenza medica}.

\paragraph{Fase 2 --- Dopo le 48 ore
      (fisioterapista/medico)}\label{fase-2-dopo-le-48-ore-fisioterapistamedico}

Si può intervenire con:

\begin{itemize}
      \tightlist
      \item
            Laser, ultrasuoni, ionoforesi, magnetoterapia, ipertermia
      \item
            I massaggi (\textbf{massoterapia}) non vanno eseguiti prima di una
            settimana dalla lesione: rischiano di aggravare la situazione
\end{itemize}

\paragraph{Fase 3 --- Riatletizzazione}\label{fase-3-riatletizzazione}

Fase molto delicata in cui l\textquotesingle atleta deve recuperare:

\begin{itemize}
      \tightlist
      \item
            \textbf{Elasticità muscolare} (ritorno alle condizioni pre-infortunio)
      \item
            \textbf{Forza muscolare}
\end{itemize}

Accorciare artificialmente questa fase aumenta significativamente il
rischio di ricaduta. Il passaggio al preparatore atletico per il
recupero funzionale completo è fondamentale prima del ritorno
all\textquotesingle attività a pieno regime.

\begin{center}\rule{0.5\linewidth}{0.5pt}\end{center}

\subsection{2. Crioterapia: Uso Corretto del
      Ghiaccio}\label{crioterapia-uso-corretto-del-ghiaccio}

Il ghiaccio produce una \textbf{vasocostrizione locale} iniziale che
riduce l\textquotesingle afflusso di sangue, limita il versamento e
blocca il riassorbimento delle sostanze infiammatorie rilasciate dal
tessuto danneggiato. Prima si applica, migliore sarà il risultato.

\subsubsection{2.1 Modalità di
      Applicazione}\label{modalituxe0-di-applicazione}

\begin{itemize}
      \tightlist
      \item
            Durata di ogni applicazione: \textbf{circa 20 minuti}
      \item
            Frequenza: \textbf{ogni 2-4 ore}
      \item
            Periodo: \textbf{3-4 giorni} (con diminuzione progressiva
            dell\textquotesingle utilizzo)
\end{itemize}

\subsubsection{2.2 Tipi di Ghiaccio a
      Confronto}\label{tipi-di-ghiaccio-a-confronto}

\begin{longtable}[]{@{}p{0.22\linewidth}p{0.22\linewidth}p{0.22\linewidth}p{0.22\linewidth}@{}}
      \toprule\noalign{}
      Tipo                                              & Descrizione                                               & Pro             & Contro \\
      \midrule\noalign{}
      \endhead
      \bottomrule\noalign{}
      \endlastfoot
      \textbf{Ghiaccio umido}                           & Cubetti del congelatore in sacchetti o borsa
      del ghiaccio                                      & Economico, riutilizzabile, ideale in impianti con bambini
                                                        & Richiede congelatore                                                                 \\
      \textbf{Ghiaccio secco}                           & Bustine chimiche monouso (si rompe la bolla
      interna e i reagenti si mescolano)                & Pratico, pronto
      all\textquotesingle uso, non richiede congelatore & Monouso; se la
      bustina si lesiona, i reagenti interni sono \textbf{tossici} e non vanno
      toccati                                                                                                                                  \\
      \textbf{Ghiaccio spray}                           & Spray refrigerante                                        & Rapido e comodo &
      Raffreddamento superficiale e non duraturo; \textbf{pericoloso} se usato
      in modo scorretto                                                                                                                        \\
\end{longtable}

\textbf{Utilizzo corretto del ghiaccio spray:}

\begin{itemize}
      \tightlist
      \item
            Spruzzare a una distanza di \textbf{10-15 cm} dalla zona interessata
      \item
            Fermarsi \textbf{non appena compare l\textquotesingle alone
                  grigiastro}: procedere oltre rischia di creare danni tissutali e
            lesioni infettive
\end{itemize}

\begin{center}\rule{0.5\linewidth}{0.5pt}\end{center}

\subsection{3. Contusioni e Distorsioni}\label{contusioni-e-distorsioni}

\subsubsection{3.1 Contusione}\label{contusione}

La \textbf{contusione} è una lesione della parete cutanea e del muscolo
sottostante causata da un trauma diretto (impatto con un corpo estraneo)
che \textbf{non provoca l\textquotesingle apertura della cute}. Può
essere associata a un ematoma sottostante.

\subsubsection{3.2 Distorsione}\label{distorsione}

La \textbf{distorsione} è l\textquotesingle allontanamento temporaneo
dei capi articolari, che poi ritornano in sede. Può essere complicata
da:

\begin{itemize}
      \tightlist
      \item
            Versamento di sangue nell\textquotesingle articolazione
      \item
            Frattura associata
      \item
            Lesioni legamentose (es. crociati del ginocchio, menischi)
\end{itemize}

Esempi frequenti: \textbf{storta alla caviglia}, distorsione del
ginocchio.

\textbf{Segni e sintomi comuni a contusioni e distorsioni:}

\begin{itemize}
      \tightlist
      \item
            Dolore vivo e costante
      \item
            Gonfiore immediato
      \item
            Ecchimosi o ematoma
\end{itemize}

\textbf{Trattamento:} metodo \textbf{RICE} (riposo, ghiaccio, bendaggio
compressivo, arto elevato).

\begin{center}\rule{0.5\linewidth}{0.5pt}\end{center}

\subsection{4. Lussazione}\label{lussazione}

La \textbf{lussazione} è la perdita permanente del contatto tra i capi
articolari: l\textquotesingle articolazione si "smonta". Il segno
caratteristico è la \textbf{deformazione anatomica} visibile e
importante.

Esempi: lussazione di spalla (la più frequente), del gomito, del
ginocchio.

\subsubsection{4.1 Possibili
      Complicazioni}\label{possibili-complicazioni}

La lussazione può comprimere strutture adiacenti:

\begin{itemize}
      \tightlist
      \item
            Arteria
      \item
            Vena
      \item
            Nervo
\end{itemize}

creando danni a valle della zona interessata.

\subsubsection{4.2 Cosa Fare e Cosa NON
      Fare}\label{cosa-fare-e-cosa-non-fare}

\textbf{NON rimettere a posto l\textquotesingle articolazione}: questa
manovra spetta esclusivamente all\textquotesingle ortopedico in ambiente
protetto. Rimettere l\textquotesingle articolazione in modo errato può
causare danni gravi e comporta \textbf{responsabilità civile} per il
soccorritore.

Nota: alcuni soggetti con lassità articolare possono rimettersi da soli
la spalla dopo episodi ripetuti. In quel caso non ci si assume alcuna
responsabilità. Se invece lo fa il soccorritore e genera un danno, ne
risponde civilmente.

\textbf{Cosa fare:}

\begin{itemize}
      \tightlist
      \item
            Lasciare l\textquotesingle articolazione nella posizione in cui si
            trova
      \item
            Bloccarla per limitare i movimenti
      \item
            Applicare il ghiaccio
      \item
            Controllare la \textbf{circolazione a valle} (polso, colorito,
            temperatura della parte distale)
      \item
            Inviare il soggetto in ambiente ospedaliero
\end{itemize}

\begin{center}\rule{0.5\linewidth}{0.5pt}\end{center}

\subsection{5. Fratture}\label{fratture}

\subsubsection{5.1 Meccanismo}\label{meccanismo}

Le ossa hanno un certo grado di elasticità: fino a una soglia si
piegano, superata quella soglia si rompono. La resistenza ossea varia
con l\textquotesingle età:

\begin{itemize}
      \tightlist
      \item
            \textbf{Bambini}: ossa molto elastiche, raramente si fratturano
      \item
            \textbf{Anziani con osteoporosi}: rischio elevato anche per traumi
            lievi
\end{itemize}

\subsubsection{5.2 Classificazione}\label{classificazione}

\paragraph{Per allineamento dei
      segmenti}\label{per-allineamento-dei-segmenti}

\begin{longtable}[]{@{}p{0.47\linewidth}p{0.47\linewidth}@{}}
      \toprule\noalign{}
      Tipo               & Caratteristiche                                     \\
      \midrule\noalign{}
      \endhead
      \bottomrule\noalign{}
      \endlastfoot
      \textbf{Composta}  & L\textquotesingle osso si rompe ma i segmenti
      restano allineati                                                        \\
      \textbf{Scomposta} & I segmenti si spostano; i monconi taglienti possono
      ledere arterie, vene o nervi                                             \\
\end{longtable}

\paragraph{Per comunicazione con
      l\textquotesingle esterno}\label{per-comunicazione-con-lesterno}

\begin{longtable}[]{@{}p{0.30\linewidth}p{0.30\linewidth}p{0.30\linewidth}@{}}
      \toprule\noalign{}
      Tipo                      & Caratteristiche                            & Rischio principale \\
      \midrule\noalign{}
      \endhead
      \bottomrule\noalign{}
      \endlastfoot
      \textbf{Chiusa}           & L\textquotesingle osso non comunica con
      l\textquotesingle esterno & Emorragia interna                                               \\
      \textbf{Aperta/esposta}   & Un frammento osseo fuoriesce oltre la cute &
      \textbf{Osteomielite} (infezione ossea)                                                     \\
\end{longtable}

\paragraph{Tipi morfologici
      principali}\label{tipi-morfologici-principali}

\begin{itemize}
      \tightlist
      \item
            \textbf{Obliqua}: linea di frattura diagonale, trattamento
            conservativo (gesso)
      \item
            \textbf{Comminuta}: più frammenti, spesso richiede placca chirurgica
      \item
            \textbf{Spiroide}: frattura a spirale, spesso richiede chiodo
            endomidollare
      \item
            \textbf{A legno verde}: tipica dei bambini, frattura parziale (come il
            legno verde che non si spezza del tutto)
      \item
            \textbf{Segmentare e trasversa}: di interesse prevalentemente
            ospedaliero
\end{itemize}

\subsubsection{5.3 Diagnosi sul Campo}\label{diagnosi-sul-campo}

Non è possibile identificare il tipo di frattura senza una radiografia.
Il soccorritore deve \textbf{sospettarla sempre} in caso di
deformazione, dolore intenso e gonfiore, e trattarla come se fosse
scomposta.

\subsubsection{5.4 Trattamento}\label{trattamento}

\textbf{Regola fondamentale:} non muovere il soggetto. Aspettare 10
minuti che arrivi l\textquotesingle ambulanza è mille volte meglio che
spostare i monconi e recidere un\textquotesingle arteria, portando il
soggetto allo shock ipovolemico.

\textbf{Cosa fare:}

\begin{enumerate}
      \tightlist
      \item
            Chiamare il 112
      \item
            \textbf{Non muovere il soggetto} se non strettamente necessario
      \item
            Immobilizzare l\textquotesingle arto nella posizione in cui si trova
            (le \emph{steccobende} --- rigide, a depressione o gonfiabili ---
            verranno usate dagli operatori del 112)
      \item
            In ambiente acquatico tagliare delicatamente i vestiti sopra e sotto
            la zona di frattura (in piscina di solito non c\textquotesingle è
            molto da tagliare)
      \item
            Tamponare eventuali emorragie associate
      \item
            \textbf{Non rimuovere mai oggetti conficcati}: fanno da tappo e la
            loro rimozione può causare emorragie massive
      \item
            Non tentare di riallineare i monconi e non forzare la zona interessata
\end{enumerate}

\textbf{Per le fratture esposte:}

\begin{itemize}
      \tightlist
      \item
            Applicare un telino sterile
      \item
            Mantenere umida la zona con soluzione fisiologica
      \item
            Effettuare disinfezione
      \item
            Inviare urgentemente in pronto soccorso
\end{itemize}

\begin{center}\rule{0.5\linewidth}{0.5pt}\end{center}

\subsection{6. Trauma Cranico}\label{trauma-cranico}

\subsubsection{6.1 Fisiopatologia}\label{fisiopatologia}

Il cranio è una scatola ossea rigida e non deformabile. Al suo interno
il cervello è una struttura morbida. Se si verifica
un\textquotesingle emorragia intracranica, il sangue --- non essendo
comprimibile --- \textbf{comprime il cervello}, che è la struttura più
delicata.

\subsubsection{6.2 Segni e Sintomi}\label{segni-e-sintomi}

\begin{itemize}
      \tightlist
      \item
            Disorientamento e sonnolenza
      \item
            Perdita di sangue da naso, bocca o orecchie
      \item
            Paralisi di un lato del corpo
      \item
            Vertigini, vomito, cefalea intensa
\end{itemize}

\subsubsection{6.3 Cosa Fare}\label{cosa-fare}

\begin{itemize}
      \tightlist
      \item
            \textbf{Monitorare} continuamente la situazione
      \item
            Se il sangue fuoriesce da naso, bocca o orecchie: si tratta di
            \textbf{emorragia interna esteriorizzata}, quindi \textbf{non va
                  tamponata} (va fatto uscire per ridurre la pressione intracranica)
      \item
            Distinguere tra emorragia del cuoio capelluto (molto vascolarizzato,
            sanguina molto ma non è pericolosa quanto sembra) e emorragia
            intracranica (invisibile, ma potenzialmente letale)
      \item
            Chiamare il 112 e garantire la sorveglianza continua
\end{itemize}

\begin{center}\rule{0.5\linewidth}{0.5pt}\end{center}

\subsection{7. BLS - Basic Life Support}\label{bls---basic-life-support}

\subsubsection{7.1 Definizione e Principi}\label{definizione-e-principi}

Il \textbf{BLS} (\emph{Basic Life Support}) è la sequenza di azioni
finalizzata a garantire un flusso di sangue ossigenato al cervello
quando cuore e polmoni non funzionano autonomamente. Si parla di
\textbf{BLSD} quando si aggiunge la lettera D, che indica
l\textquotesingle uso del defibrillatore.

Queste manovre si imparano in corsi pratico-teorici della durata di 4-6
ore. \textbf{Non vanno mai provate su persone reali} al di fuori
dell\textquotesingle emergenza: si lavora esclusivamente su manichini.

\subsubsection{7.2 Diagnosi di Arresto
      Cardiocircolatorio}\label{diagnosi-di-arresto-cardiocircolatorio}

Per diagnosticare l\textquotesingle arresto cardiocircolatorio occorre
verificare due condizioni:

\begin{enumerate}
      \tightlist
      \item
            \textbf{Assenza di coscienza}: il soggetto non risponde quando viene
            chiamato o stimolato
      \item
            \textbf{Assenza di respiro}: nessun atto respiratorio spontaneo
            visibile
\end{enumerate}

\textbf{Nota sulla ricerca del polso carotideo:} questo controllo è
stato \textbf{abolito nel 2005} dai protocolli internazionali (ILCOR).
Motivo: i non addestrati tendono a percepire il proprio polso nelle
dita, generando false informazioni. La sola assenza di respiro è
sufficiente per iniziare le manovre.

Considerare che se un soggetto non respira, il cuore andrà in arresto
entro \textbf{massimo 30 secondi}.

\subsubsection{7.3 Cosa NON Fare in Caso di
      Arresto}\label{cosa-non-fare-in-caso-di-arresto}

\begin{itemize}
      \tightlist
      \item
            \textbf{Non cercare il polso radiale}: inaffidabile in caso di shock
            (pressione troppo bassa)
      \item
            \textbf{Non misurare la pressione arteriosa}: un sistema
            cardiocircolatorio fermo ha pressione zero; il tempo perso a
            montare/smontare lo sfigmomanometro (1-1,5 minuti) causa danni
            neurologici irreversibili per anossia cerebrale
      \item
            \textbf{Non appoggiare l\textquotesingle orecchio al torace}: tecnica
            abbandonata da oltre trent\textquotesingle anni
      \item
            \textbf{Non iniettare farmaci}: la somministrazione di qualsiasi
            farmaco, anche il più comune, è di esclusiva competenza medica
\end{itemize}

\subsubsection{7.4 Quando Iniziare le
      Manovre}\label{quando-iniziare-le-manovre}

Le manovre di BLS vanno \textbf{sempre iniziate}, indipendentemente da:

\begin{itemize}
      \tightlist
      \item
            Età apparente del soggetto (anche se anziano)
      \item
            Aspetto cadaverico (tutti i soggetti in arresto cardiocircolatorio
            sembrano morti)
      \item
            Ipotermia (temperatura corporea fredda non è una controindicazione)
      \item
            Differenza nella dilatazione delle pupille
\end{itemize}

\subsubsection{7.5 Quando Sospendere le
      Manovre}\label{quando-sospendere-le-manovre}

\begin{itemize}
      \tightlist
      \item
            Arrivo di \textbf{altri soccorritori} (operatori del 112 o colleghi)
            che prendono il cambio
      \item
            \textbf{Ripresa spontanea} dell\textquotesingle attività cardiaca e
            respiratoria del soggetto
      \item
            \textbf{Esaurimento fisico} del soccorritore (se si è in più persone,
            darsi il cambio)
\end{itemize}

\subsubsection{7.6 La Catena della
      Sopravvivenza}\label{la-catena-della-sopravvivenza}

La catena della sopravvivenza è composta da più anelli collegati: ogni
anello è indispensabile e nessuno può compensare la mancanza di un
altro.

\begin{longtable}[]{@{}p{0.30\linewidth}p{0.30\linewidth}p{0.30\linewidth}@{}}
      \toprule\noalign{}
      Anello & Azione                                            & Chi                      \\
      \midrule\noalign{}
      \endhead
      \bottomrule\noalign{}
      \endlastfoot
      1      & \textbf{Prevenzione} dell\textquotesingle evento  & Istruttore               \\
      2      & \textbf{Accesso precoce} al sistema 112           & Istruttore               \\
      3      & \textbf{RCP precoce}: compressioni + ventilazioni & Istruttore               \\
      4      & \textbf{Defibrillazione precoce} (DAE)            & Istruttore (con
      abilitazione)                                                                         \\
      5      & \textbf{Soccorso avanzato} con farmaci            & Solo personale sanitario \\
\end{longtable}

\begin{center}\rule{0.5\linewidth}{0.5pt}\end{center}

\subsection{8. Il Defibrillatore (DAE)}\label{il-defibrillatore-dae}

Il \textbf{DAE} (\emph{Defibrillatore Automatico Esterno}) è
obbligatorio per legge in tutti gli impianti sportivi.

\subsubsection{8.1 Come Funziona}\label{come-funziona}

La macchina è dotata di un algoritmo interno che analizza il ritmo
cardiaco e decide autonomamente se erogare la scarica. Il
\textbf{margine di errore è quasi zero}.

\textbf{Responsabilità del soccorritore:} garantire la
\textbf{sicurezza} durante l\textquotesingle analisi e la scarica. Non
decidere se defibrillare o meno (lo fa la macchina).

\subsubsection{8.2 Procedura di Sicurezza}\label{procedura-di-sicurezza}

Durante la fase di analisi: \textbf{nessuno deve toccare il soggetto},
nemmeno sfiorarlo, per evitare interferenze che falserebbero
l\textquotesingle analisi.

Formula da applicare: \textbf{"via io, via tu, via tutti"} --- controllo
visivo che nessuno sia a contatto con il soggetto prima di confermare la
scarica.

\begin{center}\rule{0.5\linewidth}{0.5pt}\end{center}

\subsection{9. Ostruzione delle Vie Aeree e Manovra di
      Heimlich}\label{ostruzione-delle-vie-aeree-e-manovra-di-heimlich}

Un corpo estraneo nelle vie aeree blocca il passaggio
dell\textquotesingle aria, portando prima ad \textbf{arresto
      respiratorio} e poi ad \textbf{arresto cardiaco}.

\textbf{Segni di ostruzione:}

\begin{itemize}
      \tightlist
      \item
            Difficoltà a respirare
      \item
            Impossibilità a parlare o tossire
      \item
            Il soggetto si porta le mani alla gola
\end{itemize}

\subsubsection{9.1 Ostruzione Parziale}\label{ostruzione-parziale}

Il soggetto ancora respira o tossisce (un po\textquotesingle{}
d\textquotesingle aria passa ancora).

\textbf{Cosa fare:}

\begin{itemize}
      \tightlist
      \item
            Chiamare il 112
      \item
            \textbf{Incoraggiare a tossire}
      \item
            \textbf{Non eseguire manovre}: intervenire potrebbe trasformare
            l\textquotesingle ostruzione parziale in totale
\end{itemize}

\subsubsection{9.2 Ostruzione Totale --- Soggetto
      Cosciente}\label{ostruzione-totale-soggetto-cosciente}

Il soggetto non respira e non tossisce: l\textquotesingle aria non
passa.

\textbf{Sequenza da eseguire:}

\begin{enumerate}
      \tightlist
      \item
            Chiamare il 112
      \item
            Posizionarsi al fianco del soggetto con una mano davanti
            all\textquotesingle addome come fermo
      \item
            \textbf{5 pacche sulla schiena} (vigorose, in sequenza rapida) per
            favorire l\textquotesingle espulsione del corpo estraneo
\end{enumerate}

Se le pacche non bastano e il soggetto è in piedi o seduto:

\begin{enumerate}
      \setcounter{enumi}{3}
      \tightlist
      \item
            Posizionarsi dietro le spalle del soggetto
      \item
            Abbracciarlo da dietro
      \item
            Porre le mani a pugno \textbf{sotto l\textquotesingle arcata costale}
            (zona epigastrica)
      \item
            Dare colpi \textbf{dal basso verso l\textquotesingle alto} --- questa
            è la \textbf{manovra di Heimlich}: comprime il diaframma verso
            l\textquotesingle alto, "strizzando" i polmoni e creando un getto di
            aria che dovrebbe espellere il corpo estraneo
\end{enumerate}

\subsubsection{9.3 Ostruzione --- Soggetto Non Cosciente o Trovato a
      Terra}\label{ostruzione-soggetto-non-cosciente-o-trovato-a-terra}

Se il soggetto ha perso conoscenza o viene trovato già a terra, non si
eseguono le manovre di Heimlich.

\textbf{Cosa fare:}

\begin{itemize}
      \tightlist
      \item
            Avviare le manovre di \textbf{BLS} (compressioni toraciche)
      \item
            Dopo ogni serie di 30 compressioni, \textbf{verificare visivamente} se
            il corpo estraneo è diventato visibile in bocca
      \item
            Se visibile, rimuoverlo con cautela indossando i guanti, usando un
            movimento a uncino
      \item
            Se disponibile un presidio di protezione (pocket mask, pallone Ambu
            con filtro), procedere anche con le insufflazioni
\end{itemize}

\begin{center}\rule{0.5\linewidth}{0.5pt}\end{center}

\subsection{Conclusione: Punti Chiave per
      l\textquotesingle Esame}\label{conclusione-punti-chiave-per-lesame-4}

\begin{enumerate}
      \item
            \textbf{Lesioni muscolari:} il tessuto muscolare si rigenera
            completamente solo per le lesioni minori; le gravi esitano in
            cicatrice fibrosa. I fattori di rischio principali: riscaldamento
            inadeguato, disidratazione, stanchezza, squilibri idrosalini.
      \item
            \textbf{RICE:} nelle prime 48 ore --- \emph{Rest, Ice, Compression,
                  Elevation}. Prima si applica il ghiaccio, meglio è. La massoterapia
            non va eseguita prima di una settimana.
      \item
            \textbf{Ghiaccio:} applicazione di 20 minuti ogni 2-4 ore per 3-4
            giorni. Lo spray va usato a 10-15 cm; fermarsi
            all\textquotesingle alone grigiastro. La bustina chimica è monousa e i
            reagenti sono tossici se a contatto con la pelle.
      \item
            \textbf{Lussazione:} non rimettere mai l\textquotesingle articolazione
            a posto. Competenza esclusiva dell\textquotesingle ortopedico.
            Conseguenze civili per chi lo fa in modo improprio.
      \item
            \textbf{Fratture:} immobilizzare sempre, non spostare mai il soggetto.
            Non rimuovere oggetti conficcati. Fratture esposte: telino sterile e
            soluzione fisiologica per mantenere umida la zona.
      \item
            \textbf{Trauma cranico:} il sangue da naso/bocca/orecchie non va
            tamponato (è emorragia esteriorizzata: deve uscire). Monitorare il
            soggetto e chiamare il 112.
      \item
            \textbf{BLS:} diagnosi = soggetto non cosciente + assenza di respiro.
            La ricerca del polso è abolita dal 2005. Non misurare mai la pressione
            (perde tempo prezioso). Le manovre vanno sempre iniziate
            indipendentemente da età, aspetto o temperatura corporea.
      \item
            \textbf{DAE:} la macchina decide se scaricare --- il soccorritore
            garantisce solo la sicurezza. "Via io, via tu, via tutti" prima
            dell\textquotesingle analisi. Obbligatorio per legge in tutti gli
            impianti sportivi.
      \item
            \textbf{Ostruzione vie aeree:} ostruzione parziale = incoraggiare a
            tossire, non intervenire. Ostruzione totale con soggetto cosciente = 5
            pacche sulla schiena + manovra di Heimlich (spinte dal basso verso
            l\textquotesingle alto sotto l\textquotesingle arcata costale).
            Soggetto non cosciente = BLS standard.
\end{enumerate}

\section{L\textquotesingle Uomo e
  l\textquotesingle Acqua}\label{luomo-e-lacqua}

\begin{center}\rule{0.5\linewidth}{0.5pt}\end{center}

\subsection{Introduzione}\label{introduzione-9}

Questa lezione affronta le fondamenta teoriche del lavoro
dell\textquotesingle istruttore di nuoto: perché
l\textquotesingle ambiente acquatico richiede un percorso educativo
specifico, quali benefici porta l\textquotesingle attività in piscina
nelle diverse aree di sviluppo del bambino, quali competenze deve
possedere l\textquotesingle istruttore e quali sono i concetti chiave
della teoria del movimento applicata all\textquotesingle acqua. La
lezione costituisce il quadro di riferimento pedagogico e teorico sul
quale si costruisce tutto il lavoro pratico.

\begin{center}\rule{0.5\linewidth}{0.5pt}\end{center}

\subsection{1. Perché l\textquotesingle Ambiente Acquatico è Diverso da
      Tutti gli
      Altri}\label{perchuxe9-lambiente-acquatico-uxe8-diverso-da-tutti-gli-altri}

\subsubsection{1.1 Il Problema degli Schemi Motori di
      Base}\label{il-problema-degli-schemi-motori-di-base}

Gli istruttori degli sport terrestri si trovano in una condizione
privilegiata: lavorano con bambini che già possiedono le competenze
motorie di base, i cosiddetti \textbf{schemi motori di base}, ovvero
camminare, correre, saltare, arrampicarsi, rotolare, afferrare,
lanciare. Queste abilità si sviluppano autonomamente durante la
crescita, senza che nessuno le insegni.

In acqua tutto questo non esiste. Un bambino --- ma anche un adulto
principiante --- non dispone di \textbf{nessuna competenza motoria
      pregressa} idonea a muoversi in acqua.

La ragione è evolutiva: ci siamo sviluppati per decine di migliaia di
anni come animali terrestri. La natura ci ha dotato di un
\textbf{software motorio} --- gli schemi motori di base ---
perfettamente adeguato alla sopravvivenza sulla terraferma, ma del tutto
inadatto all\textquotesingle ambiente acquatico. Questo è il motivo per
cui istintivamente teniamo lontani dall\textquotesingle acqua.

\subsubsection{1.2 Conseguenze per il Percorso
      Educativo}\label{conseguenze-per-il-percorso-educativo}

Prima di poter acquisire le abilità specifiche del nuoto, della
pallanuoto, del nuoto artistico, del nuoto per salvamento o dei tuffi, è
indispensabile un \textbf{percorso educativo} che parte
dall\textquotesingle ambientamento e procede con
l\textquotesingle acquisizione delle basi della locomozione acquatica.

\begin{quote}
      La dicitura \textbf{"scuola nuoto"} è in realtà fuorviante. I quattro
      stili classici (farfalla, dorso, rana, stile libero) sono solo una
      piccola parte --- e probabilmente nemmeno la più importante --- delle
      abilità che un soggetto deve acquisire in acqua. La denominazione più
      corretta è \textbf{"scuola di abilità acquatiche"}.
\end{quote}

\subsubsection{1.3 Un Criterio di Divisione
      Diverso}\label{un-criterio-di-divisione-diverso}

In piscina i bambini non vengono divisi per \textbf{fascia
      d\textquotesingle età}, ma in base al \textbf{grado di confidenza} che
hanno o non hanno con l\textquotesingle acqua. Prima di imparare a
nuotare bisogna imparare a stare in acqua, acquisendo la consapevolezza
di come l\textquotesingle acqua agisce sul corpo e ne modifica la
percezione.

\begin{center}\rule{0.5\linewidth}{0.5pt}\end{center}

\subsection{2. I Benefici dell\textquotesingle Attività Motoria in
      Acqua}\label{i-benefici-dellattivituxe0-motoria-in-acqua}

L\textquotesingle attività in piscina influisce su tutte le principali
aree di sviluppo del bambino.

\subsubsection{2.1 Area
      Morfologico-Funzionale}\label{area-morfologico-funzionale}

\begin{itemize}
      \tightlist
      \item
            Sviluppa e potenzia le strutture anatomiche, fisiologiche e
            biochimiche che sostengono le funzioni vitali
      \item
            Previene difetti morfologici (posture scorrette) e difetti
            senso-motori
      \item
            La grande varietà di stimoli dell\textquotesingle ambiente acquatico
            favorisce il miglioramento della percezione e della consapevolezza
            dell\textquotesingle ambiente
\end{itemize}

\subsubsection{2.2 Area
      Intellettivo-Cognitiva}\label{area-intellettivo-cognitiva}

\begin{itemize}
      \tightlist
      \item
            Stimola il \textbf{sistema nervoso centrale} (cervello, cervelletto,
            tronco spinale), struttura unitaria: migliorare la funzionalità in
            ambito motorio genera benefici in tutti gli ambiti
      \item
            Migliora la senso-percezione e la capacità di immaginare il proprio
            corpo in condizioni sempre differenti
      \item
            È ampiamente dimostrato che bambini e ragazzi che praticano attività
            sportiva in modo continuativo migliorano le \textbf{prestazioni
                  scolastiche} e il funzionamento in tutti gli aspetti della vita
            quotidiana
\end{itemize}

\subsubsection{2.3 Area Affettivo-Emotiva}\label{area-affettivo-emotiva}

\begin{itemize}
      \tightlist
      \item
            Le ore in piscina (2-3 a settimana nel migliore dei casi)
            rappresentano una grande occasione di \textbf{socializzazione} in
            un\textquotesingle epoca in cui i bambini hanno vita sempre più
            strutturata, passano più tempo con gli adulti che con i coetanei
      \item
            La socializzazione previene squilibri affettivi, frustrazioni e
            complessi
      \item
            Apprendere nuove competenze in un ambiente sconosciuto favorisce
            l\textquotesingle{}\textbf{autostima}, una buona percezione di sé e il
            desiderio di autonomia
      \item
            I bambini sportivi mostrano un funzionamento sociale ed emotivo
            complessivamente migliore
\end{itemize}

\subsubsection{2.4 Area Sociale}\label{area-sociale}

\begin{itemize}
      \tightlist
      \item
            Frequentare la piscina sviluppa il \textbf{senso di appartenenza}
      \item
            Insegna il rispetto delle regole e la fiducia nella loro applicazione
      \item
            Sviluppa lo spirito di collaborazione e di emulazione
      \item
            Favorisce l\textquotesingle{}\textbf{apprendimento cooperativo} (far
            lavorare i bambini insieme è il modo migliore per farli imparare)
      \item
            I bambini sportivi sono più propensi a rispettare se stessi, gli altri
            e le regole della comunità
\end{itemize}

\begin{center}\rule{0.5\linewidth}{0.5pt}\end{center}

\subsection{3. Il Ruolo e le Competenze
      dell\textquotesingle Istruttore}\label{il-ruolo-e-le-competenze-dellistruttore}

\subsubsection{3.1 L\textquotesingle Istruttore come Educatore
      Motorio}\label{listruttore-come-educatore-motorio}

L\textquotesingle istruttore di nuoto non è un semplice tecnico
sportivo: è un \textbf{educatore motorio in ambiente acquatico}. Dalla
qualità del suo lavoro dipendono le possibilità del soggetto di
sviluppare in maniera piena il proprio potenziale motorio acquatico.

Due livelli di educazione:

\begin{itemize}
      \tightlist
      \item
            \textbf{Educazione al nuoto}: insegnare le tecniche di locomozione in
            acqua
      \item
            \textbf{Educazione attraverso il nuoto}: contribuire a uno sviluppo
            equilibrato e armonico della personalità del bambino in tutte le sue
            aree
\end{itemize}

La sicurezza è sempre la priorità assoluta. \textbf{Regola
      fondamentale:} dal momento in cui gli allievi vengono affidati
all\textquotesingle istruttore fino al momento del saluto, tutti gli
allievi devono essere contemporaneamente all\textquotesingle interno del
\textbf{campo visivo} dell\textquotesingle istruttore.

\subsubsection{3.2 I Quattro Saperi
      dell\textquotesingle Istruttore}\label{i-quattro-saperi-dellistruttore}

\begin{longtable}[]{@{}p{0.47\linewidth}p{0.47\linewidth}@{}}
      \toprule\noalign{}
      Sapere                  & Descrizione                                           \\
      \midrule\noalign{}
      \endhead
      \bottomrule\noalign{}
      \endlastfoot
      \textbf{Sapere}         & Conoscenze teoriche legate al proprio ruolo: poche ma
      da avere ben chiare                                                             \\
      \textbf{Saper fare}     & Capacità pratica di riprodurre le tecniche
      richieste agli allievi e di intervenire in caso di necessità                    \\
      \textbf{Saper far fare} & Comunicazione adeguata all\textquotesingle età
      e al grado di comprensione degli allievi per metterli in condizione di
      ottenere i risultati prefissati                                                 \\
      \textbf{Saper essere}   & Consapevolezza del proprio ruolo; essere ottimi
      osservatori e ascoltatori degli allievi per comprenderne le esigenze e
      trovare le soluzioni più adeguate per ciascuno                                  \\
\end{longtable}

\subsubsection{3.3 Le Due Regole Fondamentali della
      Comunicazione}\label{le-due-regole-fondamentali-della-comunicazione}

\begin{enumerate}
      \tightlist
      \item
            \textbf{Non si può non comunicare}: tutto ciò che si fa, si dice, gli
            atteggiamenti che si assumono lanciano messaggi chiari agli allievi
      \item
            \textbf{Il messaggio è sempre quello che arriva, non quello che
                  parte}: se gli allievi non eseguono correttamente un esercizio, il
            problema non è che hanno capito male, ma che il linguaggio, il tono o
            la postura dell\textquotesingle istruttore non erano adeguati alle
            loro capacità di comprensione in quel momento
\end{enumerate}

Nota contestuale: la piscina è un ambiente rumoroso, si sente poco e si
capisce ancora meno. La scelta del registro comunicativo è quindi di
importanza critica.

\begin{center}\rule{0.5\linewidth}{0.5pt}\end{center}

\subsection{4. I Due Concetti Chiave}\label{i-due-concetti-chiave}

\subsubsection{4.1 Multilateralità}\label{multilateralituxe0}

La \textbf{multilateralità} è la necessità di diversificare al massimo
le esperienze motorie. Si sviluppa in due direzioni:

\begin{itemize}
      \tightlist
      \item
            \textbf{Orizzontale}: fornire agli allievi il maggior numero possibile
            di esperienze motorie per stimolare il sistema nervoso centrale a
            trovare sempre strade e soluzioni differenti a qualsiasi tipo di
            problema motorio. Aumentare le abilità motorie acquatiche significa
            aumentare in modo esponenziale le tecniche che potranno essere apprese
            in futuro
      \item
            \textbf{Verticale}: fare in modo che ciascuna tecnica appresa venga
            eseguita in varie condizioni di \textbf{velocità, ritmo e ampiezza} di
            movimento
\end{itemize}

\subsubsection{4.2 Senso Percezione}\label{senso-percezione}

La \textbf{senso percezione} è il processo attraverso il quale il corpo
percepisce e interpreta le informazioni che l\textquotesingle ambiente
fornisce attraverso i cinque sensi, le integra e le traduce in
movimento. Equivale alla consapevolezza della posizione del proprio
corpo nell\textquotesingle acqua e del modo in cui
l\textquotesingle acqua reagisce al movimento.

\begin{center}\rule{0.5\linewidth}{0.5pt}\end{center}

\subsection{5. Le Proprietà Fisiche
      dell\textquotesingle Acqua}\label{le-proprietuxe0-fisiche-dellacqua}

L\textquotesingle acqua induce l\textquotesingle istruttore a operare in
maniera \textbf{controintuitiva} rispetto a ciò che avviene sulla
terraferma.

\subsubsection{5.1 Resistenza
      dell\textquotesingle Acqua}\label{resistenza-dellacqua}

\begin{itemize}
      \tightlist
      \item
            L\textquotesingle acqua ha una resistenza circa \textbf{800 volte
                  superiore} a quella dell\textquotesingle aria
      \item
            Per muoversi efficacemente in acqua bisogna prima \textbf{ridurre le
                  resistenze} e poi aumentare le propulsioni: l\textquotesingle opposto
            rispetto alla terraferma
      \item
            La resistenza è di tipo \textbf{progressivo}: aumenta con il quadrato
            della velocità, quindi il posizionamento corretto del corpo è
            fondamentale
\end{itemize}

\subsubsection{5.2 Spinta Idrostatica}\label{spinta-idrostatica}

La \textbf{spinta idrostatica} è l\textquotesingle unico elemento utile
al movimento in acqua. In base al principio di Archimede, un corpo
immerso in un liquido riceve una spinta verso l\textquotesingle alto
proporzionale alla quantità di liquido spostato.

Conseguenze pratiche:

\begin{itemize}
      \tightlist
      \item
            Per massimizzare la spinta idrostatica occorre mantenere una
            \textbf{posizione orizzontale} del corpo
      \item
            Bisogna imparare a \textbf{rilassare la muscolatura} per distendere il
            peso del corpo sulla porzione di acqua più ampia possibile
\end{itemize}

\begin{quote}
      Questo è particolarmente controintuitivo con i bambini: i bambini sono
      perennemente in movimento rapido ed esplosivo, mentre in acqua devono
      imparare a rilassarsi, distendere il corpo e rimanere in posizione
      orizzontale.
\end{quote}

\begin{center}\rule{0.5\linewidth}{0.5pt}\end{center}

\subsection{6. La Teoria del Movimento}\label{la-teoria-del-movimento}

\subsubsection{6.1 Prestazione Sportiva}\label{prestazione-sportiva}

La \textbf{prestazione} è l\textquotesingle esito misurabile di
un\textquotesingle azione sportiva completa. Non è un concetto riservato
all\textquotesingle agonismo: è una prestazione sportiva sia nuotare i
100 metri in stile libero sotto i 47 secondi, sia spostarsi dal bordo
vasca alla corsia più vicina.

I quattro fattori della prestazione:

\begin{longtable}[]{@{}p{0.30\linewidth}p{0.30\linewidth}p{0.30\linewidth}@{}}
      \toprule\noalign{}
      Fattore                                                              & Descrizione                                      & Metafora \\
      \midrule\noalign{}
      \endhead
      \bottomrule\noalign{}
      \endlastfoot
      \textbf{Costituzione}                                                & Misura e proporzione dei segmenti corporei,
      massa, composizione corporea, mobilità articolare                    & La carrozzeria                                              \\
      \textbf{Condizione}                                                  & Grado di sviluppo delle capacità condizionali
      (legato ai grandi apparati: cardiocircolatorio e cardiorespiratorio) &
      La cilindrata del motore                                                                                                           \\
      \textbf{Coordinazione}                                               & Capacità di controllare e regolare il movimento
      in modo funzionale                                                   & La guida                                                    \\
      \textbf{Controllo}                                                   & Processi emotivi, cognitivi e motivazionali (es.
      difficoltà di replicare in gara le prestazioni di allenamento)       & La
      testa del pilota                                                                                                                   \\
\end{longtable}

L\textquotesingle obiettivo del lavoro dell\textquotesingle istruttore è
trasformare le \textbf{capacità} in \textbf{abilità motorie},
presupposto della prestazione sportiva.

\subsubsection{6.2 Capacità Motorie:
      Panoramica}\label{capacituxe0-motorie-panoramica}

Le capacità motorie si dividono in due grandi categorie più una
intermedia:

\begin{longtable}[]{@{}p{0.30\linewidth}p{0.30\linewidth}p{0.30\linewidth}@{}}
      \toprule\noalign{}
      Categoria                    & Allenabilità                                    & Periodo ottimale \\
      \midrule\noalign{}
      \endhead
      \bottomrule\noalign{}
      \endlastfoot
      \textbf{Condizionali}        & Sempre, lungo tutta la vita                     & Nessuna finestra
      specifica                                                                                         \\
      \textbf{Coordinative}        & Molto limitata fuori dalla finestra di sviluppo
                                   & 6-14 anni (età scolastica)                                         \\
      \textbf{Mobilità articolare} & Intermedia (componente ereditaria forte)
                                   & Tendenzialmente in età evolutiva                                   \\
\end{longtable}

\begin{center}\rule{0.5\linewidth}{0.5pt}\end{center}

\subsection{7. Capacità Coordinative}\label{capacituxe0-coordinative-1}

Le capacità coordinative sono le capacità di organizzare e regolare il
movimento. Si dividono in tre categorie:

\begin{itemize}
      \tightlist
      \item
            \textbf{Capacità di apprendimento}: imparare un gesto tecnico
      \item
            \textbf{Capacità di controllo}: regolare il gesto in modo raffinato
      \item
            \textbf{Capacità di adattamento e trasformazione}: modificare il gesto
            in base alle necessità del momento
\end{itemize}

\subsubsection{7.1 Stadi dell\textquotesingle Apprendimento
      Motorio}\label{stadi-dellapprendimento-motorio}

\begin{longtable}[]{@{}p{0.47\linewidth}p{0.47\linewidth}@{}}
      \toprule\noalign{}
      Stadio                           & Descrizione                           \\
      \midrule\noalign{}
      \endhead
      \bottomrule\noalign{}
      \endlastfoot
      \textbf{Coordinazione grezza}    & Il movimento viene appreso per
      imitazione nella sua struttura fondamentale                              \\
      \textbf{Coordinazione fine}      & Il movimento viene raffinato          \\
      \textbf{Disponibilità variabile} & Il movimento è padroneggiato con tale
      sicurezza da poter essere modificato in base alle necessità              \\
\end{longtable}

\subsubsection{7.2 La Fase Sensibile}\label{la-fase-sensibile}

Il periodo compreso tra i \textbf{6 e i 14 anni} (età della scuola
dell\textquotesingle obbligo) è definito \textbf{fase sensibile dello
      sviluppo coordinativo}: il sistema nervoso e l\textquotesingle apparato
muscolo-scheletrico sono in una fase di sviluppo che permette
l\textquotesingle acquisizione ottimale di nuove competenze motorie.

Lo \textbf{schema di Martin} (rappresentato in diapositiva) sintetizza
per fascia d\textquotesingle età le capacità sulle quali è più opportuno
lavorare:

\begin{itemize}
      \tightlist
      \item
            \textbf{6-12 anni}: priorità alle capacità coordinative
      \item
            \textbf{Pre-adolescenza e pubertà in poi}: inizio del lavoro
            significativo sulle capacità condizionali
\end{itemize}

\begin{center}\rule{0.5\linewidth}{0.5pt}\end{center}

\subsection{8. Capacità Condizionali}\label{capacituxe0-condizionali-1}

\subsubsection{8.1 Forza}\label{forza-1}

La \textbf{forza} è la capacità di opporsi a una resistenza attraverso
la tensione muscolare. Si articola in tre forme:

\begin{longtable}[]{@{}p{0.30\linewidth}p{0.30\linewidth}p{0.30\linewidth}@{}}
      \toprule\noalign{}
      Tipo                            & Descrizione                                    & Adatta ai bambini? \\
      \midrule\noalign{}
      \endhead
      \bottomrule\noalign{}
      \endlastfoot
      \textbf{Forza massima}          & Massima potenza espressa in breve tempo        & No:
      sovraccarico su struttura muscoloscheletrica immatura; i bambini non
      possono incrementare la massa muscolare                                                               \\
      \textbf{Forza resistente}       & Sostentamento di uno sforzo di forza per un
      periodo protratto               & No per bambini e preadolescenti; da pubertà in
      poi                                                                                                   \\
      \textbf{Forza veloce/esplosiva} & Forza di elevata intensità nel più
      breve tempo possibile           & Sì: è la tipica espressione motoria infantile
      (es. tuffo dal blocco di partenza)                                                                    \\
\end{longtable}

\subsubsection{8.2 Resistenza}\label{resistenza-1}

La \textbf{resistenza} è la capacità di sopportare uno sforzo prolungato
contrastando l\textquotesingle insorgere della fatica. Ha una componente
fisiologica (gestione delle risorse energetiche) e una componente
psicologica (volontà di sopportare lo sforzo).

\begin{itemize}
      \tightlist
      \item
            \textbf{Resistenza generale}: capacità di sopportare uno sforzo
            prolungato indipendentemente dal tipo di attività
      \item
            \textbf{Resistenza specifica}: capacità legata a una determinata
            disciplina sportiva
\end{itemize}

Nota importante per il nuoto: molti nuotatori con grande resistenza in
acqua si stancano rapidamente sulla terraferma, perché le \textbf{catene
      cinetiche} attivate durante la nuotata sono completamente diverse da
quelle usate per muoversi sulla terraferma.

I lavori di resistenza sono indicati per bambini e ragazzi, ma tendono a
essere \textbf{monotoni}: è essenziale introdurre variazioni di
esercizio, ritmo e ampiezza del gesto per mantenere la motivazione del
gruppo.

\subsubsection{8.3 Velocità}\label{velocituxe0}

La velocità viene citata come terza capacità condizionale ma non
approfondita in questa lezione.

\subsubsection{8.4 Mobilità Articolare}\label{mobilituxe0-articolare}

Capacità intermedia con forte componente ereditaria e un certo grado di
allenabilità. Per le fasce di età della scuola nuoto non è necessario un
lavoro specifico: gli esercizi del nuoto richiedono già una grande
ampiezza di movimento e questo è sufficiente.

\begin{center}\rule{0.5\linewidth}{0.5pt}\end{center}

\subsection{9. Gli Analizzatori}\label{gli-analizzatori-1}

Gli \textbf{analizzatori} sono strutture anatomico-funzionali che
selezionano, codificano e trasmettono gli stimoli ambientali al sistema
nervoso centrale, il quale elabora una risposta che si traduce in
un\textquotesingle azione motoria.

\subsubsection{9.1 Analizzatori Esterocettivi (informazioni
      dall\textquotesingle esterno)}\label{analizzatori-esterocettivi-informazioni-dallesterno}

\begin{longtable}[]{@{}p{0.47\linewidth}p{0.47\linewidth}@{}}
      \toprule\noalign{}
      Analizzatore      & Funzione                                         \\
      \midrule\noalign{}
      \endhead
      \bottomrule\noalign{}
      \endlastfoot
      \textbf{Acustico} & Percepire suoni, rumori, ritmi                   \\
      \textbf{Ottico}   & Informazioni sul movimento proprio e altrui      \\
      \textbf{Tattile}  & Informazioni relative alle pressioni sulle varie
      parti del corpo                                                      \\
\end{longtable}

\subsubsection{9.2 Analizzatori Propriocettivi (informazioni
      dall\textquotesingle interno)}\label{analizzatori-propriocettivi-informazioni-dallinterno}

\begin{longtable}[]{@{}p{0.47\linewidth}p{0.47\linewidth}@{}}
      \toprule\noalign{}
      Analizzatore                            & Funzione                             \\
      \midrule\noalign{}
      \endhead
      \bottomrule\noalign{}
      \endlastfoot
      \textbf{Cinestesico} (motorio)          & Valutare posizione, velocità e altri
      parametri relativi alle parti del corpo e agli oggetti/individui in
      movimento                                                                      \\
      \textbf{Vestibolare} (statico-dinamico) & Segnali collegati al movimento
      nello spazio; rapporto tra posizione della testa e gli altri segmenti
      corporei. Sede: \textbf{orecchio interno} (motivo per cui chi soffre di
      infezione dell\textquotesingle orecchio tende a perdere
      l\textquotesingle equilibrio)                                                  \\
\end{longtable}

\begin{center}\rule{0.5\linewidth}{0.5pt}\end{center}

\subsection{10. Il Percorso di Educazione Motoria
      Acquatica}\label{il-percorso-di-educazione-motoria-acquatica}

Il percorso dell\textquotesingle allievo nella scuola di abilità
acquatiche segue una progressione precisa:

\begin{enumerate}
      \tightlist
      \item
            \textbf{Esplorazione del nuovo ambiente}: presa di confidenza con un
            contesto sconosciuto
      \item
            \textbf{Ambientamento}: stare in acqua volentieri, padroneggiare il
            proprio corpo, orientarsi
      \item
            \textbf{Galleggiamento statico}
      \item
            \textbf{Galleggiamento dinamico}
      \item
            \textbf{Prime forme propulsive}
      \item
            \textbf{Acquisizione progressiva di molteplici abilità motorie
                  acquatiche}
      \item
            \textbf{Costruzione delle tecniche specifiche} per ciascuno stile e
            disciplina
\end{enumerate}

\textbf{Obiettivi di fondo:}

\begin{itemize}
      \tightlist
      \item
            Massima \textbf{flessibilità e trasferibilità} delle abilità apprese:
            l\textquotesingle allievo deve adattarsi al cambiare delle condizioni
            esterne e alla propria evoluzione psicofisica (specialmente in
            adolescenza, quando cambia il galleggiamento e la capacità di
            scivolare in acqua)
      \item
            Strutturare una \textbf{mappa motoria} (la rappresentazione cerebrale
            del movimento) il più possibile ricca ed elastica, dalla quale
            costruire un bagaglio motorio ampio, dinamico, trasferibile e
            potenzialmente inesauribile
\end{itemize}

\begin{quote}
      Si acquisiscono tante abilità quante se ne apprendono: quanto più
      l\textquotesingle istruttore sarà proattivo, dinamico, creativo e
      attento, tanto più soddisfacente sarà il percorso di crescita dei propri
      allievi.
\end{quote}

La \textbf{competenza motoria} così sviluppata è la base sulla quale si
costruiscono le possibilità di esprimere il talento sportivo, agonistico
o meno.

\begin{center}\rule{0.5\linewidth}{0.5pt}\end{center}

\subsection{Conclusione: Punti Chiave per
      l\textquotesingle Esame}\label{conclusione-punti-chiave-per-lesame-5}

\begin{enumerate}
      \item
            \textbf{Unicità dell\textquotesingle ambiente acquatico}: non esistono
            schemi motori di base per l\textquotesingle acqua.
            L\textquotesingle uomo è un animale terrestre, l\textquotesingle acqua
            richiede un percorso educativo specifico che parte da zero.
      \item
            \textbf{Scuola di abilità acquatiche}: il termine "scuola nuoto" è
            impreciso; i quattro stili sono solo una parte
            dell\textquotesingle obiettivo. I bambini si dividono per grado di
            confidenza con l\textquotesingle acqua, non per età anagrafica.
      \item
            \textbf{Quattro aree di beneficio}: morfologico-funzionale,
            intellettivo-cognitiva, affettivo-emotiva, sociale.
            L\textquotesingle attività in piscina è un potente strumento di
            sviluppo globale della personalità.
      \item
            \textbf{Quattro saperi dell\textquotesingle istruttore}: sapere
            (teoria), saper fare (tecnica), saper far fare (comunicazione), saper
            essere (consapevolezza e osservazione). La regola fondamentale di
            sicurezza: tutti gli allievi sempre nel campo visivo.
      \item
            \textbf{Due regole della comunicazione}: non si può non comunicare; il
            messaggio è sempre quello che arriva (se non capiscono, il problema è
            della comunicazione, non degli allievi).
      \item
            \textbf{Multilateralità}: orizzontale (varietà di esperienze) e
            verticale (stessa tecnica in condizioni diverse di velocità, ritmo e
            ampiezza).
      \item
            \textbf{Proprietà dell\textquotesingle acqua}: resistenza 800 volte
            superiore all\textquotesingle aria, progressiva con il quadrato della
            velocità; spinta idrostatica = posizione orizzontale + muscolatura
            rilassata. Logica controintuitiva rispetto alla terraferma.
      \item
            \textbf{Capacità coordinative vs. condizionali}: le coordinative si
            allenano principalmente tra i 6 e i 14 anni (fase sensibile); le
            condizionali sono allenabili per tutta la vita. Con i bambini: forza
            esplosiva sì, forza massima no.
      \item
            \textbf{Analizzatori}: esterocettivi (acustico, ottico, tattile) e
            propriocettivi (cinestesico, vestibolare). Il vestibolare ha sede
            nell\textquotesingle orecchio interno.
      \item
            \textbf{Schema di Martin}: tra 6 e 12 anni priorità alle capacità
            coordinative; dalla pre-adolescenza si inizia con quelle condizionali.
            La mappa motoria ricca ed elastica è la base del talento sportivo.
\end{enumerate}

\section{Ambientamento e
  Acquaticità}\label{ambientamento-e-acquaticituxe0}

\begin{center}\rule{0.5\linewidth}{0.5pt}\end{center}

\subsection{Introduzione}\label{introduzione-10}

Questa lezione tratta il primo passo del percorso di educazione motoria
acquatica: l\textquotesingle{}\textbf{ambientamento}. I contenuti si
riferiscono alla scuola nuoto per bambini e ragazzi;
l\textquotesingle adulto principiante ha caratteristiche specifiche che
richiedono una gestione differente e che vengono affrontate in una
lezione dedicata. L\textquotesingle ambientamento non è una semplice
fase introduttiva: è il presupposto indispensabile di qualsiasi
apprendimento acquatico successivo.

\begin{center}\rule{0.5\linewidth}{0.5pt}\end{center}

\subsection{1. Definizione di
      Ambientamento}\label{definizione-di-ambientamento}

Per \textbf{ambientamento} si intende la fase preliminare
all\textquotesingle acquisizione di qualsiasi competenza motoria
acquatica. Può essere definito come il raggiungimento di uno
\textbf{stato di benessere psicofisico} conseguente
all\textquotesingle immersione e alla conoscenza del modo in cui il
proprio corpo reagisce all\textquotesingle immersione.

Raggiungere questa condizione è complesso. Come stabilito nelle lezioni
precedenti, l\textquotesingle essere umano non è un animale acquatico:
l\textquotesingle istinto lo porta a sviluppare paura e diffidenza nei
confronti dell\textquotesingle acqua.

Questa è la ragione per cui esistono corsi di \textbf{acquaticità} per
bambini da \textbf{0 a 3 anni} e da \textbf{3 a 6 anni}: non perché i
bambini siano maturi per sviluppare tecniche di locomozione acquatica,
ma perché prendano confidenza con l\textquotesingle ambiente e imparino
a considerare l\textquotesingle acqua come un elemento amico.

\begin{center}\rule{0.5\linewidth}{0.5pt}\end{center}

\subsection{2. Le Fasi dell\textquotesingle Ambientamento in Vasca
      Piccola}\label{le-fasi-dellambientamento-in-vasca-piccola}

La trattazione assume come contesto di riferimento la disponibilità di
una \textbf{vasca per l\textquotesingle ambientamento} (vasca piccola),
dove i bambini toccano il fondo. Le variazioni per il caso in cui si
disponga solo di acqua fonda vengono trattate nella sezione dedicata.

\subsubsection{2.1 Prima di Entrare in Acqua: Conoscere
      l\textquotesingle Ambiente}\label{prima-di-entrare-in-acqua-conoscere-lambiente}

Prima ancora di mettere piede in acqua, i bambini devono familiarizzare
con l\textquotesingle ambiente piscina nel suo complesso:

\begin{itemize}
      \tightlist
      \item
            L\textquotesingle istruttore si presenta e impara i \textbf{nomi di
                  tutti gli allievi}: non ricordare il nome dei propri bambini è una
            delle sconfitte più gravi per un istruttore
      \item
            I bambini esplorano gli spazi: dove cambiarsi, dove lasciare le
            ciabatte e la cappatoia
      \item
            Si stabiliscono le \textbf{regole di ingaggio}: ad esempio fare la
            doccia prima di entrare in acqua
\end{itemize}

Questa fase favorisce la familiarizzazione con
l\textquotesingle ambiente e lo sviluppo
dell\textquotesingle{}\textbf{autonomia}, che --- come ricordato nelle
lezioni precedenti --- è la base dell\textquotesingle autostima.

\subsubsection{2.2 Posizione dell\textquotesingle Istruttore e Gestione
      dello Spazio}\label{posizione-dellistruttore-e-gestione-dello-spazio}

In vasca piccola l\textquotesingle istruttore \textbf{sta sempre in
      acqua con i bambini}, per tre ragioni:

\begin{enumerate}
      \tightlist
      \item
            Rassicurare i bambini con la propria presenza fisica
      \item
            Poter intervenire immediatamente in situazioni di difficoltà
      \item
            Mantenere un contatto visivo adeguato (un bambino in acqua guarda
            l\textquotesingle adulto dal basso verso l\textquotesingle alto: se
            l\textquotesingle istruttore è sul bordo, la comunicazione è
            difficile)
\end{enumerate}

\textbf{Gestione dello spazio:} ridurre al massimo
l\textquotesingle area acqua attraverso le corsie. Un bambino
principiante non ha bisogno di 10-20 metri di spazio libero: un riquadro
di \textbf{4-5 metri di lato} è più che sufficiente per tutte le
esercitazioni propedeutiche.

\subsubsection{2.3 Il Principio Metodologico: Imparare
      Giocando}\label{il-principio-metodologico-imparare-giocando}

I bambini imparano giocando. Questo principio deve guidare assolutamente
tutto il lavoro di ambientamento:

\begin{itemize}
      \tightlist
      \item
            L\textquotesingle istruttore parla \textbf{poco} e fa \textbf{muovere
                  molto}
      \item
            Le proposte devono stimolare la \textbf{voglia di giocare}, che è il
            principale motore di apprendimento nei bambini
\end{itemize}

\subsubsection{2.4 Le Competenze da Acquisire in
      Sequenza}\label{le-competenze-da-acquisire-in-sequenza}

Il livello di ambientamento si misura osservando il \textbf{grado di
      tensione muscolare}: un bambino rigido è un bambino che non ha ancora
superato la paura. Man mano che l\textquotesingle ambientamento
progredisce, la muscolatura si distende.

\paragraph{1. Entrare in acqua}\label{entrare-in-acqua}

Per un bambino impaurito non è affatto un\textquotesingle azione banale.
Va proposta con gradualità attraverso il gioco.

\paragraph{2. Riprodurre gli schemi motori terrestri in
      acqua}\label{riprodurre-gli-schemi-motori-terrestri-in-acqua}

\begin{itemize}
      \tightlist
      \item
            Camminare avanti e indietro
      \item
            Saltare
      \item
            Muoversi con confidenza e in modo rilassato
\end{itemize}

\paragraph{3. Immersione del viso e apertura degli occhi
      sott\textquotesingle acqua}\label{immersione-del-viso-e-apertura-degli-occhi-sottacqua}

\begin{itemize}
      \tightlist
      \item
            Prima con gli occhialini, poi senza
      \item
            \textbf{Nota sugli occhialini:} non vanno demonizzati, ma non devono
            diventare uno strumento indispensabile. I bambini devono imparare a
            stare sott\textquotesingle acqua anche senza, perché se dovessero
            cadere accidentalmente in acqua fuori dalla piscina, devono saper
            agire in ogni condizione
\end{itemize}

\paragraph{4. Respirare in acqua}\label{respirare-in-acqua}

La respirazione in acqua è una competenza tutt\textquotesingle altro che
banale.

\begin{itemize}
      \tightlist
      \item
            In acqua è necessaria un\textquotesingle{}\textbf{espirazione di tipo
                  forzato}: la pressione dell\textquotesingle acqua che preme sul torace
            deve essere vinta attivamente, mentre sulla terraferma la cassa
            toracica si espande spontaneamente
      \item
            \textbf{Inspirazione}: sempre con la \textbf{bocca} (inspirare col
            naso espone al rischio di inalare acqua nelle vie respiratorie)
      \item
            \textbf{Espirazione}: con la bocca o con il naso
      \item
            Obiettivo aggiuntivo: imparare a \textbf{trattenere
                  l\textquotesingle aria nei polmoni il più a lungo possibile}, perché
            con i polmoni pieni si galleggia meglio
\end{itemize}

\paragraph{5. Tuffarsi in acqua}\label{tuffarsi-in-acqua}

\begin{itemize}
      \tightlist
      \item
            In vasca piccola (acqua bassa) ci si tuffa \textbf{sempre in piedi e
                  mai di testa}, per ragioni di sicurezza
      \item
            Il \textbf{tuffo di testa} è riservato esclusivamente alle attività in
            acqua fonda
\end{itemize}

\paragraph{6. Galleggiamento statico e
      dinamico}\label{galleggiamento-statico-e-dinamico}

\begin{itemize}
      \tightlist
      \item
            Galleggiamento sul petto (\emph{prono}) e sul dorso (\emph{supino})
      \item
            Prima in forma statica, poi in forma dinamica: gli
            \textbf{scivolamenti} (spingersi dal bordo e scivolare cercando di
            arrivare il più lontano possibile)
\end{itemize}

\begin{center}\rule{0.5\linewidth}{0.5pt}\end{center}

\subsection{3. Il Passaggio dall\textquotesingle Acqua Bassa
      all\textquotesingle Acqua
      Fonda}\label{il-passaggio-dallacqua-bassa-allacqua-fonda}

Imparare a nuotare significa imparare a nuotare dove non si tocca. Il
passaggio dall\textquotesingle acqua bassa all\textquotesingle acqua
alta rappresenta un \textbf{piccolo trauma}: tutte le acquisizioni
conquistate nelle settimane precedenti sembrano scomparire e il bambino
torna a essere rigido e spaventato.

\subsubsection{3.1 Introduzione Graduale all\textquotesingle Acqua
      Fonda}\label{introduzione-graduale-allacqua-fonda}

È fondamentale iniziare il prima possibile la conoscenza
dell\textquotesingle acqua fonda, già durante il lavoro in vasca
piccola. La "promozione" non deve mai avvenire in modo brusco, ma
attraverso un \textbf{percorso progressivo}:

\begin{longtable}[]{@{}p{0.47\linewidth}p{0.47\linewidth}@{}}
      \toprule\noalign{}
      Sessione         & Attività proposta                                 \\
      \midrule\noalign{}
      \endhead
      \bottomrule\noalign{}
      \endlastfoot
      Prima visita     & Portare i bambini a vedere la vasca grande (senza
      entrare)                                                             \\
      Seconda          & Farli sedere sul bordo e battere i piedi          \\
      Terza            & Brevi immersioni con uscita immediata             \\
      Progressivamente & Aumentare il tempo e le attività in acqua fonda   \\
\end{longtable}

\subsubsection{3.2 Priorità di Sicurezza in Acqua
      Fonda}\label{priorituxe0-di-sicurezza-in-acqua-fonda}

La prima cosa da assicurarsi prima di qualsiasi altro lavoro è che i
bambini siano in grado di \textbf{sostenersi dove non toccano}. Prima di
preoccuparsi di qualsiasi abilità tecnica, l\textquotesingle istruttore
deve insegnare loro a mantenere il \textbf{galleggiamento verticale} con
la testa fuori dall\textquotesingle acqua.

Qualsiasi tecnica è accettata:

\begin{itemize}
      \tightlist
      \item
            Nuoto a cagnolino
      \item
            Gambe a bicicletta
      \item
            Gambe a rana
\end{itemize}

L\textquotesingle importante è che, appena arrivati in vasca grande, i
bambini non rischino di affondare se staccano le mani dal bordo.

\subsubsection{3.3 L\textquotesingle Esercizio di
      Sicurezza}\label{lesercizio-di-sicurezza}

Esiste un esercizio specifico che fornisce
all\textquotesingle istruttore la \textbf{ragionevole tranquillità} che
i bambini siano sempre in grado di cavarsela in autonomia. La sequenza è
la seguente:

\begin{enumerate}
      \tightlist
      \item
            Staccarsi dal bordo
      \item
            Spostarsi per un breve tratto (\textbf{non più di 4-5 metri}) in
            posizione \textbf{prona}
      \item
            Raggiungere la posizione \textbf{verticale} e mantenerla per qualche
            secondo (suggerimento: contare fino a 10)
      \item
            Tornare al bordo in posizione \textbf{supina}
\end{enumerate}

La sequenza va poi ripetuta in senso inverso: partendo supini e tornando
pronti al bordo.

La capacità di svolgere questo esercizio in entrambi i sensi garantisce
all\textquotesingle istruttore che i bambini siano in grado di gestire
qualsiasi situazione di difficoltà in acqua fonda.

\begin{center}\rule{0.5\linewidth}{0.5pt}\end{center}

\subsection{4. Ambientamento in Assenza di Vasca Piccola (Solo Acqua
      Fonda)}\label{ambientamento-in-assenza-di-vasca-piccola-solo-acqua-fonda}

Questa è una situazione meno ottimale ma gestibile con le dovute
precauzioni.

\subsubsection{4.1 Posizione
      dell\textquotesingle Istruttore}\label{posizione-dellistruttore}

In acqua fonda --- intesa come l\textquotesingle acqua in cui neppure
l\textquotesingle istruttore tocca il fondo ---
l\textquotesingle istruttore deve \textbf{rimanere fuori
      dall\textquotesingle acqua}, per ragioni di sicurezza. Non può assistere
i bambini dall\textquotesingle interno della vasca senza mettere a
rischio la propria capacità di intervento.

\subsubsection{4.2 Inversione dell\textquotesingle Ordine delle
      Acquisizioni}\label{inversione-dellordine-delle-acquisizioni}

In vasca piccola la sequenza era: entrata in acqua → immersione del capo
→ respirazione → galleggiamento. In acqua fonda l\textquotesingle ordine
cambia radicalmente:

\begin{longtable}[]{@{}p{0.30\linewidth}p{0.30\linewidth}p{0.30\linewidth}@{}}
      \toprule\noalign{}
      Ordine       & Competenza                                         & Motivazione            \\
      \midrule\noalign{}
      \endhead
      \bottomrule\noalign{}
      \endlastfoot
      \textbf{1°}  & \textbf{Galleggiamento verticale}                  & Sicurezza assoluta: il
      bambino deve poter tenersi con la testa fuori
      dall\textquotesingle acqua                                                                 \\
      \textbf{2°}  & \textbf{Brevi spostamenti orizzontali}             & Una volta che il
      bambino si sente sicuro, si lavora sulla mobilità                                          \\
      \textbf{3°}  & Respirazione, immersione del capo, galleggiamento,
      scivolamenti & Solo dopo aver consolidato sicurezza e mobilità di
      base                                                                                       \\
\end{longtable}

La logica è semplice: una volta che il bambino è in grado di mantenersi
con la testa fuori dall\textquotesingle acqua, la sua autostima migliora
e diventa più disponibile agli apprendimenti successivi.

\subsubsection{4.3 Condizioni Operative}\label{condizioni-operative}

\begin{itemize}
      \tightlist
      \item
            Il lavoro in acqua fonda senza vasca piccola richiede \textbf{numeri
                  di allievi molto ridotti}
      \item
            Non è la situazione ideale, ma con attenzione e piccoli gruppi il
            disagio è superabile
\end{itemize}

\begin{center}\rule{0.5\linewidth}{0.5pt}\end{center}

\subsection{Conclusione: Punti Chiave per
      l\textquotesingle Esame}\label{conclusione-punti-chiave-per-lesame-6}

\begin{enumerate}
      \item
            \textbf{Definizione di ambientamento}: raggiungimento di uno stato di
            benessere psicofisico conseguente all\textquotesingle immersione e
            alla conoscenza di come il corpo reagisce all\textquotesingle acqua. È
            la fase preliminare a qualsiasi competenza motoria acquatica.
      \item
            \textbf{Acquaticità 0-3 e 3-6 anni}: non per insegnare a nuotare, ma
            per far considerare l\textquotesingle acqua un elemento amico.
      \item
            \textbf{Istruttore in vasca piccola}: sempre in acqua con i bambini.
            Spazio ridotto a 4-5 metri di lato. Approccio ludico prevalente:
            parlare poco, far muovere molto.
      \item
            \textbf{Indicatore dell\textquotesingle ambientamento}: il grado di
            tensione muscolare del bambino. Maggiore rigidità = maggiore paura.
      \item
            \textbf{Respirazione in acqua}: inspirare con la \textbf{bocca} (non
            col naso per evitare inalazione d\textquotesingle acqua), espirare con
            bocca o naso. Espirazione forzata perché la pressione
            dell\textquotesingle acqua la impedisce passivamente.
      \item
            \textbf{Tuffo in vasca piccola}: sempre in piedi, mai di testa. Il
            tuffo di testa solo in acqua fonda.
      \item
            \textbf{Passaggio all\textquotesingle acqua fonda}: avviene
            gradualmente durante le lezioni in vasca piccola. Non va mai fatto in
            modo brusco.
      \item
            \textbf{Esercizio di sicurezza in vasca grande}: staccarsi dal bordo →
            4-5 metri in posizione prona → verticalizzazione (contare fino a 10) →
            ritorno al bordo in posizione supina. Da ripetere anche in senso
            inverso.
      \item
            \textbf{Solo acqua fonda}: l\textquotesingle istruttore rimane fuori
            dall\textquotesingle acqua. L\textquotesingle ordine delle
            acquisizioni si inverte: prima il galleggiamento verticale, poi gli
            spostamenti, poi tutto il resto. Gruppi molto piccoli.
\end{enumerate}

\section{Galleggiamenti, Spostamenti, Posture e
  Propulsioni}\label{galleggiamenti-spostamenti-posture-e-propulsioni}

\textbf{Docente:} Cristiano Guerra, Settore Istruzione Tecnica della
Federazione Italiana Nuoto\\
\textbf{Riferimento dispensa:} punti 4.9 e 4.10

\begin{center}\rule{0.5\linewidth}{0.5pt}\end{center}

\subsection{Introduzione}\label{introduzione-11}

Questa lezione costituisce il ponte teorico tra
l\textquotesingle ambientamento e l\textquotesingle apprendimento
tecnico delle nuotate. Vengono trattati: la distinzione fondamentale tra
ambientamento e acquaticità, il ruolo del corpo in acqua (centro di
gravità, centro di galleggiamento, tre masse critiche), la postura
corretta, il ruolo della respirazione, i principi metodologici per
costruire qualsiasi nuotata (continuità, ampiezza, traiettoria,
posizione del corpo), il concetto di forma e contenuto, le
caratteristiche generali dei quattro stili e il cosiddetto "quinto
stile".

\begin{center}\rule{0.5\linewidth}{0.5pt}\end{center}

\subsection{1. Ambientamento e Acquaticità: una Distinzione
      Fondamentale}\label{ambientamento-e-acquaticituxe0-una-distinzione-fondamentale}

Questi due termini vengono spesso confusi ma indicano cose radicalmente
diverse.

\begin{longtable}[]{@{}p{0.30\linewidth}p{0.30\linewidth}p{0.30\linewidth}@{}}
      \toprule\noalign{}
      Concetto               & Natura                                  & Durata                     \\
      \midrule\noalign{}
      \endhead
      \bottomrule\noalign{}
      \endlastfoot
      \textbf{Ambientamento} & Una \textbf{fase}                       & Ha un inizio e una fine: è
      la fase preliminare all\textquotesingle acquisizione di qualsiasi
      competenza motoria acquatica                                                                  \\
      \textbf{Acquaticità}   & Un\textquotesingle{}\textbf{attitudine} & Inizia
      parallelamente all\textquotesingle ambientamento in forma elementare e
      \textbf{non termina mai}                                                                      \\
\end{longtable}

L\textquotesingle acquaticità è una caratteristica che si sviluppa
continuamente. Anche gli atleti di alto livello svolgono esercitazioni
sempre più fini di acquaticità. Quando un atleta si allontana troppo a
lungo dall\textquotesingle acqua perde parte di questa attitudine e deve
ripristinarla con esercizi specifici.

\subsubsection{1.1 Esperienze vs.
      Esercizi}\label{esperienze-vs.-esercizi}

\begin{itemize}
      \tightlist
      \item
            Nella fase di \textbf{ambientamento}: si propongono
            \textbf{esperienze}, non esercizi
      \item
            Nella fase di \textbf{scuola nuoto}: si propongono \textbf{esercizi}
            (e anche esperienze)
\end{itemize}

\textbf{Errore comune:} anticipare l\textquotesingle approccio
tecnico-meccanico durante la fase di ambientamento. In questa fase
bisogna privilegiare la \textbf{quantità
      dell\textquotesingle esperienza} rispetto alla qualità tecnica. La
qualità tecnica diventa progressivamente il focus man mano che
l\textquotesingle allievo avanza nel percorso.

Se non vengono costruiti correttamente i presupposti
dell\textquotesingle ambientamento, tutto
l\textquotesingle apprendimento tecnico successivo viene compromesso.

\begin{center}\rule{0.5\linewidth}{0.5pt}\end{center}

\subsection{2. Il Corpo in Acqua: Centro di Gravità e Centro di
      Galleggiamento}\label{il-corpo-in-acqua-centro-di-gravituxe0-e-centro-di-galleggiamento}

Ogni soggetto in acqua ha due punti critici da allineare:

\begin{itemize}
      \tightlist
      \item
            \textbf{Centro di gravità}: punto in cui si concentra la massa
            corporea
      \item
            \textbf{Centro di galleggiamento}: punto in cui si concentra la spinta
            idrostatica
\end{itemize}

Quando questi due centri non sono correttamente collocati, il corpo
tende a inclinarsi, con le gambe che cadono verso il fondo. Questo è
particolarmente evidente negli adulti principianti, che hanno una massa
delle gambe maggiore rispetto ai bambini.

\textbf{Obiettivo:} collocare centro di gravità e centro di
galleggiamento in modo da garantire stabilità e assumere una
\textbf{posizione orizzontale} che riduca al minimo la resistenza
all\textquotesingle avanzamento (\emph{drag}).

\begin{center}\rule{0.5\linewidth}{0.5pt}\end{center}

\subsection{3. Le Tre Masse Critiche}\label{le-tre-masse-critiche}

Il corpo del nuotatore può essere immaginato come un kayak sezionato in
tre parti: se le sezioni non sono allineate, la penetrazione
nell\textquotesingle acqua risulta disturbata.

Le tre masse critiche sono:

\begin{enumerate}
      \tightlist
      \item
            \textbf{Testa}
      \item
            \textbf{Torace}
      \item
            \textbf{Gambe}
\end{enumerate}

Queste tre masse devono essere allineate fra loro in modo da:

\begin{itemize}
      \tightlist
      \item
            Garantire stabilità in acqua
      \item
            Ridurre il \emph{drag} (resistenza all\textquotesingle avanzamento)
\end{itemize}

Se il corpo non è in equilibrio, parte dell\textquotesingle energia
prodotta da braccia e gambe verrà spesa per compensare
l\textquotesingle instabilità, invece di essere interamente destinata
all\textquotesingle avanzamento. Nei principianti questo si manifesta
con movimenti circolari e compensatori (non finalizzati
all\textquotesingle avanzamento ma al sostentamento).

\subsubsection{3.1 La Posizione della Testa: Fonte Primaria degli
      Errori}\label{la-posizione-della-testa-fonte-primaria-degli-errori}

La testa --- nei bambini proporzionalmente più grande e pesante rispetto
al corpo --- è la massa critica che più frequentemente scatena errori a
cascata:

\begin{itemize}
      \tightlist
      \item
            \textbf{Testa troppo alta} → i fianchi si abbassano → le gambe cadono
            → postura non corretta
      \item
            L\textquotesingle intervento corretto non è "solleva i fianchi" ma
            "abbassa la testa"
\end{itemize}

\begin{quote}
      La testa "regna sovrana" nello scatenare gli errori di postura.
\end{quote}

\textbf{Posizione corretta della testa:}

\begin{itemize}
      \tightlist
      \item
            \textbf{Ovale del viso appoggiato all\textquotesingle acqua}
      \item
            \textbf{Collo lungo e rilassato} (non retratto, non con pieghe nella
            parte posteriore)
      \item
            \textbf{Sguardo verso il basso}, non diretto in avanti
      \item
            La nuca deve essere piatta e liscia, senza sollevamenti né
            abbassamenti eccessivi
\end{itemize}

La tendenza naturale è guardare l\textquotesingle orizzonte (riflesso
terrestre): va corretta sin dalle prime proposte di galleggiamento e
scivolamento.

Questa posizione può essere fatta sentire anche a \textbf{secco} (fuori
dall\textquotesingle acqua) o in vaschetta, per creare la sensazione
corretta prima di trasferirla in acqua.

\begin{center}\rule{0.5\linewidth}{0.5pt}\end{center}

\subsection{4. La Respirazione come Azione
      Motoria}\label{la-respirazione-come-azione-motoria}

La respirazione in acqua non è un accessorio: il docente usa il termine
\textbf{"respirazione-azione"} per sottolineare quanto incida
sull\textquotesingle azione motoria complessiva.

\subsubsection{4.1 Il Problema della Respirazione
      Superficiale}\label{il-problema-della-respirazione-superficiale}

Accontentarsi di bambini o adulti che respirano in modo regolare ma non
profondo, autonomo e dinamico è un errore metodologico grave. Un allievo
che lavora al 75-80\% della propria capacità vitale sarà limitato nello
sviluppo atletico, e questo viene spesso erroneamente interpretato come
un limite del potenziale individuale, quando è invece un \textbf{limite
      tecnico} superabile.

\subsubsection{4.2 Esercitazioni sulla
      Respirazione}\label{esercitazioni-sulla-respirazione}

\begin{itemize}
      \tightlist
      \item
            Galleggiamenti e scivolamenti \textbf{in apnea}: amplificano il senso
            di galleggiamento (polmoni pieni = si galleggia meglio)
      \item
            Galleggiamenti e scivolamenti \textbf{in espirazione}: si percepisce
            come i polmoni senza aria fanno affondare
      \item
            \textbf{Esercizi di autoaffondamento}: posizione raccolta → apnea per
            qualche secondo → espirazione completa → raggiungimento del fondo
            della vasca
      \item
            Rapporto respiratorio di riferimento: \textbf{3:1} (3 cicli o
            bracciata su 1 respirazione), da praticare sia in forma statica che
            dinamica
\end{itemize}

Un consolidamento insufficiente della respirazione nelle prime fasi è
uno dei principali fattori limitanti a lungo termine, anche per atleti
di alto livello.

\begin{center}\rule{0.5\linewidth}{0.5pt}\end{center}

\subsection{5. Multilateralità e Trasferibilità: una Necessità
      Didattica}\label{multilateralituxe0-e-trasferibilituxe0-una-necessituxe0-didattica}

Le prime proposte nella scuola di educazione acquatica non sono elementi
tipici del nuoto. Sono più vicine al \textbf{nuoto artistico} (controllo
della postura, equilibrio in acqua, verticalità) o alla
\textbf{pallanuoto} (autosostentamento in acqua alta, recupero della
posizione verticale, posizione orizzontale).

Attingere trasversalmente a tutte le discipline acquatiche non è una
scelta metodologica: è una \textbf{necessità didattica}. La
multilateralità arricchisce il bagaglio motorio e offre possibilità di
trasferire e intrecciare abilità tra discipline diverse,
un\textquotesingle opportunità enorme nell\textquotesingle ambiente
acquatico.

\begin{center}\rule{0.5\linewidth}{0.5pt}\end{center}

\subsection{6. I Tre Stadi del Percorso
      dell\textquotesingle Allievo}\label{i-tre-stadi-del-percorso-dellallievo}

\begin{longtable}[]{@{}p{0.22\linewidth}p{0.22\linewidth}p{0.22\linewidth}p{0.22\linewidth}@{}}
      \toprule\noalign{}
      Fase                    & Nome                 & Rapporto con l\textquotesingle acqua                  & Strumento
      didattico                                                                                                          \\
      \midrule\noalign{}
      \endhead
      \bottomrule\noalign{}
      \endlastfoot
      1                       & Educazione acquatica & L\textquotesingle allievo si \textbf{muove}
      in acqua                & Esperienze                                                                               \\
      2                       & Scuola nuoto         & L\textquotesingle allievo si \textbf{sposta} in acqua
                              & Esercizi                                                                                 \\
      3                       & Agonismo             & L\textquotesingle allievo \textbf{attraversa}
      l\textquotesingle acqua & Allenamenti                                                                              \\
\end{longtable}

\begin{center}\rule{0.5\linewidth}{0.5pt}\end{center}

\subsection{7. Come si Costruisce una Nuotata: Principio
      Metodologico}\label{come-si-costruisce-una-nuotata-principio-metodologico}

\subsubsection{7.1 Togliere, non
      Aggiungere}\label{togliere-non-aggiungere}

Il principio fondamentale nella costruzione di qualsiasi nuotata è
\textbf{sottrarre, non aggiungere}:

\begin{itemize}
      \tightlist
      \item
            Lasciare solo ciò che serve
      \item
            Costruire prima una struttura di base solida, poi aggiungere i
            dettagli
      \item
            L\textquotesingle atleta di alto livello trasmette immediatamente
            semplicità ed efficacia, perché usa esclusivamente i movimenti
            finalizzati all\textquotesingle obiettivo
\end{itemize}

\subsubsection{7.2 L\textquotesingle Azione Motoria vs. il Semplice
      Movimento}\label{lazione-motoria-vs.-il-semplice-movimento}

Il \textbf{movimento} non è un\textquotesingle azione motoria. Il
processo corretto è:

\begin{enumerate}
      \tightlist
      \item
            \textbf{Comprendere il compito e lo scopo} (es. le braccia devono far
            avanzare il corpo, non solo girare)
      \item
            \textbf{Elaborare una strategia} (la mente lavora insieme al corpo)
      \item
            \textbf{Mettere in pratica} i tentativi (\emph{prova ed errori})
      \item
            \textbf{Eseguire il movimento} (ultima fase del processo)
\end{enumerate}

Proporre solo il movimento senza passare lo scopo determina
\textbf{addestramento}.\\
Proporre l\textquotesingle azione motoria completa con scopo e strategia
determina \textbf{apprendimento}.

\subsubsection{7.3 Forma vs. Contenuto}\label{forma-vs.-contenuto}

Gli istruttori meno esperti tendono a copiare l\textquotesingle aspetto
\textbf{formale} (estetico, coreografico) di quello che vedono, senza
trasmettere i \textbf{contenuti}.

\begin{itemize}
      \tightlist
      \item
            La \textbf{forma} copiata senza contenuto non garantisce propulsione
            né avanzamento
      \item
            I \textbf{contenuti} trasmessi correttamente generano forma
            personalizzata sulle caratteristiche individuali
\end{itemize}

Esempio per la bracciata a stile libero:

\begin{quote}
      "Aggancia l\textquotesingle acqua in avanti, portala forte sotto la
      pancia, spingi forte verso i piedi"
\end{quote}

Questi contenuti restituiscono una forma che ciascun allievo
personalizzerà in base alle proprie leve e proporzioni corporee. Due
allievi con leve diverse produrranno forme diverse dallo stesso
contenuto: è la tecnica sportiva corretta, non una coreografia uniforme.

\subsubsection{7.4 Feedback e Correzione}\label{feedback-e-correzione}

Dopo l\textquotesingle azione motoria avviene il \textbf{controllo e la
      regolazione del gesto}:

\begin{itemize}
      \tightlist
      \item
            \textbf{Feedback interno} (\emph{fit-make interno}): il processo di
            autoregolazione dell\textquotesingle allievo
      \item
            \textbf{Feedback esterno} (\emph{fit-make esterno}): le indicazioni
            dell\textquotesingle istruttore
\end{itemize}

\begin{quote}
      La correzione non la fa l\textquotesingle istruttore: la esegue
      l\textquotesingle allievo. L\textquotesingle istruttore fornisce
      \textbf{indicazioni} perché l\textquotesingle allievo possa autoregolare
      il gesto.
\end{quote}

Nel nuoto l\textquotesingle allievo non può vedere molti aspetti della
propria tecnica: l\textquotesingle istruttore è i suoi occhi,
restituisce informazioni che l\textquotesingle allievo usa per
autoregolarsi verso il modello ottimale.

\begin{center}\rule{0.5\linewidth}{0.5pt}\end{center}

\subsection{8. I Quattro Parametri Fondamentali di Qualsiasi
      Nuotata}\label{i-quattro-parametri-fondamentali-di-qualsiasi-nuotata}

Questi quattro elementi sono la struttura di base di tutte le nuotate.
Se uno soltanto manca, l\textquotesingle apprendimento tecnico risulta
compromesso. I dettagli vengono dopo.

\subsubsection{8.1 Posizione del Corpo in
      Acqua}\label{posizione-del-corpo-in-acqua}

Come discusso in precedenza: tre masse critiche allineate, testa in
posizione corretta, \emph{drag} ridotto. È il presupposto di tutto.

\subsubsection{8.2 Continuità
      d\textquotesingle Azione}\label{continuituxe0-dazione}

Il movimento deve essere \textbf{continuo, senza pause}. In tutte e
quattro le nuotate non esistono interruzioni del gesto (la rana ha una
pausa di scivolamento ma ha continuità d\textquotesingle azione
nell\textquotesingle insieme).

Attenzione: non proporre mai esercitazioni che prevedano pause interne
all\textquotesingle esercizio, perché questo si trasferisce nella
nuotata come abitudine negativa.

\subsubsection{8.3 Ampiezza del Movimento}\label{ampiezza-del-movimento}

I movimenti devono essere \textbf{ampi}, non corti e stretti. Bisogna
completare tutte le fasi di spinta per applicare forza per il tempo più
lungo possibile e ottenere il massimo spostamento.
Un\textquotesingle ampiezza ridotta equivale a perdita di propulsione.

\subsubsection{8.4 Traiettoria}\label{traiettoria}

La \textbf{traiettoria} è il percorso descritto dalla mano durante la
passata subacquea (\emph{underwater pull}):

\begin{itemize}
      \tightlist
      \item
            Non deve attraversare la linea mediana del corpo
      \item
            Non deve essere troppo esterna (mano troppo lontana dal corpo)
      \item
            Deve rimanere \textbf{piuttosto vicina al corpo}
      \item
            Esistono variabili soggettive e limiti di tolleranza individuali
\end{itemize}

\begin{center}\rule{0.5\linewidth}{0.5pt}\end{center}

\subsection{\texorpdfstring{9. Visione Moderna della Tecnica: il
            Concetto di \emph{Inside
                  Out}}{9. Visione Moderna della Tecnica: il Concetto di Inside Out}}\label{visione-moderna-della-tecnica-il-concetto-di-inside-out}

\subsubsection{9.1 Vecchia vs. Nuova
      Visione}\label{vecchia-vs.-nuova-visione}

\begin{longtable}[]{@{}p{0.30\linewidth}p{0.30\linewidth}p{0.30\linewidth}@{}}
      \toprule\noalign{}
      Visione                       & Approccio                               & Ruolo del corpo \\
      \midrule\noalign{}
      \endhead
      \bottomrule\noalign{}
      \endlastfoot
      \textbf{Vecchia (settoriale)} & Gambe + braccia + respirazione separati
                                    & Il corpo è semplicemente da trasportare                   \\
      \textbf{Nuova (integrata)}    & Il corpo nuota nella sua interezza      & Il
      corpo genera e trasmette la forza dal nucleo alla periferia                               \\
\end{longtable}

\subsubsection{\texorpdfstring{9.2 Il Concetto di \emph{Inside Out}
            (dall\textquotesingle Interno
            all\textquotesingle Esterno)}{9.2 Il Concetto di Inside Out (dall\textquotesingle Interno all\textquotesingle Esterno)}}\label{il-concetto-di-inside-out-dallinterno-allesterno}

La forza non si genera localmente nel braccio o nella gamba: viene
generata nel \textbf{nucleo} (\emph{core}), ovvero nella zona
addomino-lombare, e trasmessa alla periferia in due direzioni (verso
l\textquotesingle alto e verso il basso).

Come nel lancio del giavellotto: non è il lancio del braccio, è una
trasmissione di impulso dal \emph{core} all\textquotesingle estremità.

Questo è il motivo per cui costruire un \textbf{corpo stabile in acqua}
è il presupposto di tutto: senza stabilità del \emph{core}, la forza
generata si disperde in micromovimenti compensatori anziché tradursi in
avanzamento.

\begin{center}\rule{0.5\linewidth}{0.5pt}\end{center}

\subsection{10. Caratteristiche Generali delle Quattro
      Nuotate}\label{caratteristiche-generali-delle-quattro-nuotate}

\subsubsection{10.1 Movimenti Comuni a Tutti gli
      Stili}\label{movimenti-comuni-a-tutti-gli-stili}

\begin{itemize}
      \tightlist
      \item
            Tutti i movimenti sono di \textbf{flessione ed estensione} continua
            (braccia e gambe si flettono e si distendono in modo costante)
      \item
            Presenza di movimenti sia \textbf{circolari} che \textbf{rettilinei},
            che interagiscono tra loro nella passata subacquea (movimenti molto
            più ridotti rispetto agli anni \textquotesingle80)
      \item
            \textbf{Velocità non costante} all\textquotesingle interno del gesto:
            le fasi finali di spinta vengono eseguite in \textbf{forte
                  accelerazione}; le fasi iniziali di appoggio e trazione sono più lente
\end{itemize}

\subsubsection{10.2 Roll e Ondulazione}\label{roll-e-ondulazione}

\begin{longtable}[]{@{}p{0.30\linewidth}p{0.30\linewidth}p{0.30\linewidth}@{}}
      \toprule\noalign{}
      Movimento                     & Stili                                       & Descrizione \\
      \midrule\noalign{}
      \endhead
      \bottomrule\noalign{}
      \endlastfoot
      \textbf{Rollio} (\emph{roll}) & Stile libero (\emph{freestyle/crawl}),
      Dorso (\emph{backstroke})     & Movimento laterale del corpo che parte dai
      \textbf{fianchi}. A stile libero non si gira la testa per respirare:
      sono i fianchi che ruotano, la testa rimane in asse                                       \\
      \textbf{Ondulazione}          & Rana (\emph{breaststroke}), Farfalla
      (\emph{butterfly/delfino})    & Movimento del bacino che schiaccia verso il
      basso; il corpo si muove come un\textquotesingle onda                                     \\
\end{longtable}

\subsubsection{10.3 Efficienza Comparata dei Quattro
      Stili}\label{efficienza-comparata-dei-quattro-stili}

\begin{longtable}[]{@{}p{0.30\linewidth}p{0.30\linewidth}p{0.30\linewidth}@{}}
      \toprule\noalign{}
      Stile                                          & Efficienza & Motivazione              \\
      \midrule\noalign{}
      \endhead
      \bottomrule\noalign{}
      \endlastfoot
      \textbf{Stile libero} (\emph{freestyle/crawl}) & Massima    & Posizione del
      corpo favorevole, grandi masse muscolari coinvolte, resistenza quasi
      costante                                                                               \\
      \textbf{Farfalla} (\emph{butterfly/delfino})   & Alta       & Simile al dorso          \\
      \textbf{Dorso} (\emph{backstroke})             & Alta       & Simile alla farfalla     \\
      \textbf{Rana} (\emph{breaststroke})            & Minima     & Posizione del corpo crea
      maggiore \emph{drag}; forte coinvolgimento del quadricipite femorale (a
      differenza degli altri stili); variazioni di resistenza durante il
      ciclo                                                                                  \\
\end{longtable}

\subsubsection{10.4 La Funzione delle Gambe: Triplice
      Valenza}\label{la-funzione-delle-gambe-triplice-valenza}

La battuta di gambe ha tre funzioni che variano di importanza nel corso
del percorso didattico:

\begin{longtable}[]{@{}p{0.30\linewidth}p{0.30\linewidth}p{0.30\linewidth}@{}}
      \toprule\noalign{}
      Funzione               & Importanza in scuola nuoto      & Importanza nel nuoto avanzato \\
      \midrule\noalign{}
      \endhead
      \bottomrule\noalign{}
      \endlastfoot
      \textbf{Propulsiva}    & Alta (primo motore del bambino) & Minore
      (compensatoria)                                                                          \\
      \textbf{Galleggiante}  & Media                           & Alta                          \\
      \textbf{Stabilizzante} & Media                           & Alta                          \\
\end{longtable}

Man mano che la scuola nuoto avanza, le funzioni galleggiante e
stabilizzante diventano prioritarie rispetto a quella propulsiva.
\textbf{Eccezione:} nella rana il movimento delle gambe mantiene
un\textquotesingle importanza simile a quella delle braccia anche ai
livelli avanzati.

\begin{center}\rule{0.5\linewidth}{0.5pt}\end{center}

\subsection{11. Il Quinto Stile}\label{il-quinto-stile}

Il \textbf{quinto stile} comprende tutte le fasi svolte
\textbf{sott\textquotesingle acqua}: le partenze (\emph{starts}) e le
virate (\emph{turns}). Viene chiamato "quinto stile" perché può arrivare
a coprire fino al \textbf{60\% di una gara} in vasca corta (\emph{short
      course}): non è un semplice dettaglio tecnico ma a tutti gli effetti un
vero stile.

È presente in tutte e quattro le nuotate e può essere lavorato nella
scuola nuoto attraverso:

\begin{itemize}
      \tightlist
      \item
            Scivolamenti superficiali (\emph{surface glides})
      \item
            Scivolamenti subacquei (\emph{underwater glides})
      \item
            Galleggiamenti verticali
      \item
            Tuffo e scivolamento
      \item
            Bracciata subacquea a rana (\emph{underwater breaststroke pull}) nel
            dorso
\end{itemize}

\textbf{Posizione caratteristica del quinto stile:}

\begin{itemize}
      \tightlist
      \item
            Collo lungo
      \item
            Braccia distese
      \item
            Gomiti stretti
      \item
            Corpo in tenuta (\emph{streamline})
\end{itemize}

\begin{quote}
      La posizione del quinto stile è evolutivamente simile ai semplici
      scivolamenti della scuola nuoto elementare: si tratta dello stesso
      gesto, reso progressivamente più raffinato e preciso negli anni.
\end{quote}

\begin{center}\rule{0.5\linewidth}{0.5pt}\end{center}

\subsection{Conclusione: Punti Chiave per
      l\textquotesingle Esame}\label{conclusione-punti-chiave-per-lesame-7}

\begin{enumerate}
      \item
            \textbf{Ambientamento vs. acquaticità}: il primo è una fase con inizio
            e fine; la seconda è un\textquotesingle attitudine che non termina
            mai. In ambientamento si propongono esperienze, non esercizi.
      \item
            \textbf{Tre masse critiche}: testa, torace, gambe devono essere
            allineate per ridurre il \emph{drag} e garantire stabilità. La testa è
            la causa principale degli errori a cascata.
      \item
            \textbf{Posizione corretta della testa}: ovale del viso appoggiato
            all\textquotesingle acqua, collo lungo, sguardo verso il basso. Mai
            guardare in avanti (riflesso terrestre da correggere).
      \item
            \textbf{Respirazione-azione}: non va trattata superficialmente.
            Lavorare al 75-80\% della capacità vitale è un limite tecnico, non di
            potenziale. Usare esercizi in apnea e in espirazione; rapporto 3:1.
      \item
            \textbf{Tre fasi del percorso}: si muove (educazione acquatica) → si
            sposta (scuola nuoto) → attraversa (agonismo).
      \item
            \textbf{Principio metodologico}: togliere, non aggiungere. Proporre
            l\textquotesingle azione motoria (scopo + strategia + movimento), non
            il semplice movimento. Il movimento è l\textquotesingle ultima fase
            del processo.
      \item
            \textbf{Forma vs. contenuto}: trasmettere contenuti ("aggancia, tira,
            spingi"), non forme estetiche. I contenuti generano forma
            personalizzata; la forma senza contenuto è solo coreografia.
      \item
            \textbf{Quattro parametri fondamentali} di qualsiasi nuotata:
            posizione del corpo, continuità d\textquotesingle azione, ampiezza del
            movimento, traiettoria. Tutti e quattro devono essere presenti.
      \item
            \textbf{Inside out}: la forza viene generata nel \emph{core} e
            trasmessa alla periferia. Il corpo nuota nella sua interezza, non è
            solo trasportato.
      \item
            \textbf{Gambe --- triplice valenza}: propulsiva (fondamentale
            all\textquotesingle inizio), galleggiante e stabilizzante
            (fondamentali nell\textquotesingle avanzamento).
            L\textquotesingle efficienza degli stili: stile libero \textgreater{}
            farfalla $\approx$ dorso \textgreater{} rana.
      \item
            \textbf{Quinto stile}: fasi subacquee che possono coprire il 60\%
            della gara. Va costruito fin dalla scuola nuoto attraverso
            scivolamenti ed esercizi di \emph{streamline}.
\end{enumerate}

\section{13. Tecnica e Didattica delle Quattro Nuotate
  Regolamentari}\label{tecnica-e-didattica-delle-quattro-nuotate-regolamentari}

Docente: Cristiano Guerra, settore istruzione tecnica della Federazione
Italiana Nuoto.\\
Riferimento dispensa: capitoli 4.12 -- 4.16.

\begin{center}\rule{0.5\linewidth}{0.5pt}\end{center}

\subsection{Introduzione}\label{introduzione-12}

Questa lezione affronta la tecnica e la didattica delle quattro nuotate
regolamentari --- stile libero, dorso, rana e farfalla --- con
l\textquotesingle obiettivo di fornire all\textquotesingle istruttore un
modello prestativo universale di riferimento. Il modello non va inteso
come schema rigido da replicare in ogni allievo, ma come strumento di
lettura del gesto tecnico per guidare l\textquotesingle allievo verso la
forma più efficace ed efficiente in base alle sue caratteristiche
individuali.

Un concetto trasversale a tutta la lezione: ogni nuotata è sempre un
\textbf{movimento unico}, suddiviso in fasi per comodità descrittiva. Le
fasi si fondono l\textquotesingle una nell\textquotesingle altra e la
qualità di ciascuna determina la riuscita della successiva.

\begin{center}\rule{0.5\linewidth}{0.5pt}\end{center}

\subsection{\texorpdfstring{Stile Libero (\emph{Freestyle /
                  Crawl})}{Stile Libero (Freestyle / Crawl)}}\label{stile-libero-freestyle-crawl}

\subsubsection{Caratteristiche Generali}\label{caratteristiche-generali}

\begin{itemize}
      \tightlist
      \item
            Nuotata \textbf{asimmetrica}: gli arti superiori eseguono movimenti
            alternati e opposti.
      \item
            \textbf{Alto rendimento meccanico}: è la nuotata più efficiente del
            ventaglio regolamentare.
      \item
            Fra le prime nuotate proposte nella scuola nuoto.
      \item
            Offre un\textquotesingle{}\textbf{ottima posizione idrodinamica} e una
            buona stabilità in acqua, garantita dalla battuta di gambe e dal
            rollio.
      \item
            È presente il \textbf{rollio}: movimento laterale che parte dai
            fianchi e viene trasmesso al busto; la testa rimane in asse.
      \item
            \textbf{Tolleranza agli errori coordinativi} nelle fasi di
            apprendimento (a differenza di rana e farfalla, dove errori di
            coordinazione compromettono l\textquotesingle intera nuotata).
      \item
            La coordinazione interna prevede \textbf{tre ritmi} da unire: gambe,
            braccia e respirazione.
\end{itemize}

\subsubsection{Azione delle Braccia}\label{azione-delle-braccia}

La bracciata è composta da due fasi principali:

\begin{itemize}
      \tightlist
      \item
            \textbf{Fase aerea} (passiva): il recupero fuori
            dall\textquotesingle acqua.
      \item
            \textbf{Fase subacquea} (attiva): genera
            l\textquotesingle avanzamento.
\end{itemize}

La fase subacquea si articola in quattro momenti che si fondono in un
movimento continuo:

\begin{longtable}[]{@{}p{0.30\linewidth}p{0.30\linewidth}p{0.30\linewidth}@{}}
      \toprule\noalign{}
      Fase                                                    & Descrizione                                        & Velocità esecutiva \\
      \midrule\noalign{}
      \endhead
      \bottomrule\noalign{}
      \endlastfoot
      \textbf{Appoggio e presa}                               & La mano si appoggia in acqua e aggancia,
      "prende" l\textquotesingle acqua                        & Moderata                                                                \\
      \textbf{Trazione}                                       & Il corpo viene tirato in avanti; termina quando si
      perde la perpendicolarità del braccio rispetto al corpo & Moderata                                                                \\
      \textbf{Spinta}                                         & Il corpo viene spinto in avanti con \textbf{forte
      accelerazione}                                          & Elevata -- in crescendo                                                 \\
      \textbf{Recupero}                                       & Inizia già sott\textquotesingle acqua nella fase
      finale di spinta                                        & Decontratto                                                             \\
\end{longtable}

\begin{quote}
      \textbf{Punti critici per l\textquotesingle istruttore:}

      \begin{itemize}
            \tightlist
            \item
                  Il recupero \textbf{non} inizia quando il braccio esce
                  dall\textquotesingle acqua, ma già nella fase finale di spinta.
            \item
                  La \textbf{spinta finale} deve essere eseguita in fortissima
                  accelerazione: una spinta lenta causa una pausa che rompe la
                  continuità.
            \item
                  La fase finale di spinta avviene dietro l\textquotesingle allievo, che
                  non la percepisce visivamente: usare feedback verbali e dimostrazione
                  a secco.
            \item
                  Il recupero aereo si fa a \textbf{gomito alto} (è il gomito a guidare
                  il riporto del braccio, non la mano) con muscolatura decontratta.
      \end{itemize}
\end{quote}

\subsubsection{Rollio}\label{rollio}

\begin{itemize}
      \tightlist
      \item
            Movimento laterale che origina dai fianchi, si trasmette al busto; la
            testa rimane in asse agevolando la respirazione.
      \item
            È necessario per coinvolgere il nucleo corporeo (\emph{core}) nella
            generazione della forza, trasmessa poi alla periferia (analogia con il
            lanciatore di giavellotto).
      \item
            Un \textbf{rollio eccessivo} costringe braccia e gambe a funzione
            stabilizzante invece che propulsiva: attenzione alle cause
            (instabilità di base).
      \item
            Come conseguenza del rollio, la battuta di gambe non è mai puramente
            verticale ma \textbf{obliqua/diagonale}.
\end{itemize}

\subsubsection{Battuta di Gambe}\label{battuta-di-gambe}

\begin{itemize}
      \tightlist
      \item
            \textbf{Due fasi}: discendente (attiva, eseguita con \emph{frustata}
            ad alta velocità) e ascendente (passiva, muscolatura decontratta).
      \item
            \textbf{Errore comune}: applicare la stessa intensità in entrambe le
            fasi → il recupero attivo azzera l\textquotesingle avanzamento.
      \item
            \textbf{6 colpi di gambe per ciclo di bracciata}: la mano cambia tre
            volte direzione durante la passata subacquea; ciascun cambio di
            direzione viene compensato in modo inconsapevole da un colpo di gambe
            (3 cambi × 2 = 6).
      \item
            Funzione delle gambe nello stile libero: \textbf{galleggiante},
            \textbf{stabilizzante} e in misura minore propulsiva (quando le
            braccia entrano in azione, la loro importanza propulsiva supera quella
            delle gambe).
\end{itemize}

\subsubsection{Respirazione}\label{respirazione}

\begin{itemize}
      \tightlist
      \item
            Collegata al timing del rollio: la testa ruota di lato, non si
            solleva.
      \item
            La coordinazione respiratoria si affina con il progredire della
            tecnica; il principiante respira spontaneamente,
            l\textquotesingle atleta inserisce la respirazione nel timing della
            nuotata.
\end{itemize}

\begin{center}\rule{0.5\linewidth}{0.5pt}\end{center}

\subsection{\texorpdfstring{Dorso
            (\emph{Backstroke})}{Dorso (Backstroke)}}\label{dorso-backstroke}

\subsubsection{Caratteristiche
      Generali}\label{caratteristiche-generali-1}

\begin{itemize}
      \tightlist
      \item
            Nuotata \textbf{asimmetrica} (movimento alternato degli arti),
            eseguita in \textbf{posizione supina}.
      \item
            Buon rendimento meccanico; fra i primi stili insegnati.
      \item
            Il viso è fuori dall\textquotesingle acqua, ma la respirazione ha
            comunque una \textbf{ritmica precisa} collegata alla bracciata.
      \item
            Presenza del \textbf{rollio} come nello stile libero.
      \item
            È una \textbf{nuotata cieca}: la fase subacquea avviene fuori dal
            campo visivo dell\textquotesingle allievo, che deve affidarsi
            all\textquotesingle{}\textbf{analizzatore cinestesico} (senso della
            posizione segmentaria nello spazio e nel tempo).
      \item
            \textbf{Tolleranza alle sbavature} nella fase di apprendimento, simile
            allo stile libero.
\end{itemize}

\subsubsection{Azione delle Braccia}\label{azione-delle-braccia-1}

Fasi identiche allo stile libero (appoggio e presa → trazione → spinta →
recupero), con alcune specificità:

\begin{itemize}
      \tightlist
      \item
            \textbf{Appoggio e presa}: la mano entra in acqua col \textbf{mignolo}
            (non con il pollice); l\textquotesingle ascella deve essere aperta.
      \item
            \textbf{Trazione}: mano, gomito e spalla si trovano sullo stesso
            piano; la mano cerca profondità.
      \item
            \textbf{Spinta}: verso il basso, in fortissima accelerazione; sfrutta
            il principio fisico opposto a quello di un palloncino spinto in
            profondità (la reazione dell\textquotesingle acqua proietta la mano
            verso l\textquotesingle alto, avviando il recupero).
      \item
            \textbf{Recupero}: inizia molto sott\textquotesingle acqua nella
            nuotata evoluta. Nella fase aerea il braccio può uscire col pollice,
            col dorso o col mignolo (non è determinante); ciò che conta è che
            \textbf{esca a braccio disteso e vicino al corpo}, con una rotazione
            verso l\textquotesingle esterno prima dell\textquotesingle ingresso in
            acqua (sempre col mignolo).
\end{itemize}

\begin{quote}
      \textbf{Note didattiche:}

      \begin{itemize}
            \tightlist
            \item
                  L\textquotesingle istruttore deve diventare gli "occhi"
                  dell\textquotesingle allievo, trasmettendo sensazioni semplici e
                  dirette.
            \item
                  Evitare di far uscire il mignolo deliberatamente (porta a una
                  contrazione errata della spalla).
            \item
                  Le mani devono rimanere vicine al corpo durante tutta la traiettoria:
                  mani lontane indicano compensazioni laterali per instabilità.
      \end{itemize}
\end{quote}

\subsubsection{Battuta di Gambe}\label{battuta-di-gambe-1}

\begin{itemize}
      \tightlist
      \item
            Come nello stile libero: fase ascendente \textbf{attiva} (propulsiva),
            fase discendente passiva (recupero con muscolatura decontratta).
      \item
            A causa del rollio, la gambata è sempre
            \textbf{diagonale/trasversale}, non puramente verticale.
\end{itemize}

\subsubsection{Respirazione}\label{respirazione-1}

\begin{itemize}
      \tightlist
      \item
            Avviene dalla bocca (più aria in meno tempo → migliore
            galleggiamento).
      \item
            Ritmica e collegata al ciclo di bracciata; nel principiante è
            spontanea, nell\textquotesingle atleta è inserita nel timing.
      \item
            Alcuni nuotatori evoluti usano il \textbf{tappanaso} per favorire il
            galleggiamento.
\end{itemize}

\begin{center}\rule{0.5\linewidth}{0.5pt}\end{center}

\subsection{La Delfinizzazione: Movimento Ondulatorio del
      Corpo}\label{la-delfinizzazione-movimento-ondulatorio-del-corpo}

Prima di affrontare rana e farfalla, è fondamentale capire la
\textbf{delfinizzazione}, cioè il movimento ondulatorio del tronco.

\begin{itemize}
      \tightlist
      \item
            Non è un movimento delle sole gambe: l\textquotesingle onda
            \textbf{parte dal torace} e viene trasmessa al resto del corpo.
      \item
            Coinvolge \textbf{tre masse critiche} in sequenza: testa → torace →
            gambe.
      \item
            Non deve essere ampia; lo schiacciamento verso il basso della parte
            centrale del tronco genera l\textquotesingle impulso che si propaga
            alle gambe.
      \item
            Può essere proposta in scuola nuoto \textbf{fin
                  dall\textquotesingle inizio}, perché non richiede elevate competenze
            coordinative (è vicina al galleggiamento e alle posture di base).
      \item
            È presente, in misura diversa, in \textbf{tutte e quattro le nuotate}
            (più marcata in rana e farfalla, meno in stile libero e dorso).
      \item
            Nella \textbf{farfalla moderna} il movimento ondulatorio è meno
            accentuato rispetto al passato; nella \textbf{rana} è diventato più
            marcato.
\end{itemize}

\begin{center}\rule{0.5\linewidth}{0.5pt}\end{center}

\subsection{Aspetti Comuni a Rana e Farfalla (Nuotate
      Simmetriche)}\label{aspetti-comuni-a-rana-e-farfalla-nuotate-simmetriche}

\begin{longtable}[]{@{}p{0.30\linewidth}p{0.30\linewidth}p{0.30\linewidth}@{}}
      \toprule\noalign{}
      Caratteristica               & Rana                           & Farfalla                \\
      \midrule\noalign{}
      \endhead
      \bottomrule\noalign{}
      \endlastfoot
      Tipo di movimento            & Simmetrico e simultaneo        & Simmetrico e simultaneo \\
      Intervento del busto         & Ondulazione con schiacciamento & Ondulazione
      (meno accentuata nel delfino moderno)                                                   \\
      Posizione appoggio/presa     & A Y (ascelle aperte)           & A Y (ascelle
      aperte)                                                                                 \\
      Difficoltà coordinativa      & Molto elevata                  & Molto elevata           \\
      Insegnamento in scuola nuoto & Seconda parte del percorso     & Seconda
      parte (o fine) del percorso                                                             \\
\end{longtable}

\begin{itemize}
      \tightlist
      \item
            La \textbf{farfalla deriva dalla rana}: in origine i ranisti
            recuperavano le braccia fuori dall\textquotesingle acqua, scoprendo
            che era più efficiente; la nuotata fu poi separata regolamentarmente.
            Per questo il termine ufficiale è \emph{farfalla} (non delfino), anche
            se nell\textquotesingle uso comune prevale "delfino".
      \item
            \textbf{Altissima trasferabilità} tra i due stili: esercizi pensati
            per uno possono essere usati nell\textquotesingle altro.
      \item
            Entrambe \textbf{non tollerano micro-errori coordinativi}: una
            sovrapposizione di movimenti (rana) o un timing errato (farfalla)
            compromettono drasticamente la nuotata.
      \item
            Entrambe presentano una continua \textbf{instabilità} dovuta al
            movimento simmetrico, che non garantisce una posizione uniforme in
            acqua.
\end{itemize}

\begin{center}\rule{0.5\linewidth}{0.5pt}\end{center}

\subsection{\texorpdfstring{Rana
            (\emph{Breaststroke})}{Rana (Breaststroke)}}\label{rana-breaststroke}

\subsubsection{Caratteristiche
      Generali}\label{caratteristiche-generali-2}

\begin{itemize}
      \tightlist
      \item
            Nuotata \textbf{simmetrica}: movimento simultaneo degli arti.
      \item
            \textbf{Basso rendimento meccanico}: non ha la fase di spinta nelle
            braccia (la spinta è nelle gambe).
      \item
            Posizione in acqua instabile durante il ciclo (non garantisce
            continuità idrodinamica favorevole).
      \item
            Insegnata nella \textbf{seconda parte della scuola nuoto} per
            l\textquotesingle elevata valenza coordinativa.
      \item
            Movimento delle gambe circolare e molto caratteristico; può
            interferire con i pattern del \emph{crawl} e del dorso se insegnato
            troppo presto.
      \item
            \textbf{Silenziosa}: produce pochi suoni rispetto alle altre nuotate.
      \item
            Le gambe hanno un contributo propulsivo pari o superiore a quello
            delle braccia (circa il \textbf{50\% ± 10\%}).
\end{itemize}

\begin{quote}
      \textbf{Importanza della rana per l\textquotesingle istruttore:}\\
      La rana è fondamentale perché:

      \begin{enumerate}
            \tightlist
            \item
                  È parte integrante della \textbf{gara dei misti} (200 e 400 m misti).
            \item
                  Serve come base per la \textbf{pallanuoto} (il movimento rana diventa
                  bicicletta).
            \item
                  È essenziale nel \textbf{nuoto artistico} e nel \textbf{salvamento}
                  (trasporto del ferito, autosostentamento in acque libere).
            \item
                  È un \textbf{dovere dell\textquotesingle istruttore} restituire
                  all\textquotesingle allievo una rana soddisfacente.
      \end{enumerate}

      In alcuni Paesi (Africa, Nord Europa) la rana viene insegnata per prima,
      proprio per la sua utilità nel sopravvivere in acqua libera e per la
      complessità coordinativa che richiede tempi più lunghi di apprendimento.
\end{quote}

\subsubsection{Azione delle Braccia}\label{azione-delle-braccia-2}

\begin{itemize}
      \tightlist
      \item
            Movimento \textbf{circolare} e \textbf{simultaneo}, sempre sotto il
            campo visivo dell\textquotesingle allievo.

            \begin{itemize}
                  \tightlist
                  \item
                        Se le mani escono dal campo visivo, il movimento è troppo largo.
                  \item
                        Iniziare da movimenti stretti e allargare progressivamente, mai il
                        contrario.
            \end{itemize}
\end{itemize}

Fasi:

\begin{enumerate}
      \tightlist
      \item
            \textbf{Appoggio e presa} (posizione a Y, ascelle aperte).
      \item
            \textbf{Trazione} (mano verso il fondo, gomito alto, ascella aperta).
      \item
            \sout{Spinta} --- \textbf{non presente} nella bracciata rana: la spinta
            è delegata alle gambe.
      \item
            \textbf{Recupero} subacqueo (o leggermente fuori superficie nei
            nuotatori evoluti).
\end{enumerate}

\subsubsection{Azione delle Gambe}\label{azione-delle-gambe}

La gambata rana è il nucleo tecnico più complesso della nuotata.

\textbf{Sequenza meccanica:}

\begin{enumerate}
      \tightlist
      \item
            Flessione della gamba sulla coscia e della coscia sul bacino (tallone
            verso il sedere); ginocchia strette, piedi vicini alla superficie.
      \item
            \textbf{Extra-rotazione dei piedi} (i piedi devono essere più larghi
            delle ginocchia --- mai il contrario).
      \item
            Frustata verso dietro-basso con \textbf{estensione del piede} e forte
            accelerazione.
\end{enumerate}

\begin{quote}
      \textbf{Errori comuni:}

      \begin{itemize}
            \tightlist
            \item
                  Piedi in intra-rotazione o con la pianta rivolta indietro.
            \item
                  Stessa velocità in tutta la gambata (non crea propulsione).
            \item
                  Movimento a ventaglio (fuori-dentro) invece che circolare.
      \end{itemize}
\end{quote}

\textbf{Manipolazione}: la rana è l\textquotesingle unica nuotata che
richiede manipolazione diretta dell\textquotesingle allievo (in acqua o
a secco, in posizione orizzontale o verticale) per far sentire la
corretta extra-rotazione del piede. La manipolazione \textbf{deve sempre
      avvenire con resistenza} (mani dell\textquotesingle istruttore che
oppongono forza): senza resistenza non si attiva
l\textquotesingle analizzatore cinestesico e non c\textquotesingle è
trasferimento in acqua.

\textbf{Timing della gambata:} lento (recupero) → \textbf{forte}
(spinta). Esercizio utile: accosciata lenta → balzo forte, per far
interiorizzare il ritmo \emph{piano-forte}.

\subsubsection{Coordinazione e Ritmica}\label{coordinazione-e-ritmica}

La sequenza della rana è:\\
\textbf{Braccia → Respirazione → Gambe → Scivolo}

\begin{itemize}
      \tightlist
      \item
            \textbf{Uno a uno}: una bracciata, una respirazione, una gambata.
      \item
            La spinta delle gambe avviene quando il corpo è in \textbf{posizione
                  streamline} (braccia distese in avanti): è il momento di massima
            efficienza idrodinamica.
      \item
            Proporre agli allievi diverse pause di scivolo (pausa 1, pausa 2,
            pausa 3) per personalizzare la ritmica e separare i movimenti,
            evitando la sovrapposizione tipica del principiante.
\end{itemize}

\subsubsection{Respirazione}\label{respirazione-2}

\begin{itemize}
      \tightlist
      \item
            \textbf{Frontale} (non laterale).
      \item
            La testa non si solleva verticalmente: rimane il più possibile in
            linea con la schiena (collo lungo).
      \item
            Il sollevamento è \textbf{avanti-alto}, non verticale.
\end{itemize}

\begin{center}\rule{0.5\linewidth}{0.5pt}\end{center}

\subsection{\texorpdfstring{Farfalla (\emph{Butterfly /
                  Delfino})}{Farfalla (Butterfly / Delfino)}}\label{farfalla-butterfly-delfino}

\subsubsection{Caratteristiche
      Generali}\label{caratteristiche-generali-3}

\begin{itemize}
      \tightlist
      \item
            Nuotata \textbf{simmetrica}, movimento simultaneo degli arti.
      \item
            La \textbf{coordinazione} è il fattore determinante, non la forza:
            bambini esili possono nuotarla meglio di adulti più forti ma meno
            coordinati.
      \item
            Dopo lo stile libero e il dorso, è la nuotata più veloce.
      \item
            Viene insegnata \textbf{alla fine} del percorso, perché richiede molte
            competenze motorie già consolidate e una disponibilità minima di forza
            per compensare le lacune coordinative nelle fasi iniziali di
            apprendimento.
      \item
            Il movimento ondulatorio del corpo (delfinizzazione) e la gambata
            possono essere impostati \textbf{prima}
            dell\textquotesingle insegnamento della nuotata completa.
      \item
            Nella \textbf{farfalla moderna} il movimento ondulatorio è meno
            marcato rispetto agli anni passati: la posizione è sempre più
            orizzontale, vicina allo stile libero.
      \item
            Errori di coordinazione o timing compromettono gravemente efficacia,
            efficienza e costo energetico.
\end{itemize}

\subsubsection{Battuta di Gambe}\label{battuta-di-gambe-2}

\begin{itemize}
      \tightlist
      \item
            \textbf{2 colpi di gambe per ogni ciclo di bracciata}:
\end{itemize}

\begin{longtable}[]{@{}p{0.30\linewidth}p{0.30\linewidth}p{0.30\linewidth}@{}}
      \toprule\noalign{}
      Colpo    & Momento nel ciclo                                      & Funzione prevalente \\
      \midrule\noalign{}
      \endhead
      \bottomrule\noalign{}
      \endlastfoot
      1° colpo & Braccia in avanti (posizione idrodinamica vantaggiosa) &
      \textbf{Propulsivo}                                                                     \\
      2° colpo & Braccia lunghe ai fianchi (durante la respirazione)    &
      \textbf{Galleggiante e stabilizzante}                                                   \\
\end{longtable}

\begin{itemize}
      \tightlist
      \item
            Il 2° colpo compensa l\textquotesingle abbassamento del bacino causato
            dal sollevamento della testa per la respirazione.
      \item
            Funzione della gambata: \textbf{galleggiante}, \textbf{stabilizzante}
            e \textbf{propulsiva}.
      \item
            Man mano che la tecnica evolve, il movimento diventa più piatto e meno
            ondulatorio.
\end{itemize}

\subsubsection{Azione delle Braccia}\label{azione-delle-braccia-3}

Come in tutte le nuotate, le fasi della passata subacquea sono:
appoggio/presa → trazione → spinta → recupero.

\begin{itemize}
      \tightlist
      \item
            \textbf{Continuità d\textquotesingle azione}: è il requisito assoluto.
            Qualsiasi pausa lungo la bracciata compromette irrimediabilmente la
            nuotata.
      \item
            \textbf{Ampiezza}: il braccio deve completare il percorso disteso
            davanti fino in fondo ai fianchi.
      \item
            \textbf{Traiettoria}: molto personale, ma deve restare entro una
            tolleranza; mani troppo lontane dal corpo = applicazione di forza
            inefficace.
      \item
            La \textbf{fase finale di spinta} deve essere completata con forte
            accelerazione per permettere il lancio a doppio braccio nel recupero.
            Una spinta lenta rende il recupero quasi impossibile.
\end{itemize}

\subsubsection{Respirazione}\label{respirazione-3}

\begin{itemize}
      \tightlist
      \item
            \textbf{Avanti-alto}, non verticale: la proiezione è in avanti, il
            busto non si solleva sull\textquotesingle acqua (al contrario di
            quanto si vedeva negli anni \textquotesingle80-\textquotesingle90).
      \item
            La farfalla moderna tende verso la \textbf{respirazione laterale}
            (come nello stile libero) in alcuni nuotatori d\textquotesingle élite,
            per ridurre al minimo il sollevamento del capo.
      \item
            Il principio chiave: \textbf{il corpo rimane incollato
                  all\textquotesingle acqua}; nessuna verticalizzazione.
\end{itemize}

\begin{quote}
      \textbf{Nota didattica:} se si riesce a costruire una \textbf{corretta
            posizione del corpo}, la coordinazione della farfalla diventa
      relativamente semplice. Se la posizione è sbagliata, ogni difficoltà si
      moltiplica.
\end{quote}

\begin{center}\rule{0.5\linewidth}{0.5pt}\end{center}

\subsection{Principi Trasversali alle Quattro
      Nuotate}\label{principi-trasversali-alle-quattro-nuotate}

Questi tre elementi strutturali sono comuni a tutti gli stili e
rappresentano la base su cui si costruisce qualsiasi nuotata:

\begin{enumerate}
      \tightlist
      \item
            \textbf{Continuità d\textquotesingle azione}: assenza di pause nel
            movimento. Non va confusa con la velocità
            d\textquotesingle esecuzione.
      \item
            \textbf{Ampiezza del movimento}: sfruttare l\textquotesingle intero
            range motorio disponibile.
      \item
            \textbf{Traiettoria}: le mani devono rimanere vicine al corpo lungo
            tutta la passata subacquea.
\end{enumerate}

Ulteriori principi comuni:

\begin{itemize}
      \tightlist
      \item
            La \textbf{fase finale di spinta} è sempre eseguita in
            \textbf{fortissima accelerazione} in tutte le nuotate.
      \item
            L\textquotesingle{}\textbf{ascella aperta} nella fase di appoggio e
            presa è presente in tutti e quattro gli stili.
      \item
            La \textbf{qualità di una fase determina la qualità della successiva}:
            le fasi non sono compartimenti stagni ma un continuum.
\end{itemize}

\begin{center}\rule{0.5\linewidth}{0.5pt}\end{center}

\subsection{Conclusione --- Punti Chiave per
      l\textquotesingle Esame}\label{conclusione-punti-chiave-per-lesame-8}

\begin{itemize}
      \tightlist
      \item
            \textbf{Stile libero e dorso} sono nuotate asimmetriche con movimento
            alternato, rollio, e tollerano micro-errori coordinativi. Lo stile
            libero prevede 6 colpi di gambe per ciclo.
      \item
            \textbf{Rana e farfalla} sono simmetriche, derivano
            l\textquotesingle una dall\textquotesingle altra, hanno elevatissima
            valenza coordinativa e non tollerano errori di timing.
      \item
            La \textbf{delfinizzazione} (onda dal torace alle gambe) è alla base
            di rana e farfalla; può essere insegnata fin
            dall\textquotesingle inizio della scuola nuoto.
      \item
            Nella \textbf{rana} la spinta è nelle gambe (non nelle braccia); la
            gambata è circolare con piedi extra-ruotati; la ritmica è braccia →
            respirazione → gambe → scivolo.
      \item
            Nella \textbf{farfalla} ci sono 2 colpi di gambe per ciclo (1°
            propulsivo con braccia avanti, 2° stabilizzante con braccia ai
            fianchi); la continuità è il requisito imprescindibile.
      \item
            Il modello prestativo è un \textbf{riferimento universale}, non una
            coreografia da imporre: va usato per leggere il gesto
            dell\textquotesingle allievo e guidarlo verso la forma più efficace
            per lui.
      \item
            Stile libero e dorso si insegnano dall\textquotesingle inizio; rana e
            farfalla nella seconda parte (o fine) del percorso di scuola nuoto.
      \item
            La \textbf{rana} deve essere garantita a tutti gli allievi per i suoi
            sbocchi verso pallanuoto, nuoto artistico, salvamento e gara dei
            misti.
\end{itemize}

\section{14. La Pallanuoto in Corsia --- Parte
  I}\label{la-pallanuoto-in-corsia-parte-i}

Docente: Massimo Tafuro, allenatore di pallanuoto con esperienza diretta
nella gestione e nell\textquotesingle organizzazione di attività di
scuola nuoto.

\begin{center}\rule{0.5\linewidth}{0.5pt}\end{center}

\subsection{Introduzione}\label{introduzione-13}

La \textbf{pallanuoto in corsia} può sembrare un ossimoro, poiché la
pallanuoto si pratica in un campo di 30×20 m (600 m²), ben lontano dagli
spazi di una singola corsia. Eppure, proprio questo formato ridotto
rappresenta uno strumento didattico e organizzativo di grande valore per
qualsiasi istruttore che lavori in una scuola nuoto.

Questa lezione affronta il tema da due angolazioni complementari: la
\textbf{teoria dell\textquotesingle apprendimento motorio acquatico}
(che fornisce le basi concettuali) e la \textbf{strutturazione pratica}
di un settore preagonistico di pallanuoto all\textquotesingle interno di
una piscina convenzionale.

\begin{center}\rule{0.5\linewidth}{0.5pt}\end{center}

\subsection{Obiettivi della Pallanuoto in
      Corsia}\label{obiettivi-della-pallanuoto-in-corsia}

Praticare la pallanuoto anche in uno spazio ristretto come una corsia
serve a raggiungere più obiettivi contemporaneamente:

\begin{enumerate}
      \item
            \textbf{Migliorare e completare la strutturazione degli schemi motori
                  acquatici} --- arricchire il bagaglio motorio
            dell\textquotesingle allievo con stimolazioni che il solo nuoto di
            corsia non offre.
      \item
            \textbf{Predisporre la costruzione di un settore amatoriale o
                  agonistico} finalizzato alla pallanuoto --- avviare i bambini che ne
            hanno le capacità verso questo sport in modo graduale e coerente.
      \item
            \textbf{Offrire maggiore varietà di attività motorie acquatiche} ---
            contrastare il fenomeno dell\textquotesingle abbandono precoce. Molti
            bambini lasciano la piscina non appena apprendono le tre nuotate di
            base (dorso, stile libero, rana). Diversificare
            l\textquotesingle offerta motoria aumenta le possibilità di trattenere
            gli allievi nel percorso acquatico.
      \item
            \textbf{Ottimizzare l\textquotesingle organizzazione degli spazi
                  acqua} --- i preadolescenti sono disponibili ad allenarsi in orari
            pomeridiani (dopo la scuola) che i bambini più piccoli frequentano
            meno. Riempire queste fasce orarie riduce i costi di gestione
            dell\textquotesingle acqua (riscaldamento, clorazione, disinfezione).
      \item
            \textbf{Potenziare i fondamentali di gioco} --- anche per atleti di
            livello più avanzato, lavorare in corsia sui fondamentali tecnici
            (specialmente senza palla) è utile quando non si dispone del campo di
            gioco completo.
\end{enumerate}

\begin{center}\rule{0.5\linewidth}{0.5pt}\end{center}

\subsection{Fondamenti Teorici: Schemi
      Motori}\label{fondamenti-teorici-schemi-motori}

\subsubsection{Schema Corporeo e Schema
      Motorio}\label{schema-corporeo-e-schema-motorio}

Alla base di qualsiasi apprendimento motorio ci sono due concetti
fondamentali:

\begin{itemize}
      \item
            \textbf{Schema corporeo}: \emph{l\textquotesingle immagine del sé
                  nello spazio}. È la rappresentazione precisa e aggiornata che il
            sistema nervoso centrale ha della propria corporeità in ogni istante.
            Senza questa "fotografia interna" non è possibile strutturare alcun
            movimento coordinato.
      \item
            \textbf{Schema motorio}: lo schema del movimento. Non è solo la
            sequenza immaginata di azioni (come il tuffatore che ripete
            mentalmente il gesto prima di eseguirlo), ma qualcosa di molto più
            complesso: include la mappa precisa dei \textbf{muscoli da contrarre}
            e dei \textbf{muscoli antagonisti da rilassare} per realizzare ogni
            singolo gesto. Anche la postura eretta --- apparentemente banale ---
            richiede l\textquotesingle integrazione continua di informazioni
            provenienti dai pressocettori plantari, dall\textquotesingle apparato
            vestibolare e dai recettori articolari e muscolari.
\end{itemize}

\subsubsection{Strutturazione degli Schemi Motori: dal Generale
      all\textquotesingle Acquatico}\label{strutturazione-degli-schemi-motori-dal-generale-allacquatico}

Gli schemi motori si costruiscono progressivamente, a partire
dall\textquotesingle epoca intrauterina, accelerando dopo la nascita
grazie al \textbf{processo imitativo}. Sulla terraferma gli schemi
motori di base sono:

\begin{itemize}
      \tightlist
      \item
            camminare, correre, saltare, arrampicare
      \item
            lanciare, afferrare
\end{itemize}

Su questi schemi di base si costruiscono gli \textbf{schemi motori
      speciali} (es. il tiro a canestro = saltare + lanciare, in forma evoluta
e complessa).

Il medesimo percorso vale per l\textquotesingle ambiente acquatico. Gli
\textbf{schemi motori acquatici di base} sono:

\begin{itemize}
      \tightlist
      \item
            galleggiamenti
      \item
            apnee
      \item
            educazione respiratoria
      \item
            autonomia in acqua
      \item
            prime forme di propulsione
\end{itemize}

Su questi si costruiscono gli schemi motori acquatici speciali, tra cui:
le quattro nuotate regolamentari, la gambata alternata del
pallanuotista, la verticalizzazione del tuffatore, i movimenti del nuoto
artistico.

\begin{quote}
      \textbf{Problema specifico dell\textquotesingle ambiente acquatico}: a
      differenza delle competenze terrestri che si acquisiscono spontaneamente
      nei primi anni di vita per imitazione, le competenze acquatiche devono
      essere costruite ex novo in maniera sistematica. Questo è il motivo per
      cui un giovane pallanuotista deve allenarsi 3-5 volte a settimana
      (progressivamente fino a tutti i giorni): deve colmare il \emph{gap}
      rispetto all\textquotesingle enorme bagaglio motorio terrestre
      accumulato spontaneamente dal bambino nei primi anni di vita.
\end{quote}

\subsubsection{Coordinazione:
      Definizione}\label{coordinazione-definizione}

\emph{La coordinazione è un\textquotesingle esatta esecuzione di una
      esatta immagine motoria.}

Questa definizione --- attribuita al prof. Calabrese
dell\textquotesingle ISEF di Roma --- chiarisce che coordinazione non
significa semplicemente eseguire un gesto, ma che:

\begin{enumerate}
      \tightlist
      \item
            bisogna costruire prima un\textquotesingle{}\textbf{immagine motoria
                  esatta} (rappresentazione interna precisa del movimento);
      \item
            poi bisogna essere capaci di \textbf{eseguirla esattamente}
            (trasferire l\textquotesingle immagine in azione motoria).
\end{enumerate}

Le due cose non coincidono: si può avere un\textquotesingle immagine
corretta ma non riuscire ancora a realizzarla.

\subsubsection{Quantità e Qualità delle Esperienze
      Motorie}\label{quantituxe0-e-qualituxe0-delle-esperienze-motorie}

La qualità dell\textquotesingle immagine motoria dipende dalla
\textbf{quantità} e dalla \textbf{qualità} delle esperienze motorie
accumulate.

\begin{itemize}
      \tightlist
      \item
            \textbf{Quantità}: numero di ripetizioni del gesto (quante volte si
            esegue).
      \item
            \textbf{Qualità}: diversificazione dell\textquotesingle esecuzione,
            che introduce difficoltà crescenti per stimolare il miglioramento
            coordinativo (es. modulare la velocità di esecuzione, cambiare un
            parametro alla volta).
\end{itemize}

\begin{quote}
      "Non esistono bambini scoordinati, ma esistono soltanto bambini che non
      hanno maturato la giusta quantità e qualità di esperienze motorie."
      (prof. Calabrese)
\end{quote}

Più le esperienze motorie sono \textbf{precoci}, più facilmente il
bambino potrà usarle come base per costruire gesti sempre più complessi.
Questo è un argomento ulteriore per portare i bambini in acqua presto,
non solo per insegnare le nuotate, ma per arricchire il bagaglio motorio
globale.

\subsubsection{Il Ruolo delle Capacità
      Condizionali}\label{il-ruolo-delle-capacituxe0-condizionali}

Accanto alle capacità coordinative, anche le \textbf{capacità
      condizionali} (forza, velocità, resistenza) contribuiscono alla
coordinazione. La forza in particolare fa parte della crescita normale
del bambino fin dai primissimi mesi di vita (il lattante solleva la
testa, poi si siede, poi cammina grazie a un progressivo aumento di
forza). Non va demonizzata: va allenata in modo appropriato
all\textquotesingle età, senza sovraccarichi eccessivi.

\subsubsection{Benefici Trasversali della Coordinazione
      Acquatica}\label{benefici-trasversali-della-coordinazione-acquatica}

Una buona coordinazione spaziale e spazio-temporale
nell\textquotesingle acqua produce ricadute positive anche sugli
apprendimenti scolastici:

\begin{itemize}
      \tightlist
      \item
            La \textbf{coordinazione spaziale} è alla base della capacità di
            lettura e scrittura.
      \item
            La letteratura scientifica indica che la stimolazione motoria precoce
            e diversificata può \textbf{prevenire la disgrafia} e ridurre i
            disturbi specifici dell\textquotesingle apprendimento (\emph{DSA}),
            oggi più diffusi anche a causa della diminuzione
            dell\textquotesingle attività motoria spontanea dei bambini.
\end{itemize}

\begin{center}\rule{0.5\linewidth}{0.5pt}\end{center}

\subsection{Struttura del Settore
      Preagonistico}\label{struttura-del-settore-preagonistico}

\subsubsection{Organizzazione Generale}\label{organizzazione-generale}

\begin{longtable}[]{@{}p{0.47\linewidth}p{0.47\linewidth}@{}}
      \toprule\noalign{}
      Parametro         & Dettaglio                                                       \\
      \midrule\noalign{}
      \endhead
      \bottomrule\noalign{}
      \endlastfoot
      Età degli allievi & Nati dal 2011 in poi (8-10 anni circa)                          \\
      Gruppi            & 2 gruppi da 21 ragazzi ciascuno (max 7 per corsia, con 3 corsie
      disponibili)                                                                        \\
      Frequenza         & 3 volte a settimana per gruppo                                  \\
      Sessione tipo     & 45 min a secco (17:45-18:30) + 60 min in acqua
      (18:30-19:30)                                                                       \\
      Piscina           & 25 m × 12,50 m --- profondità min. 1,40 m / max. 3 m            \\
      Allenatore        & Scelto tra gli allenatori del settore agonistico
      (continuità relazionale con i ragazzi)                                              \\
\end{longtable}

\subsubsection{Attività a Secco (45
      minuti)}\label{attivituxe0-a-secco-45-minuti}

L\textquotesingle attività a secco ha come obiettivo
l\textquotesingle ampliamento del bagaglio motorio generale:

\begin{itemize}
      \tightlist
      \item
            Coordinazione spaziale e spazio-temporale
      \item
            Irrobustimento generale
      \item
            Andature ginnico-sportive
      \item
            Utilizzo di attrezzi (palloni) con spunti tratti da altri sport
            (pallavolo, basket)
\end{itemize}

\subsubsection{Attività in Acqua (60
      minuti)}\label{attivituxe0-in-acqua-60-minuti}

L\textquotesingle attività in acqua si concentra
sull\textquotesingle insegnamento dei \textbf{fondamentali della
      pallanuoto}, senza palla e con palla, in corsia.

\subsubsection{Momenti di Gioco in
      Campo}\label{momenti-di-gioco-in-campo}

Due domeniche al mese si radunano i gruppi per attività con il campo
montato (non in corsia): i ragazzi capiscono la finalità del lavoro in
corsia vedendola applicata nel contesto di gioco reale.

Partecipazione anche a \textbf{tornei e campionati federali} riservati
alla categoria \textbf{Under 11}.

\begin{center}\rule{0.5\linewidth}{0.5pt}\end{center}

\subsection{Prerequisiti per l\textquotesingle Accesso al Settore
      Preagonistico}\label{prerequisiti-per-laccesso-al-settore-preagonistico}

Prima di essere inserito nel gruppo preagonistico, il bambino deve
essere in grado di:

\begin{itemize}
      \tightlist
      \item
            Nuotare correttamente a \textbf{stile libero}
      \item
            Nuotare correttamente a \textbf{rana} con un buon colpo di gambe
      \item
            Nuotare correttamente a \textbf{farfalla} (recupero delle braccia
            fuori dall\textquotesingle acqua)
\end{itemize}

Questi tre prerequisiti sono imprescindibili. A partire da queste basi
si insegna il \textbf{movimento di gambe alternate} (\emph{eggbeater
      kick} / gambe bicicletta).

\begin{center}\rule{0.5\linewidth}{0.5pt}\end{center}

\subsection{I Fondamentali della Pallanuoto in
      Corsia}\label{i-fondamentali-della-pallanuoto-in-corsia}

\subsubsection{\texorpdfstring{Fondamentali Senza Palla
            (\emph{Without-Ball
                  Fundamentals})}{Fondamentali Senza Palla (Without-Ball Fundamentals)}}\label{fondamentali-senza-palla-without-ball-fundamentals}

Sono le azioni motorie che il giocatore compie indipendentemente dal
possesso del pallone, ma che sono alla base della partecipazione al
gioco.

\paragraph{Nuotate Orizzontali}\label{nuotate-orizzontali}

\begin{itemize}
      \tightlist
      \item
            \textbf{Stile libero testa alta} (\emph{front crawl head-up}):
            mantiene il capo fuori dall\textquotesingle acqua per osservare il
            gioco; il gomito alto è un elemento tecnico rilevante.
      \item
            \textbf{Dorso testa alta} (\emph{backstroke head-up}): permette di
            arretrare guardando in avanti.
      \item
            Il \textbf{druggio} (stile libero con traiettoria obliqua per eludere
            l\textquotesingle avversario).
\end{itemize}

\begin{quote}
      Nota: la pallanuoto si gioca in posizione orizzontale con gli occhi a
      circa 20 cm sopra il livello dell\textquotesingle acqua. La difficoltà
      di lettura del gioco da questa prospettiva è enorme, e la capacità di
      mantenere la testa alta durante la nuotata è fondamentale.
\end{quote}

\paragraph{Arresti, Avvitamenti e Cambi di
      Direzione}\label{arresti-avvitamenti-e-cambi-di-direzione}

\begin{itemize}
      \tightlist
      \item
            \textbf{Cambio di direzione}: non implica necessariamente un cambio di
            fronte (si può avanzare nella stessa direzione con la fronte rivolta
            verso il fondo o verso il soffitto).
      \item
            \textbf{Cambio di fronte}: indica una rotazione del corpo,
            indipendente dalla direzione di spostamento.
      \item
            \textbf{Avvitamenti sul dorso}: movimento classico di
            \emph{dribbling}, con cui si invita l\textquotesingle avversario ad
            attaccare su un lato per poi ripartire sul suo lato debole.
\end{itemize}

I cambi di direzione con o senza cambio di fronte, uniti alle
transizioni tra galleggiamento verticale e orizzontale, sono esercizi ad
altissimo valore coordinativo e spaziale.

\paragraph{\texorpdfstring{Nuotate Verticali (\emph{Vertical
                  Swimming})}{Nuotate Verticali (Vertical Swimming)}}\label{nuotate-verticali-vertical-swimming}

\begin{itemize}
      \tightlist
      \item
            \textbf{Gambe a bicicletta} (\emph{eggbeater kick}): alternanza delle
            gambe in posizione verticale.
      \item
            \textbf{Braccia a pressione}: azione delle braccia che mantiene il
            corpo in posizione alta sull\textquotesingle acqua.
      \item
            \textbf{Verticalità}: mantenimento della posizione verticale,
            prerequisito per il tiro e il passaggio.
      \item
            \textbf{Salti} e \textbf{scivolamenti}.
\end{itemize}

\begin{quote}
      La pallanuoto alterna circa il \textbf{50\% del tempo di gioco} in
      posizione orizzontale (spostamento) e il \textbf{50\%} in posizione
      verticale (costruzione dell\textquotesingle azione, marcatura, tiro).
\end{quote}

\subsubsection{\texorpdfstring{Fondamentali Con Palla (\emph{With-Ball
                  Fundamentals})}{Fondamentali Con Palla (With-Ball Fundamentals)}}\label{fondamentali-con-palla-with-ball-fundamentals}

\begin{itemize}
      \tightlist
      \item
            \textbf{Presa della palla e tenuta in verticale}: il bambino deve
            saper impugnare il pallone, sollevarlo dall\textquotesingle acqua e
            contestualmente trovare la posizione verticale di tiro/passaggio.
      \item
            \textbf{Impostazione del palleggio}: prerequisito per il passaggio; si
            esegue dalla posizione orizzontale fermandosi, girandosi e
            palleggiando con un compagno.
      \item
            \textbf{Impostazione del passaggio}
      \item
            \textbf{Impostazione dei vari tipi di tiro}: tiro in verticale, tiro
            in orizzontale (\emph{tiro a schizzo}, \emph{tiro colonnello}), tiro
            con palla impugnata.
\end{itemize}

\begin{quote}
      \textbf{Nota tecnica sul trasporto della palla}: la palla non deve
      essere toccata dall\textquotesingle avambraccio durante la nuotata, ma
      deve essere spinta dall\textquotesingle{}\textbf{onda creata dal viso}
      dell\textquotesingle allievo. Il contatto con
      l\textquotesingle avambraccio fa perdere il controllo
      dell\textquotesingle attrezzo.
\end{quote}

\subsubsection{\texorpdfstring{Comportamenti Tattici (\emph{Tactical
                  Behaviors})}{Comportamenti Tattici (Tactical Behaviors)}}\label{comportamenti-tattici-tactical-behaviors}

Man mano che i fondamentali vengono assimilati, si introducono i
comportamenti tattici individuali e di squadra:

\begin{itemize}
      \tightlist
      \item
            \textbf{Difesa dello spazio} (es. impedire
            all\textquotesingle avversario di entrare nella propria zona,
            segnalata da coni sul bordo vasca)
      \item
            \textbf{Marcatura dell\textquotesingle avversario}
      \item
            \textbf{Smarcarsi} dalla marcatura
      \item
            \textbf{Protezione della palla} dall\textquotesingle attacco
            dell\textquotesingle avversario
      \item
            \textbf{Scavalcamento dell\textquotesingle avversario} per ripartire
            sul suo lato debole
      \item
            \textbf{Giochi derivati} da altri sport (es. esercizio mutuato dal
            rugby: portare la palla passandosela e appoggiarla sul bordo vasca
            avversario senza tirarla)
      \item
            \textbf{Regola dei 3 secondi}: il giocatore ha 3 secondi per passare
            la palla; il compagno deve smarcarsi in quel lasso di tempo.
\end{itemize}

\begin{quote}
      Per tattica non si intende solo la difesa a zona o il \emph{pressing},
      ma più in generale la \textbf{capacità del giocatore di utilizzare la
            tecnica per raggiungere un obiettivo}: prima individuale, poi di
      squadra.
\end{quote}

\begin{center}\rule{0.5\linewidth}{0.5pt}\end{center}

\subsection{Progressione Didattica e Organizzazione
      Pratica}\label{progressione-didattica-e-organizzazione-pratica}

\subsubsection{Apprendimento del Movimento di Gambe Alternate
      (Bicicletta)}\label{apprendimento-del-movimento-di-gambe-alternate-bicicletta}

\begin{enumerate}
      \tightlist
      \item
            Esercizio addossato alla parete (apprendimento in condizioni
            controllate).
      \item
            Utilizzo di \textbf{braccioli} come ausilio al sostentamento, per
            permettere al bambino di concentrarsi sul movimento delle gambe senza
            preoccuparsi del galleggiamento.
      \item
            Combinazione di \textbf{braccia a pressione} + \textbf{gambe
                  bicicletta}: produce la postura tipica del pallanuotista in verticale.
\end{enumerate}

Gli errori iniziali (\emph{sforbiciare}, asincronie) fanno parte
dell\textquotesingle apprendimento e non devono essere motivo di
preoccupazione, a condizione che il ragazzo capisca a che cosa serve il
gesto.

\subsubsection{Lavoro in Parallelo con la Scuola
      Nuoto}\label{lavoro-in-parallelo-con-la-scuola-nuoto}

Il settore preagonistico può coesistere nelle stesse corsie con la
scuola nuoto tradizionale: nelle corsie adiacenti
l\textquotesingle istruttore insegna le nuotate regolamentari, mentre in
un\textquotesingle altra corsia si lavora sui fondamentali di
pallanuoto. Non c\textquotesingle è incompatibilità.

\subsubsection{Raduni Estivi}\label{raduni-estivi}

Nel periodo estivo i due gruppi vengono riuniti, disponendo di più
spazio. Si completa e si amplia il bagaglio tecnico con:

\begin{itemize}
      \tightlist
      \item
            Finte e movimenti di protezione della palla
      \item
            Partite 3 vs 3 in spazi ridotti (più campetti paralleli = molti
            ragazzi allenati contemporaneamente)
      \item
            Tiro con porta (anche autocostruita o gonfiabile)
      \item
            Giochi di ruolo (portiere, attaccante, difensore)
\end{itemize}

\begin{center}\rule{0.5\linewidth}{0.5pt}\end{center}

\subsection{Galleggiamento Verticale: Elemento Trasversale
      Prioritario}\label{galleggiamento-verticale-elemento-trasversale-prioritario}

Il \textbf{galleggiamento verticale} è un apprendimento che, secondo il
docente, viene stimolato meno di quanto dovrebbe nella scuola nuoto
tradizionale, pur essendo fondamentale per:

\begin{itemize}
      \tightlist
      \item
            Il gioco della pallanuoto (50\% del tempo in verticale)
      \item
            La trasversalità verso altri sport acquatici (tuffi, nuoto artistico,
            salvamento)
      \item
            La padronanza generale del corpo nell\textquotesingle acqua
\end{itemize}

Da praticare insieme ai \textbf{passaggi tra i vari tipi di
      galleggiamento}:

\begin{itemize}
      \tightlist
      \item
            Verticale $\leftrightarrow$ orizzontale ventrale
      \item
            Verticale $\leftrightarrow$ orizzontale dorsale
      \item
            Con avvitamenti e cambi di fronte
\end{itemize}

\begin{center}\rule{0.5\linewidth}{0.5pt}\end{center}

\subsection{Conclusione --- Punti Chiave per
      l\textquotesingle Esame}\label{conclusione-punti-chiave-per-lesame-9}

\begin{itemize}
      \tightlist
      \item
            La \textbf{pallanuoto in corsia} serve a diversificare le stimolazioni
            motorie, ridurre l\textquotesingle abbandono della scuola nuoto,
            ottimizzare gli spazi acqua e avviare i ragazzi alla pallanuoto.
      \item
            \textbf{Schema corporeo} = immagine del sé nello spazio;
            \textbf{schema motorio} = mappa precisa di movimenti, contrazioni e
            rilassamenti muscolari necessari per eseguire un gesto.
      \item
            \textbf{Coordinazione} = \emph{esatta esecuzione di una esatta
                  immagine motoria}. Si migliora aumentando la \textbf{quantità}
            (ripetizioni) e la \textbf{qualità} (diversificazione e difficoltà
            crescente) delle esperienze motorie.
      \item
            Le esperienze motorie acquatiche non vengono acquisite spontaneamente
            come quelle terrestri: devono essere costruite in modo sistematico, il
            che spiega la necessità di un addestramento frequente per i giovani
            pallanuotisti.
      \item
            Il settore preagonistico si rivolge a bambini di \textbf{8-10 anni},
            con \textbf{3 sedute settimanali} (45\textquotesingle{} a secco +
            60\textquotesingle{} in acqua); i prerequisiti sono: stile libero,
            rana e farfalla corretti.
      \item
            \textbf{Fondamentali senza palla}: nuotate orizzontali testa alta
            (\emph{front crawl head-up}, \emph{backstroke head-up}), gambe
            bicicletta, braccia a pressione, verticalità, arresti, avvitamenti,
            cambi di direzione.
      \item
            \textbf{Fondamentali con palla}: presa e tenuta in verticale,
            palleggio, passaggio, tiro (verticale, orizzontale, \emph{schizzo},
            \emph{colonnello}).
      \item
            La palla non si tocca con l\textquotesingle avambraccio: va spinta
            dall\textquotesingle{}\textbf{onda del viso}.
      \item
            Il \textbf{galleggiamento verticale} va stimolato molto di più
            rispetto alla prassi consueta della scuola nuoto.
      \item
            La tecnica individuale si trasforma progressivamente in
            \textbf{comportamento tattico} (individuale e poi di squadra): questo
            è il cuore dello sport della pallanuoto.
\end{itemize}

\section{15. La Pallanuoto in Corsia --- Parte
  II}\label{la-pallanuoto-in-corsia-parte-ii}

Docente: Massimo Tafuro, allenatore di pallanuoto.

\begin{center}\rule{0.5\linewidth}{0.5pt}\end{center}

\subsection{Introduzione}\label{introduzione-14}

Questa seconda parte riprende e amplia quanto trattato nella Parte I,
spostando il focus dalla scuola nuoto e dalla preagonistica verso i
livelli successivi: il \textbf{settore amatoriale} e il \textbf{settore
      agonistico}. Il filo conduttore rimane lo stesso: la pallanuoto in
corsia come strumento didattico e di allenamento a tutti i livelli, dal
principiante all\textquotesingle atleta agonista.

La lezione affronta anche i contenuti tecnici avanzati
dell\textquotesingle allenamento in corsia per atleti di livello
superiore, con particolare attenzione al passaggio dalla tecnica
individuale al \textbf{comportamento tattico}.

\begin{center}\rule{0.5\linewidth}{0.5pt}\end{center}

\subsection{Riepilogo della Parte I}\label{riepilogo-della-parte-i}

Prima di procedere, il docente richiama i concetti della Parte I:

\begin{itemize}
      \tightlist
      \item
            Importanza della pallanuoto in corsia per la strutturazione degli
            schemi motori acquatici.
      \item
            Lavoro sulla verticalità, galleggiamento, cambi di direzione e di
            fronte.
      \item
            Settore preagonistico: bambini di 8-10 anni avviati ai fondamentali
            senza palla (nuotate orizzontali e verticali) e con palla.
      \item
            I fondamentali tecnici diventano progressivamente
            \textbf{comportamenti tattici}: marcare l\textquotesingle avversario,
            favorire la circolazione della palla, realizzare un goal.
\end{itemize}

\begin{center}\rule{0.5\linewidth}{0.5pt}\end{center}

\subsection{Dal Settore Preagonistico ai Livelli
      Successivi}\label{dal-settore-preagonistico-ai-livelli-successivi}

\subsubsection{Contenuti da Ampliare}\label{contenuti-da-ampliare}

Il percorso di preagonistica non è sufficiente da solo per costruire una
squadra. I fondamentali già acquisiti devono essere ampliati con:

\begin{itemize}
      \tightlist
      \item
            \textbf{Cambi di assetto} e cambi di fronte più complessi
      \item
            Utilizzo della palla in modo sempre più specializzato: tiro,
            circolazione sul perimetro
      \item
            \textbf{Attacco a uomini pari}
      \item
            Costruzione della difesa
      \item
            Tutto ciò che attiene alla \textbf{squadra} come sistema collettivo
\end{itemize}

\subsubsection{Differenza tra Settore Amatoriale e Settore
      Agonistico}\label{differenza-tra-settore-amatoriale-e-settore-agonistico}

In partenza la differenza è quasi nulla: entrambe le categorie
provengono dalla preagonistica. La distinzione emerge in seguito, in
base alla disponibilità e alla volontà di impegno dei ragazzi.

\begin{longtable}[]{@{}p{0.30\linewidth}p{0.30\linewidth}p{0.30\linewidth}@{}}
      \toprule\noalign{}
      Aspetto                    & Pallanuoto Amatoriale                        & Pallanuoto Agonistica \\
      \midrule\noalign{}
      \endhead
      \bottomrule\noalign{}
      \endlastfoot
      Provenienza                & Preagonistica                                & Preagonistica         \\
      Frequenza allenamenti      & 2 volte/settimana + eventuali domeniche      & Da 4
      volte/settimana fino a giornaliero                                                                \\
      Competizioni               & Tornei amatoriali, attività FIN amatoriale   & Campionati
      federali regionali e nazionali                                                                    \\
      Impegno                    & Moderato --- compatibile con altri interessi & Molto elevato e
      costante                                                                                          \\
      Numero di atleti (esempio) & Circa 30 ragazzi                             & Circa 60 ragazzi (3
      sottogruppi)                                                                                      \\
\end{longtable}

\begin{center}\rule{0.5\linewidth}{0.5pt}\end{center}

\subsection{Struttura del Settore
      Amatoriale}\label{struttura-del-settore-amatoriale}

\begin{itemize}
      \tightlist
      \item
            Circa \textbf{30 ragazzi} seguiti da \textbf{un solo allenatore} (con
            supporto occasionale in occasione dei tornei).
      \item
            Vasca da \textbf{25 m × 12,50 m}.
      \item
            \textbf{Allenamenti}: mercoledì sera (20:15-21:30) e sabato
            (13:30-14:30), per un totale di circa \textbf{2 ore e 30 minuti a
                  settimana}.
      \item
            \textbf{Domenica mattina}: possibilità di attività aggiuntiva con
            campo montato, in abbinamento con la preagonistica (tornei
            amatoriali).
      \item
            Partecipazione a \textbf{tornei amatoriali} organizzati da club del
            territorio o dalla Federazione Italiana Nuoto.
      \item
            Il settore amatoriale risponde all\textquotesingle esigenza di quei
            ragazzi che vogliono continuare a praticare la pallanuoto senza
            l\textquotesingle impegno giornaliero richiesto
            dall\textquotesingle agonismo.
\end{itemize}

\begin{center}\rule{0.5\linewidth}{0.5pt}\end{center}

\subsection{Struttura del Settore
      Agonistico}\label{struttura-del-settore-agonistico}

\subsubsection{Organizzazione Generale}\label{organizzazione-generale-1}

Il settore agonistico parte dall\textquotesingle ingresso in
\textbf{prima media} (circa 11 anni), momento in cui gli orari
scolastici diventano più compatibili con allenamenti pomeridiani.

\begin{longtable}[]{@{}p{0.47\linewidth}p{0.47\linewidth}@{}}
      \toprule\noalign{}
      Parametro               & Dettaglio                        \\
      \midrule\noalign{}
      \endhead
      \bottomrule\noalign{}
      \endlastfoot
      Età degli atleti        & 11-16 anni                       \\
      Numero totale           & Circa 60 atleti                  \\
      Sottogruppi             & 3 (Under 12, Under 14, Under 16) \\
      Frequenza iniziale      & 4 volte/settimana                \\
      Frequenza massima (U16) & Allenamento giornaliero          \\
\end{longtable}

\subsubsection{Categorie e Livello
      Competitivo}\label{categorie-e-livello-competitivo}

\begin{longtable}[]{@{}p{0.47\linewidth}p{0.47\linewidth}@{}}
      \toprule\noalign{}
      Categoria        & Livello campionato                        \\
      \midrule\noalign{}
      \endhead
      \bottomrule\noalign{}
      \endlastfoot
      Under 12         & Regionale                                 \\
      Under 14         & Regionale                                 \\
      Under 16         & Nazionale                                 \\
      Under 16 (extra) & Campionato Under 18 (solo con atleti U16) \\
\end{longtable}

\begin{quote}
      Gli atleti che superano i 16 anni vengono \textbf{indirizzati ad altre
            società}, perché la struttura descritta (una sola piscina) non dispone
      delle risorse per aprire categorie senior. Questa scelta riflette una
      filosofia di esclusivo settore giovanile, coerente con le disponibilità
      di spazio.
\end{quote}

\subsubsection{Orari e Spazi}\label{orari-e-spazi}

\begin{itemize}
      \tightlist
      \item
            \textbf{Pomeriggio}: 14:00-16:15 --- tutta la piscina a disposizione;
            16:15-16:45 --- 3 corsie disponibili.
      \item
            \textbf{Mercoledì sera}: 21:30-22:45 --- spazio aggiuntivo per coprire
            le esigenze di tutti i gruppi.
      \item
            \textbf{Attività a secco} (a seguire l\textquotesingle allenamento
            pomeridiano, 16:45-17:45):

            \begin{itemize}
                  \tightlist
                  \item
                        3 volte/settimana per il gruppo più avanzato (U16)
                  \item
                        2 volte/settimana per il gruppo intermedio
                  \item
                        1 volta/settimana per i nuovi arrivati (primo anno agonistico)
            \end{itemize}
      \item
            \textbf{Spazio esterno}: 200 m² attrezzati con bacchette, coni,
            ostacoli, bilancieri, manubri, palloni medicinali ed elastici per la
            preparazione fisica generale e di forza.
\end{itemize}

\begin{center}\rule{0.5\linewidth}{0.5pt}\end{center}

\subsection{Contenuti Tecnici Avanzati in
      Corsia}\label{contenuti-tecnici-avanzati-in-corsia}

\subsubsection{Il Principio Fondamentale: Pensare e Agire
      Velocemente}\label{il-principio-fondamentale-pensare-e-agire-velocemente}

Un atleta di alto livello deve possedere due qualità inscindibili:

\begin{enumerate}
      \tightlist
      \item
            \textbf{Pensare velocemente}: riconoscere una situazione di gioco con
            anticipo rispetto all\textquotesingle avversario, grazie
            all\textquotesingle esperienza accumulata.
      \item
            \textbf{Agire velocemente}: tradurre quella decisione in
            un\textquotesingle azione motoria nei tempi giusti.
\end{enumerate}

Pensare velocemente senza agire altrettanto velocemente è inutile. Se un
atleta comprende che il compagno sta per servirgli un passaggio da
sfruttare per un tiro, ma è troppo lento nell\textquotesingle assumere
la posizione verticale corretta rispetto al compagno, alla porta e al
tempo di trasmissione della palla, fallirà comunque
l\textquotesingle azione.

\begin{quote}
      La tecnica individuale non è un fine in sé: è lo strumento per
      trasformare la decisione tattica in azione efficace nei tempi giusti.
\end{quote}

\subsubsection{Dalla Tecnica al Comportamento
      Tattico}\label{dalla-tecnica-al-comportamento-tattico}

Il percorso metodologico segue sempre la stessa progressione:

\begin{enumerate}
      \tightlist
      \item
            \textbf{Esercizi analitici} (tecnica individuale senza avversario ---
            i tempi li determina l\textquotesingle atleta).
      \item
            \textbf{Introduzione dell\textquotesingle avversario} (i tempi si
            adattano all\textquotesingle avversario --- coordinazione temporale
            esterna).
      \item
            \textbf{Introduzione della palla} (aggiunta
            dell\textquotesingle attrezzo --- ulteriore elemento di complessità).
      \item
            \textbf{Introduzione della porta} (finalizzazione --- avvicinamento al
            gioco reale).
\end{enumerate}

\begin{center}\rule{0.5\linewidth}{0.5pt}\end{center}

\subsection{Esercitazioni Specifiche in Corsia (Livello
      Avanzato)}\label{esercitazioni-specifiche-in-corsia-livello-avanzato}

\subsubsection{Nuotate Orizzontali
      Evolute}\label{nuotate-orizzontali-evolute}

\begin{itemize}
      \tightlist
      \item
            \textbf{Braccia a stile libero testa alta + gambata rana}: permette di
            occupare spazio in modo rapido rimanendo molto alti
            sull\textquotesingle acqua; favorisce l\textquotesingle attacco sul
            perimetro avversario.
      \item
            \textbf{Ricerca della lunghezza dell\textquotesingle azione}: potenzia
            gli arti inferiori e permette l\textquotesingle occupazione dello
            spazio orizzontale (diversamente dalla verticalità, che occupa lo
            spazio verticale).
      \item
            \textbf{Nuotate laterali}: fondamentale individuale che diventa poi
            controllo dell\textquotesingle avversario diretto e intercetto del
            passaggio.
\end{itemize}

\subsubsection{Modulazione delle Posizioni di Partenza e
      Arrivo}\label{modulazione-delle-posizioni-di-partenza-e-arrivo}

Variare le posizioni iniziali e finali degli esercizi costringe
l\textquotesingle atleta a \textbf{ricoordinarsi rapidamente},
stimolando la flessibilità motoria. Le rotazioni nelle varie direzioni,
le azioni aritmiche, gli arresti improvvisi hanno tutti questo
obiettivo.

\subsubsection{\texorpdfstring{Attacco Monobraccio (\emph{One-Arm
                  Attack})}{Attacco Monobraccio (One-Arm Attack)}}\label{attacco-monobraccio-one-arm-attack}

\begin{itemize}
      \tightlist
      \item
            Esercizio di simulazione dell\textquotesingle attacco al gomito
            dell\textquotesingle avversario in possesso di palla.
      \item
            È diventato fondamentale dopo la modifica regolamentare che
            \textbf{vieta di uscire con entrambe le braccia fuori
                  dall\textquotesingle acqua} (punito con fallo da espulsione).
      \item
            Tutta l\textquotesingle azione difensiva-offensiva deve quindi essere
            condotta con \textbf{un solo braccio}, lasciando
            all\textquotesingle altro e soprattutto agli \textbf{arti inferiori}
            il compito di sostenere e spingere.
      \item
            Combinazione tipica: \textbf{camminata avanti} (attacco al perimetro)
            + \textbf{attacco al gomito monobraccio} (pressione
            sull\textquotesingle avversario senza incorrere in fallo).
\end{itemize}

\subsubsection{\texorpdfstring{Azioni Combinate (\emph{Combined
                  Actions})}{Azioni Combinate (Combined Actions)}}\label{azioni-combinate-combined-actions}

Si parte da azioni semplici e analitiche per arrivare a sequenze sempre
più articolate:

\begin{itemize}
      \tightlist
      \item
            Uscita sul perimetro → attacco al gomito → recupero della verticalità
            per intercettare un eventuale passaggio.
\end{itemize}

Ogni fase dell\textquotesingle azione si innesta sulla precedente; le
azioni combinate sono la forma più prossima al comportamento tattico
reale.

\subsubsection{Copertura di Spazi
      Predeterminati}\label{copertura-di-spazi-predeterminati}

\begin{itemize}
      \tightlist
      \item
            Spazio di riferimento: \textbf{6 m × 6 m} (o simile).
      \item
            Spostamenti in tutte le direzioni su comando
            dell\textquotesingle allenatore; ad ogni arresto segue
            l\textquotesingle esecuzione di un gesto tecnico specifico (già
            acquisito in precedenza).
      \item
            Obiettivo: rapidità di reazione e adattamento posizionale.
\end{itemize}

\subsubsection{\texorpdfstring{L\textquotesingle1 contro 1
            (\emph{One-on-One})}{L\textquotesingle1 contro 1 (One-on-One)}}\label{l1-contro-1-one-on-one}

L\textquotesingle{}\textbf{1 contro 1 è il fondamentale tattico più
      importante} della pallanuoto. Ogni azione difensiva e ogni attacco ---
anche all\textquotesingle interno di sistemi di gioco collettivi --- si
risolve inevitabilmente in un duello individuale. Anche la difesa a zona
più elaborata, quando l\textquotesingle avversario si avvicina alla
porta, diventa difesa individuale.

Un giocatore che non padroneggia l\textquotesingle1 contro 1 non sarà
mai un buon pallanuotista.

\subsubsection{\texorpdfstring{Tecniche di Ruolo (\emph{Role-Specific
                  Techniques})}{Tecniche di Ruolo (Role-Specific Techniques)}}\label{tecniche-di-ruolo-role-specific-techniques}

Si introducono i movimenti tipici dei ruoli specifici:

\begin{itemize}
      \tightlist
      \item
            \textbf{Centro boa} (\emph{hole set}): posizionamento davanti alla
            porta, movimenti di smarcamento, ricezione e tiro in verticale.
      \item
            \textbf{Marcatore del centro boa}: tutti i movimenti per neutralizzare
            il giocatore avversario più pericoloso.
\end{itemize}

Questi movimenti sono la sintesi di tutto il lavoro analitico
precedente, finalizzato a situazioni di gioco reali.

\subsubsection{\texorpdfstring{Mantenimento del Possesso di Palla
            (\emph{Ball
                  Possession})}{Mantenimento del Possesso di Palla (Ball Possession)}}\label{mantenimento-del-possesso-di-palla-ball-possession}

\begin{itemize}
      \tightlist
      \item
            Esercizio con la \textbf{regola dei 3 secondi}: il giocatore in
            possesso ha 3 secondi per passare la palla a un compagno.
      \item
            Il compagno deve smarcarsi in quel lasso di tempo.
      \item
            Questo esercizio richiede l\textquotesingle integrazione di: presa
            della posizione, mantenimento della verticalità, presa
            dell\textquotesingle attrezzo, scelta e scarico del passaggio.
\end{itemize}

\subsubsection{Posizioni 2, 3 e 4: Smarcamento e
      Marcatura}\label{posizioni-2-3-e-4-smarcamento-e-marcatura}

\begin{itemize}
      \tightlist
      \item
            Disposizione \textbf{trasversale} rispetto al campo di gioco, che
            richiama le classiche posizioni di attacco (2, 3, 4).
      \item
            \textbf{Attaccanti}: azioni di smarcamento dalla marcatura.
      \item
            \textbf{Difensori}: recupero della marcatura sul proprio avversario
            diretto.
      \item
            Tutti i movimenti eseguiti sono la versione complessa e articolata dei
            gesti appresi in corsia.
\end{itemize}

\subsubsection{\texorpdfstring{Finalizzazione con la Porta
            (\emph{Shooting with
                  Goal})}{Finalizzazione con la Porta (Shooting with Goal)}}\label{finalizzazione-con-la-porta-shooting-with-goal}

L\textquotesingle obiettivo finale di tutto il percorso è
l\textquotesingle introduzione della porta e il gioco in spazio ridotto
(\textbf{3 contro 3 a una porta}), dove si finalizzano tutte le azioni
precedentemente costruite. Non è il gioco canonico 7 contro 7, ma
rappresenta la sintesi di:

\begin{itemize}
      \tightlist
      \item
            Verticalizzazione (dalla scuola nuoto)
      \item
            Galleggiamento verticale e orizzontale
      \item
            Fondamentali senza palla e con palla
      \item
            Comportamenti tattici individuali e di squadra
\end{itemize}

\begin{center}\rule{0.5\linewidth}{0.5pt}\end{center}

\subsection{La Corsia come Strumento
      Permanente}\label{la-corsia-come-strumento-permanente}

Anche per le squadre agonistiche di alto livello il lavoro in corsia
rimane utile e necessario. Già negli anni \textquotesingle60,
l\textquotesingle allenatore Mario Maglioni --- che vinse
l\textquotesingle oro olimpico a Londra nel 1948 --- descriveva come
riuscisse ad allenare efficacemente i propri giocatori in corsia nei
giorni o negli spazi in cui non aveva a disposizione il campo completo.

In corsia si lavora principalmente su:

\begin{itemize}
      \tightlist
      \item
            \textbf{Tecnica individuale} (fondamentali senza palla)
      \item
            \textbf{Nuotate orizzontali e verticali}
      \item
            \textbf{Affinamento dei gesti tecnici}
\end{itemize}

I contenuti in corsia vengono poi trasferiti e riutilizzati nelle
esercitazioni con campo completo.

\begin{center}\rule{0.5\linewidth}{0.5pt}\end{center}

\subsection{Conclusione --- Punti Chiave per
      l\textquotesingle Esame}\label{conclusione-punti-chiave-per-lesame-10}

\begin{itemize}
      \tightlist
      \item
            Il \textbf{settore amatoriale} (2 allenamenti/settimana + domeniche) e
            quello \textbf{agonistico} (da 4 allenamenti/settimana fino a
            giornaliero) originano entrambi dalla preagonistica; la distinzione è
            nell\textquotesingle impegno richiesto.
      \item
            Il settore agonistico è articolato in \textbf{Under 12}, \textbf{Under
                  14} e \textbf{Under 16}; le prime due categorie hanno campionato
            regionale, la terza nazionale. A 16 anni gli atleti vengono
            indirizzati ad altre società.
      \item
            L\textquotesingle atleta avanzato deve \textbf{pensare velocemente}
            (anticipazione tattica) e \textbf{agire velocemente} (tecnica pronta):
            entrambe le qualità sono necessarie; nessuna è sufficiente da sola.
      \item
            La progressione metodologica è sempre: analitico (senza avversario) →
            con avversario → con palla → con porta.
      \item
            L\textquotesingle{}\textbf{attacco monobraccio} è obbligatorio per via
            del regolamento che vieta l\textquotesingle uscita con entrambe le
            braccia (fallo da espulsione); l\textquotesingle azione si regge sugli
            arti inferiori.
      \item
            L\textquotesingle{}\textbf{1 contro 1} è il fondamentale tattico più
            importante della pallanuoto: qualsiasi sistema difensivo o offensivo
            si riduce sempre a un duello individuale.
      \item
            I movimenti di \textbf{ruolo} (centro boa e suo marcatore) sono la
            sintesi e la finalizzazione di tutto il percorso tecnico individuale.
      \item
            La \textbf{regola dei 3 secondi} è un esercizio chiave per integrare
            verticalità, possesso di palla, smarcamento e scelta del passaggio.
      \item
            La corsia rimane uno strumento valido a tutti i livelli, compresi i
            professionisti: permette di lavorare su tecnica individuale anche
            quando non è disponibile il campo completo.
      \item
            Il percorso complessivo --- dalla scuola nuoto alla pallanuoto --- va
            costruito \textbf{partendo da lontano}, dall\textquotesingle infanzia,
            con priorità assoluta alla tecnica individuale e
            all\textquotesingle esperienza motoria quantitativamente e
            qualitativamente ricca.
\end{itemize}

\section{16. Il Nuoto Artistico in Corsia --- Parte
  I}\label{il-nuoto-artistico-in-corsia-parte-i}

Docente: Barbara Quaranta, tecnica di nuoto artistico sincronizzato con
oltre 30 anni di esperienza, insegnante di educazione fisica.

\begin{center}\rule{0.5\linewidth}{0.5pt}\end{center}

\subsection{Introduzione}\label{introduzione-15}

Il \textbf{nuoto artistico sincronizzato} è una disciplina acquatica che
unisce tecnica natatori, espressività artistica e coordinazione. Questa
prima parte affronta la storia della disciplina, i suoi elementi
costitutivi, le tipologie di competizione e i fondamentali tecnici:
nuotate, bicicletta, remate e posizioni di base.

\begin{center}\rule{0.5\linewidth}{0.5pt}\end{center}

\subsection{Definizione della
      Disciplina}\label{definizione-della-disciplina}

Il nuoto artistico sincronizzato può essere definito in tre modi
complementari:

\begin{itemize}
      \tightlist
      \item
            \textbf{Disciplina natatoria}: il presupposto fondamentale è stare e
            spostarsi nell\textquotesingle acqua; le nuotate sono
            l\textquotesingle elemento primario.
      \item
            \textbf{Disciplina tecnico-compositoria}: gli elementi tecnici
            essenziali determinano direttamente la prestazione agonistica.
      \item
            \textbf{Disciplina artistico-espressiva}: gli elementi tecnici si
            combinano in attività artistiche e acrobatiche ad alto impatto
            spettacolare.
\end{itemize}

Attinge, per la sua componente tecnica, a discipline quali:

\begin{itemize}
      \tightlist
      \item
            Danza
      \item
            Ginnastica artistica
      \item
            Ginnastica acrobatica
      \item
            Tuffi
\end{itemize}

\begin{center}\rule{0.5\linewidth}{0.5pt}\end{center}

\subsection{Cenni Storici}\label{cenni-storici}

\begin{longtable}[]{@{}p{0.47\linewidth}p{0.47\linewidth}@{}}
      \toprule\noalign{}
      Periodo                 & Evento                                                  \\
      \midrule\noalign{}
      \endhead
      \bottomrule\noalign{}
      \endlastfoot
      Preistoria              & Pitture rupestri nella \textbf{grotta dei nuotatori del
            Sahara egiziano}: probabili riti di ringraziamento legati
      all\textquotesingle acqua                                                         \\
      Epoca imperiale romana  & Mosaici che testimoniano pratiche simili al
      nuoto artistico                                                                   \\
      Primo Novecento         & Notizie di nuotate in schieramento usate dagli
      eserciti dell\textquotesingle Europa del Nord tra le due guerre
      mondiali                                                                          \\
      Anni \textquotesingle50 & \textbf{Esther Williams} porta ad Hollywood le
      grandi coreografie in acqua, rendendo famosa la disciplina a livello
      mondiale                                                                          \\
      Anni \textquotesingle70 & Il nuoto sincronizzato nasce in
      \textbf{Italia}, creando un movimento abbastanza sviluppato da
      permettere la partecipazione olimpica                                             \\
      \textbf{1984}           & \textbf{Prima edizione olimpica} del nuoto
      sincronizzato: Giochi di Los Angeles                                              \\
\end{longtable}

\begin{center}\rule{0.5\linewidth}{0.5pt}\end{center}

\subsection{Elementi Costitutivi della
      Disciplina}\label{elementi-costitutivi-della-disciplina}

La struttura tecnica del nuoto artistico si articola in cinque
componenti fondamentali:

\begin{enumerate}
      \tightlist
      \item
            \textbf{Le nuotate}: stile libero, dorso, rana, delfino, stile sul
            fianco, bicicletta. Sono la base dello spostamento e del sostegno in
            acqua.
      \item
            \textbf{Bracciate, gambate e remate}: permettono di spostare e
            sostenere il corpo nelle diverse posizioni.
      \item
            \textbf{Posizioni fondamentali}: configurazioni del corpo, in
            emersione o in semie mersione, definite e normate dal regolamento con
            nome, coefficiente di difficoltà e criteri di esecuzione.
      \item
            \textbf{Movimenti}: collegano tra loro le posizioni di base; devono
            rispettare criteri di lentezza, precisione e controllo.
      \item
            \textbf{Figure obbligatorie}: elementi che combinano posizioni e
            movimenti di base secondo le disposizioni regolamentari.
\end{enumerate}

\begin{center}\rule{0.5\linewidth}{0.5pt}\end{center}

\subsection{Tipologie di Esercizi in
      Gara}\label{tipologie-di-esercizi-in-gara}

\subsubsection{\texorpdfstring{Esercizi Obbligatori (\emph{Compulsory
                  Figures})}{Esercizi Obbligatori (Compulsory Figures)}}\label{esercizi-obbligatori-compulsory-figures}

\begin{itemize}
      \tightlist
      \item
            Eseguiti individualmente, in forma \textbf{anonima}: costume nero,
            cuffia bianca.
      \item
            Nessuno spostamento nella vasca, nessun accompagnamento musicale.
      \item
            Di solito \textbf{quattro esercizi obbligatori} fanno parte del
            programma tecnico del quadriennio olimpico.
      \item
            Devono essere eseguiti in un \textbf{ordine di rotazione}
            prestabilito.
      \item
            Criteri di esecuzione: \textbf{massima lentezza} (salvo diversa
            indicazione regolamentare), \textbf{massima estensione},
            \textbf{massima tenuta} e \textbf{massima precisione} del gesto.
\end{itemize}

\subsubsection{Esercizi con Accompagnamento
      Musicale}\label{esercizi-con-accompagnamento-musicale}

Si dividono in:

\begin{itemize}
      \tightlist
      \item
            \textbf{Esercizi tecnici}: obbligo di includere elementi tecnici
            imposti.
      \item
            \textbf{Esercizi liberi}: composizione libera
            dell\textquotesingle atleta/squadra.
\end{itemize}

Entrambe le categorie possono essere di \textbf{solo}, \textbf{duo} o
\textbf{squadra}.

\subsubsection{\texorpdfstring{Esercizio \emph{High
                  Light}}{Esercizio High Light}}\label{esercizio-high-light}

\begin{itemize}
      \tightlist
      \item
            Tipologia di nuova generazione.
      \item
            Prevede una serie di \textbf{elementi acrobatici} (tuffi, torri,
            spinte verso l\textquotesingle alto) da eseguire in un ordine
            determinato.
      \item
            Si colloca a cavallo tra l\textquotesingle esercizio tecnico e la
            squadra combinata.
\end{itemize}

\subsubsection{Esercizio Combinato}\label{esercizio-combinato}

\begin{itemize}
      \tightlist
      \item
            Svolto da \textbf{8 a 10 atlete} contemporaneamente.
      \item
            La vasca funge da palcoscenico: le atlete si avvicendano nello
            specchio d\textquotesingle acqua alternando esercizi di solo, duo e
            squadra \textbf{senza soluzione di continuità}.
      \item
            Rappresenta una vera e propria narrazione attraverso il movimento.
\end{itemize}

\begin{center}\rule{0.5\linewidth}{0.5pt}\end{center}

\subsection{Le Nuotate nel Nuoto
      Artistico}\label{le-nuotate-nel-nuoto-artistico}

Nel nuoto artistico le nuotate convenzionali vengono adattate:
l\textquotesingle{}\textbf{assetto è più verticalizzato}, il busto è
generalmente più immerso (effetto espressivo e coreografico),
l\textquotesingle azione propulsiva delle gambe avviene in immersione e
\textbf{non si rompe la superficie dell\textquotesingle acqua} (niente
schiuma) se non per esigenze coreografiche.

\subsubsection{Stile Libero}\label{stile-libero}

\begin{itemize}
      \tightlist
      \item
            Assetto più verticale rispetto al crawl tradizionale.
      \item
            Bracciate più brevi e più veloci.
      \item
            La \textbf{fase aerea} può essere variata: braccio teso (\emph{teso})
            oppure suddiviso in più tempi (\emph{frazionato}), con movimenti della
            mano a titolo espressivo.
      \item
            Il busto più immerso lascia emergere spalle e volto, utilizzabili a
            fini espressivi.
\end{itemize}

\subsubsection{\texorpdfstring{Dorso
            (\emph{Backstroke})}{Dorso (Backstroke)}}\label{dorso-backstroke-1}

\begin{itemize}
      \tightlist
      \item
            Assetto anch\textquotesingle esso più verticalizzato.
      \item
            Maggiore libertà espressiva nella fase aerea della bracciata.
      \item
            Le braccia vengono usate su piani e con articolazioni diversificate
            (polso, gomito, spalla) per massimizzare l\textquotesingle effetto
            coreografico.
\end{itemize}

\subsubsection{\texorpdfstring{Stile sul Fianco
            (\emph{Sidestroke})}{Stile sul Fianco (Sidestroke)}}\label{stile-sul-fianco-sidestroke}

\begin{itemize}
      \tightlist
      \item
            Permette di controllare la compagna (negli esercizi di duo e squadra)
            e di spostarsi velocemente.
      \item
            Il braccio superiore è libero per riempire l\textquotesingle esercizio
            con gesti artistici.
\end{itemize}

\subsubsection{\texorpdfstring{Rana
            (\emph{Breaststroke})}{Rana (Breaststroke)}}\label{rana-breaststroke-1}

\begin{itemize}
      \tightlist
      \item
            Eseguita \textbf{senza la fase di immersione}, sfruttando al massimo
            la fase di trazione per ricercare quasi la verticalità completa.
      \item
            Tecnica caratteristica: \textbf{\emph{kick-pull-kick}}
            (spingo-tiro-spingo), usata per:

            \begin{itemize}
                  \tightlist
                  \item
                        Effetti ritmici
                  \item
                        Spostamenti rapidi nelle formazioni di squadra
                  \item
                        Riempimento ritmico dell\textquotesingle esercizio
            \end{itemize}
\end{itemize}

\subsubsection{\texorpdfstring{Delfino
            (\emph{Butterfly})}{Delfino (Butterfly)}}\label{delfino-butterfly}

\begin{itemize}
      \tightlist
      \item
            Utilizzato prevalentemente solo per la \textbf{gambata} (\emph{dolphin
                  kick}).
      \item
            Funzioni: spostamenti subacquei rapidi o effetti coreografici sulla
            superficie dell\textquotesingle acqua.
\end{itemize}

\subsubsection{\texorpdfstring{Bicicletta (\emph{Eggbeater
                  Kick})}{Bicicletta (Eggbeater Kick)}}\label{bicicletta-eggbeater-kick}

Elemento fondamentale della disciplina; i principi esecutivi sono gli
stessi della pallanuoto, con due differenze importanti:

\begin{longtable}[]{@{}p{0.30\linewidth}p{0.30\linewidth}p{0.30\linewidth}@{}}
      \toprule\noalign{}
      Parametro              & Pallanuoto & Nuoto Artistico \\
      \midrule\noalign{}
      \endhead
      \bottomrule\noalign{}
      \endlastfoot
      Frequenza              & Normale    & Più elevata     \\
      Estensione della gamba & Completa   & Mai completa    \\
\end{longtable}

\begin{itemize}
      \tightlist
      \item
            La frequenza più alta garantisce \textbf{stabilità e fermezza}:
            l\textquotesingle atleta deve dare l\textquotesingle impressione di
            avere i piedi appoggiati al fondo.
      \item
            Il busto deve essere \textbf{eretto in massima estensione}, con il
            capo ben elevato (presentazione e tenuta scenica).
      \item
            La postura corretta permette l\textquotesingle esecuzione dei
            movimenti delle braccia (una o entrambe ---
            quest\textquotesingle ultima molto più impegnativa perché il corpo è
            sostenuto da solo dalle gambe).
      \item
            La bicicletta si esegue in \textbf{tutte le direzioni}:
\end{itemize}

\begin{longtable}[]{@{}p{0.47\linewidth}p{0.47\linewidth}@{}}
      \toprule\noalign{}
      Direzione               & Posizione delle ginocchia                        \\
      \midrule\noalign{}
      \endhead
      \bottomrule\noalign{}
      \endlastfoot
      In avanzamento frontale & Leggermente più basse e intra-ruotate            \\
      In arretramento         & Leggermente più alte                             \\
      Laterale                & Arto corrispondente alla direzione più verticale \\
\end{longtable}

\begin{center}\rule{0.5\linewidth}{0.5pt}\end{center}

\subsection{\texorpdfstring{Le Remate
            (\emph{Sculling})}{Le Remate (Sculling)}}\label{le-remate-sculling}

Le \textbf{remate} sono il secondo grande elemento tecnico del nuoto
artistico. Permettono di \textbf{sostenere e spostare il corpo} quando
si trova in posizione supina o prona (scivolamenti).

\subsubsection{Meccanica del Gesto}\label{meccanica-del-gesto}

\begin{itemize}
      \tightlist
      \item
            Movimento di \textbf{flesso-estensione
                  dell\textquotesingle avambraccio sul braccio}.
      \item
            La parte del braccio dalla spalla al gomito rimane
            \textbf{generalmente ferma}, con il moncone della spalla abbassato.
      \item
            La \textbf{direzione o il sostegno} sono determinati
            dall\textquotesingle{}\textbf{angolo di inclinazione della mano
                  sull\textquotesingle avambraccio}, cioè dalla flessione del polso.
\end{itemize}

\subsubsection{Tre Tipologie di Remata}\label{tre-tipologie-di-remata}

\begin{longtable}[]{@{}p{0.30\linewidth}p{0.30\linewidth}p{0.30\linewidth}@{}}
      \toprule\noalign{}
      Tipo di remata                            & Posizione della mano                                   & Effetto \\
      \midrule\noalign{}
      \endhead
      \bottomrule\noalign{}
      \endlastfoot
      \textbf{Remata stazionaria}               & Mano sul prolungamento
      dell\textquotesingle avambraccio (neutra) & Solo sostegno: nessuno
      spostamento, pressione rivolta verso il fondo                                                                \\
      \textbf{Remata verso la testa}            & Mano in \textbf{flessione dorsale}
      (dita verso l\textquotesingle alto)       & Spostamento in direzione della
      testa (se braccia lungo i fianchi)                                                                           \\
      \textbf{Remata verso i piedi}             & Mano in \textbf{flessione palmare} (dita
      verso il fondo)                           & Spostamento in direzione dei piedi (se braccia lungo i
      fianchi)                                                                                                     \\
\end{longtable}

\begin{quote}
      \textbf{Nota}: la direzione cambia se le braccia sono sopra la testa
      (braccia a fianco delle orecchie). In quel caso la remata con mano in
      flessione dorsale sposta il corpo verso i piedi, e la remata con mano in
      flessione palmare sposta il corpo verso la testa --- la logica è sempre
      quella della reazione alla pressione esercitata
      sull\textquotesingle acqua.
\end{quote}

\begin{center}\rule{0.5\linewidth}{0.5pt}\end{center}

\subsection{Posizioni Fondamentali}\label{posizioni-fondamentali}

Le posizioni sono normate dal regolamento, che ne definisce nome,
coefficiente di difficoltà e criteri di esecuzione. Si eseguono su
supporto delle remate e della bicicletta.

\subsubsection{\texorpdfstring{Galleggiamento sul Dorso (\emph{Back
                  Layout})}{Galleggiamento sul Dorso (Back Layout)}}\label{galleggiamento-sul-dorso-back-layout}

\begin{itemize}
      \tightlist
      \item
            Corpo in completo galleggiamento, in \textbf{massima tenuta e massima
                  estensione}.
      \item
            In emersione: dorso dei piedi, cosce, petto e volto.
      \item
            È la posizione di base da cui si costruiscono tutte le altre.
\end{itemize}

\subsubsection{\texorpdfstring{Posizione di \emph{Tube} o
            Semiflessa}{Posizione di Tube o Semiflessa}}\label{posizione-di-tube-o-semiflessa}

\begin{itemize}
      \tightlist
      \item
            Dalla posizione di galleggiamento sul dorso:

            \begin{itemize}
                  \tightlist
                  \item
                        \textbf{Cosce perpendicolari alla superficie
                              dell\textquotesingle acqua}.
                  \item
                        \textbf{Gambe flesse sulle cosce} con angolo posteriore al ginocchio
                        di \textbf{90°}.
                  \item
                        Busto il più esteso possibile.
            \end{itemize}
      \item
            In emersione: gambe (dal ginocchio ai piedi) e volto; bacino
            affondato.
\end{itemize}

\subsubsection{\texorpdfstring{Posizione di
            \emph{Tuck}}{Posizione di Tuck}}\label{posizione-di-tuck}

\begin{itemize}
      \tightlist
      \item
            \textbf{Massima raccolta}: flessione ulteriore rispetto alla posizione
            di tube.
      \item
            Cosce praticamente sul busto.
      \item
            \textbf{Fronte a contatto con le ginocchia} (obiettivo da
            raggiungere).
      \item
            Gli ischi devono tendere a toccare i talloni.
      \item
            Gambe (dal ginocchio ai piedi) mantenute sulla superficie
            dell\textquotesingle acqua.
\end{itemize}

\subsubsection{\texorpdfstring{Posizione di Vela (\emph{Sailboat
                  Position})}{Posizione di Vela (Sailboat Position)}}\label{posizione-di-vela-sailboat-position}

\begin{itemize}
      \tightlist
      \item
            Si costruisce dalla posizione di galleggiamento sul dorso:

            \begin{itemize}
                  \tightlist
                  \item
                        \textbf{Una gamba si flette} in modo che la coscia sia
                        perpendicolare alla superficie dell\textquotesingle acqua.
                  \item
                        Il piede della gamba flessa deve toccare la parte \textbf{interna
                              della gamba distesa}, all\textquotesingle altezza del ginocchio.
                  \item
                        Busto in massima estensione, bacino il più possibile vicino alla
                        superficie.
                  \item
                        Il piede della gamba estesa deve stare sulla superficie
                        dell\textquotesingle acqua.
            \end{itemize}
      \item
            Il nome deriva dal \textbf{triangolo d\textquotesingle aria}
            (\emph{vela}) che dovrebbe comparire nella zona poplitea (dietro il
            ginocchio) della gamba flessa.
      \item
            Appare semplice ma richiede molto controllo posturale.
\end{itemize}

\subsubsection{\texorpdfstring{Gamba di Balletto (\emph{Ballet
                  Leg})}{Gamba di Balletto (Ballet Leg)}}\label{gamba-di-balletto-ballet-leg}

\begin{itemize}
      \tightlist
      \item
            Si costruisce dalla posizione di galleggiamento sul dorso:

            \begin{itemize}
                  \tightlist
                  \item
                        \textbf{Una gamba elevata perpendicolare alla superficie}, con il
                        piede sopra il bacino.
                  \item
                        Tutto il corpo in \textbf{massima estensione}.
                  \item
                        Ginocchia in massima estensione (anche la gamba sollevata).
                  \item
                        Bacino \textbf{sostenuto attivamente}, con le anche pari (non
                        inclinate).
            \end{itemize}
      \item
            È la posizione simbolo del nuoto artistico.
      \item
            Richiede ottime capacità di controllo del corpo e remata stazionaria
            molto efficace.
\end{itemize}

\begin{center}\rule{0.5\linewidth}{0.5pt}\end{center}

\subsection{I Movimenti di Base}\label{i-movimenti-di-base}

I movimenti \textbf{collegano tra loro le posizioni} e sono
anch\textquotesingle essi normati dal regolamento.

\subsubsection{Sviluppo della Gamba di Balletto (esempio di
      movimento)}\label{sviluppo-della-gamba-di-balletto-esempio-di-movimento}

\begin{enumerate}
      \tightlist
      \item
            Partenza dalla \textbf{posizione supina} (galleggiamento sul dorso).
      \item
            Una gamba si flette → \textbf{posizione di vela} (coscia
            perpendicolare, ginocchio sulla verticale).
      \item
            La gamba si estende verso l\textquotesingle alto → \textbf{gamba di
                  balletto} (mantenuta per alcuni istanti).
      \item
            La gamba si riflette: il \textbf{ginocchio non si sposta dalla
                  verticale} → si ritorna in posizione di vela.
      \item
            La gamba si estende verso il basso → si raggiunge
            l\textquotesingle altra gamba.
      \item
            Il bacino risale alla superficie → si torna alla \textbf{posizione di
                  galleggiamento sul dorso}.
\end{enumerate}

\subsubsection{Altri Movimenti
      Fondamentali}\label{altri-movimenti-fondamentali}

\begin{itemize}
      \tightlist
      \item
            \textbf{Capovolta indietro} e \textbf{capovolta in avanti}
      \item
            \textbf{Raggiungimento della gamba di balletto} e \textbf{chiusura
                  della gamba di balletto}
      \item
            \textbf{Raggiungimento della posizione di avanti} (\emph{front
                  layout})
      \item
            \textbf{Passeggiata avanti} e \textbf{passeggiata indietro}
\end{itemize}

Tutti i movimenti devono essere eseguiti con:

\begin{itemize}
      \tightlist
      \item
            \textbf{Massima lentezza}
      \item
            \textbf{Massima precisione} nel passaggio tra una posizione e
            l\textquotesingle altra
      \item
            \textbf{Massima tenuta del corpo}
      \item
            \textbf{Massimo controllo}
\end{itemize}

\begin{center}\rule{0.5\linewidth}{0.5pt}\end{center}

\subsection{Conclusione --- Punti Chiave per
      l\textquotesingle Esame}\label{conclusione-punti-chiave-per-lesame-11}

\begin{itemize}
      \tightlist
      \item
            Il nuoto artistico sincronizzato è contemporaneamente una disciplina
            \textbf{natatoria}, \textbf{tecnico-compositoria} e
            \textbf{artistico-espressiva}; attinge a danza, ginnastica artistica,
            ginnastica acrobatica e tuffi.
      \item
            La prima edizione olimpica è del \textbf{1984} (Los Angeles); in
            Italia la disciplina nasce negli \textbf{anni \textquotesingle70}.
      \item
            Gli elementi costitutivi sono: nuotate, bracciate/gambate/remate,
            posizioni fondamentali, movimenti e figure obbligatorie.
      \item
            Gli \textbf{esercizi obbligatori} si eseguono in anonimato (costume
            nero, cuffia bianca), senza musica né spostamenti, con massima
            lentezza, estensione, tenuta e precisione.
      \item
            Le nuotate nel nuoto artistico hanno \textbf{assetto più
                  verticalizzato} e non rompono la superficie dell\textquotesingle acqua
            (no schiuma).
      \item
            La \textbf{bicicletta} nel nuoto artistico ha frequenza più alta
            rispetto alla pallanuoto e non prevede mai la distensione completa
            della gamba: dà l\textquotesingle impressione di solidità e stabilità
            assoluta.
      \item
            Le \textbf{remate} sostengono o spostano il corpo in posizione
            supina/prona: la direzione dipende dall\textquotesingle angolo del
            polso (flessione dorsale = verso la testa; flessione palmare = verso i
            piedi; mano neutra = stazionaria).
      \item
            Le posizioni fondamentali da conoscere: \textbf{galleggiamento sul
                  dorso, tube/semiflessa, tuck, vela, gamba di balletto} --- ciascuna
            con i propri criteri tecnici precisi.
      \item
            I movimenti collegano le posizioni e devono rispettare sempre i
            criteri di \textbf{lentezza, precisione e controllo}.
\end{itemize}

\section{17. Il Nuoto Artistico in Corsia --- Parte
  II}\label{il-nuoto-artistico-in-corsia-parte-ii}

Docente: Barbara Quaranta, tecnica di nuoto artistico sincronizzato.

\begin{center}\rule{0.5\linewidth}{0.5pt}\end{center}

\subsection{Introduzione}\label{introduzione-16}

Questa seconda parte chiarisce un possibile paradosso: come può la
\textbf{corsia} essere il luogo ideale per iniziare il percorso nel
nuoto artistico sincronizzato? La risposta sta nel fatto che la corsia
offre tutto il necessario per costruire le \textbf{capacità motorie di
      base} che questa disciplina richiede. La lezione è organizzata attorno
ai pilastri fondamentali dell\textquotesingle apprendimento:
equilibramento, tenuta, respirazione, uso dello spazio, ritmo e, infine,
espressione artistica.

\begin{center}\rule{0.5\linewidth}{0.5pt}\end{center}

\subsection{Capacità Motorie Richieste dalla
      Disciplina}\label{capacituxe0-motorie-richieste-dalla-disciplina}

Il nuoto artistico sincronizzato è una disciplina ad \textbf{elevata
      complessità motoria}. Richiede:

\begin{itemize}
      \tightlist
      \item
            \textbf{Coordinazione intersegmentaria}: l\textquotesingle azione
            differenziata e dissociata degli arti superiori e inferiori è visibile
            nelle riprese subacquee delle competizioni ed è alla base
            dell\textquotesingle esecuzione della maggior parte degli elementi
            tecnici.
      \item
            \textbf{Differenziazione e aggiustamento motorio}: capacità di
            modulare finemente l\textquotesingle intensità e la forma del gesto.
      \item
            \textbf{Equilibramento e senso-percezione}: adattare, riconoscere,
            modificare e creare movimenti nuovi in un ambiente instabile come
            l\textquotesingle acqua.
      \item
            \textbf{Base condizionale}: resistenza, forza e velocità come supporto
            all\textquotesingle esecuzione tecnica.
      \item
            \textbf{Flessibilità}: capacità a cavallo tra condizionale e
            coordinativa, essenziale per l\textquotesingle esecuzione degli
            elementi tecnici.
\end{itemize}

\begin{center}\rule{0.5\linewidth}{0.5pt}\end{center}

\subsection{Ruolo dell\textquotesingle Istruttore in
      Corsia}\label{ruolo-dellistruttore-in-corsia}

L\textquotesingle istruttore in corsia può e deve \textbf{costruire il
      bagaglio motorio di base} che verrà poi speso nel nuoto artistico. Lo
strumento principale è la \textbf{variazione sul tema}: non proporre
sempre la nuotata standard, ma modificarne sistematicamente i parametri.

\subsubsection{Variazioni sulle Nuotate}\label{variazioni-sulle-nuotate}

\begin{itemize}
      \tightlist
      \item
            Stile libero con testa fuori dall\textquotesingle acqua (\emph{head-up
                  crawl})
      \item
            Braccia tese vs. braccia più flesse durante la bracciata
      \item
            Fase aerea frazionata (il braccio passa in più tempi con movimenti
            espressivi della mano)
      \item
            Gambate a diversi livelli di profondità (più a fondo o più in
            superficie)
      \item
            \textbf{Coordinazione crociata}: gamba destra + braccio sinistro a
            stile, gamba sinistra + braccio destro a stile --- abitua
            all\textquotesingle utilizzo creativo e non convenzionale degli arti.
\end{itemize}

\begin{quote}
      Il punto cruciale non è solo eseguire la variazione, ma
      \textbf{percepire} cosa cambia: cosa rende più facile o difficile
      l\textquotesingle esecuzione, dove si trova difficoltà e dove vantaggio.
      Stimolare questa consapevolezza senso-percettiva è
      l\textquotesingle obiettivo.
\end{quote}

\begin{center}\rule{0.5\linewidth}{0.5pt}\end{center}

\subsection{\texorpdfstring{La Bicicletta (\emph{Eggbeater Kick}) in
            Corsia: Approcci
            Didattici}{La Bicicletta (Eggbeater Kick) in Corsia: Approcci Didattici}}\label{la-bicicletta-eggbeater-kick-in-corsia-approcci-didattici}

La bicicletta è un elemento tecnico avanzato, ma si possono proporre
attività propedeutiche in forma giocosa sin da subito.

\subsubsection{Strumenti Didattici}\label{strumenti-didattici}

\begin{itemize}
      \tightlist
      \item
            \textbf{Ciambelle galleggianti} (una per gamba): insegnano a tenere le
            gambe divaricate e le ginocchia in intra-rotazione; se le ginocchia si
            allargano, i salvagenti sfuggono.
      \item
            \textbf{Tubi galleggianti} (\emph{pool noodles}) a cavallo: forniscono
            sostegno e permettono di concentrarsi sull\textquotesingle azione
            ciclica delle gambe.
      \item
            \textbf{Esercizi contro la parete}: permettono di percepire la
            corretta postura del busto (verticale) e la giusta flessione della
            coscia sul bacino.
\end{itemize}

\subsubsection{Giochi Propedeutici}\label{giochi-propedeutici}

\begin{itemize}
      \tightlist
      \item
            \textbf{Gioco del gallo in acqua}: in coppia, mano contro mano,
            ciascuno cerca di spostare l\textquotesingle altro usando la
            bicicletta (variante acquatica di un gioco consueto a secco).
      \item
            \textbf{Schiena contro schiena}: due allievi cercano di spingersi
            l\textquotesingle uno contro l\textquotesingle altro.
      \item
            \textbf{Spingere il bordo della piscina}: sviluppa la sensibilità
            della pressione dell\textquotesingle acqua in assetti diversificati.
\end{itemize}

Tutti questi giochi possono essere proposti senza necessariamente citare
la bicicletta del nuoto sincronizzato: l\textquotesingle obiettivo è
\textbf{formativo}, non ancora tecnico-analitico.

\subsubsection{\texorpdfstring{Attività di \emph{Boost} (Bicicletta +
            Salto)}{Attività di Boost (Bicicletta + Salto)}}\label{attivituxe0-di-boost-bicicletta-salto}

\begin{itemize}
      \tightlist
      \item
            Esercizi di salto verso oggetti tenuti dall\textquotesingle istruttore
            dal bordo vasca.
      \item
            Il \emph{boost} (bicicletta + spinta + salto) è propedeutico alle
            \textbf{spinte verticali} del nuoto sincronizzato (torri, lanci delle
            compagne) e, in pallanuoto, alle prese e ai salti del portiere.
\end{itemize}

\begin{center}\rule{0.5\linewidth}{0.5pt}\end{center}

\subsection{Trasporto e Tenuta}\label{trasporto-e-tenuta}

Le \textbf{attività di trasporto} (di oggetti o di compagni) sono
trasversali a più discipline natatorie e hanno una duplice utilità:

\begin{itemize}
      \tightlist
      \item
            Sviluppano la tenuta e il controllo del corpo.
      \item
            Sono propedeutiche alla \textbf{tenuta} (\emph{body tension}),
            elemento tecnico fondamentale nel nuoto artistico.
\end{itemize}

\begin{center}\rule{0.5\linewidth}{0.5pt}\end{center}

\subsection{Equilibramento nelle Posizioni di
      Base}\label{equilibramento-nelle-posizioni-di-base}

\subsubsection{Il Principio Chiave}\label{il-principio-chiave}

Per mantenere le posizioni fondamentali in massima emersione, gli
allievi devono imparare a \textbf{percepire e sfruttare le spinte di
      galleggiamento}.

\subsubsection{Attività di
      Equilibramento}\label{attivituxe0-di-equilibramento}

\begin{itemize}
      \tightlist
      \item
            Galleggiamenti classici (prono, supino) e galleggiamenti in posizioni
            insolite.
      \item
            Galleggiamenti sul fondo: fermarsi mimando la corsa
            sott\textquotesingle acqua, "fotografare" la posizione e percepire
            cosa succede al proprio corpo in quell\textquotesingle assetto.
      \item
            Indicazioni su come \textbf{sfruttare} le spinte (stare più in
            superficie) o \textbf{contrastarle} (affondare).
\end{itemize}

L\textquotesingle obiettivo è aumentare la \textbf{percezione di sé}
nell\textquotesingle acqua, riconoscere le spinte di galleggiamento e
poterle sfruttare quando si innesta l\textquotesingle elemento tecnico.

\begin{center}\rule{0.5\linewidth}{0.5pt}\end{center}

\subsection{Tenuta e Controllo del
      Corpo}\label{tenuta-e-controllo-del-corpo}

\subsubsection{Il Paradosso del Nuoto
      Artistico}\label{il-paradosso-del-nuoto-artistico}

Nel nuoto si insegna a rilassarsi nell\textquotesingle acqua e a
sfruttare la spinta di galleggiamento. Nel nuoto artistico questo rimane
vero, ma va affiancato a una capacità opposta: \textbf{rendere il corpo
      solido e bloccato} nelle posizioni che costituiscono gli elementi di
base.

\subsubsection{Esercizi di Tenuta}\label{esercizi-di-tenuta}

\begin{itemize}
      \tightlist
      \item
            \textbf{Scivolamento dopo spinta} (gara con un compagno: chi arriva
            più lontano = chi è più teso e controllato).
      \item
            \textbf{Essere tirati} da un compagno durante lo scivolamento.
      \item
            \textbf{Affondamento controllato}: mantenere il corpo in una posizione
            che penetra meglio l\textquotesingle acqua mentre si viene spinti
            verso il fondo.
      \item
            \textbf{Uscita come un missile}: essere spinti dalla superficie o dal
            fondo verso l\textquotesingle alto, cercando di uscire
            dall\textquotesingle acqua il più possibile con il corpo ben teso.
\end{itemize}

\begin{center}\rule{0.5\linewidth}{0.5pt}\end{center}

\subsection{\texorpdfstring{Sensibilizzazione alla Remata
            (\emph{Sculling
                  Sensitivity})}{Sensibilizzazione alla Remata (Sculling Sensitivity)}}\label{sensibilizzazione-alla-remata-sculling-sensitivity}

\subsubsection{Prerequisito: Sensibilità
      all\textquotesingle Acqua}\label{prerequisito-sensibilituxe0-allacqua}

La remata è tecnicamente precisa, ma è priva di significato se
l\textquotesingle allievo non ha sviluppato una \textbf{sensibilità
      all\textquotesingle acqua}. La sensibilità si costruisce prima che nelle
posizioni di base, in qualsiasi situazione.

\subsubsection{Esercizi di
      Sensibilizzazione}\label{esercizi-di-sensibilizzazione}

\begin{itemize}
      \tightlist
      \item
            Spostarsi in tutte le direzioni (verso la testa, verso i piedi) usando
            solo le mani, in galleggiamento supino.
      \item
            Tenersi in sospensione tra superficie e fondo.
      \item
            Tenersi contro la parete con la schiena in galleggiamento supino, con
            i piedi che spingono contro il muro.
      \item
            Variante con la testa contro la parete.
      \item
            \textbf{Schiena contro schiena}: due compagni si spingono
            l\textquotesingle uno contro l\textquotesingle altro usando le azioni
            di remata.
      \item
            \textbf{Piedi contro piedi}: simula una bozza di posizione di
            \emph{tuck} (vista nella Parte I), in galleggiamento supino.
      \item
            \textbf{Piccole coreografie} a tre o quattro con i piedi al centro
            (es. "fiore"): l\textquotesingle obiettivo non è la precisione della
            posizione, ma l\textquotesingle orientamento nello spazio e la
            sensibilizzazione delle mani per sostegno e direzionalità.
\end{itemize}

\begin{quote}
      Proporre le remate senza necessariamente chiamarle tali, con una
      finalità \textbf{formativa} (aumentare la sensibilità
      all\textquotesingle acqua) piuttosto che tecnico-analitica.
\end{quote}

\begin{center}\rule{0.5\linewidth}{0.5pt}\end{center}

\subsection{Respirazione nel Nuoto
      Artistico}\label{respirazione-nel-nuoto-artistico}

\subsubsection{Caratteristica della
      Disciplina}\label{caratteristica-della-disciplina}

Il nuoto artistico sincronizzato prevede \textbf{momenti di apnea
      prolungati} sia negli esercizi obbligatori che in quelli liberi.
Conoscere la propria respirazione è garanzia di attività sana.

\subsubsection{Effetti dell\textquotesingle Apnea sul
      Galleggiamento}\label{effetti-dellapnea-sul-galleggiamento}

\begin{longtable}[]{@{}p{0.47\linewidth}p{0.47\linewidth}@{}}
      \toprule\noalign{}
      Tipo di apnea                         & Effetto sul galleggiamento      \\
      \midrule\noalign{}
      \endhead
      \bottomrule\noalign{}
      \endlastfoot
      \textbf{Inspiratoria} (polmoni pieni) & Maggiore spinta di
      galleggiamento --- corpo più alto                                       \\
      \textbf{Espiratoria} (polmoni vuoti)  & Minore spinta di galleggiamento
      --- corpo più basso                                                     \\
\end{longtable}

Questa differenza è direttamente collegata alle fasi di equilibramento e
può essere sfruttata o gestita nell\textquotesingle esecuzione degli
elementi tecnici. Chi padroneggia bene la propria respirazione può, ad
esempio, usare l\textquotesingle espirazione per gestire esercizi che
diversamente risulterebbero difficili.

\subsubsection{Attività di Educazione
      Respiratoria}\label{attivituxe0-di-educazione-respiratoria}

\begin{itemize}
      \tightlist
      \item
            \textbf{Ritmica della respirazione}: collegare le fasi di inspirazione
            ed espirazione al ritmo della gambata o a quello della compagna.
      \item
            Stare in galleggiamento o in immersione con \textbf{apnea
                  inspiratoria} per percepire l\textquotesingle aumento della spinta
            verso il basso.
      \item
            Ritmo respiratorio scandito dalla \textbf{voce degli stessi allievi}.
\end{itemize}

\begin{center}\rule{0.5\linewidth}{0.5pt}\end{center}

\subsection{Utilizzo Tridimensionale dello Spazio in
      Corsia}\label{utilizzo-tridimensionale-dello-spazio-in-corsia}

\subsubsection{La Corsia come Spazio
      Tridimensionale}\label{la-corsia-come-spazio-tridimensionale}

Il nuoto artistico è una disciplina che si gioca sulla
\textbf{tridimensionalità}. La corsia, pur apparendo uno spazio ridotto,
contiene un \textbf{volume d\textquotesingle acqua significativo} con
altezza, larghezza e lunghezza da sfruttare in tutte le direzioni.

\subsubsection{Modalità di Utilizzo Insolito della
      Corsia}\label{modalituxe0-di-utilizzo-insolito-della-corsia}

\begin{itemize}
      \tightlist
      \item
            Avanzare e arretrare (senso tradizionale)
      \item
            Usare \textbf{tutte le pareti}: lunghezza, altezza
      \item
            Nuotare in \textbf{zig-zag} tra il muro e la corda di corsia
      \item
            Sfruttare lo \textbf{spazio subacqueo} in tutto il volume
            d\textquotesingle acqua
      \item
            \textbf{Percorsi creativi}: gli allievi disegnano su carta un percorso
            (come se la corsia fosse una vasca vista dall\textquotesingle alto),
            scambiano il foglio con un compagno e nella lezione successiva devono
            seguire il percorso inventato dall\textquotesingle altro.
\end{itemize}

Questi esercizi preparano anche all\textquotesingle uso delle
\textbf{piattaforme} e delle spinte verticali (\emph{boost}) che nel
nuoto artistico moderno permettono alle atlete di spingere le compagne
in alto con grande forza e precisione.

\begin{center}\rule{0.5\linewidth}{0.5pt}\end{center}

\subsection{Ritmo e Sincronismo}\label{ritmo-e-sincronismo}

\subsubsection{Il Ritmo nel Nuoto
      Artistico}\label{il-ritmo-nel-nuoto-artistico}

Il ritmo nel nuoto artistico è \textbf{prevalentemente esterno}:
l\textquotesingle atleta deve seguire il ritmo della musica di
accompagnamento, non solo il proprio ritmo interno.

È necessario quindi distinguere e allenare:

\begin{itemize}
      \tightlist
      \item
            \textbf{Ritmo percepito}: riconoscere il proprio ritmo esecutivo nelle
            nuotate, nelle gambate, nella bicicletta, nelle remate.
      \item
            \textbf{Ritmo esterno}: adeguarsi a un ritmo che viene da fuori
            (musica, voce dell\textquotesingle istruttore, voce dei compagni,
            filastrocca).
\end{itemize}

\begin{quote}
      All\textquotesingle interno di un brano musicale esistono molti ritmi
      sovrapposti: la cadenza di base, il tempo, le variazioni melodiche. È
      importante stimolare gli allievi a percepire questa molteplicità.
\end{quote}

\subsubsection{Fonti di Ritmo Esterno da Usare in
      Corsia}\label{fonti-di-ritmo-esterno-da-usare-in-corsia}

\begin{itemize}
      \tightlist
      \item
            Voce degli allievi stessi (es. contare ad alta voce mentre nuotano)
      \item
            Filastrocche conosciute (il ritmo della filastrocca guida la gambata)
      \item
            Musica dell\textquotesingle attività di acqua gym in corso nelle
            corsie adiacenti
      \item
            Voce o battito dell\textquotesingle istruttore
\end{itemize}

\subsubsection{Sincronismo con i
      Compagni}\label{sincronismo-con-i-compagni}

Il sincronismo consiste nel trovare con una o più compagne
\textbf{accenti ritmici comuni}. Esercizi utili:

\begin{itemize}
      \tightlist
      \item
            Eseguire una gambata consueta dando l\textquotesingle accento sulla
            \textbf{fase discendente} (tempo in battere).
      \item
            Eseguire la stessa gambata dando l\textquotesingle accento sulla
            \textbf{fase ascendente} (tempo in levare).
      \item
            Scambiare i ruoli e confrontare le sensazioni.
\end{itemize}

\begin{center}\rule{0.5\linewidth}{0.5pt}\end{center}

\subsection{Espressione Artistica}\label{espressione-artistica}

\subsubsection{Che Cos\textquotesingle è l\textquotesingle Espressione
      Artistica}\label{che-cosuxe8-lespressione-artistica}

Costruire la parte artistica non significa solo impressionare il
pubblico con figure spettacolari. Significa lavorare sulla
\textbf{sensibilizzazione della persona} e sulla sua \textbf{percezione
      emotiva}, permettendo di narrare qualcosa attraverso il gesto tecnico
--- esattamente come fa la danza.

\subsubsection{Come Costruirla in
      Corsia}\label{come-costruirla-in-corsia}

L\textquotesingle approccio più semplice e immediato è associare le
nuotate a emozioni o stati d\textquotesingle animo:

\begin{itemize}
      \tightlist
      \item
            \emph{"Nuota tre vasche come se fossi arrabbiato."}
      \item
            \emph{"Nuota come ti sentiresti dopo aver ricevuto una bella
                  notizia."}
\end{itemize}

Questo trasforma l\textquotesingle esecuzione tecnica consueta in
qualcosa di personale e filtrato da un\textquotesingle emozione. Non si
chiede di abbandonare la tecnica, ma di usarla come \textbf{mezzo
      espressivo}.

Il risultato è la costruzione della \textbf{libertà di esprimersi}:
l\textquotesingle allievo impara a narrare sé stesso o qualcosa
d\textquotesingle altro attraverso il gesto tecnico, che è la base
dell\textquotesingle arte nel nuoto artistico.

\subsubsection{Le Parole del Nuoto
      Artistico}\label{le-parole-del-nuoto-artistico}

Quando al pubblico viene chiesto cosa associa al nuoto artistico, le
risposte ricorrenti sono:

\begin{quote}
      \emph{musica, passione, squadra, sacrificio, dedizione, danza in acqua,
            emozione, alleanza, fiducia, determinazione.}
\end{quote}

Tutte parole che presuppongono un\textquotesingle{}\textbf{esperienza di
      vita}: questo è il messaggio della disciplina.

\begin{center}\rule{0.5\linewidth}{0.5pt}\end{center}

\subsection{Conclusione --- Punti Chiave per
      l\textquotesingle Esame}\label{conclusione-punti-chiave-per-lesame-12}

\begin{itemize}
      \tightlist
      \item
            Il nuoto artistico richiede \textbf{coordinazione intersegmentaria},
            differenziazione motoria, equilibramento, senso-percezione, una solida
            base condizionale (forza, resistenza, velocità) e
            \textbf{flessibilità}.
      \item
            In corsia si costruisce il bagaglio motorio di base attraverso
            \textbf{variazioni sul tema}: nuotate con testa alta, braccia
            tese/flesse, fase aerea frazionata, coordinazione crociata, gambate a
            diverse profondità.
      \item
            La \textbf{bicicletta} si introduce gradualmente con strumenti di
            supporto (ciambelle, tubi) e giochi (gioco del gallo, schiena contro
            schiena); le attività di \emph{boost} preparano ai salti verticali.
      \item
            \textbf{Equilibramento} = percepire e sfruttare le spinte di
            galleggiamento; la tenuta = saper rendere il corpo solido nelle
            posizioni, in apparente contrasto con il rilassamento insegnato in
            scuola nuoto.
      \item
            \textbf{Sensibilità alla remata}: si costruisce prima della tecnica,
            attraverso tutti i modi possibili di spostarsi
            nell\textquotesingle acqua con le mani.
      \item
            \textbf{Respirazione}: apnea inspiratoria → più galleggiamento; apnea
            espiratoria → meno galleggiamento. Conoscere la propria respirazione è
            fondamentale e garantisce anche la salute dell\textquotesingle atleta.
      \item
            La corsia va usata in modo \textbf{tridimensionale}: tutte le
            direzioni, tutte le pareti, tutto il volume subacqueo.
      \item
            Il \textbf{ritmo} nel nuoto artistico è prevalentemente esterno
            (musica); allenarsi a seguire ritmi diversi (voce, filastrocca,
            musica) e a trovare accenti ritmici comuni con le compagne costruisce
            il \textbf{sincronismo}.
      \item
            L\textquotesingle{}\textbf{espressione artistica} si costruisce
            associando emozioni alle nuotate; permette di trasformare il gesto
            tecnico in strumento narrativo, come avviene nella danza.
\end{itemize}

\section{18. Il Nuoto per Salvamento in
  Corsia}\label{il-nuoto-per-salvamento-in-corsia}

\subsection{Introduzione}\label{introduzione-17}

Il \textbf{nuoto per salvamento} è la quinta disciplina sportiva della
Federazione Italiana Nuoto (FIN). Non ancora olimpica, è una disciplina
composta da competizioni che si svolgono sia in \textbf{piscina} (in
corsia) che in \textbf{mare} (gare oceaniche). Questa lezione, tenuta da
Filippo Tassara, analizza il modello funzionale del nuotatore di
salvamento, le tecniche specifiche, le apparecchiature utilizzate e la
progressione didattica per organizzare un corso di salvamento in corsia.

\begin{center}\rule{0.5\linewidth}{0.5pt}\end{center}

\subsection{Il Modello Funzionale del Nuotatore di
      Salvamento}\label{il-modello-funzionale-del-nuotatore-di-salvamento}

Il nuotatore di salvamento può essere descritto attraverso un
\textbf{modello funzionale a tre livelli}.

\subsubsection{1. Modello Psicologico}\label{modello-psicologico}

Il nuotatore di salvamento deve possedere:

\begin{itemize}
      \tightlist
      \item
            \textbf{Motivazione} elevata e capacità di sostenere carichi di lavoro
            importanti (per gli assoluti anche fino a \textbf{10 allenamenti
                  settimanali})
      \item
            \textbf{Adattamento al lavoro} e capacità di soffrire nei momenti
            cruciali
      \item
            \textbf{Autoconoscenza e autocontrollo}
      \item
            \textbf{Capacità di lavorare in gruppo}
\end{itemize}

\subsubsection{2. Modello
      Tecnico-Coordinativo}\label{modello-tecnico-coordinativo}

Nelle gare in piscina, la tecnica natatoria è strettamente correlata
allo \textbf{stile libero} (\emph{crawl}). Le varianti tecniche, in
progressione, sono:

\begin{longtable}[]{@{}p{0.47\linewidth}p{0.47\linewidth}@{}}
      \toprule\noalign{}
      Tecnica                                  & Caratteristiche principali                     \\
      \midrule\noalign{}
      \endhead
      \bottomrule\noalign{}
      \endlastfoot
      \emph{Crawl} classico                    & Base dello stile libero                        \\
      \emph{Crawl} testa alta                  & Variante con testa sollevata per la visione in
      acqua                                                                                     \\
      \emph{Over} (\emph{overarm side stroke}) & Nuotata antica con corpo
      inclinato sull\textquotesingle avanzamento; propedeutica al
      \emph{trudgeon}                                                                           \\
      \emph{Trudgeon}                          & Azione delle gambe simmetrica a forbice (non
      alternata); arti superiori come nel \emph{crawl}                                          \\
      Rana rovesciata                          & Utilizzata in alcune fasi di gara, con o senza
      l\textquotesingle azione delle braccia; può essere associata al dorso                     \\
\end{longtable}

Nelle \textbf{gare oceaniche} si aggiungono tecniche di:

\begin{itemize}
      \tightlist
      \item
            \textbf{Corsa} (es. i 90 metri bandierine)
      \item
            \textbf{Nuoto in acque libere} (varianti del \emph{crawl})
      \item
            \textbf{Remata} con le braccia (gara di tavola, su 300-400 metri)
\end{itemize}

\subsubsection{3. Modello
      Fisico-Condizionale}\label{modello-fisico-condizionale}

Riguarda l\textquotesingle efficienza metabolica degli apparati
muscolare e scheletrico e i meccanismi energetici. Include:

\begin{itemize}
      \tightlist
      \item
            \textbf{Forza} nelle sue tipologie: forza velocità e forza resistenza
      \item
            \textbf{Velocità} e \textbf{resistenza}
      \item
            \textbf{Mobilità articolare}: fondamentale per mantenere
            l\textquotesingle economia del gesto tecnico-natatorio
\end{itemize}

Ogni gara si classifica in base ai tempi di realizzazione e alle
capacità di forza, resistenza e velocità espresse.

\begin{center}\rule{0.5\linewidth}{0.5pt}\end{center}

\subsection{Le Gare}\label{le-gare}

\subsubsection{Gare in Piscina (in
      corsia)}\label{gare-in-piscina-in-corsia}

Si distinguono per:

\begin{itemize}
      \tightlist
      \item
            Nuotate eseguite \textbf{in corsia}
      \item
            Utilizzo del \textbf{corpo} e di \textbf{attrezzi/apparecchiature}
      \item
            Stili che partono dal \emph{crawl} e si estendono ad altre discipline
\end{itemize}

\subsubsection{Gare Oceaniche}\label{gare-oceaniche}

Si svolgono "sulla spiaggia" e comprendono:

\begin{itemize}
      \tightlist
      \item
            Prove di \textbf{nuoto}
      \item
            Attività \textbf{a secco} (es. sprint con le bandierine, 90 metri)
      \item
            Prove \textbf{miste} (passaggi dall\textquotesingle acqua alla
            terraferma)
\end{itemize}

\begin{center}\rule{0.5\linewidth}{0.5pt}\end{center}

\subsection{Le Apparecchiature}\label{le-apparecchiature}

\begin{longtable}[]{@{}p{0.47\linewidth}p{0.47\linewidth}@{}}
      \toprule\noalign{}
      Attrezzo           & Descrizione e utilizzo                             \\
      \midrule\noalign{}
      \endhead
      \bottomrule\noalign{}
      \endlastfoot
      \textbf{Manichino} & Attrezzatura principale; può essere riempito o
      svuotato d\textquotesingle acqua per variarne il galleggiamento         \\
      \textbf{Ostacoli}  & Strutture sotto le quali il nuotatore si immerge a
      una profondità di \textbf{70 cm}; posizionabili tra i galleggianti delle
      corsie                                                                  \\
      \textbf{Torpedo}   & Banda galleggiante con ancoraggio che si chiude a
      cerchio intorno al nuotatore, per il trasporto del manichino            \\
      \textbf{Pinne}     & Quelle del salvamento sono specifiche, ma nelle
      attività in corsia si possono usare pinne normali                       \\
\end{longtable}

\begin{center}\rule{0.5\linewidth}{0.5pt}\end{center}

\subsection{La Scuola di Salvamento in
      Corsia}\label{la-scuola-di-salvamento-in-corsia}

L\textquotesingle istruttore ha un \textbf{ruolo fondamentale di guida e
      orientamento}. La formazione deve:

\begin{itemize}
      \tightlist
      \item
            Essere centrata sull\textquotesingle{}\textbf{individuo} e sulla sua
            \textbf{età biologica} (non cronologica)
      \item
            Rispettare gli studi di sviluppo fisico, cognitivo e affettivo
      \item
            Integrarsi con le conoscenze scientifiche: pediatria, psicologia,
            sociologia, fisiologia dell\textquotesingle esercizio
\end{itemize}

\subsubsection{Quando Iniziare un Corso di
      Salvamento}\label{quando-iniziare-un-corso-di-salvamento}

Idealmente, gli allievi dovrebbero aver \textbf{completato tutti i
      livelli della Scuola Nuoto Federale} (5° livello). Tuttavia, il percorso
può essere articolato in \textbf{due fasi}.

\begin{center}\rule{0.5\linewidth}{0.5pt}\end{center}

\subsection{Le Due Fasi del Percorso
      Formativo}\label{le-due-fasi-del-percorso-formativo}

\subsubsection{\texorpdfstring{Prima Fase - Livello Base
            (\emph{SaNuotare
                  1})}{Prima Fase - Livello Base (SaNuotare 1)}}\label{prima-fase---livello-base-sanuotare-1}

\textbf{Requisito di ingresso:} Conseguimento del \textbf{4° livello}
della Scuola Nuoto Federale.

\textbf{Frequenza:} lezioni \textbf{bisettimanali} da 40-45 minuti.

Il 4° livello corrisponde al \textbf{perfezionamento delle nuotate} e
consente al soggetto di essere autonomo in acqua (\emph{SaNuotare 1}: il
soggetto sa vivere l\textquotesingle ambiente acqua in equilibrio, senza
problemi di sopravvivenza).

\textbf{Obiettivi tecnici:}

\begin{itemize}
      \tightlist
      \item
            25 metri a rana
      \item
            25 metri a farfalla
      \item
            Sviluppo di crono nel dorso e nel \emph{crawl}
      \item
            Partenze, virate e subacque nei quattro stili (la subacqua è
            considerata il \textbf{quinto stile})
      \item
            Tuffo del salvamento
\end{itemize}

\textbf{Obiettivi per la sicurezza:}

\begin{itemize}
      \tightlist
      \item
            Immergersi a una profondità di almeno \textbf{1,60 metri}
      \item
            Spostamenti laterali in galleggiamento verticale a \textbf{bicicletta}
      \item
            25 metri di trasporto in presa da crampi
      \item
            Nuotata di avvicinamento: 25 metri a \emph{crawl} testa alta e 25
            metri a rana testa alta
      \item
            Remate stazionarie di avanzamento
      \item
            Elementi di primo soccorso
\end{itemize}

\textbf{Esercizi propedeutici al 4° livello:}

\begin{itemize}
      \tightlist
      \item
            Tuffo e 25 metri con trasporto di pericolante (con crampi)
      \item
            Nuotata di avvicinamento al pericolante
\end{itemize}

\begin{center}\rule{0.5\linewidth}{0.5pt}\end{center}

\subsubsection{\texorpdfstring{Seconda Fase - Livello Avanzato
            (\emph{SaNuotare
                  2})}{Seconda Fase - Livello Avanzato (SaNuotare 2)}}\label{seconda-fase---livello-avanzato-sanuotare-2}

\textbf{Requisito di ingresso:} Conseguimento del \textbf{5° livello}
della Scuola Nuoto Federale (\emph{competenza motoria acquatica}).

\textbf{Frequenza:} 4-5 sedute settimanali da 60 a 90 minuti.

Il 5° livello equivale al brevetto di \emph{SaNuotare 2}: il soggetto sa
nuotare, muoversi, vivere in acqua e può essere \textbf{di aiuto a
      un\textquotesingle altra persona in difficoltà}. Dà accesso, al
raggiungimento dei \textbf{16 anni}, all\textquotesingle iscrizione a un
corso di \textbf{assistente bagnante}.

\textbf{Obiettivi tecnici:}

\begin{itemize}
      \tightlist
      \item
            Esercitazioni per la \textbf{resistenza aerobica}
      \item
            Prove cronometrate nei \textbf{4 stili}
      \item
            Nuotate subacque (distinzione tra rana subacquea e avanzamento con
            flesso-estensione)
      \item
            Utilizzo di \textbf{snorkel} e \textbf{pinne}
      \item
            \textbf{Motricità acquatica verticale evoluta} (bicicletta in tutte le
            varianti, compresa la \emph{bicicletta a pressione}, fondamentale per
            la pala nuoto)
\end{itemize}

\textbf{Obiettivi per la sicurezza:}

\begin{itemize}
      \tightlist
      \item
            25 metri di trasporto ascellare
      \item
            25 metri di trasporto al capo
      \item
            25 metri vestiti
      \item
            25 metri a \emph{trudgeon}
      \item
            \emph{Over} e \emph{dorso gambe a rana} (elementi da proporre anche
            per il \emph{VLS})
\end{itemize}

\textbf{Propedeutici al 5° livello:}

\begin{itemize}
      \tightlist
      \item
            25 metri vestiti
      \item
            Nuotate subacque e utilizzo di snorkel e pinne
      \item
            25 metri di trasporto di pericolante non cosciente
\end{itemize}

\begin{center}\rule{0.5\linewidth}{0.5pt}\end{center}

\subsection{\texorpdfstring{Le Tecniche Specifiche: \emph{Over} e
            \emph{Trudgeon}}{Le Tecniche Specifiche: Over e Trudgeon}}\label{le-tecniche-specifiche-over-e-trudgeon}

\subsubsection{\texorpdfstring{Il
            \emph{Trudgeon}}{Il Trudgeon}}\label{il-trudgeon}

Il \emph{trudgeon} è un \textbf{fondamentale sia per il salvamento che
      per la pala nuoto}. Si distingue dal \emph{crawl} per:

\begin{itemize}
      \tightlist
      \item
            L\textquotesingle azione delle \textbf{gambe non è alternata}
            (flesso-estensione) ma \textbf{simmetrica e simultanea}
      \item
            Le gambe lavorano a \textbf{forbice}: una rimane
            semiflessa-semidistesa, l\textquotesingle altra esegue una flessione
            tipo rana
      \item
            L\textquotesingle azione degli \textbf{arti superiori} è simile allo
            stile libero
\end{itemize}

\subsubsection{\texorpdfstring{L\textquotesingle{}\emph{Over}
            (\emph{Overarm Side
                  Stroke})}{L\textquotesingle Over (Overarm Side Stroke)}}\label{lover-overarm-side-stroke}

L\textquotesingle{}\emph{over} è propedeutico al \emph{trudgeon}.
Caratteristiche:

\begin{itemize}
      \tightlist
      \item
            Il \textbf{corpo è inclinato} rispetto alla direzione di avanzamento
            (una parte più immersa dell\textquotesingle altra)
      \item
            Il \textbf{braccio più profondo} lavora in \textbf{trazione}
      \item
            Il \textbf{braccio più superficiale} lavora in \textbf{spinta}
      \item
            Le \textbf{gambe} agiscono con un\textquotesingle azione simile alla
            rana, adattata all\textquotesingle inclinazione del corpo
      \item
            Nell\textquotesingle{}\textbf{overarm side stroke} le braccia vengono
            recuperate \textbf{fuori dall\textquotesingle acqua}
\end{itemize}

\begin{center}\rule{0.5\linewidth}{0.5pt}\end{center}

\subsection{I Fondamentali Tecnici per il Salvamento in
      Corsia}\label{i-fondamentali-tecnici-per-il-salvamento-in-corsia}

Provenienti prevalentemente dalla \textbf{pala nuoto}:

\begin{itemize}
      \tightlist
      \item
            \emph{Crawl} classico e \emph{crawl} testa alta
      \item
            \emph{Over} e \emph{overarm side stroke}
      \item
            Sedia avanti e sedia indietro
      \item
            Passaggio libero-dorso (\emph{crawl} con transizione al dorso)
      \item
            Rana rovesciata con braccio a dorso (fondamentale nel trasporto
            dell\textquotesingle inerte dall\textquotesingle assistente bagnante)
      \item
            \emph{Trudgeon}
      \item
            \textbf{Bicicletta} e tutte le sue varianti:

            \begin{itemize}
                  \tightlist
                  \item
                        Bicicletta avanti / indietro
                  \item
                        Bicicletta a pressione
                  \item
                        Bicicletta di fronte / di fianco
                  \item
                        Dorso seduto
                  \item
                        Rana rovesciata
            \end{itemize}
\end{itemize}

\subsubsection{Fondamentali di
      Salvataggio}\label{fondamentali-di-salvataggio}

\begin{itemize}
      \tightlist
      \item
            \textbf{Tecniche di avvicinamento} al pericolante
      \item
            \textbf{Tecniche di allontanamento}
      \item
            \textbf{Nuoto subacqueo}
      \item
            \textbf{Nuoto in corrente} (favorevole e avversa, per le gare
            oceaniche)
\end{itemize}

\subsubsection{Fondamentali di
      Trasporto}\label{fondamentali-di-trasporto}

\begin{longtable}[]{@{}p{0.47\linewidth}p{0.47\linewidth}@{}}
      \toprule\noalign{}
      Presa                                                   & Contesto                                   \\
      \midrule\noalign{}
      \endhead
      \bottomrule\noalign{}
      \endlastfoot
      Presa ascellare                                         & Trasporto soggetto cosciente/collaborativo \\
      Presa al capo                                           & Trasporto soggetto non cosciente           \\
      Prese collaborative                                     & Soggetto pienamente collaborativo          \\
      Prese di emergenza (mezza \emph{Nelson}, \emph{Nelson}) & Situazioni di
      emergenza; preparazione ai corsi futuri di assistente bagnante                                       \\
\end{longtable}

\subsubsection{Fondamentali di
      Gareggiamento}\label{fondamentali-di-gareggiamento}

\begin{itemize}
      \tightlist
      \item
            Nuoto \textbf{prono}, \textbf{supino} e \textbf{in verticale} (remate)
      \item
            Posizionamenti, spostamenti, bilanciamenti (sedie avanti e indietro)
      \item
            Tecniche di \textbf{immersione}: davanti agli ostacoli e nel punto di
            recupero del manichino
      \item
            Discese in acqua controllata, in acqua profonda e discese sul fondo
\end{itemize}

\begin{center}\rule{0.5\linewidth}{0.5pt}\end{center}

\subsection{Programmi e Progressione per
      Età}\label{programmi-e-progressione-per-etuxe0}

\subsubsection{Prima Fase - Esordienti (circa 12-13 anni femmine, 13
      anni
      maschi)}\label{prima-fase---esordienti-circa-12-13-anni-femmine-13-anni-maschi}

\textbf{Obiettivo generale:} Apprendere i fondamentali in modo
divertente.

\begin{itemize}
      \tightlist
      \item
            Migliorare le \textbf{capacità motorie} (forza, velocità, resistenza,
            capacità coordinative, mobilità articolare)
      \item
            Costruire le \textbf{abilità motorie} (azioni motorie automatizzate
            tramite ripetizione)
      \item
            Privilegiare \textbf{velocità} e \textbf{resistenza di lunga durata e
                  bassa intensità}
      \item
            Sviluppo totale delle capacità coordinative
      \item
            Non dimenticare le abilità motorie di tipo diverso (corsa, salto,
            lancio) utili anche per le gare oceaniche
      \item
            Periodo critico e molto favorevole per la \textbf{mobilità
                  articolare}: va mantenuta attivamente
      \item
            Frequenza: \textbf{4-5 allenamenti settimanali}
      \item
            Obiettivo: \textbf{competere, non vincere}; organizzare piccole gare
            propedeutiche in corsia
      \item
            Sviluppo integrato delle componenti mentali, cognitive, emotive e
            affettive
      \item
            Miglioramento dell\textquotesingle{}\textbf{etica sportiva}; evitare
            periodizzazioni rigide
\end{itemize}

\subsubsection{Seconda Fase - Ragazzi (14-15 anni maschi, 13-14 anni
      femmine)}\label{seconda-fase---ragazzi-14-15-anni-maschi-13-14-anni-femmine}

\textbf{Obiettivo generale:} Allenarsi alla gara.

\begin{itemize}
      \tightlist
      \item
            Costruire una \textbf{base aerobica}
      \item
            Sviluppare \textbf{forza} e \textbf{velocità}
      \item
            \textbf{Consolidare} le abilità motorie fondamentali
      \item
            Inserire elementi \textbf{tattici}: distribuzione razionale
            dell\textquotesingle impegno fisico per ottimizzare il risultato e
            condizionare l\textquotesingle avversario
      \item
            Riservare parte dell\textquotesingle allenamento alla
            \textbf{competizione} (allenamento mentale)
      \item
            Avvio di una \textbf{periodizzazione} con gare obiettivo
      \item
            Frequenza: fino a \textbf{6 allenamenti settimanali}
\end{itemize}

\begin{center}\rule{0.5\linewidth}{0.5pt}\end{center}

\subsection{Il Collegamento con la Scuola Nuoto
      Federale}\label{il-collegamento-con-la-scuola-nuoto-federale}

Il nuoto per salvamento si integra naturalmente con la Scuola Nuoto
Federale ai \textbf{livelli 5 e 6}, che prevedono esercizi
progressivamente crescenti. Al termine del 5° livello, il soggetto ha
sviluppato un patrimonio motorio che lo rende idoneo
all\textquotesingle avviamento verso:

\begin{itemize}
      \tightlist
      \item
            \textbf{Nuoto agonistico}
      \item
            \textbf{Nuoto per salvamento}
      \item
            \textbf{Nuoto in acque libere}
      \item
            \textbf{Pala nuoto}
      \item
            \textbf{Nuoto sincronizzato} (nuoto artistico)
\end{itemize}

I \textbf{fondamentali trasversali} - come il \emph{trudgeon} e la
bicicletta - sono comuni sia al salvamento che alla pala nuoto,
confermando l\textquotesingle interconnessione tra le discipline FIN.

\begin{center}\rule{0.5\linewidth}{0.5pt}\end{center}

\subsection{Conclusione - Punti Chiave per
      l\textquotesingle Esame}\label{conclusione---punti-chiave-per-lesame-4}

\begin{itemize}
      \tightlist
      \item
            Il nuoto per salvamento è la \textbf{quinta disciplina FIN}, con gare
            in piscina e in mare.
      \item
            Il modello funzionale del nuotatore si divide in \textbf{tre livelli}:
            psicologico, tecnico-coordinativo e fisico-condizionale.
      \item
            Le \textbf{apparecchiature} principali sono: manichino, ostacoli
            (profondità 70 cm), torpedo e pinne.
      \item
            Il corso in corsia si sviluppa in \textbf{due fasi}: la prima a
            partire dal 4° livello (\emph{SaNuotare 1}), la seconda dal 5° livello
            (\emph{SaNuotare 2}).
      \item
            Le tecniche chiave sono: \emph{crawl} testa alta, \emph{over},
            \emph{overarm side stroke} e \textbf{trudgeon} (gambe a forbice
            simmetrica).
      \item
            La \textbf{bicicletta} verticale è un fondamentale sia per il
            salvamento che per la pala nuoto.
      \item
            La progressione didattica rispetta l\textquotesingle{}\textbf{età
                  biologica}, non quella cronologica, e privilegia il divertimento nelle
            fasi giovanili.
      \item
            Nei ragazzi si introduce la \textbf{periodizzazione} e la
            \textbf{tattica di gara} (distribuzione delle forze).
      \item
            Il raggiungimento dei \textbf{16 anni} consente
            l\textquotesingle accesso al corso di \textbf{assistente bagnante}.
\end{itemize}

\section{19. Cenni di Attività per
  Gestanti}\label{cenni-di-attivituxe0-per-gestanti}

\subsection{Introduzione}\label{introduzione-18}

La \textbf{gravidanza non è una malattia}, ma un fatto fisiologico che
va vissuto con gli strumenti psicofisici adatti. È un periodo durante il
quale la gestante riacquista il rapporto con se stessa e con il
nascituro. L\textquotesingle attività acquatica in gravidanza nasce
dall\textquotesingle intuizione di tre specialisti (una ginecologa, un
medico sportivo e un chirurgo vascolare) convinti che
l\textquotesingle acqua potesse prevenire e talvolta curare determinate
problematiche della gestante. Dopo un anno di studi - incluse prove con
ecografo subacqueo e collaborazioni con cattedre universitarie di
medicina dello sport e clinica ostetrica - la metodica ha ricevuto una
\textbf{validazione scientifica nel 1990}, con relativa pubblicazione
nazionale della medicina dello sport.

\begin{center}\rule{0.5\linewidth}{0.5pt}\end{center}

\subsection{Il Ruolo della FIN e
      dell\textquotesingle Istruttore}\label{il-ruolo-della-fin-e-dellistruttore}

Il settore istruzione tecnica della FIN si interessa a questa attività
proprio perché spesso viene svolta \textbf{senza nozioni specifiche} e
senza finalizzazioni precise degli esercizi. L\textquotesingle obiettivo
è che i tecnici acquisiscano una conoscenza fisiologica e medica di base
sufficiente a:

\begin{itemize}
      \tightlist
      \item
            Comprendere i meccanismi dello stato di gravidanza
      \item
            Distinguere le differenze rispetto allo stato non gravidico
      \item
            \textbf{Variare e modificare il programma} del corso con
            consapevolezza e senza paura
\end{itemize}

Questa attività \textbf{non si sostituisce} ai corsi di preparazione al
parto degli ospedali (che forniscono informazioni scientifiche e
conoscenza della struttura), ma ne è
un\textquotesingle{}\textbf{integrazione attiva}: la gestante arriva al
parto allenata - tra virgolette - e con una buona conoscenza del proprio
corpo.

\begin{center}\rule{0.5\linewidth}{0.5pt}\end{center}

\subsection{Perché l\textquotesingle Acqua}\label{perchuxe9-lacqua}

L\textquotesingle acqua ripristina l\textquotesingle equilibrio
psicofisico compromesso dai ritmi della quotidianità. Immergervisi
richiama inconsciamente le emozioni delle origini. In particolare:

\begin{itemize}
      \tightlist
      \item
            \textbf{Massaggia}, culla e rilassa
      \item
            Libera \textbf{endorfine}, aumentando il benessere
      \item
            Scarica le energie negative
      \item
            Offre \textbf{privacy} alla gestante: l\textquotesingle immersione
            copre il corpo, facilitando chi desidera muoversi senza sentirsi
            osservata (vantaggio rispetto al lavoro a secco)
\end{itemize}

\begin{center}\rule{0.5\linewidth}{0.5pt}\end{center}

\subsection{Effetti della Gravidanza per
      Trimestre}\label{effetti-della-gravidanza-per-trimestre}

\subsubsection{Primo Trimestre}\label{primo-trimestre}

\begin{itemize}
      \tightlist
      \item
            Momenti di incertezza e \textbf{ambivalenza emotiva}
      \item
            La gioia si alterna a dubbi sull\textquotesingle adeguatezza e paure
            sui cambiamenti
      \item
            Oscillazione tra accettazione e rifiuto della nuova condizione
      \item
            Paura di non tornare a come si era prima
\end{itemize}

\subsubsection{Secondo Trimestre}\label{secondo-trimestre}

\begin{itemize}
      \tightlist
      \item
            La donna vive in \textbf{simbiosi con il bambino}
      \item
            Vive piacevolmente il proprio stato
      \item
            Trova \textbf{sollievo in acqua}, con operatori competenti, in un
            ambiente tranquillo e protettivo
\end{itemize}

\subsubsection{Terzo Trimestre}\label{terzo-trimestre}

\begin{itemize}
      \tightlist
      \item
            Caratterizzato da \textbf{ansia, insofferenza e stanchezza}
      \item
            Preoccupazione per l\textquotesingle organizzazione di ogni aspetto
            materiale e non
      \item
            L\textquotesingle acqua aiuta ad acquisire \textbf{relax e benessere
                  mentale}, a scaricare le energie negative
\end{itemize}

\begin{center}\rule{0.5\linewidth}{0.5pt}\end{center}

\subsection{Benefici dell\textquotesingle Attività Fisica in
      Gravidanza}\label{benefici-dellattivituxe0-fisica-in-gravidanza}

\subsubsection{Benefici Fisici}\label{benefici-fisici}

\begin{itemize}
      \tightlist
      \item
            Mantiene il \textbf{tono muscolare}, inclusi i muscoli paravertebrali
            che sostengono il peso del bambino
      \item
            Previene gonfiori e pesantezza agli arti inferiori
      \item
            Previene rischi di traumi, tendiniti e distorsioni
      \item
            Tiene il \textbf{peso sotto controllo}
      \item
            \textbf{Migliore tolleranza allo sforzo}: in acqua
            l\textquotesingle{}\emph{emodiluizione} compensa
            l\textquotesingle{}\emph{emoconcentrazione da sforzo} che si verifica
            a secco, grazie anche all\textquotesingle aumento dei volumi
            plasmatici e della gettata sistolica
\end{itemize}

\subsubsection{Effetti Fisiologici Specifici
      dell\textquotesingle Immersione}\label{effetti-fisiologici-specifici-dellimmersione}

\begin{itemize}
      \tightlist
      \item
            Incremento del \textbf{ritorno venoso} del 60\%
      \item
            Incremento della \textbf{capacità inspiratoria}
      \item
            Buona \textbf{frequenza respiratoria}
      \item
            \textbf{Aumento della diuresi}
      \item
            Rapida espansione del \textbf{volume plasmatico}, determinata dalla
            pressione idrostatica dell\textquotesingle acqua sulle pareti del
            corpo (l\textquotesingle acqua funziona come una \emph{calza
                  elastica}): i liquidi extravascolari vengono compressi nello spazio
            vascolare
      \item
            \textbf{Stabilità della temperatura corporea}
      \item
            Assenza di variazioni di pressione arteriosa in funzione della
            posizione (che invece si verificano a secco)
      \item
            \textbf{Riduzione del peso addominale} su rachide, colonna vertebrale
            e arti inferiori
      \item
            Miglioramento del \textbf{controllo della respirazione}
\end{itemize}

\subsubsection{Benefici Psicologici e
      Generali}\label{benefici-psicologici-e-generali}

\begin{itemize}
      \tightlist
      \item
            Stimola il rilascio di \textbf{ormoni} che favoriscono il buon umore
            della madre e, di conseguenza, del feto
      \item
            Fisico più allenato e maggiore \textbf{agilità}
            nell\textquotesingle affrontare le fatiche del parto
      \item
            Miglioramento dell\textquotesingle{}\textbf{acquaticità individuale},
            utile per il parto in acqua e per la trasmissione
            dell\textquotesingle empatia acquatica al futuro bambino
      \item
            Emozioni positive dal superamento di esperienze nuove e divertenti
\end{itemize}

\begin{quote}
      Tutta l\textquotesingle attività si svolge in \textbf{modalità
            aerobica}.
\end{quote}

\begin{center}\rule{0.5\linewidth}{0.5pt}\end{center}

\subsection{Modalità di Esecuzione: Posizione Verticale vs.
      Orizzontale}\label{modalituxe0-di-esecuzione-posizione-verticale-vs.-orizzontale}

Il lavoro aerobico in acqua per le gestanti è prevalentemente in
\textbf{posizione verticale}, per sfruttare al massimo la pressione
idrostatica. La normale nuotata in posizione orizzontale produce una
minore ridistribuzione del volume sanguigno a livello centrale e una
\textbf{forza idrostatica ridotta} rispetto alla posizione verticale.

\begin{center}\rule{0.5\linewidth}{0.5pt}\end{center}

\subsection{Disturbi Comuni e
      Gestione}\label{disturbi-comuni-e-gestione}

\begin{longtable}[]{@{}p{0.47\linewidth}p{0.47\linewidth}@{}}
      \toprule\noalign{}
      Disturbo                                    & Indicazione                                           \\
      \midrule\noalign{}
      \endhead
      \bottomrule\noalign{}
      \endlastfoot
      \textbf{Lombalgia} (mal di schiena)         & Preferire esercizi di allungamento
      vertebrale o che distanzino le vertebre lombari; evitare addominali e
      gambe tese                                                                                          \\
      \textbf{Crampi}                             & Diminuire il carico di lavoro; eventuale consulto del
      ginecologo; l\textquotesingle attività può continuare                                               \\
      \textbf{Edemi agli arti inferiori / varici} & L\textquotesingle acqua è
      consigliata dall\textquotesingle angiologo come vera e propria
      \textbf{terapia}                                                                                    \\
      \textbf{Perdite vaginali di sangue}         & \textbf{Sospendere immediatamente}
      l\textquotesingle attività; richiedere nulla osta del ginecologo per
      riprendere                                                                                          \\
      \textbf{Ipertensione}                       & Attività \textbf{sconsigliata}: la gestante è
      già sottoposta a uno sforzo che potrebbe portare a ricovero
      ospedaliero                                                                                         \\
      \textbf{Contrazioni uterine}                & Evitare esercizi addominali; dopo la 32ª
      settimana sono normali gli indurimenti uterini; se non frequenti,
      l\textquotesingle attività può continuare                                                           \\
      \textbf{Problemi di crescita fetale}        & \textbf{Interrompere}
      l\textquotesingle attività: ossigeno e nutrimento devono essere
      riservati esclusivamente al bambino                                                                 \\
      \textbf{Diabete materno}                    & Attività consentita e ottimale (utilizzo del
      glucosio da parte della muscolatura); avvertire il \textbf{diabetologo}
      per dosare farmaci e dieta nei giorni di attività                                                   \\
\end{longtable}

\begin{center}\rule{0.5\linewidth}{0.5pt}\end{center}

\subsection{Finalità dell\textquotesingle Attività
      Acquatica}\label{finalituxe0-dellattivituxe0-acquatica}

\begin{itemize}
      \tightlist
      \item
            Mantenimento del \textbf{tono muscolare} a supporto della colonna
            vertebrale
      \item
            \textbf{Consapevolezza e controllo della respirazione} (non le stesse
            tecniche dell\textquotesingle ospedale, ma un aiuto a padroneggiare il
            respiro)
      \item
            \textbf{Controllo dei movimenti} e sviluppo della
            \textbf{propriocezione}
\end{itemize}

\begin{quote}
      \emph{Propriocezione}: la capacità di percepire e riconoscere la
      posizione del corpo nello spazio, contraendo e decontraendo la
      muscolatura anche senza vederla. Fondamentale per la
      \textbf{consapevolezza del movimento} durante il parto.
\end{quote}

\begin{itemize}
      \tightlist
      \item
            \textbf{Percezione e mobilitazione del bacino}: spesso la gestante non
            distingue i movimenti del bacino da quelli delle gambe
      \item
            \textbf{Controllo del perineo}
      \item
            \textbf{Controllo motorio settoriale}: arti superiori e arti inferiori
            possono eseguire azioni indipendenti
      \item
            Cura delle \emph{trombosi venose} e della \emph{stasi linfatica}
\end{itemize}

\begin{center}\rule{0.5\linewidth}{0.5pt}\end{center}

\subsection{Programmazione dell\textquotesingle Attività
      Acquatica}\label{programmazione-dellattivituxe0-acquatica}

\subsubsection{Struttura di una Seduta}\label{struttura-di-una-seduta}

\begin{itemize}
      \tightlist
      \item
            I \textbf{primi 15-20 minuti} dedicati alle \textbf{grandi masse
                  muscolari} (gambe), per poi passare ad altri obiettivi
      \item
            Per ogni serie di esercizi (durata 10-15 minuti), la \textbf{frequenza
                  cardiaca} deve aumentare gradualmente fino a circa \textbf{140
                  battiti/minuto} per garantire gli effetti benefici descritti
      \item
            Il lavoro è in \textbf{continuo aerobico}, con variazioni nella
            tipologia degli esercizi
\end{itemize}

\subsubsection{Tipologie di Esercizi}\label{tipologie-di-esercizi}

\textbf{Rilassamento progressivo e settoriale:}

\begin{itemize}
      \tightlist
      \item
            Ricerca della percezione del proprio corpo
      \item
            Utilizzo del \textbf{tubo} (attrezzo) come supporto alla gestante
      \item
            Ricerca dell\textquotesingle equilibrio e padronanza della muscolatura
\end{itemize}

\textbf{Mobilitazione del bacino:}

\begin{itemize}
      \tightlist
      \item
            Spostamenti verso l\textquotesingle alto, a sinistra, a destra
      \item
            Tutti gli esercizi di percezione del bacino
\end{itemize}

\textbf{Controllo della respirazione:}

\begin{itemize}
      \tightlist
      \item
            Inspirazioni ed espirazioni guidate
      \item
            Utilizzo del respiro per spostarsi nell\textquotesingle acqua
      \item
            Espirazione con emissione di suoni
      \item
            L\textquotesingle obiettivo non è replicare le respirazioni
            ospedaliere, ma insegnare alla gestante a usare il respiro come
            preferisce
\end{itemize}

\textbf{\emph{Stretching}:}

\begin{itemize}
      \tightlist
      \item
            Migliora l\textquotesingle estendibilità muscolare e la mobilità
            articolare
      \item
            La muscolatura da allungare deve essere sempre \textbf{vascolarizzata}
            dal movimento che precede l\textquotesingle allungamento (stesso
            gruppo muscolare)
\end{itemize}

\textbf{Esercizi al bordo (per chi non galleggia o non si sente a
      proprio agio):}

\begin{itemize}
      \tightlist
      \item
            Tutti gli esercizi possono essere adattati al bordo vasca
      \item
            Massima varietà e creatività
\end{itemize}

\textbf{Attrezzi e strumenti \emph{fitness}:}

\begin{itemize}
      \tightlist
      \item
            \emph{Tapirulan} in acqua
      \item
            Camminate (senza appoggio, con appoggio, frontali, laterali)
      \item
            Esercizi di \emph{step}
      \item
            Stretching
      \item
            Spinte dal tappeto
\end{itemize}

\textbf{Attività creative (esempio):}

\begin{itemize}
      \tightlist
      \item
            Immersioni con erogatore per l\textquotesingle adattamento
            all\textquotesingle ambiente
      \item
            Respirazioni di vario tipo per acquisire tranquillità
            sott\textquotesingle acqua
      \item
            Esercizi di \emph{Tai Chi} subacquei (movimenti lenti con respirazione
            abbinata)
\end{itemize}

\begin{center}\rule{0.5\linewidth}{0.5pt}\end{center}

\subsection{Quando Iniziare e Quando
      Fermarsi}\label{quando-iniziare-e-quando-fermarsi}

\begin{itemize}
      \tightlist
      \item
            L\textquotesingle attività acquatica in gravidanza è aperta a
            \textbf{qualsiasi epoca gestazionale}, sempre con il consenso del
            ginecologo
      \item
            Può proseguire fino al \textbf{termine del periodo gestazionale},
            compreso il \textbf{pretravaglio}
      \item
            In caso di pretravaglio (contrazioni): controllare attentamente la
            situazione
      \item
            Alcune istrici hanno continuato a lavorare anche durante la propria
            gravidanza
\end{itemize}

\begin{center}\rule{0.5\linewidth}{0.5pt}\end{center}

\subsection{Gestione Organizzativa del
      Corso}\label{gestione-organizzativa-del-corso}

\subsubsection{Prima di Avviare il
      Corso}\label{prima-di-avviare-il-corso}

\begin{itemize}
      \tightlist
      \item
            Stabilire \textbf{contatti con specialisti} di ostetricia e
            ginecologia
      \item
            Spiegare obiettivi e programma del corso
      \item
            Costruire un rapporto di \textbf{collaborazione e integrazione} senza
            invadere i campi specifici altrui
\end{itemize}

\subsubsection{Iscrizione e
      Documentazione}\label{iscrizione-e-documentazione}

\begin{itemize}
      \tightlist
      \item
            La segreteria fornisce una spiegazione chiara di: obiettivi, modalità
            e consigli di comportamento
      \item
            È obbligatorio richiedere il \textbf{certificato medico del
                  ginecologo} (non del medico curante)
\end{itemize}

\subsubsection{Prima Lezione in Vasca}\label{prima-lezione-in-vasca}

\begin{itemize}
      \tightlist
      \item
            Spiegare gli obiettivi in modo chiaro e accessibile
      \item
            \textbf{Intervistare la gestante}: conoscere la sua storia
            gestazionale e le eventuali problematiche per adattare il programma
\end{itemize}

\begin{center}\rule{0.5\linewidth}{0.5pt}\end{center}

\subsection{Il Buon Istruttore per Questa
      Attività}\label{il-buon-istruttore-per-questa-attivituxe0}

Un istruttore efficace in questo contesto deve:

\begin{itemize}
      \tightlist
      \item
            \textbf{Provare personalmente} gli esercizi (almeno ogni tre mesi) per
            comprendere realmente ciò che insegna e i limiti da rispettare
      \item
            Mantenersi in \textbf{aggiornamento continuo} attraverso corsi di
            specializzazione
      \item
            Non fossilizzarsi sui primi esercizi appresi: aggiornarsi
            costantemente
      \item
            Essere \textbf{propenso ad imparare} da ogni tipo di esperienza
      \item
            Essere sempre \textbf{disponibile} e accettare consigli
      \item
            Portare in vasca \textbf{buon umore}, rendendo l\textquotesingle ora
            di lezione un momento utile, rilassante e piacevole
\end{itemize}

\begin{center}\rule{0.5\linewidth}{0.5pt}\end{center}

\subsection{Conclusione - Punti Chiave per
      l\textquotesingle Esame}\label{conclusione---punti-chiave-per-lesame-5}

\begin{itemize}
      \tightlist
      \item
            La gravidanza è un \textbf{fatto fisiologico}, non una malattia;
            l\textquotesingle attività acquatica ne è
            un\textquotesingle integrazione, non una sostituzione dei corsi
            ospedalieri.
      \item
            La metodica ha ricevuto \textbf{validazione scientifica nel 1990} con
            pubblicazione nazionale; le prime parì partecipanti mostravano lavoro
            più breve, migliore sopportazione del dolore e parto più consapevole.
      \item
            Il lavoro si svolge prevalentemente in \textbf{posizione verticale}
            per massimizzare la pressione idrostatica.
      \item
            Il \textbf{ritorno venoso in acqua} aumenta del \textbf{60\%};
            l\textquotesingle acqua agisce come una \emph{calza elastica} sui
            liquidi extravascolari.
      \item
            La \textbf{frequenza cardiaca target} durante gli esercizi è circa
            \textbf{140 battiti/minuto}.
      \item
            I disturbi che impongono la \textbf{sospensione}
            dell\textquotesingle attività sono: perdite vaginali di sangue,
            ipertensione e problemi di crescita fetale.
      \item
            Il \textbf{diabete materno} non controindica
            l\textquotesingle attività (ottimale per il metabolismo del glucosio),
            ma richiede il coordinamento con il diabetologo.
      \item
            Fondamentale è la \textbf{propriocezione}: la capacità di percepire e
            controllare il corpo senza vederlo, utile al momento del parto.
      \item
            Il certificato richiesto è quello del \textbf{ginecologo}, non del
            medico curante.
      \item
            La seduta inizia sempre con i \textbf{grandi gruppi muscolari} (gambe)
            per i primi 15-20 minuti.
\end{itemize}

\section{20. Cenni di Attività
  Neonatale}\label{cenni-di-attivituxe0-neonatale}

\subsection{Introduzione}\label{introduzione-19}

L\textquotesingle attività neonatale in acqua è spesso denominata
\emph{nuoto baby}, ma la definizione più corretta è \textbf{corsi di
      acquaticità per neonati}. Il termine \emph{nuoto} implica un
insegnamento tecnico che non appartiene a questa fascia
d\textquotesingle età: si tratta invece di stimolare e coltivare
un\textquotesingle{}\textbf{acquaticità innata}, intesa come il
complesso delle azioni e reazioni che indicano una condizione di
serenità e agiatezza nell\textquotesingle acqua.

Questa acquaticità è naturalmente presente nel neonato, ma va stimolata
continuamente, tenendo conto che \textbf{ogni bambino è diverso}: le
differenze dipendono da fattori ambientali (prenatali, neonatali,
familiari, sociali) e genetici.

\begin{center}\rule{0.5\linewidth}{0.5pt}\end{center}

\subsection{Il Gioco come Strumento
      Fondamentale}\label{il-gioco-come-strumento-fondamentale}

Il gioco è l\textquotesingle elemento centrale di tutta
l\textquotesingle attività neonatale. Può essere utilizzato in tre modi:

\begin{itemize}
      \tightlist
      \item
            \textbf{Meccanismo di maturazione e sviluppo}: sfrutta le potenzialità
            psicologiche
      \item
            \textbf{Fonte di piacere}: trasforma le cariche energetiche e sviluppa
            le potenzialità affettive
      \item
            \textbf{Fonte di conoscenza}: favorisce l\textquotesingle interazione
            con altri bambini e sviluppa le potenzialità cognitive
\end{itemize}

Attraverso il gioco il bambino:

\begin{itemize}
      \tightlist
      \item
            Esplora l\textquotesingle ambiente e agisce in esso
      \item
            Sperimenta emozioni positive e negative
      \item
            Cerca il proprio modo di essere libero in acqua
      \item
            Socializza con altri bambini
\end{itemize}

\begin{quote}
      Il gioco deve avere delle \textbf{regole}, altrimenti non è più gioco ma
      caos.
\end{quote}

\begin{center}\rule{0.5\linewidth}{0.5pt}\end{center}

\subsection{Importanza dell\textquotesingle Acquaticità
      Neonatale}\label{importanza-dellacquaticituxe0-neonatale}

L\textquotesingle ambiente acquatico offre al bambino vantaggi unici:

\begin{itemize}
      \tightlist
      \item
            \textbf{Quasi totale assenza di gravità}: il bambino si sente libero e
            può compiere movimenti più complessi
      \item
            \textbf{Maggiore percezione del proprio corpo}: sviluppa più
            velocemente la consapevolezza degli arti e acquisisce schemi motori
            basilari
      \item
            \textbf{Stimolazione sensoriale innata}: riscoprire
            l\textquotesingle acqua attraverso il gioco attiva percezioni
            sensoriali naturali
      \item
            \textbf{Sviluppo emotivo e corporeo}: ogni stimolazione motoria
            contribuisce a una maggiore consapevolezza, favorendo una relazione
            con il mondo esterno più serena ed equilibrata
      \item
            \textbf{Scarico delle tensioni}: i giochi con l\textquotesingle acqua
            (lanciare, versare, mescolare, immergere) vengono utilizzati anche per
            scaricare le tensioni del bambino
\end{itemize}

L\textquotesingle acqua culla, avvolge, carezza e porta gioia,
rilassamento e appagamento. Il bambino la conosce prima di tutto con la
\textbf{bocca}, primo strumento di esplorazione.

La \textbf{motricità fine} cresce contemporaneamente alla maturazione e
al collegamento causa-effetto. In acqua si sviluppano:

\begin{itemize}
      \tightlist
      \item
            Capacità di equilibrio e controllo posturale
      \item
            Organizzazione sequenziale delle azioni verso un obiettivo
\end{itemize}

\begin{center}\rule{0.5\linewidth}{0.5pt}\end{center}

\subsection{Quando Iniziare}\label{quando-iniziare}

\begin{itemize}
      \tightlist
      \item
            L\textquotesingle inizio dell\textquotesingle attività è \textbf{a
                  discrezione del neonatologo}
      \item
            Corrisponde generalmente \textbf{alle prime vaccinazioni}, intorno ai
            \textbf{3 mesi} circa
      \item
            Più è precoce l\textquotesingle inizio, più diventa piacevole: si
            ripristinano meccanismi corporei e neurovegetativi innati
      \item
            Più si lascia passare il tempo, maggiore sarà lo sforzo necessario per
            riacquistare capacità respiratorie e adattative già presenti alla
            nascita
\end{itemize}

\begin{quote}
      Non esiste una divisione netta per età: c\textquotesingle è
      \textbf{grande variabilità} nelle modalità e nei tempi delle
      acquisizioni motorie, indipendente dall\textquotesingle età anagrafica.
\end{quote}

\begin{center}\rule{0.5\linewidth}{0.5pt}\end{center}

\subsection{Fasi di Sviluppo per Fascia
      d\textquotesingle Età}\label{fasi-di-sviluppo-per-fascia-detuxe0}

\subsubsection{3-10 Mesi}\label{mesi}

\begin{itemize}
      \tightlist
      \item
            Il bambino \textbf{non galleggia ancora} in modo autonomo (galleggiare
            = trovare un punto di appoggio nell\textquotesingle acqua per ogni
            parte del corpo)
      \item
            Le \textbf{apnee} inizialmente sono guidate; il \emph{diving reflex}
            si perde gradualmente
      \item
            Il movimento degli arti inferiori è inizialmente \textbf{passivo o
                  riflesso}
      \item
            La presa è inizialmente \emph{cubito-palmare} (a tutto il palmo), poi
            diventa mirata al recupero di piccoli oggetti
      \item
            La \textbf{posizione dorsale} è un buon punto di partenza: favorisce
            la respirazione (non sempre accettata dal bambino)
      \item
            Il bambino è in costante ricerca della \textbf{base sicura} (mamma,
            nonno, zia, persona familiare)
      \item
            Curiosità intensa per tutto ciò che lo circonda
\end{itemize}

\subsubsection{10-18 Mesi}\label{mesi-1}

\begin{itemize}
      \tightlist
      \item
            Capisce quando deve affrontare l\textquotesingle apnea
      \item
            Si avvia l\textquotesingle{}\textbf{educazione
                  all\textquotesingle immersione} con la progressiva perdita del
            \emph{diving reflex}
      \item
            I movimenti degli arti superiori e inferiori diventano
            \textbf{coscienti e volontari}, orientati a uno scopo
      \item
            La presa evolve verso la \emph{presa radio-digitale} (tre dita)
      \item
            Risponde se chiamato, comprende i divieti, obbedisce agli ordini
      \item
            Può momentaneamente \textbf{allontanarsi dal genitore}
      \item
            Inizia a razionalizzare le proprie paure
      \item
            Soffia nelle vaschette, esplora ogni spazio, apre e rompe oggetti
            (apprendimento per \emph{gioco euristico}: "aiutami a fare da me")
      \item
            Inizia il \textbf{gioco simbolico} ("far finta di"), imita azioni
      \item
            Inizia a \textbf{galleggiare autonomamente}
      \item
            Gioca vicino agli altri bambini, ne imita il comportamento, ma spesso
            non interagisce ancora
      \item
            Inizia la \textbf{lateralizzazione}
\end{itemize}

\subsubsection{24-36 Mesi}\label{mesi-2}

\begin{itemize}
      \tightlist
      \item
            Momento delicato: il bambino assorbe ogni proposta ma la rifiuta se
            non presentata nel modo giusto al momento giusto
      \item
            \textbf{Tensione di breve durata}; attenzione alle tempistiche
      \item
            La \textbf{base sicura} (accompagnatore) si sposta progressivamente
            dall\textquotesingle acqua al bordo vasca e poi allo spogliatoio
            (passaggio graduale, mai forzato)
      \item
            Prime prove di \textbf{autonomia}
      \item
            Giochi in parallelo con altri bambini; tende a imitare
            l\textquotesingle adulto in modo ancora grossolano
      \item
            Iniziano i \textbf{giochi di ruolo} ("giochiamo a...", "facciamo finta
            di...")
      \item
            Atteggiamento ancora \textbf{egocentrico}: crede che tutti pensino
            come lui
      \item
            L\textquotesingle aspetto ludico è \textbf{fondamentale} in questo
            periodo
\end{itemize}

\begin{center}\rule{0.5\linewidth}{0.5pt}\end{center}

\subsection{\texorpdfstring{Il \emph{Diving
                  Reflex}}{Il Diving Reflex}}\label{il-diving-reflex}

Il \textbf{\emph{diving reflex}} (o \emph{apnea riflessa}) è un riflesso
fisiologico del periodo neonatale:

\begin{itemize}
      \tightlist
      \item
            Si attiva quando l\textquotesingle acqua entra in contatto con i
            \textbf{recettori cutanei delle zone periorali e perinasali}
      \item
            Questo stimolo attiva un \textbf{riflesso vagale} che arresta
            l\textquotesingle attività respiratoria
      \item
            Quando il bambino riemerge, mancando il riflesso, ricomincia a
            respirare normalmente
      \item
            \textbf{Scompare gradualmente} dopo la nascita, in tempi variabili da
            bambino a bambino
      \item
            Si trasforma poi in \textbf{apnea volontaria}
\end{itemize}

Le fasi legate alla respirazione prevedono lavoro su:

\begin{itemize}
      \tightlist
      \item
            Esercizi di immersione
      \item
            Spostamenti in superficie e subacquei
      \item
            Tuffi
      \item
            Controllo del galleggiamento prono e supino
      \item
            Spostamento con o senza supporti
\end{itemize}

\begin{center}\rule{0.5\linewidth}{0.5pt}\end{center}

\subsection{Educazione all\textquotesingle Acqua vs. Educazione
      attraverso
      l\textquotesingle Acqua}\label{educazione-allacqua-vs.-educazione-attraverso-lacqua}

Nell\textquotesingle attività neonatale esistono due approcci distinti e
complementari:

\begin{longtable}[]{@{}p{0.47\linewidth}p{0.47\linewidth}@{}}
      \toprule\noalign{}
      Tipo                                                   & Contenuto      \\
      \midrule\noalign{}
      \endhead
      \bottomrule\noalign{}
      \endlastfoot
      \textbf{Educazione all\textquotesingle acqua}          & Adattamento
      all\textquotesingle ambiente acquatico; sviluppo della sfera
      morfologico-funzionale attraverso spostamenti, scivolamenti e
      galleggiamenti. Non è insegnamento delle tecniche di nuoto.             \\
      \textbf{Educazione attraverso l\textquotesingle acqua} & Sviluppo degli
      aspetti \textbf{cognitivi} (conoscenza, esplorazione),
      \textbf{psicologici} (gestione di ansie e paure, socializzazione),
      \textbf{affettivi} (ambientamento, comportamento, sentimenti),
      \textbf{motori e sensoriali} (controllo dei movimenti, perfezionamento,
      mettersi alla prova)                                                    \\
\end{longtable}

\begin{center}\rule{0.5\linewidth}{0.5pt}\end{center}

\subsection{Il Metodo: Nessun Metodo}\label{il-metodo-nessun-metodo}

Non esiste un metodo univoco per svolgere questa attività. Bisogna
avere:

\begin{itemize}
      \tightlist
      \item
            \textbf{Attenzione e capacità di osservazione}
      \item
            \textbf{Nessuna forzatura}: è il bambino a indicare la strada giusta
\end{itemize}

Il principio fondamentale è \textbf{benessere e felicità}: solo in
questo stato il bambino riesce a esplorare, fare esperienza e crescere
nell\textquotesingle ambiente acquatico.

Un bambino non impara perché qualcuno spiega, ma perché
\textbf{sperimenta e prova}.

\begin{center}\rule{0.5\linewidth}{0.5pt}\end{center}

\subsection{Il Ruolo
      dell\textquotesingle Istruttore}\label{il-ruolo-dellistruttore}

L\textquotesingle istruttore è:

\begin{itemize}
      \tightlist
      \item
            Il \textbf{tramite} per proporre tecniche adeguate
            all\textquotesingle età
      \item
            L\textquotesingle{}\textbf{anello di congiunzione} tra bambino,
            genitore e acqua
\end{itemize}

Cosa osservare:

\begin{itemize}
      \tightlist
      \item
            Come si meraviglia, si muove, scopre, si diverte
      \item
            Come conosce e si adegua ai giochi proposti
      \item
            Che emozioni prova e quando
      \item
            Come gioca, quando è teso, quando si rilassa
\end{itemize}

Il compito dell\textquotesingle istruttore è:

\begin{itemize}
      \tightlist
      \item
            \textbf{Imparare a reggere il risultato} ottenuto
      \item
            Mantenere e variare il proprio comportamento per ottenere il risultato
            voluto (sia con il genitore che con il bambino)
      \item
            Suggerire prese, movimenti e \textbf{modo di parlare} al genitore
      \item
            Motivare con obiettivi raggiungibili e raccomandare sempre
            \textbf{pazienza e progressività}
\end{itemize}

\begin{center}\rule{0.5\linewidth}{0.5pt}\end{center}

\subsection{\texorpdfstring{Il
            \emph{Baby-Talking}}{Il Baby-Talking}}\label{il-baby-talking}

Il \textbf{\emph{baby-talking}} è il modo consigliato di parlare al
bambino in acqua:

\begin{itemize}
      \tightlist
      \item
            \textbf{Tono di voce più alto} del normale
      \item
            \textbf{Ritmo lento}
      \item
            \textbf{Musicalità delle parole}
\end{itemize}

\begin{center}\rule{0.5\linewidth}{0.5pt}\end{center}

\subsection{L\textquotesingle Ambiente}\label{lambiente}

\begin{longtable}[]{@{}p{0.47\linewidth}p{0.47\linewidth}@{}}
      \toprule\noalign{}
      Parametro                                     & Valore consigliato              \\
      \midrule\noalign{}
      \endhead
      \bottomrule\noalign{}
      \endlastfoot
      Temperatura dell\textquotesingle acqua        & \textbf{30°C} (troppo calda
      causerebbe eccessivo rilassamento e impedirebbe il raggiungimento degli
      obiettivi)                                                                      \\
      Temperatura dell\textquotesingle aria esterna & \textbf{26-27°C}                \\
      Qualità dell\textquotesingle acqua            & Aspetto di \textbf{fondamentale
            importanza}: il sistema immunitario dei neonati non è in grado di
      compensare carenze igieniche                                                    \\
\end{longtable}

\begin{center}\rule{0.5\linewidth}{0.5pt}\end{center}

\subsection{Struttura della Lezione}\label{struttura-della-lezione}

Una lezione tipo dura \textbf{45 minuti}, suddivisi in:

\begin{itemize}
      \tightlist
      \item
            \textbf{5-10 minuti iniziali}: studio del bambino e termoregolazione
      \item
            \textbf{30 minuti di gioco finalizzato}: raggiungimento degli
            obiettivi prefissati
      \item
            \textbf{5 minuti finali di gioco libero}: tranquillizzarsi e
            prepararsi all\textquotesingle uscita
\end{itemize}

\begin{center}\rule{0.5\linewidth}{0.5pt}\end{center}

\subsection{Il Genitore: Ruolo e
      Gestione}\label{il-genitore-ruolo-e-gestione}

\subsubsection{Al Momento
      dell\textquotesingle Iscrizione}\label{al-momento-delliscrizione}

\begin{itemize}
      \tightlist
      \item
            I genitori devono avere chiari gli \textbf{obiettivi del corso}
            (dedicare una lezione a questa spiegazione non è tempo perso)
      \item
            Fornire \textbf{consigli comportamentali} basati
            sull\textquotesingle esperienza
\end{itemize}

\subsubsection{Colloquio Preliminare}\label{colloquio-preliminare}

\begin{itemize}
      \tightlist
      \item
            Conoscere gli \textbf{aspetti del bambino e la sua storia} (anche
            tramite una scheda preimpostata)
      \item
            Valutare le \textbf{ansie e le paure dei genitori} per metterli a loro
            agio
\end{itemize}

\subsubsection{Principi Chiave nel Rapporto con i
      Genitori}\label{principi-chiave-nel-rapporto-con-i-genitori}

\begin{itemize}
      \tightlist
      \item
            I genitori \textbf{trasmettono ogni stato d\textquotesingle animo} al
            bambino, comprese ansie e paure
      \item
            Valutare e considerare l\textquotesingle{}\textbf{acquaticità del
                  genitore} stesso: preferire come accompagnatore il più tranquillo tra
            i due, senza offendere l\textquotesingle altro
      \item
            La \textbf{prima esperienza in piscina rimarrà indelebile}: massima
            attenzione
      \item
            Parlare di \textbf{diade} (genitore-bambino): il bambino deve avere
            un\textquotesingle unica base sicura di riferimento
      \item
            Non è necessaria la presenza di entrambi i genitori
            contemporaneamente; quando entrambi sono presenti,
            l\textquotesingle obiettivo è solo \textbf{divertirsi}, non
            raggiungere obiettivi specifici
\end{itemize}

\begin{center}\rule{0.5\linewidth}{0.5pt}\end{center}

\subsection{Conclusione - Punti Chiave per
      l\textquotesingle Esame}\label{conclusione---punti-chiave-per-lesame-6}

\begin{itemize}
      \tightlist
      \item
            L\textquotesingle attività neonatale è correttamente definita
            \textbf{acquaticità per neonati}, non \emph{nuoto baby}.
      \item
            \emph{Acquaticità}: complesso di azioni e reazioni che indicano
            serenità e agiatezza in acqua; è \textbf{innata} e va stimolata.
      \item
            L\textquotesingle attività inizia indicativamente a \textbf{3 mesi},
            dopo le prime vaccinazioni, su indicazione del neonatologo.
      \item
            Il \emph{diving reflex} è un riflesso fisiologico neonatale che
            arresta la respirazione al contatto dell\textquotesingle acqua con le
            zone periorali/perinasali; scompare gradualmente.
      \item
            Esistono due approcci distinti: \textbf{educazione
                  all\textquotesingle acqua} (adattamento) e \textbf{educazione
                  attraverso l\textquotesingle acqua} (sviluppo cognitivo, affettivo,
            motorio).
      \item
            \textbf{Non esiste un metodo} fisso: è il bambino a indicare la
            strada; mai forzare.
      \item
            Il principio guida è \textbf{benessere e felicità}.
      \item
            La temperatura dell\textquotesingle acqua è \textbf{30°C}; quella
            dell\textquotesingle aria \textbf{26-27°C}.
      \item
            La lezione dura \textbf{45 minuti}: 5-10 di termoregolazione, 30 di
            gioco finalizzato, 5 di gioco libero.
      \item
            Il genitore trasmette le proprie ansie al bambino: va gestito e messo
            a proprio agio dall\textquotesingle istruttore.
      \item
            Il \textbf{\emph{baby-talking}} prevede tono alto, ritmo lento e
            musicalità.
      \item
            La \textbf{prima esperienza in piscina} è indelebile: va curata con la
            massima attenzione.
      \item
            Il \emph{gioco euristico} (10-18 mesi) si basa sul principio "aiutami
            a fare da me": lasciare che il bambino esplori e sperimenti
            autonomamente.
\end{itemize}

\section{21. Cenni di Attività
  Prescolare}\label{cenni-di-attivituxe0-prescolare}

\subsection{Introduzione}\label{introduzione-20}

L\textquotesingle attività prescolare in piscina riguarda la fascia
d\textquotesingle età compresa tra i \textbf{3 e i 6 anni}, periodo in
cui le scuole nuoto generalmente accettano i bambini. In questa fase il
bambino ha forti motivazioni allo sport, soprattutto attraverso il gioco
e il divertimento. I due fattori primari che agiscono da molla sono il
\textbf{gioco} e l\textquotesingle{}\textbf{agonismo} (inteso come
desiderio di mettersi alla prova), oltre a fattori secondari legati al
contesto sociale e familiare.

Un aspetto fondamentale da comunicare subito ai genitori: il bambino
deve \textbf{divertirsi} in ciò che fa. Iscriverlo a un corso
prestigioso dove non si sente a proprio agio è \textbf{deleterio}.

\begin{center}\rule{0.5\linewidth}{0.5pt}\end{center}

\subsection{Caratteristiche delle Fasce
      d\textquotesingle Età}\label{caratteristiche-delle-fasce-detuxe0}

\subsubsection{3 anni e mezzo - Sviluppo
      Motorio}\label{anni-e-mezzo---sviluppo-motorio}

\begin{itemize}
      \tightlist
      \item
            Salta e riesce a stare su un \textbf{piede solo} per qualche secondo
      \item
            Afferra una palla \textbf{al volo}
      \item
            Manipola il proprio cibo (miscolare, tagliare, versare) in modo
            abbastanza corretto
      \item
            Comincia a usare bene le \textbf{forbici}
\end{itemize}

\subsubsection{3 anni e mezzo - Sviluppo Relazionale e
      Cognitivo}\label{anni-e-mezzo---sviluppo-relazionale-e-cognitivo}

\begin{itemize}
      \tightlist
      \item
            Comincia a piacergli fare cose nuove
      \item
            Gioca a mamma e papà; crescente inventiva nei giochi di immaginazione
      \item
            Non sempre distingue fantasia dalla realtà
      \item
            Gli piace giocare \textbf{in compagnia} degli altri bambini
      \item
            Comprende l\textquotesingle idea del contare e acquisisce un primo
            senso del tempo
      \item
            Anticipa gli eventi di una storia che conosce; presta attenzione al
            racconto
\end{itemize}

\subsubsection{5 anni e mezzo - 6 anni: l\textquotesingle Età
      d\textquotesingle Oro della
      Motricità}\label{anni-e-mezzo---6-anni-letuxe0-doro-della-motricituxe0}

A partire dai 5 anni e mezzo ha inizio la cosiddetta \textbf{età
      d\textquotesingle oro della motricità}:

\begin{itemize}
      \tightlist
      \item
            Riconosce e denomina le varie parti del corpo su di sé e sugli altri
      \item
            Riconosce, classifica e rielabora le informazioni provenienti dagli
            organi di senso: sensazioni \textbf{visive, uditive, tattili e
                  cinestesiche}
      \item
            Sviluppa un ottimo senso dell\textquotesingle{}\textbf{equilibrio}
      \item
            Crescente destrezza nei giochi con la palla (tirare, prendere, far
            rimbalzare, calciare)
\end{itemize}

\subsubsection{5 anni e mezzo - 6 anni: Sviluppo Relazionale e
      Cognitivo}\label{anni-e-mezzo---6-anni-sviluppo-relazionale-e-cognitivo}

\begin{itemize}
      \tightlist
      \item
            Entra nella \textbf{fase del pensiero intuitivo}: comprende, spiega e
            cerca di controllare il mondo circostante
      \item
            Inizia a distinguere \textbf{fantasia da realtà}
      \item
            Aumenta la curiosità per le differenze
      \item
            Migliora la capacità di \textbf{collaborare e seguire le regole}
      \item
            Disegna una figura umana con maggiori dettagli
      \item
            Migliora la comprensione di combinazioni di forme, colori e sequenze
\end{itemize}

\begin{quote}
      In questa fascia d\textquotesingle età il bambino è un vulcano:
      \textbf{tutto è gioco e movimento}, gli unici due strumenti che ha per
      scoprire il mondo.
\end{quote}

\begin{center}\rule{0.5\linewidth}{0.5pt}\end{center}

\subsection{I Corsi: Obiettivi e
      Organizzazione}\label{i-corsi-obiettivi-e-organizzazione}

\subsubsection{Obiettivi Primari}\label{obiettivi-primari}

\begin{itemize}
      \tightlist
      \item
            Conoscenza dell\textquotesingle{}\textbf{ambiente piscina}
      \item
            Conoscenza dell\textquotesingle{}\textbf{acqua}
      \item
            \textbf{Adattamento degli schemi motori terrestri}
            all\textquotesingle ambiente acquatico: correre, saltare, rotolare,
            equilibrarsi in acqua
      \item
            \textbf{Creazione di nuovi schemi motori acquatici}: immergersi,
            galleggiare, scivolare in superficie, scivolare
            sott\textquotesingle acqua
      \item
            Rafforzamento dell\textquotesingle{}\textbf{autonomia personale}
\end{itemize}

\subsubsection{Organizzazione Pratica}\label{organizzazione-pratica}

\begin{itemize}
      \tightlist
      \item
            Durata della lezione: \textbf{30-40 minuti}
      \item
            Rapporto istruttore/bambini: \textbf{1 a 6}
      \item
            Il corso inizia già con l\textquotesingle{}\textbf{ingresso
                  nell\textquotesingle impianto}: per il bambino
            l\textquotesingle esperienza è globale, non si limita alla vasca
\end{itemize}

\begin{center}\rule{0.5\linewidth}{0.5pt}\end{center}

\subsection{Il Ruolo
      dell\textquotesingle Istruttore}\label{il-ruolo-dellistruttore-1}

L\textquotesingle istruttore è la \textbf{figura di riferimento adulta}
che sostituisce il genitore per tutta la durata della lezione. Deve:

\begin{itemize}
      \tightlist
      \item
            Stare \textbf{in acqua con i bambini} per l\textquotesingle intera
            durata della lezione (sia con i più piccoli che con i più grandi)
      \item
            Essere \textbf{fantasioso} e \textbf{immediato}: le spiegazioni troppo
            lunghe non sono efficaci
      \item
            Essere \textbf{carismatico} e preparato
      \item
            Essere \textbf{flessibile}: l\textquotesingle improvvisazione deve
            nascondere una solida progettazione del percorso
      \item
            \textbf{Motivare l\textquotesingle apprendimento}:
            l\textquotesingle apprendimento passa attraverso le emozioni
      \item
            Saper \textbf{entusiasmare, divertire, creare fiducia}
      \item
            Essere un \textbf{buon osservatore} (non solo guardare, ma osservare):
            notare i cambiamenti nel comportamento motorio e scegliere le proposte
            didattiche adeguate
      \item
            Conoscere i \textbf{personaggi dei cartoni animati} e i messaggi che
            trasmettono, per poterli utilizzare in vasca e rendere le lezioni più
            coinvolgenti
\end{itemize}

\begin{quote}
      L\textquotesingle istruttore è una \textbf{guida educativa}:
      l\textquotesingle obiettivo è educare \emph{attraverso}
      l\textquotesingle acqua.
\end{quote}

\subsubsection{Osservare vs. Guardare}\label{osservare-vs.-guardare}

\begin{itemize}
      \tightlist
      \item
            \textbf{Guardare}: registrare passivamente ciò che i bambini fanno
      \item
            \textbf{Osservare}: rendersi conto dei cambiamenti nel comportamento e
            nel comportamento motorio, e scegliere di conseguenza le proposte
            didattiche
\end{itemize}

\begin{center}\rule{0.5\linewidth}{0.5pt}\end{center}

\subsection{I Tre Elementi Chiave del Lavoro con i
      Bambini}\label{i-tre-elementi-chiave-del-lavoro-con-i-bambini}

Il lavoro dell\textquotesingle istruttore si fonda su tre elementi che
devono essere \textbf{completamente compenetrati} tra loro e
caratterizzare l\textquotesingle intera durata della lezione.

\subsubsection{1. Il Gioco}\label{il-gioco}

Il gioco non è un mero passatempo: è uno strumento serio e fondamentale.

Attraverso il gioco il bambino:

\begin{itemize}
      \tightlist
      \item
            Esprime creatività e immaginazione
      \item
            Si cimenta in nuove conquiste, alimentando
            l\textquotesingle{}\textbf{autostima}
      \item
            Vince ansie e paure
      \item
            Scopre il mondo e le regole che lo governano
      \item
            Affronta novità e complessità in condizioni semplificate
      \item
            Mette alla prova emozioni e sentimenti
      \item
            Conosce gradualmente il mondo degli altri
\end{itemize}

\begin{quote}
      Il gioco è un \textbf{diritto sancito dalla Convenzione sui Diritti
            dell\textquotesingle Infanzia}, perché possiede una funzione centrale
      nello sviluppo cognitivo, sociale, emotivo, creativo e motorio del
      bambino.
\end{quote}

\subsubsection{Tipologie di Gioco per Fascia
      d\textquotesingle Età}\label{tipologie-di-gioco-per-fascia-detuxe0}

\begin{longtable}[]{@{}p{0.30\linewidth}p{0.30\linewidth}p{0.30\linewidth}@{}}
      \toprule\noalign{}
      Tipo di gioco                   & Età indicativa & Caratteristiche                  \\
      \midrule\noalign{}
      \endhead
      \bottomrule\noalign{}
      \endlastfoot
      \textbf{Gioco simbolico}        & Da 2 anni      & Oggetti che rappresentano cose
      diverse (una scopa = un cavallo; una banana = un telefono)                          \\
      \textbf{Gioco socio-drammatico} & 4-5 anni       & Il bambino interpreta ruoli
      tenendo conto della propria parte e di quella degli altri (es. giocare a
      mamma e papà)                                                                       \\
      \textbf{Gioco di ruolo}         & 6 anni         & Descrive e definisce i ruoli,
      progetta il gioco in anticipo, assegna ruoli ai partecipanti                        \\
      \textbf{Giochi di regole}       & 6-11 anni      & Prima regole proprie, poi quelle
      condivise dal gruppo                                                                \\
\end{longtable}

\subsubsection{2. La Fabulazione}\label{la-fabulazione}

La \textbf{fabulazione} è la capacità di creare una favola attorno a una
consegna tecnica o a un esercizio proposto.

\begin{itemize}
      \tightlist
      \item
            Gli strumenti di lavoro (giocattoli) assumono \textbf{colori e
                  funzioni fiabesche}: il tappeto diventa una barca, la tavoletta un
            vassoio
      \item
            Sviluppa il \textbf{gioco simbolico}, modalità privilegiata di
            relazione nel mondo infantile
      \item
            Permette all\textquotesingle istruttore di rendere una consegna
            tecnica più accattivante
      \item
            Gli ambienti del gioco simbolico si riferiscono principalmente a
            mestieri e attività pratiche (fare la spesa, cucinare, preparare un
            picnic, giocare al meccanico)
      \item
            Inventare storie è tra le attività che piacciono di più ai bambini
\end{itemize}

\begin{quote}
      La mente è uno strumento creatore: da un piccolo elemento può generare
      storie elaborate e originali.
\end{quote}

\subsubsection{3. L\textquotesingle Imitazione}\label{limitazione}

L\textquotesingle imitazione è una forma fondamentale di apprendimento:
il bambino \textbf{riproduce ciò che vede}, senza però ripetere
meccanicamente il comportamento dell\textquotesingle adulto.

\begin{itemize}
      \tightlist
      \item
            Fornisce la possibilità di correggere e fare esperienza
      \item
            Il bambino può eseguire gesti motori solo se ha già avuto esperienze
            motorie più semplici ed elementari
      \item
            Attraverso l\textquotesingle imitazione il bambino esprime, mette alla
            prova e comprende le \textbf{dinamiche di relazione}, sperimenta
            emozioni e sentimenti
\end{itemize}

\begin{center}\rule{0.5\linewidth}{0.5pt}\end{center}

\subsection{Le Paure dei Bambini}\label{le-paure-dei-bambini}

Le paure nascono nel momento in cui si sviluppa la
\textbf{consapevolezza del pericolo}. Elaborare una propria paura
rafforza la stima in se stessi.

Nei \textbf{primi cinque anni di vita} il bambino sperimenta
\textbf{cinque forme di paura} che lo accompagneranno per tutta la vita:

\begin{enumerate}
      \tightlist
      \item
            Perdita del contatto fisico
      \item
            Paura dell\textquotesingle estraneo
      \item
            Paura della separazione
      \item
            Paura dell\textquotesingle annegamento
      \item
            Paura della morte
\end{enumerate}

\begin{quote}
      Non bisogna \textbf{deridere} né \textbf{compiangere} il bambino per le
      sue paure: va sostenuto, compreso e accompagnato, anche nella paura
      dell\textquotesingle acqua.
\end{quote}

\begin{center}\rule{0.5\linewidth}{0.5pt}\end{center}

\subsection{I Due Principi di
      Programmazione}\label{i-due-principi-di-programmazione}

\subsubsection{Principio
      dell\textquotesingle Adeguatezza}\label{principio-delladeguatezza}

Gli esercizi devono essere sempre in rapporto alle \textbf{reali
      capacità psicomotorie} del bambino, che dipendono da:

\begin{itemize}
      \tightlist
      \item
            Età
      \item
            Sesso
      \item
            Grado di preparazione
      \item
            Vissuto motorio personale
\end{itemize}

\subsubsection{Principio della
      Sistematicità}\label{principio-della-sistematicituxe0}

Prevede una progressione sequenziale:

\begin{itemize}
      \tightlist
      \item
            Dal \textbf{semplice} al \textbf{complesso}
      \item
            Dal \textbf{facile} al \textbf{difficile}
      \item
            Dal \textbf{conosciuto} all\textquotesingle{}\textbf{ignoto}
\end{itemize}

\begin{quote}
      Questi tre elementi - gioco, fabulazione, imitazione - e i due principi
      sono strettamente legati. Nessuno può sussistere senza gli altri; il
      gioco li lega tutti.
\end{quote}

\begin{center}\rule{0.5\linewidth}{0.5pt}\end{center}

\subsection{Prove di Autonomia: dalla Vasca Piccola alla Vasca
      Grande}\label{prove-di-autonomia-dalla-vasca-piccola-alla-vasca-grande}

Per chi utilizza la vasca con acqua bassa, è importante procedere in
modo \textbf{programmato e graduale} verso la scoperta
dell\textquotesingle acqua profonda:

\begin{enumerate}
      \tightlist
      \item
            Osservare la vasca grande da fuori
      \item
            Immergere solo i piedi
      \item
            Entrare e uscire con l\textquotesingle aiuto
            dell\textquotesingle istruttore
\end{enumerate}

\begin{quote}
      Un passaggio \textbf{brusco} dalla vasca piccola alla vasca profonda
      (dove il bambino non tocca) può rovinare settimane e mesi di lavoro
      sull\textquotesingle ambientamento.
\end{quote}

\begin{center}\rule{0.5\linewidth}{0.5pt}\end{center}

\subsection{Struttura della Lezione}\label{struttura-della-lezione-1}

La lezione si articola in \textbf{tre momenti fondamentali}:

\subsubsection{1. Accoglienza}\label{accoglienza}

\begin{itemize}
      \tightlist
      \item
            Il gruppo si ritrova e ricorda le \textbf{poche regole} per giocare
            insieme
      \item
            Si ricordano i nomi dei bambini presenti e di quelli assenti
      \item
            Momento speciale di connessione e appartenenza al gruppo
\end{itemize}

\subsubsection{2. Giochiamo ad Imparare}\label{giochiamo-ad-imparare}

\begin{itemize}
      \tightlist
      \item
            Il gioco ha \textbf{sempre una finalità}
      \item
            Si gioca per imparare e per divertirsi
      \item
            La proposta didattica è strutturata ma condotta in forma ludica
\end{itemize}

\subsubsection{3. Il Saluto}\label{il-saluto}

\begin{itemize}
      \tightlist
      \item
            Ci si complimenta e ci si saluta
      \item
            Si \textbf{dà appuntamento} alla lezione successiva, creando attesa e
            aspettativa positiva
\end{itemize}

\begin{center}\rule{0.5\linewidth}{0.5pt}\end{center}

\subsection{Conclusione - Punti Chiave per
      l\textquotesingle Esame}\label{conclusione---punti-chiave-per-lesame-7}

\begin{itemize}
      \tightlist
      \item
            La fascia d\textquotesingle età prescolare va dai \textbf{3 ai 6
                  anni}; le scuole nuoto accettano bambini generalmente dai \textbf{3
                  anni in su}.
      \item
            Il bambino deve \textbf{divertirsi}: un corso inadatto è deleterio.
      \item
            A \textbf{5 anni e mezzo} ha inizio l\textquotesingle{}\textbf{età
                  d\textquotesingle oro della motricità}.
      \item
            La durata della lezione è di \textbf{30-40 minuti} con rapporto
            \textbf{1 a 6}.
      \item
            L\textquotesingle istruttore è \textbf{in acqua} con i bambini per
            tutta la lezione.
      \item
            I tre elementi chiave del lavoro sono: \textbf{gioco},
            \textbf{fabulazione} e \textbf{imitazione}.
      \item
            Le quattro tipologie di gioco per fascia d\textquotesingle età:
            \emph{simbolico} (2 anni), \emph{socio-drammatico} (4-5 anni),
            \emph{di ruolo} (6 anni), \emph{di regole} (6-11 anni).
      \item
            La \textbf{fabulazione} trasforma gli attrezzi e gli esercizi tecnici
            in elementi di una storia o favola.
      \item
            I due principi di programmazione sono: \textbf{adeguatezza} (calibrare
            sulle reali capacità) e \textbf{sistematicità} (dal semplice al
            complesso, dal facile al difficile, dal conosciuto
            all\textquotesingle ignoto).
      \item
            Le \textbf{cinque paure dei primi cinque anni}: perdita del contatto
            fisico, estraneo, separazione, annegamento, morte.
      \item
            La lezione si struttura in: \textbf{accoglienza}, \textbf{giochiamo ad
                  imparare}, \textbf{saluto}.
      \item
            Il passaggio dalla vasca piccola a quella grande deve essere
            \textbf{graduale e mai brusco}.
      \item
            L\textquotesingle istruttore deve saper \textbf{osservare} (non solo
            guardare) e conoscere i personaggi dei cartoni animati per utilizzarli
            in vasca.
\end{itemize}

\section{22. Cenni di Attività per Persone con
  Disabilità}\label{cenni-di-attivituxe0-per-persone-con-disabilituxe0}

\subsection{Introduzione}\label{introduzione-21}

La disabilità non è un mondo a parte: è \textbf{una parte del mondo}.
Secondo l\textquotesingle Organizzazione Mondiale della Sanità (2011),
il \textbf{10\% della popolazione mondiale} presenta una qualche forma
di disabilità. Questo dato rende indispensabile che le piscine diventino
ambienti \textbf{accessibili e accoglienti per tutti}, per una scelta di
civiltà verso ogni tipo di utenza.

L\textquotesingle obiettivo è costruire una comunità di istruttori
\textbf{professionisti e professionali}, capaci di guardare le
\textbf{potenzialità} di ogni soggetto - non ciò che non sa fare. In
piscina, dove l\textquotesingle acqua è un ambiente completamente
diverso da quello terrestre, la disabilità si manifesta spesso in modo
accentuato, ma altrettanto spesso l\textquotesingle acqua diventa uno
spazio di straordinaria libertà.

\begin{center}\rule{0.5\linewidth}{0.5pt}\end{center}

\subsection{La Definizione di Disabilità: Evoluzione del
      Concetto}\label{la-definizione-di-disabilituxe0-evoluzione-del-concetto}

\subsubsection{Dal Modello Medico al Modello
      Bio-Psico-Sociale}\label{dal-modello-medico-al-modello-bio-psico-sociale}

Storicamente, la disabilità è stata classificata attraverso etichette
diagnostiche. Termini come \emph{scemo}, \emph{cretino},
\emph{deficiente}, \emph{imbecille}, \emph{mongoloide} erano
originariamente terminologie mediche, oggi diventate offensive.
L\textquotesingle errore sta nel \textbf{ridurre la persona alla propria
      diagnosi}.

I testi di riferimento internazionale sono
l\textquotesingle{}\textbf{ICF} (\emph{International Classification of
      Functioning}, 2001) e l\textquotesingle{}\textbf{ICF-CY} (\emph{Children
      and Youth}, 2007), che definiscono un approccio nuovo:

\begin{quote}
      La disabilità non va interpretata come qualcosa di
      \textbf{unidirezionale e statico}, ma come qualcosa di \textbf{dinamico
            e interattivo}, dipendente dal contesto ambientale e dalla
      partecipazione del soggetto.
\end{quote}

\subsubsection{Il Nuovo Lessico
      dell\textquotesingle ICF}\label{il-nuovo-lessico-dellicf}

\begin{longtable}[]{@{}p{0.30\linewidth}p{0.30\linewidth}p{0.30\linewidth}@{}}
      \toprule\noalign{}
      Vecchia terminologia & Nuova terminologia ICF      & Significato                    \\
      \midrule\noalign{}
      \endhead
      \bottomrule\noalign{}
      \endlastfoot
      Menomazione          & \textbf{Abilità residue}    & Non ciò che il soggetto ha
      perso, ma ciò che sa ancora fare                                                    \\
      Disabilità           & \textbf{Partecipazione}     & Il soggetto non è disabile in
      assoluto, ma lo è in una determinata situazione che non gli consente di
      partecipare                                                                         \\
      Handicap             & \textbf{Fattore ambientale} & L\textquotesingle ambiente può
      essere una barriera (limitante) o un facilitatore                                   \\
\end{longtable}

\textbf{Esempio pratico:} una piscina con un \textbf{sollevatore} per
carrozzine riduce al minimo la disabilità di un utente tetraplegico. Una
piscina priva di questo strumento ne amplifica invece il disagio. La
disabilità dipende quindi anche da come è organizzato
l\textquotesingle ambiente.

L\textquotesingle ICF segna il passaggio da una \textbf{prospettiva
      medico-sanitaria} a una \textbf{prospettiva bio-psico-sociale}: anche
l\textquotesingle istruttore di nuoto ha ora voce in capitolo, perché
può osservare il soggetto nel proprio contesto di vita e valutarne la
partecipazione.

\subsubsection{L\textquotesingle Etichetta Non È la
      Persona}\label{letichetta-non-uxe8-la-persona}

La pratica migliore non è usare etichette, ma conoscere i propri allievi
\textbf{per nome}. Non "un autistico" ma "Nathan". Non "un caso grave di
cerebrolesione" ma "Edoardo". Questo non significa ignorare la diagnosi,
ma non ridurre la persona ad essa.

\begin{center}\rule{0.5\linewidth}{0.5pt}\end{center}

\subsection{Tipologie di Disabilità}\label{tipologie-di-disabilituxe0}

\subsubsection{Congenita vs. Acquisita}\label{congenita-vs.-acquisita}

Una distinzione fondamentale riguarda \textbf{quando} si è manifestata
la disabilità:

\begin{itemize}
      \tightlist
      \item
            \textbf{Congenita}: presente dalla nascita; il soggetto non ha mai
            vissuto senza quella condizione
      \item
            \textbf{Acquisita}: insorta successivamente (trauma, malattia,
            incidente)
\end{itemize}

Lavorare psicologicamente con un soggetto traumatizzato da un incidente
è molto diverso rispetto a lavorare con chi è nato con una certa
menomazione. Anche il rapporto con l\textquotesingle acqua e con
l\textquotesingle istruttore cambia radicalmente.

\subsubsection{Tipologie Principali}\label{tipologie-principali}

\begin{longtable}[]{@{}p{0.47\linewidth}p{0.47\linewidth}@{}}
      \toprule\noalign{}
      Tipo di disabilità                  & Caratteristiche                                   \\
      \midrule\noalign{}
      \endhead
      \bottomrule\noalign{}
      \endlastfoot
      \textbf{Fisica}                     & Riguarda principalmente il corpo (es. paraplegia,
      tetraplegia, spasticità)                                                                \\
      \textbf{Sensoriale}                 & Colpisce principalmente udito e vista (cecità,
      ipoacusia)                                                                              \\
      \textbf{Intellettiva / relazionale} & Include i disturbi dello spettro
      autistico, le disabilità cognitive e comportamentali                                    \\
\end{longtable}

\subsubsection{In Piscina: Tre
      Categorie}\label{in-piscina-tre-categorie}

\begin{itemize}
      \tightlist
      \item
            \textbf{Disabilità permanenti}: traumatizzati spinali, spasticità,
            cerebrolesioni, disturbi comportamentali, dislessia e disturbi
            specifici dell\textquotesingle apprendimento (che influenzano la
            parametrazione spazio-temporale)
      \item
            \textbf{Disabilità temporanee}: persone che arrivano dopo malattie
            debilitanti, interventi chirurgici o traumi ortopedici. Un corso
            ponderato può facilitare il recupero
      \item
            \textbf{Nuove disabilità}: \emph{sindrome ipocinetica} (incapacità di
            muoversi per prolungata inattività), obesità, perdita
            dell\textquotesingle autosufficienza, \emph{senescenza} (anziani che
            richiedono corsi specifici)
\end{itemize}

\begin{center}\rule{0.5\linewidth}{0.5pt}\end{center}

\subsection{Organizzazione della
      Lezione}\label{organizzazione-della-lezione}

\subsubsection{Tempistica}\label{tempistica}

La frequenza tipica è \textbf{una volta alla settimana} per circa
\textbf{45 minuti}. Sebbene possa sembrare poco sul piano quantitativo,
l\textquotesingle obiettivo è massimizzare il \textbf{coinvolgimento
      qualitativo}. Per molti soggetti la piscina diventa un punto fisso della
routine settimanale, con un forte valore motivazionale: si passa da una
\textbf{motivazione estrinseca} (il genitore che accompagna) a una
\textbf{motivazione intrinseca} (il bambino che vuole tornare).

\subsubsection{Tipologie di Lezione}\label{tipologie-di-lezione}

\begin{longtable}[]{@{}p{0.47\linewidth}p{0.47\linewidth}@{}}
      \toprule\noalign{}
      Formato                      & Quando è indicato                                 \\
      \midrule\noalign{}
      \endhead
      \bottomrule\noalign{}
      \endlastfoot
      \textbf{1 a 1} (individuale) & Disabilità che non consentono
      l\textquotesingle inserimento in gruppo; fase di ambientamento
      iniziale                                                                         \\
      \textbf{1 a 2}               & Disabilità moderate; buona alternativa al piccolo
      gruppo                                                                           \\
      \textbf{Piccolo gruppo}      & Soggetti con buon adattamento sociale             \\
      \textbf{Grande gruppo}       & Soggetti con disabilità fisiche/sensoriali;
      attività ludiche guidate                                                         \\
\end{longtable}

\begin{quote}
      Il consiglio generale è di iniziare con almeno una \textbf{prima lezione
            in rapporto 1 a 1}, per dare al soggetto il tempo di conoscere
      l\textquotesingle ambiente acquatico, e valutare successivamente se
      inserirlo in un gruppo.
\end{quote}

\subsubsection{Integrazione nel Gruppo
      Normotipico}\label{integrazione-nel-gruppo-normotipico}

L\textquotesingle integrazione è \textbf{possibile ma non obbligatoria}.
In acqua molti soggetti con disabilità mostrano risposte migliori
rispetto all\textquotesingle ambiente scolastico o familiare, grazie
alla maggiore libertà di movimento. Tuttavia,
l\textquotesingle integrazione presenta difficoltà reali:

\begin{itemize}
      \tightlist
      \item
            Soggetti con disabilità \textbf{fisiche o sensoriali} si inseriscono
            generalmente con maggiore facilità
      \item
            Soggetti con \textbf{disabilità intellettive o della sfera
                  relazionale} possono incontrare problemi più seri: le differenze nelle
            modalità di apprendimento e relazione emergono progressivamente, con
            rischio di \textbf{stigmatizzazione e marginalizzazione}
      \item
            La \textbf{discontinuità} (cambio di gruppo, progressione diversa) può
            interrompere legami amicali consolidati
      \item
            Problematiche specifiche come l\textquotesingle{}\textbf{incontinenza}
            o la \textbf{percezione distorta della diversità} da parte degli altri
            utenti richiedono un istruttore altamente sensibile
\end{itemize}

\begin{quote}
      L\textquotesingle inclusione sociale è un obiettivo importante ma è un
      \textbf{percorso lungo}, non una soluzione immediata.
\end{quote}

\begin{center}\rule{0.5\linewidth}{0.5pt}\end{center}

\subsection{L\textquotesingle Ambientamento}\label{lambientamento}

L\textquotesingle ambientamento è la \textbf{fase più importante}
dell\textquotesingle inserimento in acqua. L\textquotesingle immersione
è un\textquotesingle esperienza globale: il corpo viene coinvolto
interamente e sottoposto a un forte bombardamento sensoriale.

Le tre dimensioni dell\textquotesingle ambientamento:

\begin{enumerate}
      \tightlist
      \item
            \textbf{Percezione degli effetti dell\textquotesingle acqua} sul
            proprio corpo
      \item
            \textbf{Controllo del corpo immerso} (parametri completamente diversi
            da quelli terrestri)
      \item
            \textbf{Gestione dell\textquotesingle ansia} con tempi e ritmi propri
            del soggetto
\end{enumerate}

Le acquisizioni tecniche \textbf{non sono lineari né preordinabili}.
L\textquotesingle istruttore deve seguire il momento: anche avendo uno
schema di lezione pianificato, una volta in acqua subentra
l\textquotesingle imprevedibilità. Bisogna "seguire
l\textquotesingle onda", adattando la lezione alle esigenze del
soggetto, pur mantenendo un \textbf{range di libertà}
all\textquotesingle interno di confini definiti.

\begin{center}\rule{0.5\linewidth}{0.5pt}\end{center}

\subsection{Effetti dell\textquotesingle Immersione e Vantaggi
      Terapeutici}\label{effetti-dellimmersione-e-vantaggi-terapeutici}

L\textquotesingle acqua offre tre effetti fisici particolarmente utili
per i soggetti con disabilità:

\subsubsection{\texorpdfstring{1. Spinta di Galleggiamento
            (\emph{Buoyancy})}{1. Spinta di Galleggiamento (Buoyancy)}}\label{spinta-di-galleggiamento-buoyancy}

Consente a soggetti che non riescono a deambulare di rimanere a galla.
Per tetraplegici e paraplegici, la vasca grande può essere
l\textquotesingle ambiente ideale già dalle prime lezioni, perché il
galleggiamento stimola le \textbf{abilità residue} e restituisce libertà
di movimento.

\subsubsection{2. Pressione Idrostatica}\label{pressione-idrostatica}

Produce un effetto \textbf{massaggiante} su tutto il corpo immerso,
stimolando ulteriormente gli arti verso nuove funzionalità.

\subsubsection{3. Resistenza
      all\textquotesingle Avanzamento}\label{resistenza-allavanzamento}

Permette al soggetto di ricalibrare le capacità coordinative.
L\textquotesingle acqua "fa muro" e consente atteggiamenti in posizione
eretta che sul piano terrestre sarebbero impossibili (es. mantenimento
dell\textquotesingle equilibrio statico).

\begin{quote}
      Nota importante: la \textbf{temperatura dell\textquotesingle acqua} è
      determinante. Un soggetto spastico, per esempio, non riesce a lavorare
      con acqua troppo fredda.
\end{quote}

\begin{center}\rule{0.5\linewidth}{0.5pt}\end{center}

\subsection{Approccio Fisico e
      Comunicativo}\label{approccio-fisico-e-comunicativo}

\subsubsection{Posizionamento
      dell\textquotesingle Istruttore}\label{posizionamento-dellistruttore}

L\textquotesingle istruttore, preferibilmente, \textbf{entra in acqua}
con il soggetto. Non esiste un posizionamento unico: dipende
dall\textquotesingle attività e dalle caratteristiche
dell\textquotesingle allievo.

\begin{itemize}
      \tightlist
      \item
            \textbf{Frontale}: per il contatto visivo e la comunicazione diretta
      \item
            \textbf{Alle spalle}: per sostenere e guidare il movimento
      \item
            \textbf{A distanza}: per stimolare l\textquotesingle autonomia nelle
            fasi più avanzate
\end{itemize}

\subsubsection{Comunicazione}\label{comunicazione}

Il modo di comunicare varia in base al tipo di disabilità:

\begin{itemize}
      \tightlist
      \item
            Con soggetti \textbf{nello spettro autistico}: usare sempre "tu" al
            posto di "noi", per rendere chiaro il riferimento personale
      \item
            Con soggetti \textbf{ciechi}: essere molto dettagliati nella
            descrizione dell\textquotesingle ambiente e dei movimenti,
            completamente nuovo per loro
      \item
            Con soggetti \textbf{non verbali}: non trattarli come se avessero una
            grave carenza cognitiva. Parlare normalmente, in modo completo. Molti
            soggetti non verbali capiscono tutto quello che viene detto: esprimono
            gioia, dolore e preferenze, anche se in modi atipici
\end{itemize}

\begin{quote}
      Conoscere bene i propri allievi permette di distinguere, ad esempio,
      un\textquotesingle{}\textbf{urla di gioia} da una di dolore in un
      soggetto autistico non verbale.
\end{quote}

\subsubsection{Manipolazione}\label{manipolazione}

La \textbf{manipolazione} - il contatto fisico guidato per far assumere
al soggetto posizioni in acqua - deve sempre essere \textbf{preceduta da
      una spiegazione verbale}, anche con soggetti apparentemente non
responsivi:

\begin{itemize}
      \tightlist
      \item
            "Ti sto prendendo alle spalle"
      \item
            "Sto per sdraiarti"
      \item
            "Appoggio la testa sulla mia spalla"
\end{itemize}

Queste frasi, ripetute costantemente, creano \textbf{routine} che il
soggetto interiorizza e che facilitano successivamente la cooperazione.
La manipolazione è fondamentale per chi non è autosufficiente, perché
sviluppa le afferenze tattili acquatiche e consente di assumere
atteggiamenti corretti, ma presuppone sempre che il soggetto \textbf{si
      fidi dell\textquotesingle istruttore}.

Il contatto fisico tende a \textbf{scemare} progressivamente man mano
che il soggetto acquisisce autonomia in acqua.

\begin{center}\rule{0.5\linewidth}{0.5pt}\end{center}

\subsection{La Trasformazione del Soggetto in
      Acqua}\label{la-trasformazione-del-soggetto-in-acqua}

L\textquotesingle obiettivo del percorso in piscina è una
\textbf{trasformazione progressiva} che si articola in tre stadi:

\begin{longtable}[]{@{}p{0.47\linewidth}p{0.47\linewidth}@{}}
      \toprule\noalign{}
      Stadio                       & Descrizione                              \\
      \midrule\noalign{}
      \endhead
      \bottomrule\noalign{}
      \endlastfoot
      \textbf{Corpo di attenzione} & Potenzialità ridotte al minimo; pensiero
      e azione spesso separati; qualcuno fa le cose al posto del soggetto     \\
      \textbf{Corpo di intenzione} & Il soggetto usa le proprie abilità
      residue per essere il più autonomo possibile; ciò che impara in acqua si
      trasferisce fuori dall\textquotesingle acqua                            \\
      \textbf{Corpo di relazione}  & Il soggetto si mette in contatto con gli
      altri; sa quando chiedere aiuto; partecipa attivamente                  \\
\end{longtable}

\begin{quote}
      L\textquotesingle obiettivo finale non è l\textquotesingle autonomia
      \textbf{acquatica}, ma l\textquotesingle{}\textbf{autonomia nella vita
            terrestre}: l\textquotesingle acqua è lo strumento, non il fine. Il
      lavoro in vasca deve trasferirsi fuori, anche nello spogliatoio
      (vestirsi/svestirsi da soli, fare la doccia autonomamente).
\end{quote}

\begin{center}\rule{0.5\linewidth}{0.5pt}\end{center}

\subsection{Il Ruolo della Famiglia}\label{il-ruolo-della-famiglia}

La famiglia è:

\begin{itemize}
      \tightlist
      \item
            \textbf{Responsabile del progetto di vita} del soggetto
      \item
            \textbf{Strumento primario} di relazione con
            l\textquotesingle ambiente esterno
      \item
            \textbf{Coattrice} dei progetti del bambino
      \item
            \textbf{Collaborante} con la rete di tutti gli attori (medici,
            educatori, istruttori)
\end{itemize}

L\textquotesingle istruttore non deve mai \textbf{contestare o sminuire}
la famiglia. Essa conosce difficoltà quotidiane che
l\textquotesingle istruttore non può immaginare dai soli 45 minuti
settimanali in piscina. Il \textbf{colloquio preliminare} è fondamentale
per conoscere le esigenze della famiglia e costruire
un\textquotesingle{}\textbf{alleanza} duratura.

\begin{center}\rule{0.5\linewidth}{0.5pt}\end{center}

\subsection{Il Profilo dell\textquotesingle Istruttore
      Ideale}\label{il-profilo-dellistruttore-ideale-1}

Un buon istruttore per persone con disabilità:

\begin{itemize}
      \tightlist
      \item
            Riconosce le \textbf{potenzialità residue positive} di ogni soggetto
      \item
            È in \textbf{aggiornamento continuo}, non solo per essere utile agli
            altri, ma per se stesso
      \item
            Contribuisce al \textbf{progetto di vita} del soggetto, trasferendo il
            lavoro acquatico verso l\textquotesingle esterno
      \item
            Sostiene l\textquotesingle{}\textbf{autostima} come priorità assoluta:
            "sei bravo, sei bravo, sei bravo" è un \emph{leitmotiv} quotidiano
      \item
            Punta all\textquotesingle{}\textbf{autoaffermazione} e
            all\textquotesingle{}\textbf{autoefficacia}: non fare le cose al posto
            del soggetto, ma farle fare a lui
      \item
            Ha ampie \textbf{capacità relazionali}: sa comunicare, sa gestire le
            frustrazioni proprie e dell\textquotesingle allievo, sa essere fermo
            senza essere rigido
      \item
            Possiede una ricca \textbf{"cassetta degli attrezzi"} di competenze,
            conoscenze ed esperienze a cui attingere nei momenti di
            imprevedibilità
      \item
            Porta con sé una visione positiva: la disabilità è \textbf{una parte
                  di sé, non tutto sé}
\end{itemize}

\begin{center}\rule{0.5\linewidth}{0.5pt}\end{center}

\subsection{Conclusione - Punti Chiave per
      l\textquotesingle Esame}\label{conclusione---punti-chiave-per-lesame-8}

\begin{itemize}
      \tightlist
      \item
            Secondo l\textquotesingle OMS (2011) il \textbf{10\% della popolazione
                  mondiale} è disabile; la piscina deve essere un ambiente per tutti.
      \item
            L\textquotesingle{}\textbf{ICF} (2001) e
            l\textquotesingle{}\textbf{ICF-CY} (2007) definiscono la disabilità
            come fenomeno dinamico e interattivo: si passa da
            menomazione/disabilità/handicap a \textbf{abilità residue /
                  partecipazione / fattore ambientale}.
      \item
            La disabilità può essere \textbf{congenita o acquisita}: cambia
            profondamente l\textquotesingle approccio psicologico.
      \item
            Le tre categorie di disabilità in piscina: \textbf{permanente},
            \textbf{temporanea}, \textbf{nuove disabilità} (sindrome ipocinetica,
            obesità, senescenza).
      \item
            Le tre tipologie principali sono: \textbf{fisica},
            \textbf{sensoriale}, \textbf{intellettiva/relazionale}.
      \item
            La lezione dura circa \textbf{45 minuti}, generalmente una volta a
            settimana; si mira alla massima qualità del coinvolgimento.
      \item
            I tre effetti dell\textquotesingle acqua utilizzabili
            terapeuticamente: \textbf{galleggiamento}, \textbf{pressione
                  idrostatica}, \textbf{resistenza all\textquotesingle avanzamento}.
      \item
            L\textquotesingle{}\textbf{ambientamento} è la fase più importante:
            non può essere pianificato rigidamente; bisogna seguire il momento.
      \item
            La \textbf{manipolazione} deve essere sempre preceduta da una
            spiegazione verbale, anche con soggetti non verbali.
      \item
            La trasformazione auspicata: \textbf{corpo di attenzione} →
            \textbf{corpo di intenzione} → \textbf{corpo di relazione}.
      \item
            L\textquotesingle obiettivo finale non è l\textquotesingle autonomia
            acquatica, ma l\textquotesingle{}\textbf{autonomia nella vita
                  terrestre}.
      \item
            La \textbf{famiglia} è alleata fondamentale: non va mai contestata; il
            colloquio preliminare è obbligatorio.
      \item
            L\textquotesingle istruttore ideale punta su \textbf{autostima},
            \textbf{autoefficacia} e trasferimento delle abilità fuori
            dall\textquotesingle acqua.
      \item
            Conoscere l\textquotesingle allievo \textbf{per nome} è più importante
            che conoscere la sua etichetta diagnostica.
\end{itemize}

\section{23. Cenni di Attività per Adulti e Terza e Quarta
  Età}\label{cenni-di-attivituxe0-per-adulti-e-terza-e-quarta-etuxe0}

\subsection{Introduzione}\label{introduzione-22}

La popolazione occidentale sta invecchiando in modo esponenziale. In
Italia le nascite sono scese sotto le \textbf{400.000 unità}
l\textquotesingle anno, con un saldo negativo tra anziani e giovani che
continua a crescere. Questo fenomeno ha conseguenze dirette sul lavoro
dell\textquotesingle istruttore di nuoto: \textbf{quasi il 50\% delle
      attività nelle piscine è oggi rivolto a un pubblico adulto o anziano}.
L\textquotesingle istruttore ha quindi l\textquotesingle obbligo
professionale e morale di conoscere le strategie di comunicazione, i
comportamenti e le esigenze di questa fascia d\textquotesingle utenza,
che sarà sempre più numerosa nei prossimi decenni.

\begin{center}\rule{0.5\linewidth}{0.5pt}\end{center}

\subsection{Il Contesto Demografico e la Salute
      Percepita}\label{il-contesto-demografico-e-la-salute-percepita}

\subsubsection{Non basta vivere a lungo}\label{non-basta-vivere-a-lungo}

Non è la morte il timore principale degli anziani, ma la
\textbf{dipendenza}: dover sopravvivere invece di vivere, non essere
autonomi, dipendere da terzi per ogni azione quotidiana.
L\textquotesingle obiettivo non è allungare la vita in sé, ma garantire
una \textbf{vita attiva e autonoma}.

\subsubsection{Due Indici di Salute
      Percepita}\label{due-indici-di-salute-percepita}

\begin{enumerate}
      \tightlist
      \item
            \textbf{Stato fisico}: capacità di compiere azioni quotidiane in
            autonomia (esempio: allacciarsi le scarpe)
      \item
            \textbf{Stato psicologico}: strettamente dipendente dallo stato
            fisico; la percezione di non essere più autonomi genera una crisi
            psicologica profonda
\end{enumerate}

Esempio significativo: perdere la patente o essere impediti di muoversi
liberamente genera in circa il \textbf{50\%} dei soggetti anziani una
percezione fortemente negativa della propria condizione di vita.

\subsubsection{Dati sull\textquotesingle Aspettativa di Vita in
      Italia}\label{dati-sullaspettativa-di-vita-in-italia}

\begin{longtable}[]{@{}p{0.30\linewidth}p{0.30\linewidth}p{0.30\linewidth}@{}}
      \toprule\noalign{}
      Periodo                        & Uomini  & Donne   \\
      \midrule\noalign{}
      \endhead
      \bottomrule\noalign{}
      \endlastfoot
      Anni \textquotesingle50 (1950) & 50 anni & 51 anni \\
      Oggi (post-Covid)              & 80 anni & 83 anni \\
\end{longtable}

\begin{quote}
      La \textbf{speranza di vita in buona salute} in Italia è oggi di circa
      \textbf{58-59 anni}: da quel momento in poi il soggetto adulto potrebbe
      iniziare ad avere deficit, problematiche e patologie progressive.
\end{quote}

\begin{center}\rule{0.5\linewidth}{0.5pt}\end{center}

\subsection{I Tre Ambiti della
      Prevenzione}\label{i-tre-ambiti-della-prevenzione}

L\textquotesingle istruttore non è un medico né un fisioterapista, ma
può fornire \textbf{linee guida di massima} orientate a tre ambiti
interconnessi, che formano un \emph{loop}:

\subsubsection{1. Prevenzione Metabolica}\label{prevenzione-metabolica}

Riguarda le \textbf{abitudini alimentari}: cosa si mangia, come si
mangia, come si distribuiscono i pasti. Una cattiva alimentazione genera
conseguenze motorie e strutturali.

\subsubsection{2. Prevenzione
      Strutturale}\label{prevenzione-strutturale}

Ha due significati complementari:

\begin{itemize}
      \tightlist
      \item
            \textbf{Personale}: mantenersi aggiornati dal punto di vista
            sanitario, con esami periodici adeguati; affrontare i problemi
            strutturali prima che diventino invalidanti
      \item
            \textbf{Territoriale}: appoggiarsi a strutture sul territorio
            (piscine, centri sportivi) che forniscano servizi adeguati
\end{itemize}

\subsubsection{3. Prevenzione Motoria}\label{prevenzione-motoria}

L\textquotesingle attività motoria varia e continuativa è il terzo
pilastro indispensabile. Il corpo umano è l\textquotesingle unica
"macchina" che \textbf{migliora la propria risposta con
      l\textquotesingle utilizzo}: più viene usato, più diventa stabile e
versatile. I criteri fondamentali sono:

\begin{itemize}
      \tightlist
      \item
            \textbf{Intensità del carico}
      \item
            \textbf{Quantità}
      \item
            \textbf{Adeguatezza dello stimolo}
      \item
            \textbf{Stile di vita complessivo}
\end{itemize}

\begin{quote}
      I tre ambiti sono un \emph{loop}: sono collegati tra loro e nessuno può
      essere trascurato. Obesità, sedentarietà e cattiva alimentazione si
      rinforzano vicendevolmente.
\end{quote}

\begin{center}\rule{0.5\linewidth}{0.5pt}\end{center}

\subsection{Dati sulle Cause di Morte in Italia
      (2018)}\label{dati-sulle-cause-di-morte-in-italia-2018}

\begin{longtable}[]{@{}p{0.47\linewidth}p{0.47\linewidth}@{}}
      \toprule\noalign{}
      Causa                    & Percentuale dei decessi \\
      \midrule\noalign{}
      \endhead
      \bottomrule\noalign{}
      \endlastfoot
      Eccesso di peso          & \textbf{44,8\%}         \\
      Sedentarietà             & \textbf{35,7\%}         \\
      Fumo                     & \textbf{19,4\%}         \\
      Inadeguata alimentazione & \textbf{19,6\%}         \\
      Alcol                    & \textbf{16,7\%}         \\
\end{longtable}

Questi dati mostrano che la maggior parte delle morti premature è legata
a \textbf{stili di vita modificabili}. La sedentarietà in particolare
genera un circolo vizioso: dolore → immobilità → aumento di peso →
peggioramento della salute → ulteriore immobilità.

\begin{center}\rule{0.5\linewidth}{0.5pt}\end{center}

\subsection{L\textquotesingle Indice Funzionale di
      Autonomia}\label{lindice-funzionale-di-autonomia}

Non è sufficiente misurare la vita media. È necessario considerare
l\textquotesingle{}\textbf{indice funzionale di autonomia}: la capacità
funzionale del soggetto deve mantenersi in un \emph{range minimo}
sufficiente a garantire una vita attiva, man mano che
l\textquotesingle età avanza.

Essere attivi su tre fronti:

\begin{itemize}
      \tightlist
      \item
            \textbf{Mentalmente} (curiosità, apprendimento continuo)
      \item
            \textbf{Socialmente} (relazioni, interazioni)
      \item
            \textbf{Motoriamente} (attività fisica adeguata e costante)
\end{itemize}

\begin{center}\rule{0.5\linewidth}{0.5pt}\end{center}

\subsection{Le Quattro Capacità Motorie Correlate alla
      Salute}\label{le-quattro-capacituxe0-motorie-correlate-alla-salute}

Secondo l\textquotesingle OMS e l\textquotesingle{}\emph{American
      College of Sports Medicine}, le capacità motorie direttamente collegate
alla salute di adulti e anziani sono quattro. Tutte possono essere
sviluppate e stimolate in acqua:

\begin{longtable}[]{@{}p{0.47\linewidth}p{0.47\linewidth}@{}}
      \toprule\noalign{}
      Capacità                      & Utilità in acqua                          \\
      \midrule\noalign{}
      \endhead
      \bottomrule\noalign{}
      \endlastfoot
      \textbf{Resistenza aerobica}  & Sviluppabile attraverso nuoto, fitness in
      acqua, attività continue                                                  \\
      \textbf{Resistenza muscolare} & Stimolabile con esercizi di forza in
      ambiente acquatico                                                        \\
      \textbf{Mobilità articolare}  & Favorita dalla riduzione del carico
      gravitazionale                                                            \\
      \textbf{Equilibrio}           & Sviluppabile grazie alla resistenza
      dell\textquotesingle acqua e alla riduzione del rischio di caduta         \\
\end{longtable}

\begin{center}\rule{0.5\linewidth}{0.5pt}\end{center}

\subsection{Il Vantaggio dell\textquotesingle Acqua per gli Adulti e gli
      Anziani}\label{il-vantaggio-dellacqua-per-gli-adulti-e-gli-anziani}

L\textquotesingle acqua offre un vantaggio esclusivo rispetto ad altre
attività motorie: \textbf{in acqua non si cade}. Un anziano di 75 anni
che fa balzi o corre a ginocchia alte non rischia di farsi male se perde
l\textquotesingle equilibrio: al massimo si immerge e si rialza. Questa
caratteristica la rende \textbf{l\textquotesingle attività più sicura}
per soggetti con ridotta capacità di recupero dopo un trauma.

\begin{quote}
      Una caduta in bicicletta a 75 anni può avere conseguenze molto serie. In
      acqua questo rischio non esiste.
\end{quote}

Esiste a questo scopo l\textquotesingle{}\textbf{attività motoria
      acquatica adattata}: l\textquotesingle adattamento del servizio alla
persona specifica, sia essa disabile, anziana o adulta con patologie, è
la competenza distintiva dell\textquotesingle istruttore moderno.

\begin{center}\rule{0.5\linewidth}{0.5pt}\end{center}

\subsection{Differenze tra Adulti e Over
      65}\label{differenze-tra-adulti-e-over-65}

L\textquotesingle istruttore competente deve avere \textbf{registri
      comunicativi e didattici completamente diversi} per le due fasce.
Svolgere attività con adulti non è la stessa cosa che svolgerla con
anziani.

\subsubsection{Elementi Comuni}\label{elementi-comuni}

\begin{itemize}
      \tightlist
      \item
            Entrambi vengono in piscina per \textbf{stare in salute}
      \item
            Entrambi valorizzano la \textbf{socialità}: la piscina è un luogo di
            incontro, di scambio, di svago
\end{itemize}

\subsubsection{Differenze}\label{differenze}

\begin{longtable}[]{@{}p{0.30\linewidth}p{0.30\linewidth}p{0.30\linewidth}@{}}
      \toprule\noalign{}
      Caratteristica                    & Adulti (30-65 anni circa)                                         & Over 65 (terza e quarta
      età)                                                                                                                                \\
      \midrule\noalign{}
      \endhead
      \bottomrule\noalign{}
      \endlastfoot
      \textbf{Obiettivo principale}     & Tecnico: migliorare la nuotata,
      allenarsi                         & Autostima: sentirsi attivi, capaci, vivi                                                        \\
      \textbf{Continuità}               & Discontinua: vengono per periodi, poi spariscono e
      tornano                           & Alta fedeltà: una volta conquistata la fiducia, non lasciano
      più                                                                                                                                 \\
      \textbf{Rapporto con il training} & Vogliono allenarsi in senso proprio
                                        & Non intendono l\textquotesingle allenamento in senso tradizionale                               \\
      \textbf{Osservazione critica}     & Meno attenta ai dettagli organizzativi                            &
      Molto critici e attenti a ogni dettaglio della struttura                                                                            \\
      \textbf{Puntualità}               & Non sempre precisi                                                & Arrivano in anticipo, hanno
      bisogno di certezze e routine                                                                                                       \\
      \textbf{Feedback richiesto}       & Tecnico (correggi il mio stile)                                   &
      Emotivo-motivazionale (dimmi che sto bene, che ce la faccio)                                                                        \\
\end{longtable}

\begin{quote}
      Per gli \textbf{over 65} l\textquotesingle obiettivo non è tecnico ma di
      \textbf{autostima}: "vado in piscina due volte alla settimana perché ce
      la faccio, perché sono attivo". L\textquotesingle istruttore deve
      alimentare e rinforzare questo senso di capacità.
\end{quote}

\begin{center}\rule{0.5\linewidth}{0.5pt}\end{center}

\subsection{L\textquotesingle Aspetto Sociale: la Piscina come Centro di
      Aggregazione}\label{laspetto-sociale-la-piscina-come-centro-di-aggregazione}

Una delle evoluzioni più importanti delle scuole nuoto è la
trasformazione in \textbf{centri di aggregazione sociale}. Le strutture
più avanzate, pur mantenendo la funzione primaria di insegnamento del
nuoto e del fitness acquatico, si sono arricchite di una
\textbf{dimensione sociale extra-acquatica}:

\begin{itemize}
      \tightlist
      \item
            Gite culturali (musei, città d\textquotesingle arte)
      \item
            Escursioni in montagna, camminate
      \item
            Giornate al centro benessere
      \item
            Uscite domenicali organizzate (es. sugli sci)
\end{itemize}

Questa proposta risponde a un bisogno reale: la gente ha sempre meno
tempo per organizzare da sola le proprie attività sociali e apprezza
essere \textbf{guidata e organizzata} in un contesto di fiducia già
consolidato.

\begin{quote}
      Questa dimensione aggregativa non esula dalle competenze tecniche
      dell\textquotesingle istruttore: ne è un \textbf{complemento naturale e
            strategico} per fidelizzare l\textquotesingle utenza adulta e anziana.
\end{quote}

\begin{center}\rule{0.5\linewidth}{0.5pt}\end{center}

\subsection{\texorpdfstring{Il Movimento
            \emph{Master}}{Il Movimento Master}}\label{il-movimento-master}

Negli ultimi vent\textquotesingle anni in Italia è cresciuto
esponenzialmente il movimento del \textbf{nuoto} \textbf{\emph{master}}:
adulti e over 60 che si allenano per partecipare a competizioni
riservate a soggetti della stessa fascia d\textquotesingle età.

Esempio significativo: gli ultimi campionati europei \emph{master} a
Riccione hanno contato circa \textbf{10.000 partecipanti}.

L\textquotesingle obiettivo non è solo la vittoria: è
\textbf{confrontarsi, socializzare, viaggiare insieme, sentirsi vivi}.
Questo fenomeno è destinato a crescere ulteriormente per ragioni
demografiche e rappresenta una delle opportunità più concrete per le
piscine del futuro.

\begin{center}\rule{0.5\linewidth}{0.5pt}\end{center}

\subsection{Il Profilo dell\textquotesingle Istruttore per Adulti e
      Anziani}\label{il-profilo-dellistruttore-per-adulti-e-anziani}

L\textquotesingle istruttore competente in questo ambito deve:

\begin{itemize}
      \tightlist
      \item
            Avere \textbf{più registri comunicativi}: uno per i bambini, uno per i
            teenager, uno per gli adulti, uno per gli anziani
      \item
            Essere un \textbf{tecnico preparato} per soddisfare le richieste degli
            adulti che vogliono migliorare la propria nuotata
      \item
            Saper trasmettere \textbf{feedback positivi e motivazionali} per
            sostenere l\textquotesingle autostima degli over 65
      \item
            Comunicare non solo attraverso il \textbf{canale verbale} ma anche
            tramite quello \textbf{paraverbale}
      \item
            Avere una forte \textbf{attitudine sociale e aggregativa}
      \item
            Programmare le strutture natatorie pensando ai prossimi \textbf{20
                  anni} di evoluzione demografica
\end{itemize}

\begin{center}\rule{0.5\linewidth}{0.5pt}\end{center}

\subsection{Conclusione - Punti Chiave per
      l\textquotesingle Esame}\label{conclusione---punti-chiave-per-lesame-9}

\begin{itemize}
      \tightlist
      \item
            La popolazione occidentale invecchia: in Italia quasi il \textbf{50\%
                  delle attività in piscina} è rivolto ad adulti e anziani.
      \item
            I due indici di \textbf{salute percepita} sono lo stato fisico e lo
            stato psicologico, strettamente interconnessi.
      \item
            Speranza di vita oggi: uomini \textbf{80 anni}, donne \textbf{83
                  anni}; speranza di vita \emph{in buona salute}: circa \textbf{58-59
                  anni}.
      \item
            Le tre aree di prevenzione sono: \textbf{metabolica},
            \textbf{strutturale} e \textbf{motoria} (formano un \emph{loop}
            indissolubile).
      \item
            Dati 2018: le principali cause di morte legate allo stile di vita sono
            eccesso di peso (\textbf{44,8\%}), sedentarietà (\textbf{35,7\%}),
            fumo e alcol.
      \item
            Le quattro capacità motorie correlate alla salute (OMS /
            \emph{American College of Sports Medicine}): \textbf{resistenza
                  aerobica}, \textbf{resistenza muscolare}, \textbf{mobilità
                  articolare}, \textbf{equilibrio}.
      \item
            L\textquotesingle acqua è l\textquotesingle ambiente più sicuro per
            l\textquotesingle anziano: \textbf{in acqua non si cade}.
      \item
            Gli \textbf{adulti} hanno obiettivi prevalentemente tecnici e
            frequenza discontinua; gli \textbf{over 65} hanno obiettivi di
            autostima, sono critici, puntuali e fedelissimi una volta fidelizzati.
      \item
            L\textquotesingle{}\textbf{indice funzionale di autonomia} è il
            parametro chiave: non basta vivere a lungo, bisogna mantenere un
            livello minimo di funzionalità attiva.
      \item
            Le piscine moderne si stanno trasformando in \textbf{centri di
                  aggregazione sociale}, con attività extra-acquatiche (gite,
            escursioni, eventi culturali).
      \item
            Il movimento \textbf{\emph{master}} è in crescita costante per ragioni
            demografiche e rappresenta una grande opportunità per le strutture
            natatorie.
\end{itemize}

\section{24. Il Fitness in Acqua}\label{il-fitness-in-acqua}

\subsection{Introduzione}\label{introduzione-23}

Il \textbf{fitness in acqua} è una delle aree di attività più
significative nelle scuole nuoto federali. Il termine \emph{fitness}
affonda le radici in una concezione di salute nata nel 1978 durante una
conferenza internazionale sull\textquotesingle assistenza primaria della
popolazione ad Alma-Ata (Russia): la salute è definita come uno
\textbf{stato di completo benessere fisico e psicologico}, non
semplicemente come assenza di malattia. Stare bene è considerato un
\textbf{diritto fondamentale dell\textquotesingle uomo}.

Questa lezione è rivolta agli istruttori coordinatori, che hanno il
compito di organizzare e guidare l\textquotesingle offerta di fitness in
acqua all\textquotesingle interno delle strutture natatorie.

\begin{center}\rule{0.5\linewidth}{0.5pt}\end{center}

\subsection{Definizione e Componenti del
      Fitness}\label{definizione-e-componenti-del-fitness}

Il concetto di \emph{fitness} si articola in due ambiti complementari:

\subsubsection{1. Ambito Psicofisico (Stile di
      Vita)}\label{ambito-psicofisico-stile-di-vita}

Riguarda tutte le caratteristiche che nel loro insieme costruiscono un
quadro di buona salute:

\begin{itemize}
      \tightlist
      \item
            Peso corporeo adeguato con corretta distribuzione di \textbf{massa
                  magra} e \textbf{massa grassa}
      \item
            Alimentazione sana ed equilibrata
      \item
            Non fumare
      \item
            Non abusare di alcol
      \item
            Non abusare di medicinali
      \item
            Avere \textbf{stima di sé} e continuità nelle abitudini di vita
      \item
            Mantenersi aggiornati con esami periodici adeguati
\end{itemize}

\subsubsection{2. Ambito Pratico (Attività
      Motoria)}\label{ambito-pratico-attivituxe0-motoria}

Riguarda lo stimolo e lo sviluppo delle \textbf{capacità condizionali} e
delle \textbf{capacità coordinative}:

\begin{itemize}
      \tightlist
      \item
            \textbf{Forza}
      \item
            \textbf{Resistenza}
      \item
            \textbf{Velocità}
      \item
            \textbf{Mobilità articolare}
      \item
            \textbf{Equilibrio}
\end{itemize}

\begin{quote}
      Il \emph{fitness} si raggiunge quando questi due ambiti agiscono in
      sinergia: il soggetto si vuole bene, non eccede nelle cattive abitudini,
      si muove con costanza e raggiunge uno stato di benessere equilibrato tra
      la sfera fisica e quella psicologica.
\end{quote}

\begin{center}\rule{0.5\linewidth}{0.5pt}\end{center}

\subsection{Il Peso del Fitness in Acqua nelle Strutture
      Natatorie}\label{il-peso-del-fitness-in-acqua-nelle-strutture-natatorie}

Da un\textquotesingle analisi condotta sulla stagione 2017 in alcune
strutture del Nord Italia, emerge che:

\begin{itemize}
      \tightlist
      \item
            Le attività per \textbf{adulti} rappresentano circa il \textbf{42\%}
            degli utenti di una piscina standard
      \item
            Nell\textquotesingle ambito adulti,
            l\textquotesingle{}\textbf{attività motoria acquatica} (fitness in
            acqua) supera il \textbf{10-11\%} del totale
      \item
            Su un campione di \textbf{2.000 utenti}, circa \textbf{200-220
                  persone} praticano fitness in acqua
\end{itemize}

\begin{quote}
      Questi numeri confermano l\textquotesingle importanza strategica del
      settore fitness acquatico all\textquotesingle interno
      dell\textquotesingle offerta di una scuola nuoto federale.
\end{quote}

\begin{center}\rule{0.5\linewidth}{0.5pt}\end{center}

\subsection{Le Due Macro Categorie del
      Fitness}\label{le-due-macro-categorie-del-fitness}

Il fitness, sia in acqua che a terra, si divide in due grandi macro
categorie:

\begin{longtable}[]{@{}p{0.47\linewidth}p{0.47\linewidth}@{}}
      \toprule\noalign{}
      Macro Categoria                           & Contenuto                     \\
      \midrule\noalign{}
      \endhead
      \bottomrule\noalign{}
      \endlastfoot
      \textbf{Muscolazione / Tonificazione}     & Attività mirate allo sviluppo
      della forza muscolare e al tono corporeo                                  \\
      \textbf{Attività Aerobica / Coreografata} & Esercitazioni di tipo
      coordinativo-motorio; costruzione di coreografie su base musicale         \\
\end{longtable}

\begin{center}\rule{0.5\linewidth}{0.5pt}\end{center}

\subsection{\texorpdfstring{Le Due Famiglie FIN: \emph{Aqua Fitness} e
            \emph{Aqua
                  Training}}{Le Due Famiglie FIN: Aqua Fitness e Aqua Training}}\label{le-due-famiglie-fin-aqua-fitness-e-aqua-training}

La Federazione Italiana Nuoto distingue due famiglie principali, non
come sottocategorie ma come differenziazioni vere e proprie:

\subsubsection{\texorpdfstring{\emph{Aqua
                  Fitness}}{Aqua Fitness}}\label{aqua-fitness}

Attività proposta con \textbf{ritmo e ampiezza moderati}. Lo stesso
gesto tecnico viene eseguito a velocità ridotta e con intensità
contenuta. È adatta a soggetti con obiettivi di benessere generale,
mobilità e socialità.

\subsubsection{\texorpdfstring{\emph{Aqua
                  Training}}{Aqua Training}}\label{aqua-training}

La \textbf{stessa tipologia di gesti} dell\textquotesingle{}\emph{aqua
      fitness}, ma proposta con:

\begin{itemize}
      \tightlist
      \item
            Serie ben definite
      \item
            Intensità crescente
      \item
            Tempi di recupero stabiliti
      \item
            Lavoro \textbf{intervallato ad alta intensità}
\end{itemize}

\begin{quote}
      La differenza tra le due non riguarda gli esercizi ma
      l\textquotesingle{}\textbf{intensità e la struttura} con cui vengono
      somministrati. Questa distinzione è fondamentale anche per la
      comunicazione commerciale: chi si iscrive a un corso con la dicitura
      \emph{"acqua circuit training"} deve sapere che si tratta di
      un\textquotesingle attività ad alta intensità.
\end{quote}

\begin{center}\rule{0.5\linewidth}{0.5pt}\end{center}

\subsection{Le Tipologie di Attività di Fitness in
      Acqua}\label{le-tipologie-di-attivituxe0-di-fitness-in-acqua}

\begin{longtable}[]{@{}p{0.47\linewidth}p{0.47\linewidth}@{}}
      \toprule\noalign{}
      Attività                                                  & Caratteristiche principali                \\
      \midrule\noalign{}
      \endhead
      \bottomrule\noalign{}
      \endlastfoot
      \textbf{\emph{Acqua gym}}                                 & Ginnastica in acqua classica, la base
      storica del settore                                                                                   \\
      \textbf{\emph{Acqua aerobica}}                            & Attività aerobica strutturata su base
      musicale                                                                                              \\
      \textbf{\emph{Acqua circuito}}                            & Lavoro a stazioni con circuiti            \\
      \textbf{\emph{Acqua interval}}                            & Alternanza di fasi ad alta e bassa
      intensità                                                                                             \\
      \textbf{\emph{Pilates} in acqua}                          & Esercizi di controllo posturale e
      respirazione adattati all\textquotesingle ambiente acquatico                                          \\
      \textbf{\emph{Acqua swim training}}                       & Peculiarità FIN: il corpo è in
      posizione verticale e simula le nuotate dei 4 stili classici e degli
      sport natatori federali (nuoto, pallanuoto, ecc.) con obiettivi di
      fitness                                                                                               \\
      \textbf{\emph{Aqua kick}} / \textbf{\emph{aqua kick box}} & Gestualità
      delle arti marziali; attività intensa, apprezzata da specifiche
      tipologie di utenti                                                                                   \\
      \textbf{\emph{Step} in acqua}                             & Adattamento dello \emph{step} terrestre
      all\textquotesingle ambiente acquatico                                                                \\
      \textbf{\emph{Bike} in acqua}                             & Biciclette sommerse; attività con grande
      attrezzo                                                                                              \\
      \textbf{\emph{Aqua treadmill}}                            & \emph{Treadmill} (tappeto elastico)
      acquatico                                                                                             \\
      \textbf{\emph{Aqua jump}}                                 & Attività di salti e balzi                 \\
      \textbf{Palo acquatico}                                   & Una delle ultime evoluzioni, basata sulla
      gestualità del palo                                                                                   \\
      \textbf{\emph{Aquazumba}}                                 & Attività precoreografata su musiche
      latineggianti; brani di 4-5 minuti con gestualità ripetute fino
      all\textquotesingle apprendimento; coordinativamente impegnativa ma
      divertente                                                                                            \\
\end{longtable}

\begin{center}\rule{0.5\linewidth}{0.5pt}\end{center}

\subsection{Le Proprietà Fisiche dell\textquotesingle Acqua come
      Attrezzo}\label{le-proprietuxe0-fisiche-dellacqua-come-attrezzo}

L\textquotesingle istruttore FIN di fitness in acqua deve conoscere le
proprietà dell\textquotesingle acqua per aumentare o diminuire le
resistenze \textbf{senza necessariamente ricorrere agli attrezzi}:

\begin{itemize}
      \tightlist
      \item
            L\textquotesingle acqua è \textbf{800 volte più densa}
            dell\textquotesingle aria
      \item
            Offre una \textbf{resistenza all\textquotesingle avanzamento e
                  all\textquotesingle attrito di 12 volte superiore} a quella
            dell\textquotesingle aria
      \item
            La \textbf{spinta idrostatica} contrasta la forza di gravità: ciò che
            è difficile a terra può risultare facile in acqua e viceversa
\end{itemize}

\begin{quote}
      Variare la \textbf{forma del corpo} (superficie esposta
      all\textquotesingle acqua, velocità del movimento, ampiezza) è
      sufficiente per ottenere risultati significativi, senza aggiungere pesi
      o attrezzi.
\end{quote}

\subsubsection{Uso degli Attrezzi}\label{uso-degli-attrezzi}

Gli attrezzi non sono obbligatori ma sempre più diffusi:

\begin{itemize}
      \tightlist
      \item
            \textbf{Manubri galleggianti} e manubri ad attrito
      \item
            \textbf{Bilanceri} acquatici
      \item
            \textbf{Biciclette} in acqua
      \item
            \textbf{Palo acquatico}
      \item
            \textbf{Cinture di galleggiamento} (indispensabili se
            l\textquotesingle utente non ha appoggio podalico)
\end{itemize}

Il trend degli attrezzi acquatici è cresciuto in risposta alla ricerca
di stimoli nuovi, sebbene in parte snaturino il principio originale
dell\textquotesingle{}\textbf{acqua come attrezzo principale}.

\begin{center}\rule{0.5\linewidth}{0.5pt}\end{center}

\subsection{Organizzazione dell\textquotesingle Attività: Dove e
      Quando}\label{organizzazione-dellattivituxe0-dove-e-quando}

\subsubsection{Dove Fare Fitness in
      Acqua}\label{dove-fare-fitness-in-acqua}

Il fitness in acqua si può fare in qualsiasi struttura con acqua, con o
senza appoggio podalico. La variabile principale è la \textbf{profondità
      dell\textquotesingle acqua}:

\begin{itemize}
      \tightlist
      \item
            La cosiddetta \textbf{\emph{acqua media}} si attesta tra i \textbf{130
                  e i 155 cm}
      \item
            Con acqua fino al collo, gli spostamenti in avanzamento diventano
            difficili per chi non è un nuotatore esperto
\end{itemize}

La FIN gestisce prevalentemente \textbf{piscine comunali} (25, 33 o 50
metri), non sempre ideali per il fitness ma adattabili:

\begin{itemize}
      \tightlist
      \item
            Non è necessaria la musica: si può fare una lezione classica di
            ginnastica in acqua
      \item
            In una corsia larga 2 corsie si possono svolgere attività tonificanti,
            divertenti e socializzanti \textbf{senza attrezzi}
      \item
            L\textquotesingle assenza di musica risolve il problema della
            convivenza con corsi di nuoto bambini nelle stesse fasce orarie
\end{itemize}

\subsubsection{Quando Proporre il Fitness in
      Acqua}\label{quando-proporre-il-fitness-in-acqua}

\begin{longtable}[]{@{}p{0.30\linewidth}p{0.30\linewidth}p{0.30\linewidth}@{}}
      \toprule\noalign{}
      Fascia oraria            & Utenza prevalente                                & Tipo di attività consigliata      \\
      \midrule\noalign{}
      \endhead
      \bottomrule\noalign{}
      \endlastfoot
      \textbf{Mattino}         & Anziani, persone non più attive lavorativamente,
      turnisti                 & Attività moderate, \emph{aqua fitness}, mobilità                                     \\
      \textbf{Pausa pranzo}    & Lavoratori                                       & Attività più intensa, \emph{acqua
      training}                                                                                                       \\
      \textbf{Pomeriggio-sera} & Lavoratori, famiglie                             & Attività intensa;
      attenzione alla convivenza con bambini                                                                          \\
\end{longtable}

\begin{quote}
      Il coordinatore deve \textbf{evitare} l\textquotesingle abbinamento di
      corsi fitness per anziani con corsi di nuoto bambini (rumore,
      confusione, incompatibilità di esigenze).
\end{quote}

\begin{center}\rule{0.5\linewidth}{0.5pt}\end{center}

\subsection{Le Motivazioni dell\textquotesingle Utente: Gestire la
      Diversità del
      Gruppo}\label{le-motivazioni-dellutente-gestire-la-diversituxe0-del-gruppo}

Una delle sfide principali del fitness in acqua è che \textbf{nello
      stesso gruppo convivono motivazioni diverse}:

\begin{itemize}
      \tightlist
      \item
            Chi vuole allenarsi seriamente (migliorare tono, resistenza, forza)
      \item
            Chi vuole dimagrire
      \item
            Chi viene principalmente per la \textbf{socialità} (ritrovare amici,
            fare compagnia)
      \item
            Chi partecipa per praticità (l\textquotesingle unico corso
            nell\textquotesingle orario disponibile)
\end{itemize}

\begin{quote}
      Se durante la lezione alcuni partecipanti chiacchierano tra loro, non
      necessariamente la lezione è fallita: probabilmente
      quell\textquotesingle utente \textbf{non aveva come obiettivo primario
            il fitness}, ma la socializzazione. L\textquotesingle istruttore non
      deve frustrarsene, ma capire il proprio pubblico.
\end{quote}

\subsubsection{Implicazioni per il
      Coordinatore}\label{implicazioni-per-il-coordinatore}

Il coordinatore e la \textbf{segreteria} devono avere chiarezza su cosa
stanno vendendo:

\begin{itemize}
      \tightlist
      \item
            Denominare correttamente le attività (\emph{aqua circuit training} vs
            \emph{aqua fitness}) per orientare correttamente
            l\textquotesingle utente
      \item
            Formare il personale di desk affinché sappia spiegare la differenza
            tra le proposte
\end{itemize}

\begin{center}\rule{0.5\linewidth}{0.5pt}\end{center}

\subsection{Il Ruolo dell\textquotesingle Istruttore
      Coordinatore}\label{il-ruolo-dellistruttore-coordinatore}

L\textquotesingle istruttore coordinatore ha il compito di:

\begin{itemize}
      \tightlist
      \item
            Guidare i \textbf{gestori e i direttori tecnici} nella scelta di come,
            dove e quando inserire il fitness acquatico nel palinsesto
      \item
            Conoscere le caratteristiche specifiche della struttura e adattare
            l\textquotesingle offerta
      \item
            Formarsi e \textbf{aggiornarsi continuamente}: visitare altre piscine,
            documentarsi, seguire le evoluzioni del settore (il fitness è un mondo
            che segue le mode e si aggiorna rapidamente)
      \item
            Saper proporre \textbf{registri diversi} a gruppi con esigenze diverse
      \item
            Non considerare la musica come elemento imprescindibile
\end{itemize}

\begin{center}\rule{0.5\linewidth}{0.5pt}\end{center}

\subsection{Il Trend del Fitness in
      Acqua}\label{il-trend-del-fitness-in-acqua}

Il settore ha registrato:

\begin{itemize}
      \tightlist
      \item
            Una \textbf{crescita costante} dal 2000 fino al picco del
            \textbf{2013-2014}
      \item
            Successivamente un \textbf{trend in calo}, ormai abbastanza stabile
      \item
            La ricerca di nuove nicchie (attrezzi specifici, nuove proposte
            coreografiche) come risposta alla stabilizzazione
\end{itemize}

Il fitness acquatico tende a seguire le mode del fitness terrestre, a
sua volta influenzato dai trend americani. L\textquotesingle istruttore
competente sa leggere queste tendenze e proporre attività innovative
senza perdere il fondamento della \textbf{conoscenza
      dell\textquotesingle elemento acqua}.

\begin{center}\rule{0.5\linewidth}{0.5pt}\end{center}

\subsection{Conclusione - Punti Chiave per
      l\textquotesingle Esame}\label{conclusione---punti-chiave-per-lesame-10}

\begin{itemize}
      \tightlist
      \item
            Il \textbf{fitness} nasce come concetto globale di salute (1978,
            Alma-Ata): benessere fisico \textbf{e} psicologico, non solo assenza
            di malattia.
      \item
            Il fitness si articola in due ambiti: \textbf{psicofisico} (stile di
            vita) e \textbf{pratico} (capacità condizionali e coordinative).
      \item
            Le attività per adulti rappresentano circa il \textbf{42\%} degli
            utenti di una piscina; il fitness in acqua copre
            l\textquotesingle{}\textbf{11\%} circa.
      \item
            Le due macro categorie del fitness:
            \textbf{muscolazione/tonificazione} e \textbf{attività
                  aerobica/coreografata}.
      \item
            Le due famiglie FIN: \textbf{\emph{aqua fitness}} (intensità moderata)
            e \textbf{\emph{aqua training}} (alta intensità, serie definite,
            recuperi strutturati). Stessi gesti, intensità diversa.
      \item
            L\textquotesingle acqua è \textbf{800 volte più densa}
            dell\textquotesingle aria e offre \textbf{12 volte} più resistenza:
            variare la forma del corpo basta per modulare lo sforzo.
      \item
            La \textbf{\emph{acqua media}} si attesta tra \textbf{130 e 155 cm}.
      \item
            Il fitness in acqua non richiede necessariamente musica né attrezzi:
            l\textquotesingle{}\textbf{acqua è l\textquotesingle attrezzo
                  principale}.
      \item
            Le motivazioni degli utenti sono diverse (allenamento, dimagrimento,
            socialità, praticità): il coordinatore deve saperle riconoscere e
            gestire.
      \item
            Il coordinatore deve formare anche il personale di \textbf{segreteria}
            sulla corretta denominazione e vendita delle attività.
      \item
            La \textbf{cintura di galleggiamento} è indispensabile quando
            l\textquotesingle utente non ha appoggio podalico.
      \item
            Il trend del settore: crescita fino al 2013-2014, poi stabilizzazione;
            il mercato degli attrezzi acquatici è in espansione.
\end{itemize}

\section{25. Ginnastica Educativa e
  Prenatatoria}\label{ginnastica-educativa-e-prenatatoria}

\subsection{Introduzione}\label{introduzione-24}

La \textbf{ginnastica educativa e prenatatoria} è un ambito che negli
ultimi anni ha perso importanza nelle piscine italiane, pur mantenendo
un valore scientifico e didattico solido. La lezione analizza le ragioni
di questo declino, i benefici oggettivi dell\textquotesingle attività a
secco come supporto all\textquotesingle apprendimento motorio in acqua e
le strategie che il coordinatore di vasca può adottare per reintegrarla
nell\textquotesingle offerta della scuola nuoto federale.

\begin{center}\rule{0.5\linewidth}{0.5pt}\end{center}

\subsection{Il Contesto: La Piscina come Luogo di
      Cultura}\label{il-contesto-la-piscina-come-luogo-di-cultura}

Le piscine non sono più semplici servizi pubblici dove si va a imparare
a nuotare. Oggi rappresentano un \textbf{luogo di cultura trasversale}
con quattro dimensioni:

\begin{itemize}
      \tightlist
      \item
            \textbf{Cultura sportiva}: lo sport come motore della vita del
            soggetto
      \item
            \textbf{Cultura sociale}: condividere momenti di piacere con altre
            persone
      \item
            \textbf{Etica}: chi si mantiene in salute riduce i costi per il
            sistema sanitario nazionale e contribuisce al benessere collettivo
      \item
            \textbf{Aspetto ludico}: il divertimento come componente legittima e
            valorizzata dell\textquotesingle attività
\end{itemize}

In questo contesto più articolato, i bambini hanno \textbf{tempi
      ridotti} da dedicare al nuoto (molte attività extrascolastiche in
parallelo) e oltre il \textbf{90\% dei frequentatori} dei corsi di nuoto
per minori partecipa con frequenza \textbf{monosettimanale} (50 minuti
una volta a settimana). Questo rende la ginnastica prenatatoria ancora
più preziosa come strumento per massimizzare
l\textquotesingle apprendimento.

\begin{center}\rule{0.5\linewidth}{0.5pt}\end{center}

\subsection{Definizione di Ginnastica
      Prenatatoria}\label{definizione-di-ginnastica-prenatatoria}

La FIN definisce la ginnastica prenatatoria come:

\begin{quote}
      Un \textbf{percorso fondamentale} per costruire una motricità più
      consapevole e più ricca.
\end{quote}

Non è \textbf{indispensabile} per imparare a nuotare o per padroneggiare
l\textquotesingle acqua, ma è \textbf{fondamentale} perché:

\begin{itemize}
      \tightlist
      \item
            Chi si muove con coerenza fuori dall\textquotesingle acqua ha
            \textbf{maggiore facilità di apprendimento} anche in acqua
      \item
            Riduce il \textbf{tempo necessario} per apprendere nuovi gesti motori
      \item
            Migliora la \textbf{qualità} dell\textquotesingle apprendimento
\end{itemize}

\begin{center}\rule{0.5\linewidth}{0.5pt}\end{center}

\subsection{Il Problema del Movimento Spontaneo
      Scomparso}\label{il-problema-del-movimento-spontaneo-scomparso}

Una delle ragioni più importanti per cui la ginnastica prenatatoria è
oggi più necessaria che in passato è la \textbf{scomparsa del gioco
      spontaneo all\textquotesingle aperto}. I bambini delle generazioni
precedenti costruivano la propria motricità di base giocando liberamente
per ore in cortile o per strada. Oggi questo non accade più, nemmeno
nelle periferie.

Conseguenza diretta: molti bambini di 8 anni faticano a compiere gesti
elementari come alzare il braccio destro mantenendo quello sinistro in
avanti, o a distinguere la lateralità. La \textbf{motricità di base è
      impoverita} e il comportamento motorio spontaneo è ridotto.

Dato ISTAT (2017-2018) sui ragazzi tra 3 e 17 anni:

\begin{itemize}
      \tightlist
      \item
            Il \textbf{52\%} pratica sport in modo continuativo (almeno una volta
            a settimana)
      \item
            Il \textbf{23\%} non pratica alcuna attività motoria
      \item
            Il \textbf{29\%} pratica meno di una volta a settimana
\end{itemize}

Tra il 2010 e il 2017-2018 c\textquotesingle è stato però un incremento
medio del \textbf{5\%} di praticanti, con un trend positivo.

\begin{center}\rule{0.5\linewidth}{0.5pt}\end{center}

\subsection{Apprendimento Motorio in Acqua: Il Ruolo della Motricità
      Terrestre}\label{apprendimento-motorio-in-acqua-il-ruolo-della-motricituxe0-terrestre}

La FIN persegue da oltre dieci anni la \textbf{cultura
      dell\textquotesingle acqua}, fondata sul principio che nulla
nell\textquotesingle acqua è predefinito filogeneticamente:
l\textquotesingle apprendimento acquatico dipende \textbf{esclusivamente
      dalla qualità e dalla quantità degli stimoli motori} che il soggetto
accumula nel tempo, attraverso l\textquotesingle esperienza di fare,
sbagliare e correggersi.

Tuttavia, la motricità terrestre agevola l\textquotesingle apprendimento
acquatico perché:

\begin{itemize}
      \tightlist
      \item
            Migliora il \textbf{controllo, la regolazione e
                  l\textquotesingle organizzazione del movimento} (capacità
            coordinative)
      \item
            Stimola gli \textbf{analizzatori del sistema nervoso centrale}:
            analizzatore \emph{cinestesico}, analizzatore \emph{vestibolare},
            \emph{esterocettori}
      \item
            Chi ha più esperienza motoria generale impara più facilmente a gestire
            situazioni motorie nuove, inclusa quella acquatica
\end{itemize}

\begin{center}\rule{0.5\linewidth}{0.5pt}\end{center}

\subsection{Obiettivi della Ginnastica
      Prenatatoria}\label{obiettivi-della-ginnastica-prenatatoria}

\subsubsection{1. Sviluppo e Consolidamento dello Schema
      Corporeo}\label{sviluppo-e-consolidamento-dello-schema-corporeo}

Lo \textbf{schema corporeo} è la percezione che
l\textquotesingle individuo ha della posizione e degli atteggiamenti del
proprio corpo nello spazio. Permette di:

\begin{itemize}
      \tightlist
      \item
            Percepire la posizione degli arti senza vederli
      \item
            Correggersi autonomamente
      \item
            Sviluppare la \textbf{propriocezione}: fondamentale in acqua, dove la
            visione è limitata e la sensazione tattile è predominante
\end{itemize}

\subsubsection{2. Sviluppo degli Schemi Motori di Base
      Terrestri}\label{sviluppo-degli-schemi-motori-di-base-terrestri}

Gli schemi motori di base (camminare, correre, saltare, rotolare,
lanciare) sono ereditati filogeneticamente ma devono essere stimolati e
consolidati. Nel contesto attuale, dove il gioco spontaneo è scomparso,
è compito dell\textquotesingle istruttore recuperare e stimolare questi
schemi di base, su cui si costruisce tutta la motricità più complessa.

\subsubsection{3. Mobilità Articolare}\label{mobilituxe0-articolare-1}

La \textbf{mobilità articolare} non è una capacità condizionale né
coordinativa: è la qualità che permette di compiere i gesti in modo
\textbf{ampio e fluido}. Il nuoto è tra le attività che ne richiedono di
più. Va allenata a tutte le età:

\begin{itemize}
      \tightlist
      \item
            Bambini
      \item
            Teenager
      \item
            Adulti
      \item
            Anziani (\emph{over} 65)
\end{itemize}

\subsubsection{4. Avviamento Motorio
      (Riscaldamento)}\label{avviamento-motorio-riscaldamento}

La ginnastica prenatatoria ha anche una funzione di
\textbf{riscaldamento} pre-acqua:

\begin{itemize}
      \tightlist
      \item
            Attiva la ventilazione polmonare
      \item
            Aumenta la frequenza cardiaca
      \item
            Prepara l\textquotesingle organismo all\textquotesingle immersione e
            all\textquotesingle attività in acqua
\end{itemize}

\subsubsection{5. Sviluppo delle Capacità Condizionali e
      Coordinative}\label{sviluppo-delle-capacituxe0-condizionali-e-coordinative}

\begin{itemize}
      \tightlist
      \item
            \textbf{Capacità condizionali}: forza, resistenza, velocità
      \item
            \textbf{Capacità coordinative}: destrezza, indipendenza degli arti,
            ritmo, controllo della respirazione
\end{itemize}

\subsubsection{6. Controllo della
      Respirazione}\label{controllo-della-respirazione}

La respirazione in acqua non è naturale come sulla terra. A secco si può
\textbf{condizionare e far percepire} la respirazione forzata che verrà
poi utilizzata durante le nuotate, senza dover ricorrere
all\textquotesingle acqua per farlo.

\subsubsection{7. Sviluppo Sociale e
      Relazionale}\label{sviluppo-sociale-e-relazionale}

In palestra la comunicazione tra istruttore e allievo è più facile che
in acqua (niente rumore, distanze ridotte, niente testa immersa). La
ginnastica prenatatoria è un contesto privilegiato per:

\begin{itemize}
      \tightlist
      \item
            Costruire la \textbf{relazione} con l\textquotesingle allievo
      \item
            Conoscere meglio carattere, aspettative e lacune di ogni soggetto
      \item
            Sviluppare l\textquotesingle aspetto \textbf{sociale} del gruppo
\end{itemize}

\subsubsection{8. Crescita e Sviluppo
      Fisico}\label{crescita-e-sviluppo-fisico}

La pratica regolare di attività motoria a secco contribuisce a una
\textbf{crescita migliore e più stimolata} del corpo in via di sviluppo.

\begin{center}\rule{0.5\linewidth}{0.5pt}\end{center}

\subsection{Perché la Ginnastica Prenatatoria è in Declino: i Punti
      Deboli}\label{perchuxe9-la-ginnastica-prenatatoria-uxe8-in-declino-i-punti-deboli}

\subsubsection{1. Ambienti Non Idonei}\label{ambienti-non-idonei}

Poche piscine dispongono di \textbf{palestre adeguate} adiacenti alla
vasca. Il bordo piscina è spesso bagnato, scivoloso e insufficiente come
superficie.

\subsubsection{2. Vincoli di Tempo}\label{vincoli-di-tempo}

I bambini sono disponibili principalmente tra le \textbf{16:30 e le
      19:30-20:00}. In questa finestra temporale limitata, aggiungere 15
minuti di ginnastica a secco allunga il corso (15 min ginnastica +
doccia + 40-45 min acqua = circa 70 minuti totali), che può risultare
troppo per le famiglie con tempi ridotti.

\subsubsection{3. Resistenza degli
      Utenti}\label{resistenza-degli-utenti}

Alcuni genitori e allievi percepiscono la ginnastica prenatatoria come
una \textbf{perdita di tempo} rispetto all\textquotesingle attività in
acqua. Il coordinatore deve essere in grado di spiegare il valore
propedeutico dell\textquotesingle attività a secco.

\subsubsection{4. Scarso Interesse degli
      Istruttori}\label{scarso-interesse-degli-istruttori}

Se l\textquotesingle istruttore stesso non crede nel valore della
ginnastica prenatatoria, la proporrà male o non la proporrà. È compito
del coordinatore \textbf{formare e motivare} i propri collaboratori.

\subsubsection{5. Repertorio di Esercizi
      Insufficiente}\label{repertorio-di-esercizi-insufficiente}

Conseguenza diretta dello scarso interesse: gli istruttori ripropongono
sempre gli stessi esercizi, che diventano noiosi e poco stimolanti.

\subsubsection{6. Gruppi Eterogenei}\label{gruppi-eterogenei}

I gruppi in piscina sono spesso eterogenei per età. La ginnastica
prenatatoria \textbf{non può essere uguale} per bambini, adulti e
\emph{over} 65: gli obiettivi sono completamente diversi e richiedono
proposte differenziate.

\begin{center}\rule{0.5\linewidth}{0.5pt}\end{center}

\subsection{Strategie per il Coordinatore di
      Vasca}\label{strategie-per-il-coordinatore-di-vasca}

\subsubsection{Strategia Didattica 1: Ginnastica in Acqua per gli
      Adulti}\label{strategia-didattica-1-ginnastica-in-acqua-per-gli-adulti}

Dove mancano spazi, gli adulti possono fare i 15 minuti di ginnastica
prenatatoria \textbf{direttamente in acqua} (con piedi in appoggio
quando possibile), mentre i bambini svolgono la loro attività a secco
sul bordo vasca. I due gruppi operano separatamente nello stesso
ambiente.

\subsubsection{Strategia Didattica 2: Conoscenza
      dell\textquotesingle Allievo}\label{strategia-didattica-2-conoscenza-dellallievo}

Il momento della ginnastica prenatatoria è un\textquotesingle occasione
preziosa per \textbf{conoscere meglio ogni allievo}: carattere,
aspettative, lacune motorie. Questo migliora la qualità del servizio
complessivo.

\subsubsection{Strategia Economica (Non
      Primaria)}\label{strategia-economica-non-primaria}

Un beneficio organizzativo secondario: i 15 minuti di ginnastica a secco
permettono una \textbf{sovrapposizione di turni} nella fascia oraria
pomeridiana (16:30-19:30), consentendo di aumentare il numero di accessi
giornalieri e il gettito economico per la struttura. Tuttavia, questo
non deve essere l\textquotesingle obiettivo primario della ginnastica
prenatatoria, ma solo un effetto collaterale positivo.

\subsubsection{Diversificazione per Fasce
      d\textquotesingle Età}\label{diversificazione-per-fasce-detuxe0}

Le proposte didattiche devono essere calibrate su quattro categorie:

\begin{longtable}[]{@{}p{0.47\linewidth}p{0.47\linewidth}@{}}
      \toprule\noalign{}
      Fascia                    & Caratteristiche della proposta                      \\
      \midrule\noalign{}
      \endhead
      \bottomrule\noalign{}
      \endlastfoot
      \textbf{Bambini}          & Circuiti colorati e divertenti che stimolano tutti
      gli obiettivi in modo ludico                                                    \\
      \textbf{Ragazzi/Teenager} & Sviluppo sistematico degli schemi motori,
      mobilità, coordinazione                                                         \\
      \textbf{Adulti}           & Conoscenza e professionalità; prevenzione; mobilità
      articolare                                                                      \\
      \textbf{Over 65}          & Socialità, tranquillità, mobilità articolare;
      monitoraggio di problematiche specifiche                                        \\
\end{longtable}

\begin{center}\rule{0.5\linewidth}{0.5pt}\end{center}

\subsection{Durata Minima Consigliata}\label{durata-minima-consigliata}

Il \textbf{tempo minimo} da dedicare alla ginnastica prenatatoria per
ottenere effetti significativi è di \textbf{15 minuti}.

La struttura indicativa di un corso completo diventa quindi:

\begin{itemize}
      \tightlist
      \item
            15 minuti di ginnastica prenatatoria a secco (o in acqua)
      \item
            Doccia obbligatoria
      \item
            40-45 minuti di attività in acqua
      \item
            Totale: circa \textbf{70 minuti}
\end{itemize}

\begin{center}\rule{0.5\linewidth}{0.5pt}\end{center}

\subsection{Conclusione - Punti Chiave per
      l\textquotesingle Esame}\label{conclusione---punti-chiave-per-lesame-11}

\begin{itemize}
      \tightlist
      \item
            La ginnastica prenatatoria è \textbf{fondamentale ma non
                  indispensabile} per imparare a nuotare; è fondamentale per costruire
            una motricità più consapevole e ricca.
      \item
            In acqua nulla è predefinito filogeneticamente:
            l\textquotesingle apprendimento acquatico dipende dagli
            \textbf{stimoli motori accumulati} attraverso
            l\textquotesingle esperienza.
      \item
            La ginnastica prenatatoria compensa la \textbf{scomparsa del gioco
                  spontaneo all\textquotesingle aperto} nei bambini.
      \item
            Dato ISTAT 2017-2018: il \textbf{23\%} dei ragazzi tra 3 e 17 anni non
            pratica alcuna attività motoria.
      \item
            Gli \textbf{otto obiettivi} principali: schema corporeo, schemi motori
            di base, mobilità articolare, avviamento/riscaldamento, capacità
            condizionali e coordinative, controllo della respirazione, socialità,
            crescita e sviluppo.
      \item
            La \textbf{mobilità articolare} non è né condizionale né coordinativa;
            è fondamentale per l\textquotesingle ampiezza e la fluidità del gesto
            natatorio.
      \item
            Il tempo minimo per una sessione efficace è di \textbf{15 minuti}.
      \item
            I principali ostacoli sono: ambienti non idonei, vincoli di tempo,
            resistenza degli utenti, scarso interesse degli istruttori, repertorio
            insufficiente, gruppi eterogenei.
      \item
            La ginnastica prenatatoria \textbf{non può essere uguale} per bambini,
            adulti e \emph{over} 65: ogni fascia richiede obiettivi e proposte
            differenziati.
      \item
            Il coordinatore deve \textbf{formare e motivare} gli istruttori,
            affinché credano nel valore di questa attività.
      \item
            Strategia pratica: gli adulti possono fare la ginnastica \textbf{in
                  acqua} mentre i bambini la fanno \textbf{a secco sul bordo vasca},
            risolvendo il problema degli spazi.
\end{itemize}

\section{26. L\textquotesingle Organizzazione del Centro
  Natatorio}\label{lorganizzazione-del-centro-natatorio}

\subsection{Introduzione}\label{introduzione-25}

Il \textbf{centro natatorio} è un impianto sportivo dotato di piscine
che deve essere interpretato e gestito come una \textbf{vera e propria
      azienda}. Questa visione aziendale non è accessoria ma fondamentale:
senza un\textquotesingle organizzazione strutturata, il centro non può
funzionare correttamente né garantire i servizi che i propri utenti si
aspettano.

I centri natatori gestiscono il \textbf{tempo libero, il divertimento e
      il benessere} delle persone, e per questo richiedono un livello di
attenzione e responsabilità particolarmente elevato.

\begin{center}\rule{0.5\linewidth}{0.5pt}\end{center}

\subsection{Il Centro Natatorio come
      Azienda}\label{il-centro-natatorio-come-azienda}

Un\textquotesingle{}\textbf{azienda} è un sistema in cui un insieme di
persone, con le proprie competenze, partecipa allo svolgimento di
un\textquotesingle attività secondo i principi della \textbf{divisione
      del lavoro} e del \textbf{coordinamento}, al fine di soddisfare bisogni
umani attraverso la produzione, la fornitura o il consumo di beni e
servizi verso terzi.

Il centro natatorio è un\textquotesingle azienda che:

\begin{itemize}
      \tightlist
      \item
            Fornisce \textbf{servizi verso terzi}
      \item
            Opera in un ambiente a \textbf{rischio specifico} (utenti in costume,
            ambiente umido o bagnato)
      \item
            Deve avere un \textbf{programma} chiaro di attività e obiettivi
      \item
            Necessita di un \textbf{organigramma} che assegni le competenze alle
            persone giuste
\end{itemize}

\begin{center}\rule{0.5\linewidth}{0.5pt}\end{center}

\subsection{La Sicurezza in Piscina}\label{la-sicurezza-in-piscina}

\subsubsection{Sicurezza Attiva e Sicurezza
      Passiva}\label{sicurezza-attiva-e-sicurezza-passiva}

La sicurezza in piscina si distingue in due categorie complementari:

\begin{longtable}[]{@{}p{0.47\linewidth}p{0.47\linewidth}@{}}
      \toprule\noalign{}
      Tipo                       & Descrizione                              \\
      \midrule\noalign{}
      \endhead
      \bottomrule\noalign{}
      \endlastfoot
      \textbf{Sicurezza attiva}  & L\textquotesingle insegnamento del nuoto
      stesso: formare persone capaci di gestire situazioni di rischio in
      qualsiasi specchio d\textquotesingle acqua (piscine, mare, laghi,
      fiumi)                                                                \\
      \textbf{Sicurezza passiva} & La presenza degli \textbf{assistenti
            bagnanti} preposti alla sorveglianza continua; le manutenzioni
      periodiche; la sorveglianza igienico-sanitaria                        \\
\end{longtable}

\subsubsection{Il Rischio Accettabile:
      Zero}\label{il-rischio-accettabile-zero}

A differenza di sport di contatto come il rugby o le arti marziali, dove
esiste una percentuale di rischio accettabile, nelle piscine il
\textbf{rischio accettabile deve essere zero}.

Gli incidenti in piscina sono rari, ma quando si verificano tendono ad
avere un \textbf{grado di gravità elevato}. Per questo devono essere
prevenuti con protocolli rigidi.

\subsubsection{Principali Tipologie di
      Incidente}\label{principali-tipologie-di-incidente}

\begin{enumerate}
      \item
            \textbf{Caduta sul bordo vasca}: causata principalmente da bambini che
            corrono su superfici bagnate o umide, a piedi nudi o con ciabattine.
            Può avere conseguenze da medie a gravi. Gli spostamenti lungo il bordo
            vasca devono essere costantemente monitorati.
      \item
            \textbf{Urto contro il bordo o il fondo vasca} per tuffi eseguiti in
            modo errato: bambini, adulti o utenti poco esperti possono colpire il
            bordo durante il tuffo o, se la piscina è poco profonda, toccare il
            fondo.
      \item
            \textbf{Annegamento}: il rischio più grave, con possibilità di
            conseguenze fatali. Richiede attenzione continua da parte di
            istruttori e assistenti bagnanti.
\end{enumerate}

\subsubsection{Regola Fondamentale di Sicurezza per
      l\textquotesingle Istruttore}\label{regola-fondamentale-di-sicurezza-per-listruttore}

\begin{quote}
      L\textquotesingle istruttore deve tenere \textbf{tutti gli allievi sotto
            controllo}, costantemente e tutti contemporaneamente.
\end{quote}

\begin{center}\rule{0.5\linewidth}{0.5pt}\end{center}

\subsection{Parametri Tecnici
      dell\textquotesingle Acqua}\label{parametri-tecnici-dellacqua}

L\textquotesingle istruttore deve conoscere i parametri di legge
relativi alla qualità dell\textquotesingle acqua, anche se la loro
gestione è generalmente affidata a tecnici e manutentori. Questa
conoscenza è utile per segnalare eventuali anomalie.

\begin{longtable}[]{@{}p{0.47\linewidth}p{0.47\linewidth}@{}}
      \toprule\noalign{}
      Parametro                                       & Valore di riferimento               \\
      \midrule\noalign{}
      \endhead
      \bottomrule\noalign{}
      \endlastfoot
      \textbf{pH dell\textquotesingle acqua}          & Tra \textbf{6,5} e
      \textbf{7,5}                                                                          \\
      \textbf{Cloro attivo libero}                    & Tra \textbf{0,7} e \textbf{1,5
      mg/litro}                                                                             \\
      \textbf{Cloro combinato}                        & Non superiore a \textbf{4 mg/litro} \\
      \textbf{Temperatura dell\textquotesingle acqua} & Non inferiore a
      \textbf{24°C}, non superiore a \textbf{30°C}                                          \\
\end{longtable}

\begin{center}\rule{0.5\linewidth}{0.5pt}\end{center}

\subsection{Gestione Sicura dei Tuffi}\label{gestione-sicura-dei-tuffi}

Per evitare incidenti durante i tuffi, la FIN propone uno schema di
flussi differenziato per corsia:

\begin{itemize}
      \tightlist
      \item
            \textbf{Corsia laterale}: l\textquotesingle allievo si tuffa e
            \textbf{esce immediatamente} sul bordo laterale
      \item
            \textbf{Corsie vicine alle laterali} (allievi con abilità intermedie):
            dopo il tuffo, l\textquotesingle allievo nuota un tratto verso il
            centro e poi \textbf{taglia verso il bordo}
      \item
            \textbf{Corsie centrali} (allievi più abili): il tuffo viene eseguito
            \textbf{leggermente spostato a destra} della propria corsia, con
            uscita sul \textbf{lato sinistro} della stessa corsia
\end{itemize}

Questo sistema di flussi evita collisioni tra chi si tuffa e chi esce
dall\textquotesingle acqua.

\begin{center}\rule{0.5\linewidth}{0.5pt}\end{center}

\subsection{Le Quattro Aree di Attività del Centro
      Natatorio}\label{le-quattro-aree-di-attivituxe0-del-centro-natatorio}

Un centro natatorio ben organizzato struttura la propria offerta in
\textbf{quattro grandi aree}:

\subsubsection{1. Area Nuoto}\label{area-nuoto}

Comprende tutte le attività in cui l\textquotesingle obiettivo è
\textbf{vincere la resistenza dell\textquotesingle acqua}:

\begin{itemize}
      \tightlist
      \item
            \emph{Baby nuoto} e attività neonatale
      \item
            Scuola nuoto prescolare
      \item
            Scuola nuoto bambini (tutti i livelli)
      \item
            Scuola nuoto adulti
      \item
            Preagonismo e agonismo
      \item
            Nuoto \emph{master}
\end{itemize}

\subsubsection{2. Area Fitness}\label{area-fitness}

Comprende tutte le attività in cui si \textbf{ricerca la resistenza
      dell\textquotesingle acqua} per allenare e migliorare lo stato fisico:

\begin{itemize}
      \tightlist
      \item
            \emph{Aqua gym}
      \item
            \emph{Acqua aerobica}
      \item
            \emph{Acqua training}
      \item
            Attività con grandi e piccoli attrezzi
\end{itemize}

\begin{quote}
      \textbf{Distinzione tecnica fondamentale}: nell\textquotesingle area
      \textbf{nuoto} si insegna a \emph{vincere} la resistenza
      dell\textquotesingle acqua; nell\textquotesingle area \textbf{fitness}
      si \emph{cerca} la resistenza dell\textquotesingle acqua come strumento
      di allenamento.
\end{quote}

\subsubsection{3. Area Salute}\label{area-salute}

Comprende le attività riabilitative e terapeutiche:

\begin{itemize}
      \tightlist
      \item
            \emph{Idrochinesi terapia}
      \item
            Attività antalgica
      \item
            Altre attività riabilitative
\end{itemize}

Queste attività si svolgono prevalentemente in \textbf{piscine
      didattiche} o vasche di dimensioni ridotte.

\begin{quote}
      \textbf{Attenzione}: un istruttore di nuoto con la sola qualifica di
      istruttore \textbf{non può operare autonomamente
            nell\textquotesingle area salute}. È necessario avere qualifiche e
      conoscenze specifiche aggiuntive.
\end{quote}

\subsubsection{4. Area Benessere}\label{area-benessere}

Comprende le attività complementari che contribuiscono al benessere
della persona:

\begin{itemize}
      \tightlist
      \item
            Nuoto libero
      \item
            \emph{Campus} estivi
      \item
            Attività integrative alle altre aree
\end{itemize}

\begin{center}\rule{0.5\linewidth}{0.5pt}\end{center}

\subsection{L\textquotesingle Organigramma del Centro
      Natatorio}\label{lorganigramma-del-centro-natatorio}

L\textquotesingle{}\textbf{organigramma} definisce la struttura di
competenze e ruoli all\textquotesingle interno
dell\textquotesingle azienda-piscina. Si trova nella \textbf{Dispensa n.
      4} del manuale FIN.

\subsubsection{Struttura Gerarchica}\label{struttura-gerarchica}

\begin{verbatim}
Proprieta (pubblica o privata)
    +-- Direzione
            +-- Coordinatore
            |       +-- Istruttori
            |       +-- Allenatori
            +-- Reception
            +-- Assistenti bagnanti
            +-- Manutentori / Ditte esterne
\end{verbatim}

\subsubsection{Competenze della
      Direzione}\label{competenze-della-direzione}

\begin{itemize}
      \tightlist
      \item
            Rappresenta e amministra la società
      \item
            Definisce le \textbf{linee principali della gestione}
      \item
            Seleziona e coordina: coordinatore, personale amministrativo,
            assistenti bagnanti, ditte esterne e professionisti
\end{itemize}

\subsubsection{Competenze del
      Coordinatore}\label{competenze-del-coordinatore}

Il coordinatore è la figura di riferimento diretta per gli istruttori.

\begin{itemize}
      \tightlist
      \item
            Seleziona e coordina il \textbf{personale tecnico} (istruttori)
      \item
            Compone i \textbf{gruppi} e gestisce gli spostamenti degli allievi tra
            i livelli
      \item
            Tiene i \textbf{rapporti con i genitori}
      \item
            Gestisce l\textquotesingle{}\textbf{amministrazione a piano vasca}
      \item
            Mantiene il raccordo con la \textbf{reception}
      \item
            Pianifica e promuove l\textquotesingle{}\textbf{aggiornamento tecnico}
            degli istruttori
      \item
            Supervisiona il \emph{piano vasca} e tutto ciò che avviene
\end{itemize}

\begin{quote}
      Reception e coordinatore devono lavorare in \textbf{grandissima
            sinergia}, interfacciandosi continuamente.
\end{quote}

\subsubsection{Competenze
      dell\textquotesingle Istruttore}\label{competenze-dellistruttore}

\begin{itemize}
      \tightlist
      \item
            Conduce le \textbf{lezioni} dopo averle pianificate con il
            coordinatore
      \item
            Definisce e persegue gli \textbf{obiettivi intermedi} di apprendimento
      \item
            Mantiene la \textbf{sicurezza} su vasca (controllo costante di tutti
            gli allievi)
\end{itemize}

\subsubsection{Competenze
      dell\textquotesingle Allenatore}\label{competenze-dellallenatore}

\begin{itemize}
      \tightlist
      \item
            Seleziona, avvia e prepara gli \textbf{atleti}
      \item
            Gli atleti provengono dalla scuola nuoto: il percorso nasce
            dall\textquotesingle istruttore, passa per il coordinatore e arriva
            all\textquotesingle allenatore
\end{itemize}

\subsubsection{Competenze della
      Reception}\label{competenze-della-reception}

\begin{itemize}
      \tightlist
      \item
            Accoglie la clientela
      \item
            Fornisce informazioni \textbf{precise e puntuali}
      \item
            Indirizza gli utenti verso le discipline adatte (sotto supervisione
            del coordinatore)
      \item
            Gestisce gli \textbf{adempimenti amministrativi} e le operazioni di
            cassa
\end{itemize}

\begin{center}\rule{0.5\linewidth}{0.5pt}\end{center}

\subsection{Distribuzione delle Attività nella
      Vasca}\label{distribuzione-delle-attivituxe0-nella-vasca}

La pianificazione del \emph{piano vasca} dipende dalla configurazione
dell\textquotesingle impianto:

\textbf{Impianto con vasca didattica e vasca grande:}

\begin{itemize}
      \tightlist
      \item
            Vasca didattica: corsi prescolare e bambini piccoli
      \item
            Vasca grande (es. 6 corsie): corsia laterale per ambientamento, corsie
            centrali per livelli progressivamente più avanzati
\end{itemize}

\textbf{Impianto con sola vasca grande:}

\begin{itemize}
      \tightlist
      \item
            Corsia laterale per ambientamento
      \item
            Ultima corsia eventualmente dedicata al fitness in acqua in
            contemporanea
      \item
            Corsie intermedie per i vari livelli
\end{itemize}

\begin{center}\rule{0.5\linewidth}{0.5pt}\end{center}

\subsection{Conclusione - Punti Chiave per
      l\textquotesingle Esame}\label{conclusione---punti-chiave-per-lesame-12}

\begin{itemize}
      \tightlist
      \item
            Il centro natatorio va gestito come
            un\textquotesingle{}\textbf{azienda} che fornisce servizi verso terzi,
            con un\textquotesingle organizzazione basata su divisione del lavoro e
            coordinamento.
      \item
            La \textbf{sicurezza} è la priorità assoluta: il rischio accettabile
            in piscina è \textbf{zero}.
      \item
            Sicurezza \textbf{attiva} = scuola nuoto; sicurezza \textbf{passiva} =
            assistenti bagnanti, manutenzioni, sorveglianza igienico-sanitaria.
      \item
            Le tre tipologie principali di incidente: \textbf{caduta sul bordo},
            \textbf{urto da tuffo}, \textbf{annegamento}.
      \item
            Parametri tecnici dell\textquotesingle acqua: pH \textbf{6,5-7,5};
            cloro libero \textbf{0,7-1,5 mg/L}; cloro combinato \textbf{max 4
                  mg/L}; temperatura \textbf{24-30°C}.
      \item
            Regola fondamentale: l\textquotesingle istruttore deve tenere
            \textbf{tutti gli allievi sotto controllo} costantemente e
            contemporaneamente.
      \item
            Le quattro aree di attività: \textbf{nuoto} (vincere la resistenza),
            \textbf{fitness} (cercare la resistenza), \textbf{salute}
            (riabilitazione, richiede qualifiche specifiche), \textbf{benessere}
            (nuoto libero, campus, attività integrative).
      \item
            L\textquotesingle organigramma: Proprietà → Direzione → Coordinatore →
            Istruttori/Allenatori; Reception e coordinatore in sinergia continua.
      \item
            Il coordinatore seleziona i tecnici, compone i gruppi, tiene i
            rapporti con i genitori, pianifica l\textquotesingle aggiornamento.
      \item
            Un istruttore \textbf{non può operare autonomamente
                  nell\textquotesingle area salute} con la sola qualifica di istruttore
            di nuoto.
      \item
            La filiera dell\textquotesingle atleta: nasce
            dall\textquotesingle{}\textbf{istruttore}, passa per il
            \textbf{coordinatore}, arriva
            all\textquotesingle{}\textbf{allenatore}.
\end{itemize}

\section{27. L\textquotesingle Avviamento all\textquotesingle Attività
  Sportiva}\label{lavviamento-allattivituxe0-sportiva}

\subsection{Introduzione}\label{introduzione-26}

L\textquotesingle{}\textbf{avviamento all\textquotesingle attività
      sportiva} in piscina è il percorso che collega la scuola nuoto al
preagonismo e all\textquotesingle agonismo. Quando la scuola nuoto è
avviata e alcuni allievi dimostrano caratteristiche e potenzialità
rilevanti, il coordinatore ha il compito di costruire o consolidare un
\textbf{vivaio}, ovvero una struttura di preagonismo che alimenti le
discipline federali presenti nell\textquotesingle impianto.

Le \textbf{discipline della Federazione Italiana Nuoto} presenti nelle
piscine sono:

\begin{itemize}
      \tightlist
      \item
            \textbf{Nuoto} (comprende anche il nuoto di fondo)
      \item
            \textbf{Salvamento}
      \item
            \textbf{Nuoto artistico}
      \item
            \textbf{Tuffi}
      \item
            \textbf{Pallanuoto}
\end{itemize}

Tutte queste discipline devono avere un percorso che nasce dalla scuola
nuoto, transita attraverso il vivaio e il preagonismo, e richiede una
\textbf{selezione} basata su obiettivi, criteri e strumenti precisi.

\begin{center}\rule{0.5\linewidth}{0.5pt}\end{center}

\subsection{Obiettivi e Finalità della
      Selezione}\label{obiettivi-e-finalituxe0-della-selezione}

\subsubsection{Definire Disciplina e
      Livello}\label{definire-disciplina-e-livello}

Il primo passo è stabilire \textbf{quale disciplina} e \textbf{a quale
      livello} ci si può permettere di sviluppare all\textquotesingle interno
dell\textquotesingle impianto. I fattori determinanti sono:

\begin{itemize}
      \tightlist
      \item
            \textbf{Competenza dei tecnici} che seguiranno i giovani atleti
      \item
            \textbf{Possibilità dell\textquotesingle impianto}: spazi acqua, spazi
            per attività a secco, disponibilità economiche, visione strategica
\end{itemize}

Esistono strutture che accompagnano i propri atleti fino a un certo
livello e poi li cedono a società più grandi. Questo sistema, lungi
dall\textquotesingle essere un fallimento, rappresenta un
\textbf{meccanismo virtuoso}: si garantisce ai giovani più talentuosi la
possibilità di continuare il proprio percorso in contesti più
attrezzati.

\subsubsection{Il Vivaio Comune e la Multilateralità
      Specifica}\label{il-vivaio-comune-e-la-multilateralituxe0-specifica}

La proposta più importante della FIN è quella del \textbf{vivaio
      comune}, fondato sul principio della \textbf{multilateralità specifica}.

La \emph{multilateralità specifica} in acqua consiste nel proporre agli
allievi selezionati esperienze motorie riconducibili a \textbf{tutte le
      discipline federali} (nuoto, salvamento, pallanuoto, nuoto artistico),
senza destinarli fin da subito a una sola specialità. I bambini svolgono
le loro esperienze acquatiche in modo integrato, attingendo alla
ricchezza motoria di ciascuna disciplina.

\textbf{Vantaggi del vivaio comune:}

\begin{itemize}
      \tightlist
      \item
            Fornisce ai bambini un \textbf{bagaglio motorio ampio} e diversificato
      \item
            Permette di \textbf{rimandare la scelta della disciplina}
            all\textquotesingle età preadolescenziale, quando il bambino ha
            maturato esperienza e può scegliere in base alle proprie attitudini
      \item
            Al momento della specializzazione, gli atleti saranno già tecnicamente
            più preparati
      \item
            Si evita la precoce specializzazione unilaterale, contraria ai
            principi dello sviluppo armonico
\end{itemize}

\begin{quote}
      In età preadolescenziale, dopo anni di esperienze multilaterali
      specifiche, il bambino sceglie la propria disciplina insieme agli
      allenatori. Lo farà con un bagaglio motorio talmente ampio da
      rappresentare una risorsa straordinaria per qualsiasi specialità.
\end{quote}

\subsubsection{Costruire il Gruppo e Comunicare con i
      Genitori}\label{costruire-il-gruppo-e-comunicare-con-i-genitori}

Una volta definita la strada didattica, occorre \textbf{costruire il
      gruppo} attraverso:

\begin{itemize}
      \tightlist
      \item
            Definizione di \textbf{obiettivi comuni}
      \item
            Stabilire \textbf{regole condivise}
      \item
            Favorire la \textbf{collaborazione} tra i membri fino a formare una
            vera squadra
\end{itemize}

La \textbf{comunicazione con i genitori} è fondamentale: vanno informati
chiaramente sugli obiettivi, sul rispetto delle fasi sensibili del
bambino, sull\textquotesingle orientamento a lungo termine del percorso.
Un genitore ben informato diventa una \textbf{risorsa} per il progetto,
non un elemento di pressione.

\begin{center}\rule{0.5\linewidth}{0.5pt}\end{center}

\subsection{Criteri di Selezione}\label{criteri-di-selezione}

Per individuare i bambini da inserire nel vivaio,
l\textquotesingle istruttore e il coordinatore devono osservare i
seguenti \textbf{requisiti}:

\subsubsection{1. Fattori Antropometrici ed
      Ereditari}\label{fattori-antropometrici-ed-ereditari}

Caratteristiche fisiche come altezza e struttura corporea. Importante: i
bambini più alti tendono spesso a essere \textbf{più scoordinati}, ma
non per questo devono essere esclusi. Anzi, frequentemente sono quelli
che a lungo termine offrono le maggiori soddisfazioni. Nessuna
discriminazione: questi fattori vanno tenuti presenti, non usati come
filtro esclusivo.

\subsubsection{2. Aspetti Coordinativi e
      Cognitivi}\label{aspetti-coordinativi-e-cognitivi}

La qualità dei prerequisiti coordinativi già acquisiti: capacità di
organizzare il movimento, di percepire il proprio corpo nello spazio, di
rispondere con precisione agli stimoli.

\subsubsection{3. Capacità di Apprendimento
      Rapido}\label{capacituxe0-di-apprendimento-rapido}

L\textquotesingle{}\textbf{intelligenza motoria} --- distinta
dall\textquotesingle intelligenza accademica --- si manifesta nella
capacità di acquisire esercitazioni nuove con estrema facilità. Un
bambino che assimila rapidamente nuovi gesti è un soggetto da
valorizzare.

\subsubsection{4. Senso del Ritmo}\label{senso-del-ritmo}

Il nuoto e le attività acquatiche richiedono \textbf{ritmo} e
\textbf{tempistica} precisi tra le coordinazioni. Il senso del ritmo è
un indicatore predittivo di buone potenzialità natatorie.

\subsubsection{5. Fattori Psicologici e
      Motivazionali}\label{fattori-psicologici-e-motivazionali}

Bambini che mostrano \textbf{piacere di stare in acqua} e
\textbf{piacere di stare in gruppo} sono naturalmente predisposti a un
percorso preagonistico. La motivazione intrinseca è un fattore
irrinunciabile.

\begin{center}\rule{0.5\linewidth}{0.5pt}\end{center}

\subsection{Strumenti di Selezione}\label{strumenti-di-selezione}

La FIN propone tre strumenti principali per identificare i talenti:

\begin{longtable}[]{@{}p{0.47\linewidth}p{0.47\linewidth}@{}}
      \toprule\noalign{}
      Strumento                  & Caratteristiche                                 \\
      \midrule\noalign{}
      \endhead
      \bottomrule\noalign{}
      \endlastfoot
      \textbf{Esami di brevetto} & Valutazione del miglioramento tecnico;
      proposti una o due volte l\textquotesingle anno; gratificano il bambino
      e i genitori; permettono al coordinatore di monitorare il lavoro in
      acqua                                                                        \\
      \textbf{Gare sociali}      & Momenti agonistici (non solo tecnici); rivelano
      bambini che nell\textquotesingle anonimato del corso non si
      distinguevano ma che emergono quando si mette in palio un premio;
      strumento fondamentale per individuare la \emph{verve} agonistica            \\
      \textbf{Leve agonistiche}  & Convocazione di bambini dello stesso anno di
      nascita per allenamenti extra; momento di osservazione collettiva;
      complementano gli altri due strumenti                                        \\
\end{longtable}

\begin{quote}
      Il coordinamento tra istruttori, coordinatore, allenatori delle varie
      discipline e la direzione tecnica è indispensabile per rendere virtuoso
      questo meccanismo di selezione.
\end{quote}

\begin{center}\rule{0.5\linewidth}{0.5pt}\end{center}

\subsection{Programmazione del Vivaio}\label{programmazione-del-vivaio}

Una volta selezionato il gruppo, la programmazione si articola in cinque
fasi:

\begin{enumerate}
      \tightlist
      \item
            \textbf{Valutazione delle risorse esistenti}: spazi acqua, tempo
            disponibile, possibilità di attività a secco, attrezzature
      \item
            \textbf{Definizione di obiettivi e tempi}: obiettivi tecnici,
            metodologici e di altro tipo, con momenti di verifica pianificati e
            possibilità di aggiustamenti \emph{in itinere}
      \item
            \textbf{Pianificazione di strumenti e metodi}: scelta di cosa serve e
            con quale metodo raggiungere gli obiettivi
      \item
            \textbf{Programmazione dei contenuti}: se si sceglie il vivaio comune,
            i contenuti devono includere gli aspetti multilaterali specifici di
            tutte le discipline
      \item
            \textbf{Verifiche e aggiustamenti}: stabilire quando e come verificare
            i risultati, con flessibilità di revisione
\end{enumerate}

Il programma di sviluppo deve tenere costantemente presenti \textbf{tre
      aspetti}:

\begin{itemize}
      \tightlist
      \item
            \textbf{Tecnico}
      \item
            \textbf{Fisiologico}
      \item
            \textbf{Psicologico}
\end{itemize}

\begin{center}\rule{0.5\linewidth}{0.5pt}\end{center}

\subsection{Sviluppo Tecnico nell\textquotesingle Età
      Prepuberale}\label{sviluppo-tecnico-nelletuxe0-prepuberale}

L\textquotesingle età prepuberale è il momento in cui si concentra la
formazione tecnica di base:

\begin{itemize}
      \tightlist
      \item
            \textbf{Massimo sviluppo della coordinazione}: finestra primaria tra i
            \textbf{4 e i 9 anni}, con il massimo del sistema nervoso centrale a
            \textbf{10-11 anni}
      \item
            I movimenti devono essere acquisiti in modo \textbf{grezzo, plastico e
                  non rigido}, perché i bambini in crescita dovranno riadattare tutto
            ciò che hanno imparato
      \item
            In questa fase è possibile apprendere la tecnica in forma
            \textbf{globale e variabile}
\end{itemize}

\begin{quote}
      L\textquotesingle immaturità del sistema nervoso centrale non supporta
      ancora l\textquotesingle apprendimento di gesti particolarmente
      raffinati: è controproducente forzare una tecnica troppo precisa in
      questa fase.
\end{quote}

\begin{center}\rule{0.5\linewidth}{0.5pt}\end{center}

\subsection{Sviluppo Fisiologico nell\textquotesingle Età
      Prepuberale}\label{sviluppo-fisiologico-nelletuxe0-prepuberale}

\subsubsection{I Codici di Allenamento
      FIN}\label{i-codici-di-allenamento-fin}

La Federazione Italiana Nuoto ha definito \textbf{codici di allenamento}
riconducibili ai meccanismi energetici:

\begin{longtable}[]{@{}p{0.30\linewidth}p{0.30\linewidth}p{0.30\linewidth}@{}}
      \toprule\noalign{}
      Codice      & Meccanismo            & Caratteristiche                           \\
      \midrule\noalign{}
      \endhead
      \bottomrule\noalign{}
      \endlastfoot
      \textbf{A1} & Aerobico              & Riscaldamento, recupero tra le serie; non
      adattativo                                                                      \\
      \textbf{A2} & Aerobico adattativo   & Resistenza aerobica con obiettivo di
      adattamento                                                                     \\
      \textbf{B1} & Soglia anaerobica     & Intensità maggiore;
      l\textquotesingle acido lattico è presente ma non si accumula
      (soglia)                                                                        \\
      \textbf{B2} & Potenza aerobica      & Intensità elevata; il meccanismo
      anaerobico lattacido entra in supporto                                          \\
      \textbf{C1} & Anaerobico lattacido  & Altissima intensità; il meccanismo
      anaerobico lattacido è protagonista                                             \\
      \textbf{C2} & Anaerobico lattacido  & Altissima intensità; variante del
      C1                                                                              \\
      \textbf{C3} & Anaerobico alattacido & Rapidità e altissima velocità;
      durata brevissima (massimo 8-10 secondi)                                        \\
\end{longtable}

\begin{quote}
      \textbf{Attenzione}: i codici FIN non vanno confusi con i meccanismi
      energetici, di cui sono espressione metodologica.
\end{quote}

\subsubsection{Codici Appropriati per i Bambini
      Prepuberi}\label{codici-appropriati-per-i-bambini-prepuberi}

Per i bambini prepuberali sono adatti \textbf{esclusivamente}:

\begin{itemize}
      \tightlist
      \item
            \textbf{A1} e \textbf{A2}: lavoro nel mondo aerobico
      \item
            \textbf{C3}: rapidità e alta velocità, naturale nei bambini (che
            vivono "a due velocità: piano o forte")
\end{itemize}

I codici B1, B2, C1 e C2 \textbf{non sono adatti} ai bambini prepuberi.

\subsubsection{Capacità Condizionali da
      Sviluppare}\label{capacituxe0-condizionali-da-sviluppare}

\begin{itemize}
      \tightlist
      \item
            \textbf{Resistenza}: incrementata attraverso lavori progressivamente
            più lunghi
      \item
            \textbf{Forza specifica}: a carico naturale, come supporto ai lavori
            di resistenza
      \item
            \textbf{Rapidità e velocità}: elementi naturali
            dell\textquotesingle età infantile, da supportare e rispettare
\end{itemize}

\begin{center}\rule{0.5\linewidth}{0.5pt}\end{center}

\subsection{Sviluppo Psicologico nell\textquotesingle Età
      Prepuberale}\label{sviluppo-psicologico-nelletuxe0-prepuberale}

L\textquotesingle età prepuberale presenta caratteristiche psicologiche
che il tecnico deve conoscere e sfruttare:

\begin{itemize}
      \tightlist
      \item
            \textbf{Grande capacità di apprendimento}: il bambino è un
            "assorbitore" naturale di esperienze motorie
      \item
            \textbf{Capacità di concentrazione}: già sviluppata e stimolabile in
            modo efficace
      \item
            \textbf{Percezione globale e analitica}: il bambino riesce già a
            cogliere sia il gesto nella sua interezza sia i dettagli
      \item
            \textbf{Motivazione basata sull\textquotesingle appartenenza al
                  gruppo}: sentirsi parte di una squadra è uno dei motori motivazionali
            più potenti in questa fascia d\textquotesingle età
      \item
            \textbf{Identità di gruppo}: i bambini si identificano nel gruppo e
            nelle sue regole condivise
\end{itemize}

\begin{quote}
      Ricordare sempre che si tratta di \textbf{individui in crescita}: la
      priorità assoluta è rispettare le fasi sensibili del loro sviluppo, non
      forzare la prestazione.
\end{quote}

\begin{center}\rule{0.5\linewidth}{0.5pt}\end{center}

\subsection{Conclusione - Punti Chiave per
      l\textquotesingle Esame}\label{conclusione---punti-chiave-per-lesame-13}

\begin{itemize}
      \tightlist
      \item
            Le discipline FIN presenti nelle piscine sono: \textbf{nuoto} (e nuoto
            di fondo), \textbf{salvamento}, \textbf{nuoto artistico},
            \textbf{tuffi}, \textbf{pallanuoto}.
      \item
            Il percorso va dalla scuola nuoto → vivaio → preagonismo → agonismo,
            con una \textbf{selezione} basata su obiettivi, criteri e strumenti.
      \item
            Il \textbf{vivaio comune} applica la \textbf{multilateralità
                  specifica}: bambini di discipline diverse si allenano insieme,
            sperimentando tutte le esperienze motorie acquatiche federali.
      \item
            La scelta della disciplina avviene in età \textbf{preadolescenziale},
            dopo anni di esperienze multilaterali.
      \item
            I \textbf{cinque criteri di selezione}: fattori antropometrici ed
            ereditari, aspetti coordinativi e cognitivi, capacità di apprendimento
            rapido, senso del ritmo, fattori psicologici e motivazionali.
      \item
            I \textbf{tre strumenti di selezione}: esami di brevetto (aspetto
            tecnico), gare sociali (aspetto agonistico), leve agonistiche
            (osservazione collettiva).
      \item
            La programmazione deve sempre tenere presenti tre aspetti:
            \textbf{tecnico}, \textbf{fisiologico}, \textbf{psicologico}.
      \item
            Il sistema nervoso centrale dei bambini si sviluppa tra i \textbf{4 e
                  i 9 anni} (finestra principale) con il massimo a \textbf{10-11 anni}.
      \item
            I movimenti vanno acquisiti in forma \textbf{grezza e variabile}, non
            rigida.
      \item
            I \textbf{codici FIN}: A1 (riscaldamento/recupero aerobico), A2
            (adattamento aerobico), B1 (soglia), B2 (potenza aerobica), C1/C2
            (anaerobico lattacido), C3 (anaerobico alattacido, brevissimo e ad
            alta velocità).
      \item
            Per i bambini prepuberi: \textbf{solo A1, A2 e C3}. I codici B1, B2,
            C1, C2 sono controindicati.
      \item
            Psicologicamente, il bambino prepubere ha già \textbf{capacità di
                  apprendimento, concentrazione, percezione globale e analitica}; la
            motivazione principale è l\textquotesingle{}\textbf{appartenenza al
                  gruppo}.
      \item
            I genitori vanno informati chiaramente: diventano \textbf{risorse} del
            progetto se la comunicazione è trasparente e orientata al lungo
            termine.
\end{itemize}

\section{28. Le Scuole Nuoto Federali}\label{le-scuole-nuoto-federali}

\subsection{Introduzione}\label{introduzione-27}

Il progetto \textbf{Scuole Nuoto Federali} (SNF) è uno dei progetti più
significativi della Federazione Italiana Nuoto. Nasce
all\textquotesingle indomani delle Olimpiadi di Sydney 2000, quando dopo
104 anni di tentativi infruttuosi il nuoto italiano conquista
\textbf{tre medaglie d\textquotesingle oro}, un argento e tre bronzi. Da
un giorno all\textquotesingle altro l\textquotesingle Italia si scopre
paese di nuotatori. La FIN studia il percorso didattico di quegli atleti
per estrarne un modello replicabile, con l\textquotesingle obiettivo di
creare un sistema di \emph{total quality} capace di favorire lo sviluppo
del talento sportivo in tutto il territorio nazionale.

\begin{quote}
      Se è vero che non necessariamente in una scuola nuoto ben organizzata
      sboccerà un grande talento sportivo, è assolutamente vero che un talento
      sportivo non può svilupparsi in una scuola nuoto \textbf{male
            organizzata}.
\end{quote}

Il progetto coinvolge oggi circa il \textbf{50\% delle società}
affiliate alla Federazione Italiana Nuoto.

\begin{center}\rule{0.5\linewidth}{0.5pt}\end{center}

\subsection{Il Progetto Scuole Nuoto
      Federali}\label{il-progetto-scuole-nuoto-federali}

\subsubsection{Elementi Fondanti}\label{elementi-fondanti}

Le società che aderiscono al progetto condividono quattro elementi
comuni:

\begin{enumerate}
      \tightlist
      \item
            Un \textbf{organigramma} con figure dirigenziali e tecniche formate
            presso la FIN
      \item
            Un \textbf{percorso didattico uniforme} organizzato su un minimo di 5
            livelli tecnici
      \item
            La \textbf{certificazione} dei livelli tecnici raggiunti dagli allievi
            (\emph{brevetti})
      \item
            Il mantenimento di \textbf{standard adeguati di sicurezza e qualità
                  ambientale}
\end{enumerate}

Aderire alle SNF non è un obbligo: molte società scelgono di rimanere
fuori dal progetto. Chi sceglie di aderire si riconosce in un modello
comune che ha generato benefici sia sportivi sia organizzativi.

\subsubsection{Benefici
      dell\textquotesingle Adesione}\label{benefici-delladesione}

Il beneficio più significativo è la possibilità di organizzare
autonomamente all\textquotesingle interno della propria struttura
\textbf{corsi di formazione per assistenti bagnanti e istruttori}.
Questo crea una \textbf{filiera formativa completa}:

\begin{itemize}
      \tightlist
      \item
            Corsi per gestanti
      \item
            Attività \emph{baby} e prescolare
      \item
            Scuola nuoto per bambini e ragazzi
      \item
            Attività agonistica
      \item
            Attività \emph{master}
      \item
            Formazione tecnica (istruttore a partire dai 16 anni, assistente
            bagnante dai 18)
\end{itemize}

Questa filiera permette di trattenere i ragazzi nella piscina anche
nella fascia d\textquotesingle età più critica (\textbf{14-25 anni}), in
cui chi non pratica agonismo tende ad abbandonare
l\textquotesingle attività acquatica.

\begin{center}\rule{0.5\linewidth}{0.5pt}\end{center}

\subsection{Il Percorso Didattico: i 5
      Livelli}\label{il-percorso-didattico-i-5-livelli}

Il percorso didattico delle SNF è organizzato su \textbf{5 macro livelli
      tecnici}:

\begin{longtable}[]{@{}p{0.30\linewidth}p{0.30\linewidth}p{0.30\linewidth}@{}}
      \toprule\noalign{}
      Livello    & Nome                         & Contenuto principale                   \\
      \midrule\noalign{}
      \endhead
      \bottomrule\noalign{}
      \endlastfoot
      \textbf{1} & Ambientamento                & Superamento della paura
      dell\textquotesingle acqua; prima relazione con
      l\textquotesingle ambiente acquatico                                               \\
      \textbf{2} & Autonomia in acqua fonda     & Trasferimento
      dell\textquotesingle ambientamento all\textquotesingle acqua profonda              \\
      \textbf{3} & Impostazione delle nuotate   & Apprendimento delle nuotate
      continue (\emph{crawl}, dorso, farfalla) e discontinue (rana)                      \\
      \textbf{4} & Perfezionamento              & Miglioramento tecnico delle nuotate in
      tutte le direzioni                                                                 \\
      \textbf{5} & Competenza motoria acquatica & Padronanza completa;
      inserimento di elementi da salvamento, nuoto artistico, pallanuoto                 \\
\end{longtable}

\begin{quote}
      Il principio di fondo: prima \textbf{ambientamento}, poi
      \textbf{continuità nelle nuotate}, poi \textbf{perfezionamento in tutte
            le direzioni}.
\end{quote}

Le società possono articolare ulteriormente questi 5 livelli (passando a
6, 7, 10 livelli interni), purché la logica progressiva resti invariata.
Diventare SNF non significa semplicemente assegnare cinque brevetti: è
il \textbf{concetto} che conta.

\subsubsection{Le Nuotate}\label{le-nuotate}

\begin{itemize}
      \tightlist
      \item
            \textbf{Nuotate continue} (\emph{continuous strokes}): \emph{crawl}
            (stile libero), dorso, farfalla
      \item
            \textbf{Nuotate discontinue} (\emph{discontinuous strokes}): rana
\end{itemize}

La rana è l\textquotesingle unico stile discontinuo e per questo nella
scuola nuoto bambini viene sempre insegnata \textbf{per ultima}: la
discontinuità del movimento è il principale ostacolo alla fluidità della
nuotata.

\begin{center}\rule{0.5\linewidth}{0.5pt}\end{center}

\subsection{La Certificazione degli
      Allievi}\label{la-certificazione-degli-allievi}

\subsubsection{Certificazione in Entrata: Gruppi Omogenei per
      Livello}\label{certificazione-in-entrata-gruppi-omogenei-per-livello}

Le SNF organizzano i gruppi per \textbf{livello tecnico omogeneo}, non
per anno di nascita (a differenza di quasi tutte le altre discipline
sportive terrestri). Gli allievi si spostano \textbf{orizzontalmente}
tra i gruppi:

\begin{itemize}
      \tightlist
      \item
            Quando un allievo migliora, viene spostato nel gruppo del livello
            superiore
      \item
            Dal gruppo precedente arriveranno nuovi allievi che hanno raggiunto
            quel livello
      \item
            L\textquotesingle istruttore lavora sempre con un gruppo
            corrispondente a un determinato livello tecnico
\end{itemize}

Questo sistema garantisce \textbf{efficacia didattica} e consente un
insegnamento più mirato.

\subsubsection{Certificazione in Uscita: il
      Brevetto}\label{certificazione-in-uscita-il-brevetto}

Il \textbf{brevetto} è l\textquotesingle attestato che certifica il
livello tecnico raggiunto dall\textquotesingle allievo. Ha due funzioni
pratiche:

\begin{enumerate}
      \tightlist
      \item
            Permette di reinserire correttamente l\textquotesingle allievo dopo la
            pausa estiva
      \item
            Se l\textquotesingle allievo si trasferisce, la piscina che lo
            accoglie (se è SNF) ha indicazioni utili per inserirlo nel gruppo
            giusto
\end{enumerate}

\textbf{Raccomandazioni per il rilascio:}

\begin{itemize}
      \tightlist
      \item
            Richiedere un \textbf{livello minimo di frequenza}: alcune famiglie
            portano i bambini solo nell\textquotesingle ultimo mese della stagione
            per far conquistare il brevetto; questa pratica snatura il brevetto,
            che deve essere un \emph{mezzo} per strutturare
            l\textquotesingle apprendimento, non un \emph{fine} in sé
      \item
            Le \textbf{modalità di rilascio} variano: spesso si organizzano saggi
            di fine anno in cui i bambini si esibiscono davanti alle famiglie
\end{itemize}

\begin{quote}
      I saggi e le prove di brevetto hanno \textbf{funzione dimostrativa}, non
      di esame. Se l\textquotesingle istruttore convoca gli allievi per il
      brevetto di un certo livello, è perché ha già la certezza che quel
      livello sia stato raggiunto. Anche se il bambino non performa al meglio
      quel giorno, il brevetto va consegnato ugualmente.
\end{quote}

\begin{center}\rule{0.5\linewidth}{0.5pt}\end{center}

\subsection{La Scuola Nuoto Adulti}\label{la-scuola-nuoto-adulti}

\subsubsection{Differenze Rispetto alla Scuola Nuoto
      Bambini}\label{differenze-rispetto-alla-scuola-nuoto-bambini}

La scuola nuoto per adulti richiede un approccio \textbf{completamente
      diverso}. I 5 macro livelli tecnici previsti per i bambini non si
applicano allo stesso modo, per una ragione fondamentale:

\begin{itemize}
      \tightlist
      \item
            Il bambino \textbf{supera} la paura dell\textquotesingle acqua: prende
            confidenza con l\textquotesingle ambiente e lo vive come
            un\textquotesingle estensione dello spazio terrestre
      \item
            L\textquotesingle adulto principiante ha avuto troppo tempo per
            \textbf{consolidare la paura}; non può superarla, ma può imparare a
            \textbf{controllarla}
\end{itemize}

L\textquotesingle adulto principiante, inoltre, non ha più la
disponibilità coordinativa necessaria per acquisire in modo fluido
movimenti complessi come quelli del nuoto.

\begin{quote}
      Con i bambini la lezione di gruppo è una lezione di gruppo. Con gli
      adulti una lezione di gruppo è un \textbf{insieme di tante lezioni
            individuali}: l\textquotesingle istruttore deve adattare il programma
      alle esigenze specifiche di ciascun allievo.
\end{quote}

\subsubsection{Percorso Didattico con gli
      Adulti}\label{percorso-didattico-con-gli-adulti}

Il percorso non è rigidamente strutturato come con i bambini, ma segue
questa sequenza pratica:

\begin{enumerate}
      \tightlist
      \item
            Imparare a \textbf{controllare la paura dell\textquotesingle acqua}
      \item
            \textbf{Stare a galla}
      \item
            \textbf{Ridurre le resistenze} all\textquotesingle avanzamento
      \item
            \textbf{Aumentare la propulsione}
\end{enumerate}

\subsubsection{La Rana come Primo Stile per gli
      Adulti}\label{la-rana-come-primo-stile-per-gli-adulti}

A differenza di quanto avviene con i bambini, con gli adulti lo
strumento più efficace è spesso partire dalla \textbf{rana}:

\begin{itemize}
      \tightlist
      \item
            È lo stile relativamente più istintivo
      \item
            È lo stile che un principiante impara più rapidamente
      \item
            Con gli adulti l\textquotesingle obiettivo non è costruire una
            motricità acquatica completa come con i bambini, quindi la
            discontinuità è un problema secondario
\end{itemize}

Nella scuola nuoto bambini, invece, la rana viene tenuta sempre per
\textbf{ultima}, proprio perché la discontinuità del movimento è il
principale ostacolo a una nuotata fluida.

\subsubsection{La Ginnastica in Acqua come Strumento di
      Ambientamento}\label{la-ginnastica-in-acqua-come-strumento-di-ambientamento}

Per adulti molto impauriti o con forti limitazioni fisiche, la
\textbf{ginnastica in acqua} può rappresentare un percorso alternativo
di ambientamento, prima di affrontare la scuola nuoto tradizionale. Il
progetto SNF include esplicitamente il \textbf{fitness in acqua} tra le
attività previste dal manuale.

\begin{center}\rule{0.5\linewidth}{0.5pt}\end{center}

\subsection{Il Manuale della Scuola Nuoto
      Federale}\label{il-manuale-della-scuola-nuoto-federale}

Tutti i concetti e le linee guida del progetto SNF sono disponibili nel
\textbf{manuale della Scuola Nuoto Federale}, consultabile sul sito
\textbf{federnuoto.it}. Il manuale contiene:

\begin{itemize}
      \tightlist
      \item
            Gli esercizi minimi richiesti per il passaggio tra i livelli
      \item
            I criteri di certificazione
      \item
            Le tipologie di attività di fitness in acqua e relative esercitazioni
      \item
            Le modalità operative per aderire al progetto (soggette a modifiche
            annuali)
\end{itemize}

\begin{center}\rule{0.5\linewidth}{0.5pt}\end{center}

\subsection{Conclusione - Punti Chiave per
      l\textquotesingle Esame}\label{conclusione---punti-chiave-per-lesame-14}

\begin{itemize}
      \tightlist
      \item
            Il progetto SNF nasce dopo le \textbf{Olimpiadi di Sydney 2000}: 3
            ori, 1 argento, 3 bronzi; obiettivo: replicare il modello didattico
            dei campioni.
      \item
            Coinvolge circa il \textbf{50\% delle società FIN};
            l\textquotesingle adesione non è obbligatoria.
      \item
            Quattro elementi comuni: \textbf{organigramma FIN-formato},
            \textbf{percorso didattico uniforme} (min. 5 livelli),
            \textbf{certificazione degli allievi}, \textbf{standard di sicurezza e
                  qualità}.
      \item
            I \textbf{5 macro livelli}: ambientamento → autonomia in acqua fonda →
            impostazione nuotate → perfezionamento → competenza motoria acquatica.
      \item
            Nuotate \textbf{continue}: \emph{crawl}, dorso, farfalla. Nuotata
            \textbf{discontinua}: rana (sempre per ultima con i bambini).
      \item
            I gruppi nelle SNF sono organizzati per \textbf{livello tecnico
                  omogeneo}, non per anno di nascita. Gli spostamenti sono
            \textbf{orizzontali}.
      \item
            Il \textbf{brevetto} è mezzo, non fine: certifica il livello
            raggiunto; utile al rientro e al trasferimento tra piscine.
      \item
            I saggi di fine anno sono \textbf{dimostrativi}: il brevetto va dato
            se il livello è stato raggiunto, indipendentemente dalla performance
            del giorno.
      \item
            Con gli \textbf{adulti}: non si supera la paura
            dell\textquotesingle acqua, si impara a \textbf{controllarla}.
            Approccio \textbf{personalizzato}; la lezione di gruppo è un insieme
            di lezioni individuali.
      \item
            Con gli adulti la \textbf{rana} è spesso il primo stile da insegnare
            (più istintiva); al contrario, con i bambini va per ultima.
      \item
            Beneficio SNF: possibilità di organizzare internamente \textbf{corsi
                  per istruttori} (dai 16 anni) e \textbf{assistenti bagnanti} (dai 18
            anni), contrastando l\textquotesingle abbandono nella fascia 14-25
            anni.
      \item
            Il manuale SNF è disponibile su \textbf{federnuoto.it}.
\end{itemize}

\end{document}
